% 本文件由 LaTeX 排版 Agent 根据有效 translation.json 自动生成。
% author=Cyprian of Carthage; work_id=0507; title=论集

\chapter{论集}

% structure: p0001 source=h1 target=omitted reason=page-title-already-rendered-by-parent

\Emph{论一}\newline{}\Emph{论二}\newline{}\Emph{论三}\newline{}\Emph{论四}\newline{}\Emph{论五}\newline{}\Emph{论六}\newline{}\Emph{论七}\newline{}\Emph{论八}\newline{}\Emph{论九}\newline{}\Emph{论十}\newline{}\Emph{论十一}\newline{}\Emph{论十二（卷一）}\newline{}\Emph{论十二（卷二）}\newline{}\Emph{论十二（卷三）}\par

\SourceRecord{Translated by Robert Ernest Wallis.；Ante-Nicene Fathers；Vol. 5.；Edited by Alexander Roberts, James Donaldson, and A. Cleveland Coxe.；Buffalo, NY: Christian Literature Publishing Co.,；1886.；<http://www.newadvent.org/fathers/0507.htm>.}

% structure: p0001 source=h1 target=section reason=page-title

\section{论文一}

\Emph{论教会的合一。}\par

论点：——鉴于诺瓦西安的分裂，为阻止迦太基人对他的倾向（他们因诺瓦图斯及教会中几位长老——他们是整个骚动的起源——而已经不反对他），西彼廉撰写了这篇论文。首先，坚固他们以对抗这些人的欺骗，劝勉他们持守恒心，并教导他们：异端之所以存在，是因为未仰望教会的元首基督，那首先委托给伯多禄的共同委托被藐视，唯一的教会和唯一的主教职被离弃。然后，他藉圣经以及旧约和新约的预像证明教会的合一。\par

1. 既然主警告我们说：\BibleQuote{你们是地上的盐，}\BibleRef{玛 5:13}；又命令我们要纯朴以致无害，却要在纯朴中谨慎，那么，亲爱的弟兄们，还有什么比以忧心忡忡的远见和警醒，去察觉并提防狡猾仇敌的诡计更合宜的呢？使我们这些已穿上基督——天主圣父的智慧——的人，在维护自己救恩的事上，不致显得缺乏智慧。因为不仅迫害是可怕的，那些公开攻击要压倒和推翻天主仆人的事，也不仅如此。危险显明时更加容易警惕，当敌人显身时，心志已预先准备好应战。但当敌人秘密潜入，以和平的外表欺骗，暗地里悄悄前行时，则更可怕，也更需防备——它也由此得名为蛇。这总是它的诡诈；这是它暗中诡计，用以欺骗人类。从世界起初它就如此欺骗；用谄媚的谎言，以不慎的轻信迷惑了无经验的灵魂。它也曾这样试探主自己：它秘密地接近他，好像要再次潜伏并欺骗；但它被识破并击退，因此被推翻，因为它被辨认和察觉了。\par

2. 由此我们得到一个榜样，要避免旧人的道路，立足于得胜的基督的足迹，免得我们再次不慎转回死亡的罗网，而要预见危险，拥有我们已经领受的永生。但我们怎能拥有永生，除非我们遵守基督那驱除并战胜死亡的命令？他自己警告我们说：\BibleQuote{你若要进入生命，就该遵守诫命。}\BibleRef{玛 19:17}又说：\BibleQuote{如果你们做我所命令的事，从此我不再称你们为仆人，而是朋友。}\BibleRef{若 14:15}最后，他称这些人为坚强和稳固的；他宣告他们建立在磐石上，拥有稳固的安全，以不可动摇、不可震动的坚固，对抗世上一切暴风和飓风。他说：\BibleQuote{凡听到我的话而实行的，我要把他比作一个聪明人，把自己的房屋建在磐石上：雨淋，水冲，风吹，袭击那房屋，它却不倒塌，因为它建基在磐石上。}\BibleRef{玛 7:24}因此我们必须坚守他的话，学习和实践他所教导和所做的一切。但一个人怎能说他相信基督，却不做基督命令他做的事？或者他从哪里获得信德的赏报，如果他不遵守诫命的信德？他必然摇摆不定，迷失方向，被错误之灵像风吹起的灰尘一样卷走；他在走向救恩的旅程中毫无进展，因为他没有持守救恩之途的真理。\par

3. 但是，亲爱的弟兄们，我们不仅必须提防公开而明显的事，也要提防以狡猾欺骗来迷惑人的事。还有什么比这更狡猾、更诡诈的呢？当基督来临时，这仇敌被察觉并推翻，光明已来到万民中，救恩的光线已照射出来以保全人类，使聋子得以领受属神恩宠的听力，盲者得以睁眼看见天主，软弱者得以因永恒的健康而重新强健，跛者得以跑向教会，哑者得以用清晰的声音和祈祷来祈祷——而他看见他的偶像被遗弃，他的小巷和庙宇因众多信众的聚集而荒废——于是他就发明了一种新的欺诈，假借基督徒之名的外衣来欺骗不慎之人。他发明了异端和分裂，借此颠覆信仰，败坏真理，分裂合一。那些他无法留在旧路的黑暗中的人，他就用新路的错误来包围和欺骗。他从教会本身夺走人；当他们似乎已经接近光明，逃脱了世界的黑夜时，他却在不知不觉中再次向他们倾泻新的黑暗；以致虽然他们不坚定地持守基督的福音，以及基督的观察和律法，他们仍称自己为基督徒，行走在黑暗中，却以为有光，而仇敌正在谄媚和欺骗，他按照宗徒的话，把自己装扮成光明的天使，装备他的仆役，好像他们是正义的仆役；他们谨守黑夜而非白昼，死亡而非救恩，绝望而非希望，背信弃义而非信德，假基督而非基督之名；以致他们假装貌似真理之事，却以诡诈使真理落空。亲爱的弟兄们，这之所以发生，是因为我们没有回归真理的源头，没有寻求元首，也没有持守天上师傅的教导。\par

4. 若有人思考并查考这些事，便无需冗长的讨论与论证。在真理的简短摘要中，便有信德的明证。主对\Person{伯多禄}说：\BibleQuote{我告诉你，你是伯多禄，在这磐石上，我要建立我的教会，地狱的门不能胜过她。我要把天国的钥匙交给你：凡你在地上所捆绑的，在天上也要被捆绑；凡你在地上所释放的，在天上也要被释放。}又对同一人，在祂复活后说：\BibleQuote{喂养我的羊。}虽然祂在复活后给所有宗徒同等的权柄，并说：\BibleQuote{就如父派遣了我，我也同样派遣你们：你们领受圣神！你们赦免谁的罪，谁的罪就得赦免；你们保留谁的罪，谁的罪就被保留。}\BibleRef{若 20:21}然而，为彰显合一，祂以权威安排了合一的起源，使之始于一人。当然，其余宗徒与伯多禄相同，享有同等的尊荣与权能；但开端来自合一。圣神也在\BibleBook{雅}{歌}中以吾主的名义指称这同一教会，说：\BibleQuote{我的鸽子，我的完全人，只有一个；她是她母亲的独一者，是生养她者的宠儿。}\BibleRef{歌 6:9}那不持守教会这合一的人，还自以为持守信德吗？那与教会争斗并反抗教会的人，还自信在教会内吗？更何况蒙福的宗徒\Person{保禄}也教导同样的事，并阐明合一的圣事，说：\BibleQuote{一个身体和一个圣神，你们蒙召同有一个希望；一主，一信，一洗，一天主？}\BibleRef{弗 4:4}\par

5. 这合一我们应当坚定持守并主张，尤其是我们这些在教会中担任主教的，好使我们也能证明主教职本身是合一而不可分的。不要让任何人用谎言欺骗弟兄团体；不要让任何人用背信的推诿败坏信德的真理。主教职是一，其中每一部分由每人为了整体而持守。教会也是一，她因丰产的扩展而散布于广大的众多。就如太阳有许多光线，却只有一个光；树有许多枝条，但力量基于坚固的根而为一；又如同从一个源头流出许多溪流，虽然多样性似乎因丰沛的慷慨而散布，然而合一仍保存于源头。将光线与光的本体分离，光的合一不允许分裂；从树上折下枝条——折下后，它便不能发芽；将溪流与泉源截断，被截断的部分便干涸。同样，教会，被主的光照耀，向全世界射出她的光线，却仍是那同一道光遍照各处，身体的合一并未分离。她丰产的丰富将枝条散布于全世界。她广阔地扩展她的河流，慷慨流淌，然而她的头是一，源头是一；她是一位母亲，在丰产的果效中丰盛：从她的子宫我们出生，由她的奶我们滋养，由她的圣神我们得生命。\par

6. 基督的新娘不可能是淫乱的；她是无玷而纯洁的。她只认识一个家；她以贞洁的谦逊守护一张床榻的神圣。她为我们保留给天主。她为天国指派她所生的儿子。谁若与教会分离而与淫妇结合，就是与教会的许诺分离；离弃基督教会的人也不能获得基督的赏报。他是外人，他是亵渎者，他是敌人。谁不以教会为母亲，便不能再以天主为父亲。如果任何人曾在\Person{诺厄}方舟之外得以逃脱，那么那在教会之外的人也将可以逃脱。主警告说：\BibleQuote{那不与我相合的，就是反对我；不与我聚集的，就是分散。}\BibleRef{玛 12:30}那破坏基督的平安与和谐的人，是在反对基督；那在教会之外聚集的人，分散基督的教会。主说：\BibleQuote{我与父原是一体。}\BibleRef{若 10:30}又论及圣父、圣子及圣神，经上记载：\BibleQuote{这三位是一体。}\BibleRef{若一 5:7}岂有人相信这来自神圣能力并在天上圣事内凝聚的合一，能在教会内被分裂，能因对立意志的分裂而被分离吗？那不持守这合一的人，不持守天主的法律，不持守圣父和圣子的信德，不持守生命与救恩。\par

7. 这合一的圣事，这不可分离地凝聚一致的纽带，在福音中揭示出来：主耶稣基督的外衣完全没有被撕裂或剪开，而是作为一整件衣服被接受，并被那些为基督的衣服抽签的人（他们更应该穿上基督）当作一件无损伤、无分裂的长袍拥有。圣经说：\BibleQuote{但是关于那件外衣，因为它不是缝制，而是从上到下纺织而成，他们彼此说：我们不要撕裂它，还是抽签看是谁的。}\BibleRef{若 19:23} 那件外衣带来了从上而来的合一，即从天主圣父而来，这合一绝不能被接受者和拥有者撕裂；相反，我们获得完整而实在的全体，毫无分离。谁分裂和分割基督的教会，就不能拥有基督的外衣。另一方面，在\Person{撒罗满}死后，他的王国和人民分裂时，先知\Person{阿彼雅}在田野遇见\Person{雅洛贝罕}王，将自己的外衣撕成十二片，说：\BibleQuote{你拿十片；因为上主这样说：看，我要从\Person{撒罗满}手中夺过王国，将十个权杖赐给你；两个权杖要归他，因我仆人\Person{达味}的缘故，并为\Place{耶路撒冷}，我所拣选安置我名的城。}\BibleRef{列上 11:31} 正如以色列十二支派被分裂，先知\Person{阿彼雅}撕碎了他的外衣。但因基督的人民不能被撕碎，他的长袍——从上到下交织合一——不被拥有它的人分裂；它不分裂、合一、连结，显示了我们穿上基督的人民的凝聚和谐。藉着他外衣的圣事和标记，他宣告了教会的合一。\par

8. 那么，谁是如此邪恶无信，谁是如此被不和的疯狂所迷，竟相信天主的合一可以被分裂，或竟敢撕裂主的衣裳——基督的教会？主自己在福音中警告我们并教导说：\BibleQuote{将有一个羊群，一个牧人。}\BibleRef{若 10:16} 难道有人相信在一个地方可以有许多牧人或许多羊群吗？此外，宗徒\Person{保禄}敦促我们这同一的合一，恳求和劝勉说：\BibleQuote{弟兄们，我因我们的主耶稣基督之名，恳求你们大家言谈一致，不要有分裂，但要同心同德，同一意向。}\BibleRef{格前 1:10} 他又说：\BibleQuote{彼此以爱德包容，努力保持圣神的合一，以和平的纽带相连。}\BibleRef{弗 4:3} 你若从教会退出，为自己建造别的家室和不同的住所，你以为自己能站立存活吗？那时曾对\Person{辣哈布}（她预表教会）说：\BibleQuote{你的父亲、母亲、兄弟和你父亲的全家，你都要聚集到你家中；凡走出你家门外的人，他的血就归在他自己头上。}\BibleRef{苏 2:19} 此外，在出谷纪的法律中，逾越节的圣事无非表明那预表基督而被杀的羔羊应在一家之内吃。天主说：\BibleQuote{应在一家之内吃，不可将肉带到屋外。}\BibleRef{出 12:46} 基督的肉和主的圣物不能被带到屋外，除了一所教会，信徒没有别的家。这一家，这同心合意的家庭，圣神在圣咏集中指明说：\Quote{天主使人们同心合意地住在一家。}在天主的家，在基督的教会中，人们同心合意地居住，继续在和谐与纯朴中：\par

9. 因此圣神也以鸽子的形态降临，这是一种纯朴而喜乐的受造物，不含苦胆的苦涩，不因咬噬而残忍，不以爪撕裂而暴烈；它喜爱人类的居所，认识同一家园的团聚；当它们有幼雏时，一起生育；当它们飞行时，也彼此相伴不离，共同度过一生，以喙的亲吻承认和平的和谐，在一切事上履行一致的法律。这就是在教会中应当认识的纯朴，这就是应当达到的爱德，使兄弟之爱效法鸽子，使他们的温良和顺像羔羊和绵羊。在基督徒胸中，狼的凶残有何作为？狗的野蛮、蛇的致命毒液和野兽的血腥残忍呢？当这样的人与教会分离时，我们应当庆幸，免得他们以残忍和有毒的感染毁灭基督的鸽子和羊群。苦涩不能与甘甜并存，黑暗不能与光明共存，雨水不能与晴朗共存，战争不能与和平共存，荒芜不能与丰收共存，干旱不能与泉源共存，风暴不能与宁静共存。谁也休想好人会离开教会。风不卷走麦子，飓风不拔掉扎根稳固的树。轻浮的稻草被暴风抛掷，软弱的树木被旋风的猛力推倒。宗徒\Person{若望}痛斥并严厉攻击这些人，他说：\BibleQuote{他们从我们中间出去，却不属于我们；如果属于我们，他们必会留在我们中间。}\BibleRef{若一 2:19}\par

10. 因此，异端不仅屡次产生，而且还在继续产生；因为扭曲的心不得安宁，而不和谐的背信者也不保持合一。但主允许并容忍这些事发生，同时保留人自由选择的权能，以便在真理的辨别考验我们的心灵和思想时，那些被认可者的纯正信德能以显明的光辉闪耀。圣神藉宗徒预先警告说：\BibleQuote{在你们中间也不免有异端，好让那些被认可的人在你们中间显明出来。}\BibleRef{格前 11:19} 如此，信友被认可，背信者被揭露；如此，甚至在此世，在审判之日以前，义人与不义者的灵魂已经分开，糠秕从麦子中分离出来。这些人擅自集会，没有任何神圣的安排，就在胆敢聚集的陌生人中自居为首，他们不凭任何授职律法自封为高级神长，自称为主教，却无人授予他们主教职；圣神在\WorkTitle{圣咏集}中指出他们坐在瘟疫的座位上，是信仰的瘟疫和污点，用蛇舌欺骗，诡诈地败坏真理，从瘟疫的舌头吐出致命的毒液；他们的言语像毒癌般蔓延，他们的谈论在每个人的心和胸中形成致命的毒药。\par

11. 主向这种人呼喊；祂阻止并召回祂迷途的子民，说：\Quote{不要听假先知的话，因为他们心中的幻象欺骗他们。他们说，但不是出于上主的口。他们向那些抛弃天主话的人说：你们必得平安，凡随自己私意行走的人，凡在心中的错误中行走的人，灾祸不会临于他。我没有对他们说话，他们却预言。如果他们曾站在我的基础上（\OriginalTerm{实质}{substantia}, \OriginalTerm{实质}{\GreekText{ὑποστα}}），……}\par

\SourceRecord{Translated by Robert Ernest Wallis.；Ante-Nicene Fathers；Vol. 5.；Edited by Alexander Roberts, James Donaldson, and A. Cleveland Coxe.；Buffalo, NY: Christian Literature Publishing Co.,；1886.；<http://www.newadvent.org/fathers/050701.htm>.}

% structure: p0001 source=h1 target=section reason=page-title

\section{论述二}

论童贞女的服饰\par

西彼廉赞扬纪律的益处，并从圣经证明其效用。接着，描述童贞及其向基督誓许并奉献的童贞的荣耀、尊荣与功绩，教导绝不只是在于肉身的洁净，亦在于服饰和妆饰的得体，甚至财富也不能成为那些已弃绝世俗者过度讲究衣着的借口。相反，既然宗徒甚至对已婚妇女规定衣着须合宜适度，那么童贞女更应遵守节制。因此，即使她富有，也应当考虑如何使用财富，但须用于善事，即天主所命令的事，也就是用于周济穷人。此外，他还禁止童贞女那些已轻率成为习惯的事，如参加婚礼以及前往混杂的浴场。最后，在简短的结语中，他宣告节德带来何种益处，缺乏它又带来何种恶果，以此结束本书。\par

1. \OriginalTerm{纪律}{Discipline}，希望之保障、信德之纽带、救恩之路的向导、善良品行的激励与滋养、德行之师，使我们常存于基督，不断为天主而活，并获得天上的恩许与神圣的赏报。追随她是有益的，背离并忽视她是致命的。圣神在圣咏中说：\Quote{你们应守纪律，免得主发怒，你们便从正途丧亡，因为祂的怒气很快就要向你们发作。}又说：但天主对恶人说：\Quote{你怎敢传述我的法律，将我的盟约挂在口边？你既然恼恨纪律，将我的话抛在脑后。}我们又读到：\BibleQuote{抛弃纪律的是可怜人。}\BibleRef{智 3:11}。从撒罗满我们领受了智慧的命令，警告说：\BibleQuote{我儿，不要轻看主的纪律，受祂谴责时也不要灰心；因为主所爱的，祂必惩戒。}\BibleRef{箴 3:11}。但如果天主惩戒祂所爱的人，并且正是为了改正他们而惩戒，那么弟兄们，尤其是司铎们，也不是憎恨，而是爱那些他们惩戒的人，以便改正他们；因为天主也曾藉耶肋米亚预言，并指向我们的时代，说：\BibleQuote{我要赐给你们合我心的牧者，以纪律的食物牧养你们。}\BibleRef{耶 3:15}\par

2. 然而，如果在圣经中纪律屡屡处处被规定，宗教与信德的全部基础来自服从与敬畏，那么，我们更应热切渴望、更加希望并持守什么，莫过于将根深深扎稳，将房屋建在磐石上，以坚固的基石抵御世界的风暴与旋风，好能藉神圣的诫命获得天主的赏报？我们考虑并知道，我们的肢体藉着生命之水的圣化，从旧染的一切污秽中洗净后，成为天主的宫殿，不可被侵犯或玷污，因为侵犯它们的人反使自己受害。我们是这些宫殿的朝拜者和司铎；让我们服从祂，我们已开始属于祂。保禄在他的书信中教导我们生活之道，说：\BibleQuote{你们不是自己的，因为你们是重价买来的；要在你们身上光荣并承载天主。}\BibleRef{格前 6:14}。让我们以纯洁、贞洁的身体，以更完全的服从，光荣并承载天主；既然我们已被基督的血所救赎，让我们以一切服事的服从，顺从并推进我们救赎者的统治，以免任何不洁或亵渎之物被带入天主的宫殿，免得祂发怒，离开祂所居住的宫殿。主赐予健康并教导的话，既治愈又警告：\BibleQuote{看，你已痊愈了：不要再犯罪，免得你遭遇更坏的事。}\BibleRef{若 5:14}。祂在赐予健康之后，又赐予生活的规范、无罪的法则，也不容许人此后放纵无羁地游荡，而是更严厉地警告那重又被同样的事所奴役的人，因为他曾从这些事中被治愈；因为无疑，在你还未认识天主的纪律之前犯罪是较小的过错；但认识天主之后犯罪，便不再有宽恕。的确，无论男女、男孩女孩，每个性别和年龄都应遵守这一点，并在这方面小心谨慎，按照他们对天主所欠的宗教与信德，以同样焦虑的敬畏保存从主的恩赐中圣洁领受的一切。\par

3. 我现在向童贞女发言，她们的荣耀更卓越，因而激起更大的兴趣。这是教会种子的花朵，神恩的优雅与装饰，喜乐的性情，赞美与尊荣的健康无瑕的工作，反映主之圣德的天主肖像，基督羊群中更光辉的部分。慈母教会凭借她们而欢欣于荣耀的丰饶，在她们身上繁荣昌盛；且童贞女数量越增多，母亲的喜乐就越增长。我向她们说话，以情感而非权力劝勉她们；并非我——最末最微，且深知自己的卑微——会要求任何训斥的权利，而是因为我始终谨慎甚至忧虑，更害怕撒殚的攻击。\par

4. 因为那种为救恩之道筹谋、遵守主与生命的诫命的谨慎并非徒劳，也不是虚妄的恐惧；这样，那些将自己奉献给基督、远离肉欲、并在肉体与精神上都已向天主誓许的人，可以完成他们的工作——这工作注定要获得巨大的赏报——并且不再致力于装扮或取悦任何人，只取悦他们的主，他们也从主那里期待贞洁的赏报；正如主亲自说：\BibleQuote{这话不是人人都能领悟的，只有那些得了恩赐的人才能领悟。因为有些阉人从母腹中生来就是这样；有些阉人是被人阉的；也有些阉人是为了天国的缘故自阉的。}\BibleRef{玛 19:11}同样，藉着天使的这句话，节制的恩赐也被展示出来，贞洁被宣讲：\BibleQuote{这些人没有与女人玷污自己，因为他们保持了童贞；这些人无论羔羊往哪里去，都跟随祂。}\BibleRef{默 14:4}因为主不仅这样向人应许节制的恩宠，而忽略了女人；但因为女人是男人的一部分，从男人取出并形成，天主在圣经中几乎总是对\OriginalTerm{原祖}{Protoplast}、即首先被造者说话，因为他们两人是一体，在男人身上同时也标志着女人。\par

5. 但如果节制跟随基督，贞洁是为天主的国预备的，那么她们与世俗的服饰和装饰有什么相干呢？当她们试图藉此取悦人时，反而得罪了天主。她们没有考虑到圣经宣告说：\BibleQuote{取悦人的人被羞辱，因为天主轻视了他们；}而保禄也荣耀而崇高地说：\BibleQuote{如果我还取悦人，我就不是基督的仆人了。}\BibleRef{迦 1:10}然而，节操与谦逊不仅在于肉身的纯洁，也在于衣着的端庄和装饰的朴实；这样，按照宗徒的话，未婚的女子可以在身体和灵魂上都成为圣洁的。保禄教导我们说：\BibleQuote{没有结婚的，挂虑主的事，怎样悦乐天主；但结了婚的，挂虑世俗的事，怎样悦乐自己的妻子。所以，童贞女和未婚的女子挂虑主的事，好使身心圣洁。}\BibleRef{格前 7:14}一个童贞女不仅应当是贞洁的，也应当被看见并被相信是贞洁的：任何人在看见一位童贞女时，不应怀疑她是否是贞洁的。完美应在一切事上显出相等的品质；身体的服饰不应诋毁心灵的善。她为什么打扮着出去呢？为什么装饰着头发，好像她有或寻求一个丈夫似的？她如果是一位童贞女，倒要害怕取悦人；如果她将自己保留给更崇高和神圣的事，就不要招惹自己的危险。那些没有丈夫、并声称要取悦天主的人，应当在身体和灵魂上都保持健全和纯洁。因为童贞女为了美貌的容貌而编发，或以肉身及其美丽自夸，是不合适的，因为她没有比抗击肉身更大的争战，也没有比征服和制服身体更顽强的战斗了。\par

6. 保禄宏亮而高声地宣布说：\BibleQuote{但我绝不夸耀，只夸耀我们主耶稣基督的十字架，藉着祂，世界于我已被钉在十字架上，我于世界也被钉在十字架上。}\BibleRef{迦 6:14}然而，教会中的一位童贞女却以她肉身的容貌和身体的美丽自夸！保禄接着说：\BibleQuote{凡属于基督耶稣的人，已把肉身连同邪情和私欲钉在十字架上了。}\BibleRef{迦 5:24}而那位宣称已弃绝肉身私欲和恶习的人，却恰恰被发现置身于她所弃绝的事物之中！童贞女，你被抓住了，你被暴露了，你声称一件事，却做出另一件事。你沾染了肉欲的污点，尽管你是贞洁和谦逊的候选人。\BibleQuote{呼喊吧，}主对依撒意亚说：\BibleQuote{凡有血肉的都像草，它的一切光荣都像田野的花；草会枯萎，花会凋谢；但天主的话永远存留。}\BibleRef{依 40:6}对任何基督徒，尤其是对童贞女，都不应以任何肉身的荣耀和尊荣为夸耀，而只应渴望天主的话，拥抱那永远持续的好处。或者，如果她必须在肉身内夸耀，那么毫无疑问，当她因宣认主的名而受酷刑时；当一个女人显示得比酷刑更坚强时；当她为得冠冕而忍受火刑、十字架、刀剑或野兽时，那时她才该夸耀。这些才是肉身的珍贵宝石，这些才是身体更好的装饰。\par

7. 但是有些妇女富有，且财产丰裕，她们偏爱自己的财富，主张她们应当享用这些福祉。要让她们首先知道，那在天主内富有的才是富有；那在基督内富有的才是富足；那些福祉是属灵的、属天主的、属天的，它们引导我们到天主面前，它们与我们同在，永远在天主内拥有。但是凡属地上的事物，在现世所领受的、将随世界留在这里的，都应当被轻视，如同世界本身被轻视一样；世界的浮华和享乐我们已藉着有福的经过来到天主时放弃了，\Person{若望}激励并劝勉我们，以属灵和属天的声音作证说：\BibleQuote{不要爱世界，也不要爱世界上的事物。若有人爱世界，圣父的爱就不在他内。因为凡在世界上的，就是肉身的欲望、眼目的欲望、生活的骄傲，都不是出于圣父，而是出于世界的欲望。世界和它的欲望都要过去；但那履行天主旨意的，永远存留，如同天主也永远存留一样。}\BibleRef{若一 2:15} 因此，应当追求永恒和属天主的事物，一切都必须按天主的旨意行事，好使我们追随我们主的神圣足迹和教导，他曾警告我们，说：\BibleQuote{我从天降下，不是为承行我的旨意，而是为承行派遣我来者的旨意。}\BibleRef{若 6:38} 但如果仆人不大于他的主人，得自由的人应当服从他的解救者，那么我们这些愿意作基督徒的人就应当效法基督所言所行。经上记载、诵读、听闻，并由教会之口宣扬以为我们的榜样，\BibleQuote{那说自己住在基督内的，就该照样去行，如同祂行过的一样。}\BibleRef{若一 2:6} 因此，我们必须步调一致地行走；我们必须以竞相赛跑的步履行走。这样，对真理的追随就与对我们名号的信仰相符，并且相信者将得到赏报，如果所信的也实践出来。\par

8. 你自称富有和富裕；但是\Person{保禄}以他的声音规定了妳的服装和装饰应适度，限定在合理的范围内。\BibleQuote{女人应以羞怯和节制来装饰自己，不可用编发、金饰、珍珠或昂贵的衣服，而要以善行作为自称保持贞洁的妇女所应有的。}\BibleRef{弟前 2:9} \Person{伯多禄}也同意这些训诫，说：\BibleQuote{女人的装饰不应是外面的发型、金饰或衣服，而是内心的装饰。}\BibleRef{伯前 3:3} 但如果这些也警告我们，那些惯于以丈夫为由为她们的穿着作辩解的女人，应当因宗教的遵守而受教会的纪律约束和限制，那么，处女更应遵守这一纪律，因为她没有理由装饰自己，也不能把她的过错的欺骗性归咎于别人，她本人仍处于罪的过犯中！\par

9. 你说你富有和富裕。但并非所有能做的事都应当做；那些出于世俗骄傲的广阔欲望也不应超越处女的高贵和谦逊；因为经上记载：\BibleQuote{凡事都可行，但并非都有益处；凡事都可行，但并非都能建树。}\BibleRef{格前 10:23} 此外，如果你奢侈地梳理头发，走路时引人注目，吸引年轻人的目光，引起青年男子的叹息，滋养贪欲的私欲，点燃叹息的燃料，以致虽然你自己没有丧亡，却使别人丧亡，你把自己像刀剑或毒药一样提供给观看者；你不能以内心纯洁谦逊为借口而得到宽恕。你可耻的衣着和放荡的装饰控告你；你再也不能被算在基督的贞女和处女中，因为你活着的方式使自己成为欲望的对象。\par

10. 你说你富有和富裕；但处女不应夸耀自己的财富，因为圣经说：\BibleQuote{骄傲为我们有什么用处？财富的夸耀为我们带来什么好处？这一切都像影子一样过去了。}\BibleRef{智 5:8} 宗徒又警告我们说：\BibleQuote{购买的，要像没有买；拥有的，要像没有拥有；使用这世界的，要像没有使用。因为这世界的局面正在逝去。}\BibleRef{格前 7:30} \Person{伯多禄}，主曾将祂的羊托付给他喂养和守护，并将教会安置建立在他身上，他确然说自己没有金银，但他说他在基督的恩宠中是富有的——他在信德和德行上是富足的——他以此行了许多伟大的奇迹，以此在属灵的祝福上丰富，以致于荣耀的恩宠。这些财富，这种富足，她不能拥有，因为她宁愿在这世界富有，而不愿在基督内富有。\par

11. 你说你富有和富裕，并认为你应该使用天主愿意你拥有的那些事物。当然可以使用，但要用于救恩之事；可以使用，但要用于善目的；可以使用，但要用于天主所命令、主所规定的事。让穷人感觉到你是富有的；让贫困者感觉到你是富足的。把你的产业借给天主；给基督食物。藉许多人的祈祷打动\Emph{祂}，赐你完成贞洁的荣耀，并得以领受主的赏报。在那里托付你的财宝，那里没有贼挖掘，没有阴险的强盗闯入。为你自己准备财产；但让它们最好是天上的，那里既不锈蚀，也不被冰雹损伤，不被太阳晒坏，不被雨水毁坏你的果实，常存而持久，完全不受世俗伤害的沾染。因为就在这件事上，你在得罪天主，如果你以为祂赐给你财富是为了让你尽情享受，而不顾救恩。因为天主也赐给人声音；然而并不因此就可以唱情歌和下流的东西。天主准许铁用于耕种土地，但不可因此就进行谋杀。或者因为天主规定了乳香、酒和火，我们就因此向偶像献祭吗？或者因为你的田里牲畜成群，你就应当向诸神献祭牺牲和供物吗？否则，大产业就是诱惑，除非财富用于善举；这样，每个人按他的财富，应当藉他的遗产赎罪，而不是增加罪过。\par

12. 装饰、服装和美貌的诱惑，只适合妓女和不端庄的妇人；没有什么人的衣服比那些谦卑者的衣服更珍贵。因此，在圣经中，主愿意我们藉以得到教诲和劝诫，那淫荡的城市被描述为装饰华丽、佩戴饰物，而且正是因为那些饰物将要灭亡。经上说：\BibleQuote{有一位天使，拿着七个盂的，来同我说话，说：你来，我要指给你看那坐在众水之上的大淫妇的审判，地上的君王都同她行淫。他在神魂中带我出去；我看见一个女人骑在一只兽上，那女人穿着紫红色和朱红色的衣服，饰有金、宝石和珍珠，手里拿着金杯，盛满了可憎之物和污秽以及全地的淫荡。} \BibleRef{默 17:1} 贞洁和端庄的贞女要避免不贞洁之人的服装、不端庄之人的举止、妓院的标志、娼妓的饰物。\par

13. 此外，充满圣神的依撒意亚大声呼喊并责斥熙雍的女子，她们被金、银和衣服所败坏，斥责她们沉溺于有害的财富，因世俗的享乐而离开天主。他说：\BibleQuote{熙雍的女子高傲，挺颈而行，抛媚眼，曳步而行，脚上叮当作响。因此，主将降卑熙雍的女子，上主要暴露她们的衣裳；主必除掉她们华美的服饰、项链、发网、月牙圈、耳环、手镯、面纱、发饰、腿饰、华带、香盒、符囊、戒指、鼻环、礼服、披肩、斗篷、手袋、镜子、细麻衣、头巾、蒙身帕。必有臭烂代替馨香，绳子代替腰带，光秃代替美发。} \BibleRef{依 3:16} 天主责备这些，指出这些；由此祂宣告贞女败坏了；由此，她们离开了真正的、神圣的敬拜。她们高傲，便跌倒了；她们头戴装饰，却得到了耻辱和羞耻。她们穿上丝绸和紫色，就不能穿上基督；饰有金、珍珠和项链，就失去了心灵和精神的装饰。谁不憎恶并躲避那导致别人毁灭的东西呢？谁愿意拿取并接受那曾作为杀死别人的刀剑和武器呢？如果一个人因喝干杯中的酒而死，你就会知道他所喝的是毒药；如果一个人吃了食物就死了，你就会知道他吃的是致命的食物；你就不会吃或喝那曾使别人丧命的东西。然而，对真理何等无知，何等疯狂的思想，竟渴望那既伤害过并将永远伤害的东西，并以为你自己不会因那你知道已使别人丧命的方式而灭亡！\par

14. 因为天主既没有把羊毛染成猩红或紫色，也没有教导用草汁和贝壳染羊毛并着色，更没有安排用金子镶嵌宝石、用珍珠串成连续的串或众多的簇来装饰项链，让你遮盖祂所造的颈项；这样，天主在人身上所造的可以被掩盖，而魔鬼额外发明的东西却显现在上面。天主曾愿意在耳朵上穿孔，使尚在襁褓、对世俗之恶无知的婴孩遭受痛苦，以致后来从耳朵上的伤疤和孔洞里挂上珍贵的珠子，尽管不重，但其价值却很重？所有这些，都是犯罪和堕落的天使用他们的技艺所发明的，当他们降落到地上的污染中时，他们就舍弃了属天的活力。他们还教给人用黑色描画眼圈，用欺骗的红色染颊，用虚假的颜色改变头发，用他们自己的腐败的猛攻逐出脸上和头上的一切真实。\par

15. 事实上，就此事而言，为了信德向我提出的敬畏，为了弟兄之爱所要求的爱德，我认为不仅贞女和寡妇，而且已婚妇女以及所有女性，都应受到劝诫：天主的化工及其塑造与形成，绝不应被掺假，无论是用黄色颜料、黑粉、胭脂，还是任何能败坏天生容貌的化妆品。天主说：\BibleQuote{让我们照我们的肖像，按我们的模样造人；}\BibleRef{创 1:26} 难道有人敢改变和变更天主所造的吗？他们设法重新塑造祂所塑造的，改变其形貌，却不知道一切受造之物都是天主的化工，凡是被改变的便是魔鬼的。如果任何一位画家在绘画时，以嫉妒的色彩描绘某人的面容、肖像和身体外貌；当这幅肖像画成并完成后，另一个人却动手修改它，仿佛它已成形并已画好，而此人更加高明，能够修正它，那么前一位画家自然会感到严重的侮辱和正当的愤慨。而你自认为可以免于惩罚地犯下如此邪恶妄为之罪，冒犯作为工匠的天主吗？因为即使你在人前并非不贞，你的诱惑性染料并未让你不贞，但是当天主所造之物被败坏和侵犯时，你正犯下更恶劣的奸淫。你认为自己被打扮，你认为你的头发被梳理，这是对天主之工的侵犯，是对真理的歪曲。\par

16. 宗徒警告的声音说：\BibleQuote{你们应把旧酵母除净，好使你们成为新和的面团，正如你们原是无酵饼一样；因为我们的逾越节羔羊基督，已被祭杀作了牺牲。所以我们过节，不可用旧酵母，也不可用奸诈和邪恶的酵母，而只可用纯洁和真诚的无酵饼。}\BibleRef{格前 5:7} 但是，当天真的东西被奸淫的颜色所污染，真实的东西被欺骗性的药物染料变为谎言时，真诚和真理还能保持吗？你的主说：\BibleQuote{你不能使一根头发变白或变黑；}\BibleRef{玛 5:36} 而你，为了克服你的主的话，竟比祂更有力，以大胆的尝试和亵渎的轻蔑染你的头发。怀着对未来的凶兆，你已为自己开始拥有火焰色的头发；并且用你的头——即你身体最高贵的部分——犯罪（哦，罪恶！）！虽然关于主写着：\BibleQuote{他的头与头发洁白如羊毛或雪，}\BibleRef{默 1:14} 你却诅咒那洁白，憎恨那类似主之头的白发。\par

17. 我恳求你，你既是如此，难道不怕当复活之日来临时，你的造物主不再认识你，当你来到祂的赏报和应许前时，祂将你赶走，把你排除在外，以监察官和法官的严厉斥责你说：\Quote{这不是我的工作，也不是我们的肖像。你以虚假的化妆品污染了你的皮肤，以奸淫的颜色改变了你的头发，你的脸被谎言强行占有，你的容貌败坏了，你的面貌是别人的。你不能看见天主，因为你的眼睛不是天主所造，而是魔鬼所破坏的。你跟随了他，你模仿了蛇的红色和涂画的眼睛。既然你以你仇敌的样式装饰自己，不久你也要与他一同焚烧。} 我请问，这些难道不是天主的仆人该反省的事吗？难道不是常常日夜该畏惧的吗？让已婚妇女自己看看，她们如何在讨好丈夫的安慰中自满，并以此作为借口，却使丈夫成了有罪同意之联合的伙伴。至于贞女，这篇讲话正是要向她们呼吁，那些用这类技巧装饰自己的贞女，我认为不应被算在贞女之中，而应如被感染的羊和病畜，从神圣纯洁的贞女群中赶出，免得她们因共处而以其传染病污染其余的人；免得她们败坏他人，就像她们已经自取灭亡一样。\par

18. 既然我们追求贞洁的益处，也让我们避免一切对它有害和敌对的事。我不会忽略那些因疏忽而渐渐流行、并已为自己窃取了许可的事，这些事违反端庄和清醒的举止。有些人不知羞耻地出席婚宴，在那淫荡谈吐的自由中，混入不洁的交谈，听到不当听的，说出不合法的，暴露自己，置身于可耻的言语和醉酒的宴席之中，由此欲火被点燃，新娘被鼓动去忍受，新郎去大胆行淫。在婚礼中，对于一位心不在婚姻的人，能有什么位置呢？或者，在那些与她的愿望和意愿都不同的场合中，她又能感到什么愉快或喜乐呢？在那里学到什么——看见什么？一位贞女多么严重地背离了她的决心，当她来时端庄，却走时不知羞耻！虽然她在身体和心灵上或许仍保持为贞女，但在眼、耳、舌上，她已削弱了她所持有的德行。\par

19. 但那些混迹于混杂浴室的人，他们将自己奉献给贞洁和端庄的身体，向好奇的眼目卖弄情欲，这又是怎样的人呢？他们可耻地注视裸体男人，也被男人看见裸体，难道他们自己不提供罪恶的诱惑吗？难道他们不勾引和邀请在场者的欲望，来败坏和伤害自己吗？\Quote{让每一个人，} 你说，\Quote{留意他来这里的心态：我只关心让我可怜的身体得到休息和清洗。} 这种辩护不能洗净你，也不能免除你放荡和淫乱的罪。这样的清洗是污秽的；它不能净化或清洁四肢，反而玷污它们。你没有不端庄地注视任何人，但你自己却被不端庄地注视。你没有以可耻的愉悦污染眼睛，但在愉悦他人时，你自己却被污染了。你使浴场变成一个展览；你聚集的地方比剧场更污秽。在那里，所有端庄都随着衣服的脱去而被抛弃，身体的尊贵和端庄被搁置；童贞被暴露，受人指点、被人触摸。现在，那么，想一想，当你穿衣时，你在男人中是否端庄，当裸体的放肆导致了不端庄时。\par

20. 因此，教会常常为她的贞女哀悼；为此，她为她们可耻可憎的传闻而叹息；为此，贞女之花熄灭，节操的尊贵和端庄受损，其一切荣耀和尊严被亵渎。这样，敌对的围攻者用他的诡计潜入；这样，借欺骗的罗网，借隐秘的道路，魔鬼爬了进来。这样，当贞女们希望被更精心地装饰，并更自由地漫游时，她们就不再是贞女，而被隐秘的耻辱败坏；在结婚之前就成了寡妇，对基督而非对丈夫犯了奸淫。她们原本作为贞女注定得到巨大的赏报，如今将因失去贞女身份而遭受巨大的惩罚。\par

21. 因此，贞女们，请听我，如同听一位父亲；请听，我恳求你们，一个在警告时怀着恐惧的人；请听一个忠实地为你们的利益和益处谋划的人。要成为天主造物主所造的那个样子；要成为你们天父的手所指定的那个样子。愿你们的容颜在你们内保持无瑕，你们的颈项不加装饰，你们的形体朴素；愿耳朵上不要穿孔，愿昂贵的手链和项链的链子不要环绕你们的手臂或颈项；愿你们的脚上没有金链，你们的头发不被染染，你们的眼睛配得看见天主。愿你们的沐浴只在女人中间进行，在她们中你们的沐浴是端庄的。要避免婚姻中的无耻盛宴和放荡筵席，其传染是危险的。克服衣着，因为你们是贞女；克服黄金，因为你们克服了肉体和世界。不能为大事所胜，却在小事上显为无能，这是不相称的。通往生命的道路是狭窄而艰难的；通往荣耀的路径是艰苦而困难的。殉道者沿此路前进，贞女由此经过，各类义人由此前行。要避开宽阔而宽敞的道路。那里有致命的陷阱和带来死亡的享乐；那里魔鬼谄媚，为了欺骗；微笑，为了行恶；引诱，为了杀害。为殉道者，初果是一百倍的；为你们，则是六十倍的。如同殉道者不考虑肉体和世界，没有微小、琐碎和娇嫩的对抗；同样，在你们内，你们的赏报在恩宠上居次，愿你们在忍耐的力量上仅次于他们。攀登伟大之事并非易事。当我们努力攀登山丘和山峰时，我们承受何等的辛劳，何等的劳苦！那么，为攀登天堂，又当如何呢？如果你们注视应许的赏报，你们的辛劳就轻省了。永生赐予那坚持到底的人，永恒的生命摆在面前；主应许了天国。\par

22. 坚持吧，贞女们！坚持你们所开始成为的；坚持你们将要成为的。巨大的赏报在等待你们，巨大的德行回报，贞洁的巨大益处。你们想知道节操的德行避免什么邪恶，拥有什么美善吗？\BibleQuote{我要加增，} 天主对女人说，\BibleQuote{你的痛苦和你的呻吟；你要在痛苦中生子；你的欲望应归于你的丈夫，他要管辖你。} \BibleRef{创 3:16} 你们摆脱了这判决。你们不惧怕女人的痛苦和呻吟。你们不害怕生育；也没有丈夫管辖你们；但你们的主和元首是基督，按照男人的形象并取代男人；\Emph{与男人相同}，你们的命运和状况是平等的。主的话说：\BibleQuote{这世界的儿女嫁娶；但那堪得来世和从死者中复活的人，不娶也不嫁；因为他们不能再死，他们相似天使，成为复活之子。} \BibleRef{路 20:35} 我们将要成为的，你们已经开始成为。你们已经在这世界拥有了复活的荣耀。你们经过这世界而不受世界的沾染；你们保持贞洁和童贞，就相似天主的使者。只要你们的童贞持续存在，坚固无伤；并且如它勇敢地开始，让它持续不断，不求项链和衣饰的装饰，而是求行为的装饰。让它仰望天主和天堂，而不垂下那仰望上界事物的眼睛，去注视肉体和世界的欲望，或专注于世俗之事。\par

23. 第一条命令是生育繁殖；第二条命令规定了节欲。当世界尚是混沌空虚之时，我们藉著繁育的果效而生养众多，并增加以扩展人类种族。如今，世界已被充满，大地已得供应，那些能领受节欲、像阉人那样生活的人，为了天国而成了阉人。主并不命令这事，而是劝勉；祂也不强加必然的轭，因为意志的自由选择被保留。但当祂说在祂父的家里有许多住所时，祂指出了更好居所的住处。你们正在寻求那些更好的居所；切断肉身的欲望，你们在天上的家乡获得更大恩宠的赏报。所有藉著圣洗圣事的圣化获得神圣恩赐和产业的人，在那拯救的水泉之恩宠中脱去旧人，并由圣神从旧传染的污秽中得以更新，藉第二次诞生而被净化。但那重复诞生的更大圣洁和真理属于你们，你们不再有肉身和身体的欲望。只有属于德行和圣神的事物留在你们内，以致荣耀。这是宗徒的话，主称他为拣选的器皿，天主派遣他宣布天主的命令：\BibleQuote{第一个人出自地，属于土；第二个人出自天。那属于土的怎样，凡属于土的也怎样；那属于天的怎样，凡属于天的也怎样。我们既然带了那属于土的肖像，也要带那属于天的肖像。} \BibleRef{格前 15:47} 童贞拥有这形像，纯洁拥有它，圣洁拥有它，真理拥有它。那些思念天主的纪律拥有它，持守正义与宗教，信心稳固，敬畏谦卑，面对一切苦难勇敢，忍受冤屈温柔，施予怜悯宽厚，在兄弟般的平安中同心合意。\par

24. 你们这些好贞女，应当遵守、热爱、实践这一切，你们将自己献给天主和基督，在更高更好的部分向主前进，你们已将自己奉献给祂。你们年长的，给年轻的提供教导。你们年轻的，激励你们的同龄人。以相互劝勉激发彼此；以德行的竞争激起荣耀。勇敢地忍受，属灵地前行，幸福地抵达。只在那时记念我们，当童贞开始在你们内获得赏报之时。\par

\SourceRecord{Translated by Robert Ernest Wallis.；Ante-Nicene Fathers；Vol. 5.；Edited by Alexander Roberts, James Donaldson, and A. Cleveland Coxe.；Buffalo, NY: Christian Literature Publishing Co.,；1886.；<http://www.newadvent.org/fathers/050702.htm>.}

% structure: p0001 source=h1 target=section reason=page-title

\section{论著三}

\Emph{论背教者。}\par

论点：——首先论及教会意外的平安，以及\OriginalTerm{宣信者}{Confessors}和那些在信仰中坚定不移者的坚毅；然后以极度悲伤指出\OriginalTerm{背教者}{Lapsed}的堕落，并揭示过去迫害的原因，即纪律的疏忽和信徒的罪过；我们的作者严厉谴责背教者，他们在敌人威胁的第一句话就已向偶像献祭，而没有按照基督的劝告退避。最后，他警告读者要躲避\OriginalTerm{诺瓦提安派}{Novatians}，用许多圣经经文驳斥他们的异端。\par

1. 看，亲爱的弟兄们，教会的平安已恢复；尽管它不久前对不信者而言似乎困难，对背教者而言不可能，但我们的安稳藉着天主的助佑和报应得以重建。我们的心回归喜乐；痛苦和阴云的时节消散后，宁静和晴朗再次照耀。必须赞美天主，必须以感恩的心庆祝祂的恩惠和恩赐，即使在迫害时期，我们的声音也没有停止感恩。因为仇敌没有足够的力量阻止我们——我们全心全意、全生命、全力爱主的人——始终在荣耀中到处宣扬祂的祝福和赞美。众人恳切渴望的日子已经来临；在漫长黑夜那可怕而可憎的黑暗之后，世界因主的光照耀而放射光芒。\par

2. 我们以喜悦的面容注视那些因美名的传报而显赫、因德行和信德之赞美而荣耀的宣信者；以神圣的亲吻依偎他们，以难以满足的热切拥抱长久渴望的他们。基督士兵的白衣军团在此，他们在激烈的冲突中打破了紧迫迫害的凶猛动荡，已为监禁的痛苦做好准备，为忍受死亡武装起来。你们勇敢地抵抗了世界：你们在天主面前提供了光荣的景观；你们成了追随你们的弟兄的榜样。那虔诚的口舌曾宣认它所相信的基督之名；那曾习惯从事神圣工作的光荣双手，抵抗了亵渎的祭献；那经过天上食粮、主的圣体圣血而圣化的嘴唇，拒绝了偶像的污秽接触和残余。你们的头摆脱了那用以遮盖献祭者被掳之头的邪恶恶毒的帕子；你们额上纯洁的天主印记，不能忍受魔鬼的冠冕，而是为主自己的冠冕保留。当你们从战场归来时，慈母教会多么喜乐地在她的怀中接纳你们！她多么有福、多么欣悦地敞开大门，好让你们结成队伍进入，带着从被击倒的仇敌得来的战利品！伴随着得胜的男性，也有女性前来，她们在与世界争战时也胜过了自己的性别；还有\OriginalTerm{贞女}{Virgins}带着双重战斗的荣耀前来，以及超越年龄的德行而成长的男孩。此外，其余那些坚守的人群也追随你们的荣耀，带着赞美的标志陪伴你们的脚步，与你们的非常接近，几乎与你们的相联。在他们身上也有同样真诚的心，同样持守信仰的健全。他们安息在天上诫命的不可动摇的根基上，藉着福音的传统得到坚固，指定的流放、注定的折磨、财产的损失、身体的惩罚，都没有惊吓他们。他们证明信仰的日子是预先限定的；但凡记得自己弃绝了世界的人，不知道世界的指定日子；凡从天主期望永生的，也不再计算地上的季节。\par

3. 我亲爱的弟兄们，愿没有人，没有人轻视这种荣耀；愿没有人以恶意的诽谤减损那些站立者不可玷污的坚定。当指定否认的日子过去后，凡在那期限内没有宣称自己不是基督徒的人，都承认自己是基督徒。胜利的第一头衔是，在暴力的\OriginalTerm{外邦人}{Gentiles}手下宣认主。荣耀的第二步是，以谨慎的退避被撤离，为主保留。前者是公开的宣认，后者是私下的宣认。前者战胜这世界的审判者；后者以天主为其审判者，在内心的完整性中保持纯洁的良心。前者有更现成的勇气，后者则有更安全的谨慎。前者，当他的时辰临近时，已显为成熟；后者也许被延迟了，他离开了自己的产业，暂时退避，因为他不会否认，但如果他也被捉拿，他一定会承认。\par

4. 一种悲哀的原因，使这些殉道者天上的冠冕、这些荣耀的属灵宣认、这些站立的弟兄们极其伟大的显赫美德感到忧伤，那就是：敌意的暴力已撕裂了我们自己内脏的一部分，并在其残忍的毁灭性中将其抛弃。亲爱的弟兄们，在这事上我该做什么？在种种情感的浪潮中摇摆不定，我该说什么，又该怎样说？我需要眼泪甚于言语，来表达我们身体的创伤应当被哀悼的悲痛，来表达一个曾经众多之民的多重损失应当被哀悼的悲痛。因为谁的心如此刚硬或残忍，谁如此不顾弟兄之爱，以至于在朋友的各种废墟和因一切卑污而毁容的哀伤遗物中，能够站立并保持眼睛干涸，而不在悲痛爆发时以眼泪而非言语表达他的呻吟？我悲伤，弟兄们，我与你们一同悲伤；我自身的完整和个人的健康并未欺骗我，使我的忧伤得以安抚，因为在羊群的创伤中，牧羊人首先受伤。我与每个人心连心，我分担忧伤和哀悼的沉重负担。我与哭泣者同哭，我与哀悼者同哀，我视自己与那些倒下的人一同倒下。我的肢体同时被狂怒敌人的箭矢击中；他们残忍的剑刺穿了我的心肠；我的心灵无法保持不受触动，不受迫害入侵的影响，在倒下的弟兄中；同情心也使我倒下。\par

5. 然而，亲爱的弟兄们，必须着眼真理的原因；可怕的迫害的幽暗漆黑不应如此蒙蔽心灵和情感，以至于再没有光亮和启迪，从中可以看见神圣的诫命。如果认清了灾难的原因，就立刻找到了创伤的补救。主愿意祂的家庭受到考验；因为长期的和平败坏了那神授给我们的纪律，天上的斥责唤醒了我们那正在减弱、我几乎要说沉睡的信德；尽管我们因自己的罪该受更重的惩罚，但最仁慈的主却如此缓和了一切，以至于所发生的一切与其说像是迫害，不如说像是考验。\par

6. 每个人都渴望增加自己的产业；忘记了信徒在宗徒时代所行或始终应当行的事，他们以无法满足的贪婪热忱，致力于增加自己的财产。在司铎中，没有宗教的热忱；在执事中，没有健全的信德：在他们的行为中没有怜悯，在他们的品行中没有纪律。男人，胡须被毁容；女人，肤色被染色：眼睛被伪造，不同于天主亲手所造；头发被虚假染色。狡诈的欺诈被用来欺骗单纯人的心，隐晦的含义被用来欺骗弟兄们。他们与不信者缔结婚姻盟约；他们将基督的肢体出卖给外邦人。他们不仅轻率地发誓，甚至更甚，发假誓；他们以傲慢的膨胀藐视那些管辖他们的人，以毒舌彼此诽谤，以顽固的仇恨彼此争吵。不少主教——他们本应为他人提供劝诫和榜样——轻视他们神圣的职分，成了世俗事务的代理人，离开了他们的宝座，遗弃了他们的子民，在外省游荡，追逐市场以谋取有利可图的商品，而弟兄们在教会中挨饿。他们寻求囤积金钱，以狡诈的欺骗夺取地产，通过多重高利贷增加收益。我们这样的人，因这类罪该受怎样的苦楚呢？甚至当天上的斥责已预先警告我们，并说：\Quote{如果他们离弃我的法律，不遵行我的典章；如果他们亵渎我的律例，不遵守我的诫命，我就要用杖责罚他们的过犯，用鞭责罚他们的罪孽？}\par

7. 这些事早已向我们宣告并预言了。但我们，忘记了所要求的法律和服从，因我们的罪行事，以致当我们藐视主的诫命时，我们藉着更严厉的补救措施来纠正我们的罪和考验我们的信德。而且我们确实最终没有归向对主的敬畏，以耐心和勇敢地承受这我们的纠正和神圣的考验。在威胁的敌人的最初言语之下，大多弟兄就背叛了他们的信德，并且跌倒，不是被迫害的冲击打倒，而是因自愿的失足而自己跌倒。我问你们，有什么闻所未闻的事，什么新事发生了，以至于好像遇到未知和意外的事，对基督的义务就应该以鲁莽的仓促解除？难道先知们从前和随后的宗徒们没有讲过这些事吗？他们充满圣神，不是预言了义人的苦难，以及外邦人常有的伤害吗？那永远以武器武装我们的信德、以来自天上的声音坚固天主仆人的圣经，不是说了：\BibleQuote{你要朝拜主，你的天主，并且只事奉祂一个？}\BibleRef{申 6:13}它不又显示了神圣愤怒的怒气，并预先警告惩罚的恐惧，当它说：\BibleQuote{他们敬拜了他们手指所造的；卑贱人俯首，尊贵人自卑，我必不宽恕他们？}\BibleRef{依 2:8}再者，天主发言说：\BibleQuote{那向任何神献祭，唯独不向主献祭的，应被消灭。}\BibleRef{出 22:20}随后在福音中，主用祂的言语教导，用祂的行为实现，教导应做的事，并实行祂所教导的一切，祂不是预先告诫我们一切现在所行和将来所行的事吗？祂不是预先为那些否认祂的人规定了永罚，为那些宣认祂的人规定了救恩的赏报吗？\par

8. 可叹啊！从这些人身上，——啊，悲惨！——这一切都消失了，并已从记忆中逝去。\newline{}他们确实没有等到被抓住就升天了，也没有等到被审问就否认了。\newline{}许多人在战斗之前就被征服，在攻击之前就倒下了。\newline{}他们甚至没有留下任何可为他们辩护的话，表明他们似乎是不情愿地向偶像献祭。\newline{}他们自愿跑向市场；自由地奔向死亡，仿佛他们以前就渴望如此，仿佛要抓住这个现在给予的机会，这是他们一直想要的。\newline{}当时有多少人被长官在傍晚时分推迟执行；又有多少人甚至要求不要延误他们的毁灭！\newline{}这样的人能引述什么暴力作为借口？\newline{}他如何能洗脱自己的罪行，当正是他自己使用暴力导致了自己的毁灭？\newline{}当他们自愿来到卡皮托利神庙，——当他们自由地接近那可怕邪恶的服从——他们的脚步没有蹒跚吗？\newline{}他们的视线没有昏暗吗？他们的心没有颤抖吗？他们的手臂没有无力地垂落吗？\newline{}他们的感官没有失灵吗？他们的舌头没有贴在上颚吗？他们的言语没有变得微弱吗？\newline{}天主的仆人怎能站在那里，说话并否认基督，而他已经否认了魔鬼和世界？\newline{}那祭台，他走近去灭亡，对他而言难道不是火葬柴堆吗？\newline{}他难道不该因魔鬼的祭台而战栗并逃离吗？他看见那祭台冒烟，散发着恶臭，如同他生命的葬礼和坟墓？\newline{}不幸的人啊，你为什么带着祭品前来？\newline{}为什么献祭牺牲？\newline{}你自己来到祭台成为供品；你自己成为牺牲：在那里你祭献了你的救恩、你的望德；你的信德在致命的火焰中焚烧了。\par

9. 然而，对他们许多人来说，自己的毁灭还不够。\newline{}人们相互激励，用劝言催促对方走向毁灭；死亡在致命的杯中轮流被许诺。\newline{}为了使罪行无所不包，连婴儿也被父母抱在怀里，或被领着或抱着，在幼小的时候就失去了他们在出生之初所获得的一切。\newline{}当审判之日到来时，他们岂不会说：\Quote{我们什么也没做；我们也没有离弃主的饼和杯，去自由地奔向亵渎的接触；是别人的背信弃义毁灭了我们。我们发现父母是我们的谋杀者；他们否认了教会为我们母亲；他们否认了天主为我们父亲：以致于当我们年幼、不能预见、且不知晓如此罪行时，我们被他人牵连到邪恶的合伙中，被别人的诡计所陷害？}\par

10. 可叹啊，并没有任何正当而有力的理由能开脱这样的罪行。\newline{}需要离开自己的国家，遭受财产的损失。\newline{}然而，有哪个生而必死的人不曾在某个时候必须离开他的国家，遭受财产的损失呢？\newline{}但不可离弃基督，以致于害怕失去救恩和永恒的家园。\newline{}看，圣神藉先知呼喊说：\BibleQuote{离开，离开，从那里出去，不要触摸不洁之物；从她中间出来，你们这些携带主器皿的人，要分别出来。}\BibleRef{依 52:11}\newline{}然而，那些是主的器皿和天主的殿宇的人，却不从中出来，也不离开，以免被迫触摸不洁之物，并被致命的食物玷污和败坏。\newline{}在别处也听到从天上来的声音，预先警告天主的仆人该做什么，说：\BibleQuote{我的子民，从她那里出来，免得你们分担她的罪过，也免得你们受她的灾祸。}\BibleRef{默 18:4}\newline{}那出去并离开的人不会成为罪过的分担者；但那被发现是罪行的同伴的人，必被灾祸击伤。\newline{}因此，主在迫害中命令我们离开并逃跑；祂既教导我们这样做，自己也这样做了。\newline{}因为如同冠冕是出于天主的俯就，除非接受它的时刻来到，否则不能得到，凡留在基督里的人，暂时离开并不否认他的信德，而是等待时机；但那些拒绝离开而留下来的人，跌倒后就是否认了。\par

11. 弟兄们，真相不可掩盖；我们伤口的根源和原因也不可隐藏。\newline{}对自己财产的盲目爱恋欺骗了许多人；当财富像锁链一样束缚他们时，他们无法准备好或从容地离开。\newline{}那些就是使他们留下的锁链——那些就是束缚，使德行受阻，信德受压，精神受绑，灵魂受羁绊；因此，那些被世俗之事缠住的人，就成为那按天主的判决以尘土为食的蛇的猎物和食物。\newline{}因此，主，善事的导师，为将来作警告，说：\BibleQuote{你若愿意成为成全的，去，变卖你所有的一切，施舍给穷人，你必有宝藏在天上；然后来跟随我。}\BibleRef{玛 19:21}\newline{}如果富人这样做了，他们就不会因财富而灭亡；如果他们积攒财宝在天上，就不会有一个内在的敌人和攻击者。\newline{}心、思想和感觉会在天堂，如果宝藏在天上；那在世上无物能使他被征服的人，就不会被世界征服。\newline{}他将自由无羁地跟随主，如同宗徒们，以及宗徒时代的许多人，以及许多舍弃了财物和亲属，以不可分割的纽带与基督结合的人。\par

12. 然而，那些被财富的锁链束缚的人，怎能跟随基督？或者，那些被世俗欲望压垮的人，怎能寻求天堂，攀登崇高且高远之处？他们以为自己拥有，实则反被拥有；作为利益的奴隶，而非自己金钱的主人，实为金钱的奴仆。宗徒所指的正是这时代和这些人，他说：\BibleQuote{但那想发财的人，就陷于诱惑，落入罗网，陷入许多愚蠢且有害的欲望，这欲望使人沉沦于败坏和灭亡。因为贪爱钱财是万恶的根源；有些人因贪恋钱财，就偏离了信德，使自己受了许多痛苦。}（见：\BibleRef{弟前 6:9}）然而，主以怎样的赏报邀请我们蔑视世俗财富？以怎样的补偿来弥补这现世微小且微不足道的损失？祂说：\BibleQuote{没有人为了天国的缘故，舍弃了房屋、田地、父母、兄弟、妻子或儿女，而不在现世得到百倍，并在来世获得永生。}（见：\BibleRef{谷 10:29}）如果我们知道这些事，并从那位应许的主的真理中领悟出来，那么不仅不应害怕这种损失，反而应渴望它；正如主再次宣告并警告我们：\BibleQuote{当人们为了人子的缘故，迫害你们，把你们从他们的团体中分离出来，将你们赶出去，并以恶言辱骂你们时，你们是有福的！在那一天，你们应当欢喜踊跃，因为看，你们在天上的赏报是丰厚的。}（见：\BibleRef{路 6:22}）\par

13. 但（他们说）随后酷刑来临，严厉的苦难威胁那些反抗的人。被酷刑打败的人才会抱怨酷刑；在苦难中被征服的人才会以苦难为借口。这样的人可以质问并说：\Quote{我确实希望勇敢奋战，并牢记我的誓言，我拿起了虔敬与信德的武器；但当我正在战斗中挣扎时，各种酷刑和无休止的苦难压倒了我。我的心思坚定，我的信德坚强，我的灵魂长久奋斗，在折磨的痛苦中毫不动摇；但是，当最残忍的法官以更新的野蛮手段，使我疲乏不堪时，鞭子撕裂我，棍棒击打我，刑架拉伸我，铁爪挖掘我，烈火炙烤我；我的肉体在战斗中弃我而去，我身体的软弱屈服了——不是我的心思，而是我的身体在受苦中让步。}这样的恳求很容易获得宽恕；这种解释可以激起怜悯。因此，主曾一时宽恕了\Person{Castus}和\Person{Aemilius}；他们虽在第一次交锋中被征服，却在第二次战斗中成为胜利者。因此，昔日屈服于烈火的人变得比烈火更强，而他们在其中被战胜的事上成了征服者。他们乞求怜悯，不是因眼泪，而是因伤口；不仅以哀哭的声音，更以身体的撕裂和痛苦。血代替了哭泣，从半焦的内脏中流出的不是眼泪，而是血液。\par

14. 但现在，那些被征服的人能展示什么伤口？什么敞开的脏腑裂口，什么四肢的酷刑？在那些不是信德在交锋中跌倒，而是不信提前发起斗争的情况下？犯罪的必要性并不能为被迫者开脱，因为犯罪是出于自由意志。我这样说，并非要加重兄弟们的负担，而是要激励弟兄们去作补赎的祈祷。因为经上记载：\BibleQuote{那些称你们为有福的人，使你们迷惑，毁坏你们脚下的路。}（见：\BibleRef{依 3:12}）以谄媚的甜言蜜语安慰罪人的人，正是给罪提供刺激；他没有阻止，反而滋养了恶行。但以更勇敢的劝告，同时责备并教导兄弟的人，则把他推向救恩。主说：\BibleQuote{凡我所爱的，我都要斥责和管教。}（见：\BibleRef{默 3:19}）因此，主的司铎也不应以欺骗性的让步误导人，而应提供有益的补救。他是一名拙劣的医生，用温柔的手处理伤口肿胀的边缘，并将毒液封闭在身体深处，从而使之加重。伤口必须切开、割除，用切除腐坏部分的更强补救来治愈。病人可能会因不耐痛苦而喊叫、喧嚷、抱怨；但当他感到自己痊愈时，他将会感激。\par

15. 再者，亲爱的弟兄们，一种新型的毁灭出现了；仿佛迫害的风暴还不够猛烈，在仁慈的名义下，又加上了一堆欺骗性的祸患和貌似美好的灾难。违背福音的刚毅，违反主和天主的法律，有些人轻率地给粗心大意的人放松了共融——这是一种虚妄而虚假的平安，对授予者有害，对接受者也毫无益处。他们不寻求健康所需的耐心，也不寻求来自补赎的真正药物。补赎被赶出了他们的胸怀，他们对极其严重和极端的罪的记忆也被抹去。垂死者的伤口被掩盖，深藏在内里的致命打击被伪装的痛苦所隐瞒。从魔鬼的祭台回来，他们带着肮脏且臭气熏天的手，几乎还散发着瘟疫般的祭偶像之肉的气味，靠近主的圣地；甚至他们的口还呼出罪恶，沾染着致命的接触，就闯入主的身体，尽管圣经挡在他们面前，呼喊说：\BibleQuote{凡是洁净的人都可以吃这肉；若有什么人带着不洁吃这拯救的祭肉，就是主的肉，这人必从民中剪除。}\BibleRef{肋 7:20} 此外，宗徒也作证说：\BibleQuote{你们不能喝主的杯，又喝魔鬼的杯；你们不能共享主的筵席，又共享魔鬼的筵席。}\BibleRef{格前 10:21} 他还威胁那些顽固和悖逆的人，谴责他们说：\BibleQuote{无论何人，不配地吃主的饼，喝主的杯，就是干犯主的身和主的血。}\BibleRef{格前 11:27}\par

16. 这一切警告都被蔑视和轻看——在他们的罪未得赎之前，在他们尚未宣认罪行之前，在他们尚未藉着祭献和司铎之手洁净良心之前，在激怒和威胁之主的冒犯未得平息之前，他们对主的身和血施以暴力；他们现在用手和口犯罪反对他们的主，比他们否认主时更甚。他们以为那是平安，而有些人用欺骗的话语正在大肆宣扬；那不是平安，而是战争；那离开福音的人并没有与教会联合。为什么他们把伤害称为仁慈？为什么他们把不敬称为虔诚？为什么他们阻止那些应当不断哭泣并恳求主的人进行悔罪的哀痛，却假装接纳他们共融？这对失足者而言，如同冰雹对庄稼、风暴之星对树木、瘟疫的毁灭对畜群、狂怒的暴风对航船。他们夺去了永恒希望的安慰；他们连根拔起了树；他们用瘟疫般的话爬向致命的传染；他们将船撞在岩石上，使它无法到达港湾。这样的轻易并不赐予平安，反而夺走平安；也不给予共融，反而阻碍救恩。这是另一种迫害，另一种试探，狡猾的敌人藉此进一步攻击失足者；以隐秘的腐败袭击他们，使他们的哀歌沉寂，使他们的悲伤沉默，使罪的记忆消失，使他们心中的呻吟被抑制，使他们眼中的泪水被擦去；长久而圆满的补赎也不再向如此严重冒犯的主祈求，尽管经上写着：\BibleQuote{你要回想你是从哪里跌落的，并要悔改。}\BibleRef{默 2:5}\par

17. 愿无人欺骗自己，愿无人自欺。唯有主能施仁慈。唯有祂能赦免得罪祂的罪，祂背负了我们的罪，为我们忧伤，天主将祂为我们的罪交出。人不能大于天主，仆人也无法凭自己的宽恕免除或赦免那对主犯下的更大罪行，免得这失足者罪上加罪，如果他不知道经上所宣布的：\Quote{凡依靠人的人，是可咒骂的。}必须恳求主。必须藉我们的补赎平息主，祂曾说，凡否认祂的，祂也要否认他；祂唯独从祂的父领受了审判的全权。我们确实相信，殉道者的功德和义人的行为在审判者面前大有功效；但那是在审判之日来临时，在此生和世界结束之后，祂的子民将站在基督的审判台前。\par

18. 但如果有人因过度急躁的匆忙，轻率地认为他能给所有人赦罪，或者胆敢废除主的诫命，这不仅对失足者毫无益处，反而伤害他们。不遵守祂的审判就是激起祂的义怒，并且认为不必首先恳求天主的仁慈，反而是藐视主，僭越祂的大能。在祭台之下，被杀的殉道者的灵魂大声呼喊说：\BibleQuote{圣洁真实的主啊，祢不审判世上的人，为我们伸流血的冤，要等到几时呢？}\BibleRef{默 6:10} 他们被告知要安息，并仍要忍耐。难道有人以为，与审判者相反，一个人能为普遍的罪赦和赦免而变得有用，或者在自己获得申冤之前就能庇护他人吗？殉道者命令做某事；但只有当这件事公正合法，且可以在不反对主自己的情况下由天主的司铎施行，服从者的同意容易而顺从，请求者的节制虔诚时，才能去做。殉道者命令做某事；但如果他们所命令的没有写在主的律法中，我们必须首先知道他们已从天主那里得到了他们所求的，然后才做他们所命的事。因为人承诺所能获得的神圣威严，未必总是立即应允。\par

19. 梅瑟也曾为百姓的罪求情；然而，当他为这些罪人恳求赦免时，却没有获得应允。他说：\Quote{主，我恳求祢，这百姓犯了大罪，为自己造了金神像。如今，求祢赦免他们的罪，不然，就从祢所写的册上涂抹我的名罢。} 主对梅瑟说：\Quote{谁得罪了我，我就要从我的册上涂抹谁。} \BibleRef{出 32:31} 他是天主的朋友，常与主面对面说话，却不能得到他所求的，也无法以恳求平息天主的义怒。天主赞扬耶肋米亚，宣告说：\Quote{我尚未在母腹中形成你以前，我已认识了你；你未出母胎以前，我已圣化了你，立你作万民的先知。} \BibleRef{耶 1:5} 当这人多次为百姓的罪恳求祈祷时，主对他说：\Quote{不要为这百姓祈祷，不要为他们呼号或祈祷，因为他们呼求我的时候，在他们的灾祸之时，我必不俯听。} \BibleRef{耶 7:16} 然而，谁比诺厄更正义？当大地充满罪恶时，只有他在世上被找到是正义的。谁比达尼尔更荣耀？谁在信仰的坚忍中更勇于殉道，在天主的屈尊俯就中更蒙福？他多次在争战中得胜，得胜后又存活。谁比约伯更乐于行善，在试探中更勇敢，在苦难中更忍耐，在敬畏中更顺服，在信德中更真诚？然而天主说，即使他们求情，他也不应允。当厄则克耳先知为百姓的罪恳求时，他说：\Quote{若一地因犯极大罪过而得罪我，我必伸手攻击它，折断其粮杖，使饥荒临到它，将人与兽从其中剪除。即使诺厄、达尼尔和约伯这三人在其中，他们也不能救自己的子女，只救自己。} \BibleRef{则 14:13} 因此，所祈求的并非全在于求者的预判，而在于施予者的自由意志；除非天主的喜悦许可，人的判断不能自行主张或篡夺什么。\par

20. 在福音中，主说：\Quote{凡在人前承认我的，我也必在我天上的父前承认他；但否认我的，我也必否认他。} \BibleRef{路 12:8} 若他不否认那否认他的人，他也不承认那承认他的人；福音不能在一处稳固，在另一处动摇。两者必须皆稳固，否则两者皆失真理之力。若否认的人不构成罪恶，那么承认的人也不接受德行的赏报。再者，若得胜的信德得冠冕，那么被击败的不信必然受惩罚。因此，殉道者若福音可被破坏，则一无所能；或若福音不可破坏，则他们不能反对福音，因为他们正是因为福音才成为殉道者。亲爱的弟兄们，但愿无人贬低殉道者的尊严，无人贬损他们的荣耀和冠冕。他们无玷信德的力量稳固完好；他们的一切盼望、信德、德能和光荣都在基督内，他们不能对基督说什么或做什么反对之事。那些自己遵行了天主命令的人，不能成为主教们违背天主命令的权威。难道有人比天主更大，或比天主的仁慈更仁慈，以至于他希望取消天主所允许成就的事，或好像天主能力不足保护祂的教会，而认为我们可以靠他的帮助得以保存吗？\par

21. 除非，或许这些事是在天主不知道的情况下发生的，或所有这些都是未经祂许可而发生的；尽管圣经教导愚顽的人，提醒健忘的人，它说：\Quote{谁将雅各伯交出成为掠物，将以色列交给劫掠者？岂不是那得罪了主、不肯行在祂道中、不服从祂法律的主吗？祂将祂的怒气倾注在他们身上。} \BibleRef{依 13:24} 在别处也见证说：\Quote{难道主的手缩短了，不能拯救？或祂的耳朵沉重，不能听见？但你们的罪过使你们与你们的天主隔绝；因你们的罪恶，祂掩面不看你们，以免施恩。} \BibleRef{依 59:1} 我们不如反省自己的过犯，回顾我们的行为和我们内心的隐密；让我们衡量良心的功过；让我们心中回想，我们没有行在主的道上，离开了天主的法律，从未愿意遵守祂的诫命和救恩的劝谕。\par

22. 你能认为他有什么好处？你能假设他有什么敬畏或信德？就是那连恐惧也不能纠正、迫害也不能改造的人？他高傲僵硬的颈项，即使跌倒也不弯曲；他膨胀傲慢的灵魂，即使被征服也不破碎。倒下的人却威胁站着的人，受伤的人却威胁健全的人。因为他不能立刻用污秽的手领受主的身体，亵渎者便对司铎发怒。噢，癫狂的人啊，你过度疯狂！你向那努力将天主的愤怒从你身上转开的人发怒；你威胁那为你恳求天主仁慈的人，他感受你的创伤，你自己却不感受；他为你流泪，或许你自己从未流过泪。你仍在加重和增加你的罪行；而当你自己对天主的仆人和司铎毫不留情时，你以为主能对你息怒吗？\par

23. 你们倒要接受并采纳我们所说的。为何你们聋聩的耳朵听不进我们警告你们的救恩训诫？为何你们昏盲的眼睛看不见我们所指出的悔改之路？为何你们受伤而疏离的心智察觉不到我们从天上圣经所学并教导的活泼疗法？或者，若有些不信者对未来之事信心不足，就让当下之事令他们战栗吧。看啊，我们目睹了那些背弃者所受的何等惩罚！我们哀悼他们何等悲惨的死亡！即在此世，他们也不能免受惩罚，尽管惩罚之日尚未到来。有些人暂时受罚，为叫他人得改正。少数人的痛苦是所有人的警戒。\par

24. 有一个自愿登上\Place{卡比托利欧}去背弃的人，在他否认基督之后，变成了哑巴。惩罚从罪行开始之处开始；以致他如今无法祈求，因为他没有言语来恳求怜悯。另一个女子，当时在公共浴场（因为她的罪行和不幸尚缺这一点：她竟在失去生命之层的恩宠后，立刻去了浴场）；在那里，不洁的她被一个不洁之神附身，并用牙齿撕裂了舌头，那舌头她曾用以不虔地吃或说话。在吃下邪恶的食物之后，口的疯狂转而攻击自身。她成了自己的行刑者，此后不久便不再存活：因腹肠绞痛而气绝身亡。\par

25. 请听我在场时并亲眼所见发生的事。有些父母因恐惧而仓促逃离，不太谨慎，将一个小女儿托付给奶妈照料。奶妈将遗弃的孩子交给了官长。官长在人们聚集的偶像面前（因为她年纪尚幼，还未能吃肉），给了她掺了酒的面包，而这面包本身是那些丧命者所献祭的剩余。后来母亲找回了孩子。但女孩已不能说话，也不能表明所犯的罪行，正如她先前不能理解或阻止它一样。因此，在我们不知情的情况下，当我们正在献祭时，母亲带着她进来了。而且，那女孩与圣人们混在一起，竟不耐我们的祈祷和恳求，时而因哭泣而颤抖，时而因心灵的强烈激动如海涛般翻滚；仿佛在刑讯的逼迫下，那仍然幼小孩子的灵魂，以它所能的迹象，承认了对事实的知晓。然而，当庄严礼仪结束，执事开始向在场的人递送杯爵，当其余人领受后，轮到那女孩时，她由于神圣尊严的本能，转过脸去，用抗拒的嘴唇紧闭嘴巴，拒绝杯爵。但执事坚持，尽管她抗拒，仍强行将一些杯爵的圣事灌入她口中。随后，她便抽泣和呕吐。圣体在亵渎的身体和口中不能停留；因主宝血而圣化的饮液从污秽的胃中迸发出来。主的权能何其大，主的威严何其大！黑暗的秘密在祂的光明下被揭露，连隐藏的罪行也欺骗不了天主的司铎。\par

26. 关于一个婴儿，就说这么多；她还没有到能言说别人对她所犯罪行的年龄。但有一个年事已高、年纪较大的女子，在我们献祭时偷偷混入我们中间，她所得的并非食物，而是一把利剑；她仿佛将某种致命毒药吞入口中和体内，立刻开始受折磨，因疯狂而僵硬；她所受的不再是迫害的痛苦，而是她罪行的痛苦；她颤抖战栗，仆倒在地。她那伪装良心的罪行并没有长久不受惩罚或隐匿。她欺骗了人，却感到天主在报复。又有一个女子，当她试图以不洁的手打开她盛有主的圣体的盒子时，被从盒中升起的火焰所阻，不敢触摸。还有一个人，他本身已被玷污，却胆敢与其他人一起秘密接受司铎所献之祭的一部分；他既不能吃也不能触摸主的圣物，而是发现当他打开手时，手中只有一块灰烬。这样，通过一人的经历表明：主在被否认时便会离去；领受之物对不配之人无益于救恩，因为救恩的恩宠因圣洁的离去而变成了灰烬。每日有多少人不悔改、不承认自己罪行的知觉，而被不洁之神充满！有多少人因疯狂的暴怒而摇动，甚至心智不健全、愚痴！无需逐一列举个人的死亡，因为世上发生多种多样的过失，罪人的惩罚正如其数量繁多。让每个人思考的不是别人受了什么苦，而是他自己应受什么苦；也不要以为如果惩罚延迟一时，他就逃脱了，因为他更应惧怕天主审判者的义怒为自己保留了惩罚。\par

27. 那些未曾用可憎的祭献玷污双手，却用证书玷污良心的人，不要自欺欺人地以为自己需要悔改得少。否认者的声明，就是一个基督徒否认自己从前身份的见证。他说自己做了别人实际所犯的事；尽管经上记载：\Quote{你们不能服事两个主人}\BibleRef{玛 6:24}，他却服事了地上的主人，因为他服从了他的谕令；他更服从人的权威而非天主。无论他将自己所作的事以较少的羞耻或罪过公布于人前，都无关紧要。不管怎样，他都不能逃脱和避开天主——他的审判者，因为圣神在圣咏中说：\Quote{你的眼睛看见了我的形体，它还不完美；在你的书上，人人都被记录下来。}又说：\Quote{人看外貌，天主却看人心。}\BibleRef{撒上 16:7} 主自己也预先警告并预备我们说：\Quote{所有教会都要知道，我是那洞察肺腑和心肠的。}\BibleRef{默 2:23} 祂察看隐藏和秘密的事，顾虑那些被掩盖的事；没有人能躲避主的眼目，祂说：\Quote{我是近处的天主，不是远处的天主。人若藏在隐密处，我岂不看见他？我岂不充满天地？}\BibleRef{耶 23:23} 祂看见每个人的心和意念；祂不仅审判我们的行为，也审判我们的言语和思想。祂洞察所有人的心思、意志和构想，就在那仍然紧闭的心灵的隐密处。\par

28. 然而，那些虽然没有被祭献\Emph{偶像}的罪行或证书所束缚，却甚至想过这些事的人，他们以忧伤和率直向天主的司铎告明这事，并作良心的宣认，卸下心中的重担，为轻微和适度的伤口寻求救治之药，知道经上记载：\Quote{天主不是可以戏弄的}\BibleRef{迦 6:7}，他们在信德上何等地更大，在敬畏上何等地更好呢！天主不可被戏弄，不可被欺骗，也不可被任何诡诈所蒙蔽。是的，谁若以为天主像人，相信自己若不公开承认罪行便能逃避刑罚，这人的罪更大。基督在祂的训诫中说：\Quote{凡以我为羞耻的，人子也要以他为羞耻。}\BibleRef{谷 6:83} 那以做基督徒为羞耻或害怕的人，难道他以为自己是个基督徒吗？那因属于基督而脸红或害怕的人，怎能与基督合一呢？他没有看见偶像，没有在围观和侮辱的众人眼前亵渎信仰的圣洁，没有用致命的祭献玷污双手，没有用邪恶的食物玷污嘴唇，他的罪必然较小。这有益之处在于：过失较小，而非良心无罪。他更容易获得罪过的赦免，然而他并非无罪；但愿他不要停止做补赎，并恳求主的仁慈，免得那在过失性质上看似较小的事，因他忽略补赎而增大。\par

29. 恳求你们，亲爱的弟兄们，每个人都应告明自己的罪，趁犯罪者尚在此世，趁他的告解还可以被接纳，趁由司铎所作的补赎与宽恕在主前蒙悦纳之时。让我们全心归向主，并以真忧伤表达对罪的悔改，让我们恳求天主的仁慈。让我们的灵魂在祂面前谦卑。让我们的哀悼赔补祂。让我们的一切希望依靠祂。祂亲自告诉我们当如何祈求。祂说：\Quote{归向我，全心归向我，同时要禁食、哭泣、哀悼；要撕裂你们的心，而不是你们的衣服。}\BibleRef{岳 2:12} 让我们全心归向主。让我们用禁食、哭泣、哀悼来平息祂的愤怒和义怒，正如祂亲自告诫我们的。\par

30. 我们真相信一个人会全心哀恸、以禁食、哭泣和悲伤恳求主吗？这人从犯罪第一天起就每日与妇女出入浴场；他享用丰盛的宴席，饱食珍馐，次日仍打嗝消化不良，并不分施其饮食以救济穷人之急需。他步履欢快、面带喜悦，又如何为自己的死亡哀恸？虽然经上记载：\BibleQuote{不可毁损你胡须的容貌，} \BibleRef{肋 19:27} 他却拔胡须、理头发；现在他竟想取悦那令天主不悦的人吗？或者，那有闲暇穿戴贵重衣饰、而不思念自己失去的基督外袍的女子，会呻吟哀叹吗？她接受贵重饰品和精致项链，却不哀恸神性和天上饰品的丧失？纵然你穿上异国服装和丝绸袍子，你仍是赤身；纵然你以珍珠、宝石和黄金过度妆饰，若无基督的妆饰，你仍是丑陋的。你染了头发，如今至少在悲哀中停手吧；你用黑粉描画眼圈边缘，如今至少用眼泪洗眼吧。你若因可朽的死亡失去任何一位亲友，必会深深呻吟，面容憔悴、衣着改变、头发散乱、脸色阴沉、神色沮丧，显出悲伤的记号。可怜的人，你失去了灵魂；在此灵性上已死，却仍为自己活着；虽然自己行走，却已开始背负自己的死亡。你还不痛苦哀号？你还不不断呻吟？你不因罪的羞耻或哀恸的延续而躲藏吗？看啊，这些是更严重的罪恶创伤；看啊，这是更大的罪行——犯罪而不做补赎，行恶而不哀恸自己的恶行。\par

31. 阿纳尼雅、阿匝黎雅和米沙耳，这些杰出而高贵的青年，甚至在烈焰和熊熊火窑的炽热中，也没有停止向天主公开认罪。他们虽然良心无亏，并常以信德和敬畏的服从堪受主恩，却仍不放弃保持谦卑，并在他们光荣的德性殉道中向主做补赎。圣经说：\BibleQuote{阿匝黎雅站定，开口祈祷，在火中与同伴一起向天主认罪。} 达尼尔也如此，在他信德和纯洁的多种恩宠之后，在主的屈尊俯就屡次针对他的德行和赞美之后，仍以禁食更堪当主的恩待，披上麻衣和灰烬，悲伤地认罪说：\BibleQuote{我主天主，伟大而可畏，为爱祢和遵守祢诫命的人持守盟约和仁慈，我们犯了罪，行了不义，作了恶；我们悖逆，偏离了祢的诫命和典章；没有听从祢仆人先知们奉祢的名对我们君王、列国和全地所说的话。我主啊，正义属于祢，而羞愧属于我们。} \BibleRef{达 9:4}\par

32. 这些事是那些温良、纯朴、无辜的人，为堪当天主威严而行的；而如今那些否认主的人却拒绝向主做补赎、恳求祂。我恳求你们，弟兄们，接受有益健康的忠告，服从更好的劝诫，以你们的眼泪连合我们的眼泪，以你们的呻吟联同我们的呻吟；我们恳求你们，为使我们能恳求天主宽恕你们；我们首先向你们献上我们的祈祷——我们为你们祈祷天主，求祂怜悯你们的祈祷。请丰盛地悔改，证明你们忧伤哀恸的心灵。\par

33. 也不让某些人愚昧的错误或虚妄的麻木动摇你们，这些人虽然陷入如此严重的罪，却被心盲所打击，以致既不理解也不哀恸自己的罪恶。这乃是天主愤怒更重的惩罚；正如经上记载：\BibleQuote{天主给他们昏迷的心。} 又说：\BibleQuote{他们不领受爱真理的心，以致得救。为此，天主就要给他们一种错谬的作用，使他们相信谎言，使一切不信真理而喜爱不义的人都被定罪。} \BibleRef{得后 2:10} 他们不义地自喜，为硬心之疏离而疯狂，藐视主的诫命，疏忽伤口之药，不愿悔改。他们在罪被承认前愚昧，犯罪后执拗；既非先前坚定，又非事后恳求；本应站稳时却跌倒，本应俯伏在天主前时却自以为站立。他们未得准许却擅自为自己取得平安；被虚假应许所惑，与背教者和不信者连结，以错误取代真理；他们视与不领圣体者共融为有效；他们信人胜过信天主，虽然他们不信天主胜过人。\par

34. 要尽可能逃避这种人；以有益的谨慎躲开那些紧贴其邪恶接触的人。他们的话像癌疮般侵蚀；他们的交谈像传染病般蔓延；他们有害而恶毒的劝说比迫害本身更致命。在此情形下，唯一能补赎的是悔罪。而那些夺去罪的悔改的人，关闭了补赎之路。这样，因某些人的鲁莽，虚假的安全被许诺或信赖，真正安全之望就被夺去了。\par

35. 但是你们，亲爱的弟兄们，你们对天主的敬畏是现成的，你们的心思虽然处在跌倒之中，却记得自己的悲惨，你们要在悔改和悲伤中省察自己的罪过；承认你们良心中极其严重的罪；睁开你们心灵的眼睛，以明白你们的罪，既不要对主的仁慈绝望，也不要立刻要求祂的宽恕。天主，正如以慈父的情怀总是宽厚和良善，同样，以审判者的威严也是可敬畏的。我们既然大大地犯了罪，就让我们大大地哀哭。深重的创伤不可缺少长久而仔细的治疗；悔改不可小于罪过。你们以为主能迅速被平息吗？你们用不信的话否认了祂，你们将世俗的产业置于祂之上，你们用亵渎的接触玷污了祂的圣殿。你们以为祂会轻易怜悯你们吗？你们曾宣告祂不是你们的天主。你们必须更热切地祈祷和恳求；你们必须在哀伤中度过白日；用守夜和哭泣消磨黑夜；用所有的时光在哀痛的哭号中；躺卧在地上，紧贴灰尘，被麻衣和污秽包围；失去基督的衣裳后，现在要甘愿无衣可穿；在魔鬼的食物之后，要喜好禁食；热切从事正义的作为，罪过借此得以涤除；常常致力于施舍，灵魂借此从死亡中获得释放。仇敌从你们那里夺去的，让基督接收；你们现在也不可再持有或爱慕那曾欺骗并征服了你们的财产。财富必须被当作仇敌回避；必须像强盗一样逃离；必须被拥有者当作剑和毒药来畏惧。为此，所剩下的财产只应被用于赎罪和过犯。善工要及时而大量地施行；你们全部的财产都应用来医治你们的创伤；让我们将我们的财富和资产借给那位将要审判我们的主。如此，信德在宗徒时代兴盛；如此，第一批信众遵守了基督的诫命：他们迅速，他们慷慨，他们将一切交给宗徒分配；然而他们当时所赎的罪并不是这种性质的。\par

36. 若有人全心祈祷，若他以真悔改的哀哭和泪水叹息，若他以正义和不断的作为使主赦免他的罪，那位以这些话表达了祂的仁慈的主可能会怜悯这样的人：\BibleQuote{当你转变并哀哭时，你便要得救，并要知道你曾在何处。}\BibleRef{依 30:51} 又说：\BibleQuote{我不喜爱那死者的死亡，主说，却喜爱他回头而活。}\BibleRef{则 33:11} 岳厄尔先知在主自己的告诫中宣告了主的仁慈，他说：\BibleQuote{你们归向上主你们的天主，因为祂是宽仁慈悲的，缓于发怒，富于慈爱，对祂降下的灾祸常感后悔。}\BibleRef{岳 2:13} 祂能显示仁慈；祂能撤回祂的审判。祂能仁慈地宽恕悔改、劳苦、恳求的罪人。祂能视殉道者所恳求的、司铎所行的为有效，以助益这样的人。若有人以自己的补赎更打动祂，若他平息了祂的愤怒，若他以正义的祈求平息了恼怒的上主的义怒，祂就再次赐予武器，使被击败者可以武装起来；祂恢复并坚固力量，使更新的信德得以强壮。士兵将重新寻求战斗；他将重演战斗，他挑战仇敌，而且确实因受苦而变得更加勇敢。这样向天主赎了罪的人，因对自己的行为悔改，因对自己的罪羞愧，从跌倒的悲痛中获得了更多的德行和信德，蒙主垂听和帮助，将使他不久前使哀伤过的教会喜乐，并且如今将不仅从主那里得到宽恕，还要得到华冠。\par

\SourceRecord{Translated by Robert Ernest Wallis.；Ante-Nicene Fathers；Vol. 5.；Edited by Alexander Roberts, James Donaldson, and A. Cleveland Coxe.；Buffalo, NY: Christian Literature Publishing Co.,；1886.；<http://www.newadvent.org/fathers/050703.htm>.}

% structure: p0001 source=h1 target=section reason=page-title

\section{第四论}

\Emph{论主祷文}\par

\Emph{论证。——\Person{西彼廉}论主祷文的论文包含三部分，在此划分中他效法\Person{戴尔都良}在其\WorkTitle{论祈祷}一书中的做法：在第一部分，他指出主祷文是所有祈祷中最卓越、最深具灵性、且最有效地获得我们祈求的祈祷。在第二部分，他着手解释主祷文；并仍跟随\Person{戴尔都良}的足迹，逐一讲解其七项主要祈求。最后，在第三部分，他考虑祈祷的条件，告诉我们祈祷应当如何。} ——\par

1. 福音的训诫，亲爱的弟兄们，无非是天主的教导——是建造望德的基础，增强信德的支柱，欢愉心灵的滋养，引导道路的舵柄，获得救恩的守卫——这些训诫，当它们在世上教导信徒温顺的心灵时，引导他们到达天国。此外，天主意愿许多事情藉着祂的仆人先知们说出和听到；但圣子所说的话，那曾在先知们内的天主圣言亲自以祂自己的声音作证的话，是何等更伟大啊！现在祂不是吩咐预备祂来临的道路，而是亲自来临，为我们开启并指明道路，好使我们这些先前在死亡的黑暗中漂泊、没有远见、瞎眼的人，蒙恩宠的光照，得以持守生命之道，以主为我们的统治者和引导！\par

2. 祂在其他有益的劝诫和神圣的训诫中，以此劝告祂的子民获得救恩，祂自己也赐予了一种祈祷的形式——祂亲自建议并教导我们应当祈求什么。祂赐予我们生命的那一位，也教导我们祈祷，且是以同样的仁慈，即祂屈尊赐予并赋予一切其他事物的仁慈；好使我们向天父说出圣子所教导我们的祈祷和恳求时，更容易被垂听。祂早已预言了时候将到，\Quote{那时真正的朝拜者将以心神和真理朝拜父；} \BibleRef{若 4:23}祂如此实现了祂先前所应许的，使我们这些藉着祂的圣化而领受了圣神和真理的人，也因着祂的教导能以真理和心神朝拜。因为还有什么祈祷比那由基督——祂也将圣神赐予了我们——所赐予我们的更属圣神呢？还有什么对天父的祈祷比那由子——祂就是真理——亲口传授给我们的更真实呢？以致于以祂所教导以外的方式祈祷，不仅是无知，而且是罪；因为祂亲自确立并说：\Quote{你们废弃天主的诫命，为要遵守你们的传统。}\par

3. 因此，让我们，亲爱的弟兄们，按照天主我们的导师所教导的来祈祷。以祂自己的话语恳求天主，在基督的祈祷中达到祂的耳中，这是一种充满爱和友情的祈祷。当我们祈祷时，愿圣父认出祂圣子的话语，愿那位内住在我们胸中的祂也内住在我们的声音中。既然祂是我们在天父面前为我们的罪而辩护的辩护者，那么当我们作为罪人为我们的罪祈求时，让我们提出我们辩护者的话语。因为祂说：\Quote{你们因我的名无论向父求什么，祂必赐给你们；} \BibleRef{若 16:23}如果我们以祂自己的祈祷来求，那么以基督的名求我们所求的，岂不更有效吗？\par

4. 但当我们祈祷时，我们的言语和祈求应受约束，保持安静和谦逊。让我们思考我们是站在天主的眼前。我们必须以身体的姿态和声音的尺度取悦天主的眼睛。因为正如无羞耻之人的特征是大声喊叫，同样，谦逊之人则适合用温和的祈求来祈祷。此外，主在祂的教导中命令我们在暗中祈祷——在隐蔽和偏远的地方，在我们的卧室里——这最适合信德，好使我们知道天主无所不在，祂听到并看到一切，并以其威严的丰满甚至渗透到隐藏和秘密的地方，正如经上所记：\Quote{我是一名近在咫尺的天主，而非遥远的天主。人若隐藏在秘密处，难道我看不见他吗？我岂不充满天地？} \BibleRef{耶 23:23}又说：\Quote{主的眼睛处处都在，观察善恶。} \BibleRef{箴 15:3}当我们与弟兄们聚集在一处，并与天主的司铎一起庆祝神圣的祭献时，我们应当谨记谦逊和纪律——不要用不受抑制的声音任意抛洒我们的祈祷，也不要将本应以谦逊向天主托付的祈求以喧闹的冗词抛向天主；因为天主是听众，不是声音的听众，而是心灵的听众。祂也不需要被大声提醒，因为祂看见人的思想，正如主向我们证明时所说：\Quote{你们为什么心里思念恶事？} \BibleRef{玛 9:4}在另一处：\Quote{所有教会都要知道，我是那洞察人心和肺腑的。} \BibleRef{默 2:23}\par

5. 而\WorkTitle{撒慕尔纪上}中的\Person{亚纳}，她是教会的预像，保持并遵守这一点，因为她不是以喧闹的祈求，而是默默地、谦逊地在内心深处向天主祈祷。她以隐秘的祈祷说话，却以明显的信德说话。她不是以声音说话，而是以心灵说话，因为她知道天主如此垂听；她有效地获得了所求的，因为她以信德祈求。圣经证实了这一点，当它说：\Quote{她心里说话，只嘴唇动，却听不到她的声音，天主听见了她。} \BibleRef{撒上 1:13}我们在\WorkTitle{圣咏}中也读到：\Quote{你们当在你们心里说话，在你们床上，并要痛悔。}此外，圣神藉着\Person{耶肋米亚}也提示同样的事，并教导说：\Quote{但你们应当在心里朝拜天主。}\par

6. 亲爱的弟兄们，但愿祈祷者不要不知道，那税吏在圣殿里与法利塞人一同祈祷时是如何做的。他并不大胆地举目向天，也不骄傲地举手，而是捶着胸膛，承认心中隐秘的罪恶，恳求天主仁慈的帮助。当法利塞人自满时，这个如此恳求的人反倒更堪当被圣化，因为他将救恩的希望不是放在自己无罪的自信上——因为没有人是无罪的——而是承认自己的罪过，谦卑地祈祷，那赦免谦卑者的主垂听了恳求者。主在祂的福音中记载了这些事，说：\Quote{有两个人上圣殿去祈祷：一个是法利塞人，另一个是税吏。法利塞人站着，心里这样祈祷：天主，我感谢祢，因为我不像别人一样，勒索、不义、奸淫，也不像这个税吏。我每周禁食两次，凡我所得的，都捐献十分之一。但那税吏远远地站着，连举目望天也不敢，只是捶着胸膛说：天主，可怜我这个罪人罢！我告诉你们：这人下去，到他家里，成了正义的，而不是那法利塞人；因为凡高举自己的，必被贬抑；凡自己谦卑的，必被高举。}\newline{}\BibleRef{路 18:10}\par

7. 亲爱的弟兄们，当我们从神圣的诵读中学习了这些事，并明白了应当如何接近祈祷时，让我们也从主的教导中知道我们该祈祷什么。祂说：\Quote{你们要这样祈祷：}——\par

\Quote{我们在天的父！愿祢的名被尊为圣，愿祢的国来临，愿祢的旨意承行于地，如在天上一样。求祢今天赏给我们日用的食粮；宽免我们的罪债，如同我们也宽免得罪我们的人；不要让我们陷入诱惑，但救我们免于凶恶。阿们。}\newline{}\BibleRef{玛 6:9}\par

8. 首先，和平的导师与合一的师傅不愿意祈祷是单独和个别的，好像一个人祈祷只为他自己。因为我们不说\Quote{我在天的父}，也不说\Quote{求祢今天赏给我日用的食粮}；也不各自只求自己的罪债得宽免；也不只为自己恳求不陷入诱惑、得救免凶恶。我们的祈祷是公开而共同的；当我们祈祷时，我们不是为一人祈祷，而是为全体子民祈祷，因为我们全体子民是一体。和平的天主与和谐的导师，祂教导合一，愿意一人如此为众人祈祷，正如祂自己将我们众人合而为一。这祈祷的法则，那三个青年人在被关进烈火窑中时遵守了，他们在祈祷中异口同声，因精神的一致而同心合意；神圣的圣经信实向我们保证这一点，在告诉我们这些人如何祈祷时，它给出了一个我们应在祈祷中效法的榜样，好使我们能像他们一样：\Quote{那时，这三个人好像从同一口中唱出赞美诗，颂扬了主。}他们好像从同一口中说话，虽然基督还没有教导他们如何祈祷。因此，他们的祈祷有功效且有效，因为和平、真诚且属神的祈祷堪当主的垂顾。同样，我们发现宗徒们与门徒们在主升天后祈祷：\Quote{这些人同一些妇女及耶稣的母亲玛利亚，并祂的弟兄们，都同心合意地专务祈祷。}\newline{}\BibleRef{宗 1:14}他们同心合意地专务祈祷，既以他们祈祷的迫切，又以他们祈祷的一致表明，那\Quote{使众人同心合意住在家中}的天主，只接纳那些祈祷同心的人进入神圣而永恒的家园。\par

9. 主的祈祷包含着多么深奥的道理！多么多！多么伟大！寥寥数语中凝聚着如此丰富的属灵大能！以至于绝对没有任何事物被遗漏，不被涵盖在我们的祈祷和恳求中，如同在天上教义的撮要中一般。祂说：\Quote{你们要这样祈祷：我们在天的父。}那新人，由天主恩宠重生并恢复与天主关系的人，首先说\Quote{父}，因为他现在开始成为儿子。\Quote{祂来到自己的领域，自己的人却没有接受祂；但是，凡接受祂的，祂给他们权能成为天主的子女，就是那些信祂名字的人。}\newline{}\BibleRef{若 1:11}因此，凡信祂名字并成为天主儿子的人，应当从此开始感恩并承认自己是天主的儿子，宣告天主是他在天的父；并且，在他重生的最初话语中，就要见证他已弃绝了属世和属血肉的父亲，开始认识并拥有那唯一在天的父，正如经上所记载：\Quote{他们对自己的父亲和母亲说：我没有见过你们；又不承认自己的子女；这些遵守了祢的诫命，持守了祢的盟约。}\newline{}\BibleRef{申 33:9}主在福音中也命令我们不要\Quote{称地上的人为父}，因为只有一位在天上的父。\newline{}\BibleRef{玛 23:9}对于那提到他已死父亲的门徒，祂回答说：\Quote{任凭死人去埋葬他们的死人！}\newline{}\BibleRef{玛 8:22}因为那门徒说他的父亲死了，而信徒的父是永生的。\par

10. 我们不应只观察和明白，亲爱的弟兄们，我们应称那在天上的为\Emph{父}；但我们还加上，说\Quote{\Emph{我们的}父}，即相信者的父——那些因祂而圣化、因超性恩宠的重生而恢复、开始成为天主的子女的人。这句话还斥责和定罪犹太人，他们不仅不信地蔑视基督（祂曾由先知们向他们宣告，并先被派遣到他们那里），还残酷地将祂处死；他们现在不能称天主为他们的父，因为主谴责和驳斥他们说：\BibleQuote{你们是出于你们的父魔鬼，你们父的私欲你们偏要行。他从起初就是杀人的，不站在真理上，因为他心里没有真理。}\BibleRef{若 8:44} 又由依撒意亚先知，天主在愤怒中呼喊：\BibleQuote{我养育了儿女，使他们长大，但他们竟背叛了我。牛认识自己的主人，驴认识自己主人的槽；以色列却不认识我，我的百姓却不理解我。祸哉，犯罪的民族，恶贯满盈的人民，邪恶的种子，败坏的儿女！你们离弃了上主，激怒了以色列的圣者。}\BibleRef{依 1:3} 弃绝这些，我们基督徒祈祷时说：\Quote{我们的父}；因为祂已成为我们的父，而不再是犹太人的父，因为他们离弃了祂。犯罪的人民不能成为儿子；而儿子的名称归于那些获得罪赦、并重新被应许永生的人，正如我主亲自的话：\BibleQuote{凡犯罪的，就是罪恶的奴隶。奴隶不会永远住在家里，儿子却永远住下。}\BibleRef{若 8:34}\par

11. 然而，主的宽容何其大！祂对我们的屈尊和丰富的良善何其大，竟愿意我们这样在天主面前祈祷，称天主为父，称我们自己为天主的子女，如同基督是天主子一样——这个称呼我们中没有一个人敢在祈祷中冒昧使用，除非祂亲自允许我们这样祈祷！为此，亲爱的弟兄们，我们应当牢记并知道，当我们称天主为父时，我们的行为应当相称于天主的子女；这样，我们以怎样的程度喜爱视天主为父，祂也以怎样的程度喜爱我们。让我们作为天主的圣殿而交谈，好显明天主住在我们内。愿我们的行为不偏离圣神，好使我们这些开始成为属天和属灵的人，只思考和从事属灵和属天的事；因为主天主亲自说过：\BibleQuote{尊重我的，我必尊重他；藐视我的，他必被轻视。}\BibleRef{撒上 2:30} 蒙福的宗徒也在他的书信中写道：\BibleQuote{你们不是属于自己的，因为你们是用重价买来的。因此，要在你们的身上光荣天主，并承载天主。}\BibleRef{格前 6:20}\par

12. 在此之后，我们说：\Quote{愿你的名被尊为圣；}这不是我们愿望天主藉我们的祈祷而被圣化，而是我们恳求祂使祂的名在我们内被尊为圣。但谁能圣化天主呢？因为祂自己就是圣化者。好吧，因为祂说：\BibleQuote{你们应是圣的，因为我是圣的，}\BibleRef{肋 20:7} 我们祈求并恳请，使我们这些在圣洗中被圣化的人，能继续留在我们已经开始成为的状态中。这是我们每日所祈祷的；因为我们每日需要圣化，好使我们每日跌倒的人，藉着不断的圣化洗净自己的罪。这藉天主屈尊而赐予我们的圣化是什么，宗徒予以宣告，他说：\BibleQuote{难道你们不知道，不义的人不会继承天主的国吗？不要自欺：无论是淫乱的、拜偶像的、奸淫的、作娈童的、好男色的、偷窃的、贪婪的、醉酒的、辱骂的、勒索的，都不能继承天主的国。你们中从前也有人是这样；但你们已经洗净了，已经成义了，已经因我们的主耶稣基督之名并因我们的天主之神而圣化了。}\BibleRef{格前 6:9} 他说我们因我们的主耶稣基督之名并因我们的天主之神而圣化。我们祈祷这圣化常在我们内；因为我们的主和审判者警告那个被祂治好并复活的人，不要再犯罪，以免更糟的事临到他，我们在不断的祈祷中献上恳请，我们日夜请求，那从恩宠中接受的圣化和复活，能藉祂的保护得以保存。\par

13. 在祈祷中接着有这样的祈求：\BibleQuote{愿祢的国来临}。我们祈求天主的国得以向我们彰显，正如我们也祈求祂的名在我们内被尊为圣。因为天主何时不统治呢？或者，那常是祂所拥有、且永不停止的统治，又何时才开始呢？我们祈求的是天主所应许我们的、我们的国度来临，就是由基督的血和苦难所获得的；使我们这些先在世上作祂臣民的人，将来在基督为王时与祂一同为王，正如祂亲自应许和所说的话：\BibleQuote{你们蒙我父降福的，来吧！来承受从创世以来给你们预备了的国度}（\BibleRef{玛 25:34}）。最亲爱的弟兄们，然而基督自己可能就是天主的国，我们日复一日地渴望祂来临，我们热切盼望祂的来临快快向我们显示。因为既然祂自己是复活，既然我们在祂内复活，那么天主的国也可以理解为祂自己，因为我们在祂内将要为王。但我们寻求天主的国，即天国，是很好的，因为也有一个地上的国。但已经弃绝世界的人，更是超越了世界的尊荣和国度。因此，那将自己奉献给天主和基督的人，不贪求地上的国，而贪求天上的国。然而我们需要不断的祈祷和恳求，免得我们失去天国，如同那最初蒙赐这应许的犹太人失去了一样；正如主所指出和证明的：\BibleQuote{有许多人要从东方和西方来，同亚巴郎、依撒格、雅各伯在天国里一起坐席；本国的子民，却被驱逐到外边黑暗里去；那里要有哀号和切齿}（\BibleRef{玛 8:11}）。祂表明犹太人先前是天国的子民，只要他们继续作天主的儿女；但自从他们不再承认\Quote{父}的名号时，国度也就停止了。因此我们基督徒，在祈祷中开始称天主为我们的父，也祈求天主的国来临到我们。\par

14. 此外，我们还加上说：\BibleQuote{愿祢的旨意承行于地，如于天焉}；这不是说，要天主成就祂愿意的事，而是使我们能够成就天主所愿意的事。因为谁能抗拒天主，而不使祂成就祂愿意的事呢？但由于魔鬼阻碍我们在一切事上用心思和行为服从天主的旨意，我们就祈祷并祈求天主的旨意承行在我们身上；并且为使它承行在我们身上，我们需要天主的善意，即祂的帮助和保护，因为没有人凭自己的力量是坚强的，而是靠天主的恩宠和仁慈才安全。再者，主显示祂所背负的人性的软弱，说：\BibleQuote{父啊！如果可能，请把这杯离开我罢！}并且给祂的门徒们一个榜样，叫他们不要行自己的旨意，而应行天主的旨意，祂接着说：\BibleQuote{但不要照我的意愿，而要照祢的意愿}（\BibleRef{玛 26:39}）。在另一处祂说：\BibleQuote{我从天降下，不是为执行我的旨意，而是为执行那派遣我来者的旨意}（\BibleRef{若 6:38}）。既然圣子服从了圣父的旨意，那么仆人要服从主人的旨意，岂不更应当吗？正如若望在他的书信中也劝勉和教导我们要承行天主的旨意，说：\BibleQuote{你们不要爱世界，也不要爱世界上的事；谁若爱世界，父的爱就不在他内。因为凡世界上的事，就是肉身的贪欲、眼目的贪欲，以及人生的骄矜，都不是出于父，而是出于世界。世界和它的贪欲都要过去；但那承行天主意旨的人，却永远存在，正如天主也永远存在一样}（\BibleRef{若一 2:15}）。我们这些渴望永远存在的人，应当承行那永远存在的天主的旨意。\par

15. 如今，那正是基督所实行和教导的天主旨意：交谈时谦卑，信仰上坚定，言语温和，行事正义，工作仁慈，品行端正，不亏负人，而能忍受不公平；与弟兄保持和平；全心爱天主，爱祂为父，敬畏祂为神；凡事先把基督放在首位，因为基督未曾将任何事物置于我们之上；与祂的爱不可分离；勇敢而忠实地坚守祂的十字架；当为了祂的名和光荣而有任何争战时，在言论中表现我们宣认时的坚定，在酷刑中表现我们战斗时的信心，在死亡中表现我们获得冠冕的忍耐——这就是渴望与基督同为嗣子；这就是承行天主的诫命；这就是履行圣父的旨意。\par

16. 此外，我们求天主的旨意能行在天上，也在地上，这两者都关乎我们安全与救恩的圆满完成。因为我们从大地拥有身体，从天上拥有神魂，所以我们自己就是地又是天；并且在两者之中——即在身体和神魂中——我们祈求天主的旨意得以承行。因为肉身与神魂之间有争斗；彼此不和时便有一场日常的争执，以致我们不能行我们所愿的，因为神魂寻求天上的和神圣的事，而肉身却贪恋地上的和暂时的事；因此我们求天主帮助和扶助，使这两种性情之间达致和谐，如此，当天主的旨意在神魂中也必在肉身中承行时，那由祂重生的灵魂就能得以保存。这正是宗徒保禄在他话语中所公开明示的：\Quote{肉身与神魂交战，神魂与肉身交战：这两者彼此对立，以致你们不能行你们所愿的事。肉身的作为是显而易见的：奸淫、淫乱、污秽、淫荡、崇拜偶像、行邪术、仇恨、争执、嫉妒、忿怒、争吵、不睦、分裂、异端、妒忌、酗酒、狂欢宴饮之类的事。我预先告诉你们，如同我从前说过一样：行这样事的人，不能继承天主的国。但是神魂的果实是仁爱、喜乐、平安、宽宏、良善、信实、温和、节制、贞洁。} \BibleRef{迦 5:17} 因此，我们日复一日，不，是连续不断地恳求，求天主的旨意在我们身上能承行在天上，也在地上；因为这是天主的旨意：地上的事物应为天上的让路，属神和神圣的事物应当得胜。\par

17. 亲爱的弟兄们，也可以这样理解：既然主命令并劝诫我们甚至要爱仇敌，并为迫害我们的人祈祷，那么我们还应当为那些仍是属地、尚未开始属天的人祈求，好使天主的旨意也成就在他们身上，这事基督在保存和更新人性时已经完成了。因为门徒如今不再被祂称为地，而是称为地上的盐；宗徒又说第一个人是出于地上的尘土，第二个人是出于天上。我们理当像天父一样，祂使太阳升起照耀善人和恶人，降雨给义人和不义的人；因此我们受基督的劝告，如此祈祷和祈求，使我们为所有人的得救而祈祷：正如在天上——即因着我们的信德，在我等内——天主的旨意已经承行，使我们能属天；同样也在地上——即在那些不信的人身上——天主的旨意得以承行，好叫那些仍在第一次出生中属地的人，能藉着水和圣神重生，开始成为天上的。\par

18. 当祈祷继续进行时，我们祈求并说：\Quote{求祢今日赐给我们日用的食粮。}这可以按精神和字面两种方式理解，因为任何一种理解对于我们得救都具有丰富的神圣益处。因为基督是生命之粮；这食粮不属于众人，而是属于我们的。正如我们说\Quote{我们的天父}，因为祂是认识和相信祂的人之父；同样我们也称它为\Quote{我们的食粮}，因为基督是与祂身体合一之人的食粮。我们求这食粮每日赐给我们，好使我们这些在基督内，并每日领受圣体作为救恩食粮的人，不致因某一大罪的中阻，因被阻止、扣留而不领受，从而无法分沾天上的食粮，与基督的身体分离。正如祂亲自预言并警告说：\Quote{我是从天上降下的生命之粮。谁若吃我的粮，将永远活着；我要赐给的粮，就是我的肉，是为世界的生命而赐的。} \BibleRef{若 6:58} 因此，当祂说，谁若吃祂的粮将永远活着，既然显而易见，那些领受祂的身体并藉共融的权利领受圣体的人是生活的，那么另一方面，我们必须畏惧并祈求，免得有人被排除于共融之外，与基督的身体分离，从而远离救恩。正如祂亲自警告说：\Quote{你们若不吃人子的肉，不喝他的血，在你们内就没有生命。} \BibleRef{若 6:53} 因此我们求我们的食粮——即基督——每日赐给我们，好使我们这些在基督内居住并生活的人，不致离开祂的圣化和身体。\par

19. 但也可如此理解：我们这些弃绝了世界、因信属神恩宠而抛弃了世间的财富与虚荣的人，只应为自身祈求食物与支持，因为主教导我们说：\Quote{凡不撇下一切所有的，不能作我的门徒。}\newline{}\BibleRef{路 14:33} 然而，那已开始作基督门徒、既按师傅的话舍弃一切的人，理应为日用粮祈求，而不应将祈求的欲望延伸至长久之期，正如主又训诫说：\Quote{不要为明天忧虑，因为明天自有明天的忧虑。一天的难处一天当就够了。}\newline{}\BibleRef{玛 6:34} 因此，基督的门徒为自己祈求日用的食粮，原是合理的，因为他已被禁止为明天忧虑；因为我们祈求天主的国快快来临，却想在世上活长久，这成了矛盾与抵触之事。有福的宗徒也以此告诫我们，为我们望德与信德的坚定奠定实质与力量：\Quote{我们什么都没有带，}他说，\Quote{到世上来，也不能带什么去。只要有衣有食，就当知足。但那些想要发财的人，就陷在诱惑、罗网和许多无知有害的私欲里，叫人沉在败坏和灭亡中。贪财是万恶之根；有人贪恋钱财，就被引诱离了信德，用许多愁苦把自己刺透了。}\newline{}\BibleRef{弟前 6:7}\par

20. 祂教导我们，财富不仅应当轻看，而且充满危险；其中藏着诱人作恶的根，以隐秘的欺骗蒙蔽人灵的盲目。为此，天主也斥责那富有的愚人，他思念地上的财富，因丰收充裕而自夸，说：\Quote{无知的人哪，今夜必要你的灵魂；你所预备的要归谁呢？}\newline{}\BibleRef{路 12:20} 那当夜就要死的愚人，竟为他的积蓄欢喜；生命即将消逝的他，却想着粮食的丰裕。反之，主告诉我们，那变卖一切所有、分给穷人使用，从而在天上积蓄财宝的人，才成为完全完备的。祂说，那摆脱了阻碍、腰束带子、不受属世产业纠缠的人，能跟随祂，效法主苦难的光荣；他自由无羁，伴随那早已送往天主的财产。为此，我们每个人为能预备自己，就当这样学着祈祷，并从祈祷的性质知道应当成为怎样的人。\par

21. 因为义人的日用粮必不缺乏，正如经上记载：\BibleQuote{主必不使义人的灵魂忍饥受饿}\newline{}\BibleRef{箴 10:3}；又说：\BibleQuote{我从前年幼，现在年老，却未见过义人被弃，也未见过他的后裔讨饭。}主更应许说：\Quote{不要忧虑说：吃什么？喝什么？穿什么？这都是外邦人所求的。你们需用的这一切东西，你们的天父是知道的。你们要先求天主的国和祂的义德，这一切都要加给你们了。}\newline{}\BibleRef{玛 6:31} 对于那些寻求天主国和祂义德的人，祂应许一切都要加添。因为万物都属于天主，那拥有天主的人，若天主本身不离开他，他就一无所缺。这样，达尼尔曾得天主赐食物：当他奉王命被囚在狮子坑中，在饥饿而仍饶恕他的野兽中间，天主的人得了喂养。同样，厄里亚在逃亡时，既有乌鸦在旷野中伺候他，又有飞鸟在他受迫害时为他带来食物。可恨的人之恶毒！野兽饶恕，飞鸟喂养，人却设网罗、施暴怒！\par

22. 此后，我们也为自己的罪恳求，说：\Quote{免我们的债，如同我们免了人的债。}在求了食物之后，也求罪的赦免，使那由天主喂养的人能活于天主，不仅今世暂时的生命得到供应，而且能藉着罪赦达到永生；主称这些罪为债，如祂在福音中所说：\Quote{因为你求我，我免了你那全部的债。}\newline{}\BibleRef{玛 18:32} 这是何等必要、何等明智而有益地提醒我们：我们是罪人，因为我们被迫为自己的罪恳求，且在向天主求赦免时，灵魂回想起自己\Emph{有罪}的自觉！免得有人自以为是义人，因自高而更陷沉沦，他受训诲和教导，得知自己每日犯罪，因为他奉命每日为自己的罪恳求。此外，若望在他的书信中也警告我们说：\Quote{我们若说自己没有罪，便是自欺，真理不在我们心里了；但若我们认自己的罪，主是信实的、公义的，必要赦免我们的罪。}他在书信中将两者结合：我们既当为自己的罪恳求，又当在祈求时获得赦免。因此他说主是信实的，必赦免罪过，持守祂应许的信实；因为那教导我们为罪债祈求的主，已应许祂的慈父怜悯和赦免必会跟随。\par

23. 祂明显将法律与此联系在一起，并附加了条件，以某种条件和约定约束我们：我们应请求宽恕我们的罪债，如同我们宽恕我们的债务人一样，知道我们为自己罪行所求的，除非我们自己也以类似的方式对待我们的债务人，否则无法获得。因此，祂在另一处也说：\BibleQuote{你们用什么量器量，也必用什么量器量给你们。}那仆人在主人免了他一切债务之后，却不愿宽恕他的同伴，就被扔回监狱；因为他不愿宽恕同伴，就失去了主人曾施予他的恩惠。基督在祂的诫命中，以更大责备的力量更迫切地申明这些事。祂说：\BibleQuote{你们站着祈祷时，若你们有什么怨怼任何人，就宽恕吧，好叫你们在天之父也宽恕你们的过犯。但若你们不宽恕，你们在天之父也不会宽恕你们的过犯。}在审判之日，你们将按自己的判决受审判，毫无借口；你们做什么，也将受什么。因为天主命令我们在祂的家中成为和平之子、同心合意、一心一意；祂使我们藉第二次诞生成为这样的人，祂愿我们新生后持续如此，使我们这些已开始成为天主子民的人，能安住于天主的平安中，并拥有同一精神，也拥有同一心志。因此，天主不接受不和睦之人的祭献，反而命令他从祭台回去，先与弟兄和好，这样天主也因和平之人的祈祷而息怒。我们的和平与兄弟和睦，对天主是更大的祭献——一个在圣父、圣子、圣神的合一中合一的子民。\par

24. 因为即使在亚伯尔和加音最初献的祭中，天主看的不是他们的礼物，而是他们的心，以致那在心中蒙受悦纳的，他的礼物也蒙悦纳。亚伯尔是和平与正义的，在清白无罪中向天主献祭，也教导其他人在带礼物到祭台时，要怀着敬畏天主的心、纯朴的心、正义的法律、和睦的平安而来。有理地，他在天主的祭献中如此行，后来他本人也成为献给天主的祭品；以致那首先彰显殉道、并以他血的光荣开创主受难的人，同时拥有主的正义与和平。最后，这样的人被主加冕，这样的人将在审判之日与主一同受报复；但好争吵、不和睦、与弟兄没有平安的人，按蒙福的宗徒和圣经的见证，即使他为基督的名被杀害，也不能逃脱兄弟不和的罪，因为经上记载：\BibleQuote{凡恨自己弟兄的，便是凶手}\BibleRef{若一 3:15}，没有凶手能进入天国，也不能与天主同活。他不能与基督同在，宁愿效法犹达斯而不效法基督。这罪如此之大，连血的洗礼也不能洗去；这罪如此严重，连殉道也不能补赎！\par

25. 再者，主必要地劝诫我们在祈祷中说：\BibleQuote{不要让我们陷于诱惑。}这些话表明，除非天主事先允许，仇敌不能对我们做什么；因此我们所有的畏惧、热忱和服从都应转向天主，因为在我们的诱惑中，除非从祂得到权力，否则邪恶什么也不被允许。这由圣经证明，经上说：\BibleQuote{巴比伦王拿步高来到耶路撒冷，围攻它；上主将城交在他手中。}\BibleRef{列下 24:11}但权力因我们的罪而赐给邪恶攻击我们，正如经上记载：\BibleQuote{谁将雅各伯交出作为掠物，将以色列交给劫掠者？岂不是上主，他们得罪了祂，不肯行在祂的路上，也不听从祂的法律？祂就将祂的怒火降在他们身上。}\BibleRef{依 13:24}此外，当撒罗满犯罪，偏离了上主的诫命和道路时，经上记载：\BibleQuote{上主就兴起撒殚攻击撒罗满本人。}\BibleRef{列上 11:14}\par

26. 权力赐下攻击我们有两种方式：或是为惩罚我们犯罪时，或是为光荣我们受考验时，正如我们在约伯身上所见；天主亲自说明：\BibleQuote{看，凡他所有的，我都交在你手中；只是不要伸手害他。}\BibleRef{约 1:12}主在祂受难时的福音中说：\BibleQuote{若不是从上头赐给你的，你对我毫无权柄。}\BibleRef{若 19:11}但当我们祈求不陷于诱惑时，我们被提醒自己的软弱和脆弱，因为我们如此祈求，免得有人傲慢自夸，免得有人骄傲自大地自以为是，免得有人将宣认或受苦的光荣归为己有，因为主亲自教导谦卑，说：\BibleQuote{你们醒寤祈祷罢！免得陷于诱惑；心神固然切愿，但肉体却软弱。}\BibleRef{谷 14:38}这样，先有谦卑顺服的承认，一切归于天主，那么凡以敬畏和尊崇天主之情恳切祈求的，都因祂的慈爱而获得。\par

27. 在所有这一切之后，在祈祷的结尾有一句简短的句子，它简短而全面地总结了我们的所有恳求和祈祷。因为我们以这句话结束：\Quote{但救我们免于凶恶，}包含了所有敌人在这个世界中企图加害我们的不利之事，我们若得天主拯救，若祂向我们祈求并呼求的人施以援助，就能获得忠信而可靠的保护。当我们说\Quote{救我们免于凶恶}时，再也没有什么应当求的了。一旦我们祈求天主的保护免于凶恶，并获得了它，那么面对魔鬼和世界所施于我们的一切，我们就已安全稳固。因为今生有什么可惧怕的呢？对那个今生以天主为护守者的人而言。\par

28. 亲爱的弟兄们，如果这就是天主所教导的祈祷，这有什么可奇怪的呢？因为祂将我们所有的祈祷浓缩在一个救恩的句子里。这早已由\Person{依撒意亚}先知预言过，当他充满圣神，论及天主的威严和慈爱时说：\BibleQuote{完成并缩短祂的话，}祂说，\BibleQuote{以正义，因为主将在普世实行一句缩短的话。}\BibleRef{依 10:22} 因为当天主圣言、我们的主\Person{耶稣基督}来到所有人中间，聚集了有学问和没有学问的人，向每个性别和年龄的人颁布救恩的规诫时，祂将祂的规诫作了一个大的摘要，这样，学者的记忆不致在属天的学问中负担过重，而是能快速学会一个简单信德所必要的东西。因此，当祂教导永生时，祂以一种宏大而神圣的简洁包含了生命的圣事，说：\BibleQuote{永生就是：认识祢，唯一的真天主，和祢所派遣的耶稣基督。}\BibleRef{若 17:3} 同样，当祂要从法律和先知中提取第一条也是最大的诫命时，祂说：\BibleQuote{以色列，请听！主我们的天主是唯一的天主：你当全心、全灵、全力爱上主你的天主。这是第一条诫命。第二条与此相似：你当爱近人如你自己。}\BibleRef{玛 12:29} \BibleQuote{全部法律和先知都系于这两条诫命。}\BibleRef{玛 22:40} 又说：\BibleQuote{凡你们愿意人给你们做的，你们也要照样给人做。法律和先知即在于此。}\BibleRef{玛 7:12}\par

29. 主不仅以言语，也以行动教导我们祈祷，祂自己常常祈祷和祈求，从而用祂榜样的见证指示我们应当做什么，正如经上所记：\BibleQuote{祂却退到荒野的地方去祈祷。}\BibleRef{路 5:16} 又说：\Quote{祂出去，上山祈祷，彻夜向天主祈祷。}如果那没有罪的主尚且祈祷，那么罪人岂不更应该祈祷？如果祂不断祈祷，整夜不眠地持续祈求，那么我们岂不更应该每夜不断地反复祈祷呢？\par

30. 但主祈祷和祈求并非为自己——因为那无罪者为何要为自己祈祷呢？——而是为我们的罪，正如祂自己所说的，当祂对\Person{伯多禄}说：\BibleQuote{看，\Person{撒殚}想要筛你们像筛麦子一样。但我已为你祈祷，使你的信德不致丧失。}\BibleRef{路 13:31} 随后祂为众人向父祈求，说：\Quote{我不但为他们祈求，也为那些因他们的话而信从我的人祈求，愿众人都合而为一；父啊！就如祢在我内，我在祢内，使他们也在我们内合而为一。}主对我们的救恩的慈爱和怜悯何等伟大！祂不仅以祂的血救赎了我们，还额外为我们祈祷。请看，祂祈求的愿望是什么：要我们如同父与子合而为一一样，也保持在绝对的合一中；由此可以明白，那分裂合一与和平的人犯了多大的罪，因为连主也为此祈求，显然祂希望祂的子民这样得救并生活在平安中，因为祂知道纷争不能进入天主的国。\par

31. 此外，亲爱的弟兄们，当我们站着祈祷时，应当全神贯注，热切地投入我们的祈祷。让一切肉性和世俗的思想都消失，那时灵魂不应思想别的，只应思想祈祷的唯一对象。因此，司铎在祈祷之前，先以\Quote{请举心向上}这句话预备弟兄们的心灵，以便在民众回应\Quote{我们全心向上主}时，他得以被提醒：他自己也只应思想主。胸膛应关闭给仇敌，只向天主敞开；在祈祷时，不让天主的敌人接近它。因为他常常潜入我们中间，渗透进来，以狡猾的欺骗使我们祈祷偏离天主，使我们心中一事，口中另一事，而祈祷时，不是声音的响声，而是灵魂和心灵应以单纯的意向祈祷主。但如果你向主祈祷时，被愚昧和亵渎的思想所分心、所带走，那是何等的疏忽呢？仿佛有什么比与天主交谈更值得你思想！当你自己也不听你自己时，你怎能要求天主听你呢？你若自己也不记得自己，你怎能希望天主在你祈求时记得你呢？这绝对是毫不防备敌人；这是当你向天主祈祷时，以你祈祷的疏忽冒犯天主的尊威；这是眼睛醒着，心却睡着；而基督徒即使眼睛睡了，心也应当醒着，正如在\BibleBook{}{雅歌}中教会发言所记载的：\BibleQuote{我睡了，但我的心却醒着。}\BibleRef{歌 5:2}因此，宗徒焦虑而谨慎地警告我们说：\BibleQuote{恒心祈祷，并在此事上警醒；}\BibleRef{哥 1:2}即教导并表明，那些在天主眼中在祈祷中警醒的人，才能从天主那里获得他们所求的。\par

32. 此外，祈祷的人不应以空洞或没有果实的祈祷来到天主面前。当恳求是贫瘠的，祈求天主时是无效的。因为正如每棵不结果子的树被砍下投入火中；同样，不结果实的言语也不能从天主那里获得什么，因为它们没有任何成果。因此，圣经教导我们说：\Quote{祈祷与斋戒和施舍是好的。}\EditorialNote{多俾亚传 12:8}因为那位在审判之日将因我们的劳苦和施舍而赏报我们的主，在此生也是仁慈的听者，倾听那以善工相伴而祈祷的人。例如，百夫长科尔乃略祈祷时，理当被垂听。因为他常对百姓行许多施舍，并不断向天主祈祷。当他在第九时辰祈祷时，有一位天使显现，为他作证说：\Quote{科尔乃略，你的祈祷和你的施舍已达到天主面前，获得纪念。}\par

33. 那些祈祷迅速上升到天主面前，因为我们的劳苦功劳促使天主倾听。同样，拉斐尔天使也为多俾亚的恒常祈祷和善工作证说：\BibleQuote{显露和承认天主的事工是光荣的。因为当你和撒辣祈祷时，我把你们的祈祷带到天主的圣洁面前。当你单纯地埋葬死者，并且没有迟延起身离开晚餐，却出去遮盖死者时，我被派来试探你；天主又派遣我来治愈你和你的儿媳撒辣。我是拉斐尔，七位圣天使之一，常在天主荣耀前侍立出入。}\BibleRef{多 12:12}依撒意亚也提醒我们类似的事，说：\BibleQuote{解开一切不义的绳索，废除毫无能力的契约压迫，让困苦人安息，折断一切不义的契约。把你的饼分给饥饿的人，将无家可归的穷人领进你的房屋。见到赤身露体的，给他衣服穿；不要轻视同族同血脉的人。那时，你的光明要在季节中迸发，你的衣服要迅速生长；公义要走在你的前面，天主的荣耀要围住你。那时你呼求，天主必应允；你还在说话时，祂必说：我在这里。}\BibleRef{依 58:6}祂应许要亲近，说祂必垂听并保护那些从内心解开不义的绳索，按照祂的命令向天主家庭成员施舍的人，甚至因听从天主命令所行的事，他们自己也得蒙天主的垂听。真福宗徒保禄在患难中因弟兄们的帮助而满足时，说所行的善事是对天主的祭献。他说：\BibleQuote{我一切都够了，}他说，\BibleQuote{我已收到从厄帕夫洛狄托来的你们所送的东西，是馨香的祭品，悦纳中悦天主的牺牲。}\BibleRef{斐 4:18}因为怜悯穷人就是借给天主；凡给最微小者施舍的就是给天主施舍——向天主献上馨香的祭品。\par

34. 在履行祈祷职务时，我们发现\Person{达尼尔}与那三少年，信德坚固，在被掳中得胜，遵守第三、第六、第九时辰，如同为那在末期必须彰显的\OriginalTerm{天主圣三的圣事}{sacrament of the Trinity}。因为第一时辰进至第三时辰，显示出天主圣三的圆满数目；第四时辰进至第六时辰，又宣告了\OriginalTerm{另一个天主圣三}{another Trinity}；当第七至第九时辰完成时，完全的天主圣三每三时辰被计数。古时敬拜天主的人凭属灵辨别选定了这些时辰的分界，用作确定且合法的祈祷时间。后来这事显明了，这些事原是古时的圣事，因为古时义人正是这样祈祷。因为在第三时辰，圣神降临于门徒身上，应验了主应许的恩宠。此外，在第六时辰，\Person{伯多禄}上到屋顶，既藉着神视也藉着天主劝戒的话受教导，要他接纳众人获得救恩的恩宠，而他先前对外邦人领受圣洗感到疑惑。从第六时辰到第九时辰，主被钉在十字架上，以自己的血洗净了我们的罪；为救赎并使我们活过来，祂藉苦难完成了胜利。\par

35. 但是，亲爱的弟兄们，不仅古时所守的祈祷时辰，而且\OriginalTerm{圣事}{sacraments}的数目如今也增多了。因为我们也必须在早晨祈祷，好使主的复活藉早晨祈祷被庆祝。圣神从前在圣咏中已指出这点，说：\BibleQuote{我的君王，我的天主，因为我向祢哀号；上主，早晨祢必听见我的声音；早晨我必侍立在祢面前，并仰望祢。}主又藉先知的口说：\BibleQuote{清晨他们必为我守候，说：\InnerQuote{来，让我们归向上主我们的天主。}}\BibleRef{欧 6:1}同样，在日落和日暮时分，我们也必须再次祈祷。因为基督是真太阳和真白日，就如世俗的太阳和世俗的白日离去，当我们祈祷并祈求光明再次回到我们中间时，我们是在祈求基督的来临，祂将赐给我们永远之光的恩宠。此外，圣神在圣咏中显示基督被称为白日：祂说：\BibleQuote{匠人所弃的石头，已成了屋角的基石；这是主的作为，在我们眼中奇妙莫测。这是主所造的日子，我们要在这天行走并欢喜。}玛拉基亚先知也证明祂被称为太阳，他说：\BibleQuote{但为你们敬畏主名的人，正义的太阳将要升起，其翅膀有医治之能。}\BibleRef{拉 4:2}然而，如果在圣经中真太阳和真白日就是基督，那么没有一个时辰是基督徒不应该时常恒常地敬拜天主的；因此，我们这些在基督内——即在真太阳和真白日内——的人，应当全天候恳切祈求并祈祷；当按世俗的规律，旋转的黑夜以其交替变化来临，对祈祷的人不会造成任何损害，因为光明之子即使在黑夜也拥有白日。因为谁心里有光，会没有光呢？谁以基督为太阳和白日，会没有太阳和白日呢？\par

36. 那么，我们这些在基督内——即常在光明中——的人，即使在夜间也不可停止祈祷。如福音所记：\BibleQuote{她不离开}，经上说，\BibleQuote{圣殿，昼夜以禁食和祈祷事奉。}让那些尚未蒙光照的外邦人思想这事，或因离弃光明而仍留在黑暗中的犹太人。亲爱的弟兄们，我们这些常在主的光明中，铭记并持守因恩宠所领受而开始成为的，要将黑夜视作白日；要相信我们总在光明中行走，不要被我们已经逃脱的黑暗所阻碍。愿夜间祈祷的时辰没有缺失，不要闲散鲁莽地浪费祈祷的机会。我们已藉天主的仁慈在圣神内新造新生，要效法我们将来要成为的。既然在王国中我们将只拥有白日，而无黑夜介入，就让我们在夜间醒寤如同在白昼。既然我们要永远祈祷感恩，就让我们在此生也不停止祈祷感恩。\par

\SourceRecord{Translated by Robert Ernest Wallis.；Ante-Nicene Fathers；Vol. 5.；Edited by Alexander Roberts, James Donaldson, and A. Cleveland Coxe.；Buffalo, NY: Christian Literature Publishing Co.,；1886.；<http://www.newadvent.org/fathers/050704.htm>.}

% structure: p0001 source=h1 target=section reason=page-title

\section{第五篇论文}

\Emph{致\Person{德梅特里亚努斯}书。}\par

\Emph{论点：\Person{西彼廉}答复非洲总督\Person{德梅特里亚努斯}，后者声称当时的战争、饥荒和瘟疫应归咎于基督徒，因为他们不敬拜诸神；西彼廉公平地论证（首先指出万物随世界衰老而逐渐败坏），这些祸患的根源其实是外教人自己，因为他们不敬拜天主，并且不公义地迫害基督徒。}\par

1. \Person{德梅特里亚努斯}，我时常以轻蔑的态度对待你用亵渎的口和邪恶的话语反对独一真天主所发出的辱骂和喧嚣，认为更为谦逊和更好的方式是默默藐视一个错误之人的无知，而不是通过发言来激怒一个愚蠢之人的狂怒。我这样做并非没有神圣教导的权威，因为经上记载：\Quote{不要向愚人耳语，免得他轻视你话语的智慧。}\BibleRef{箴 23:9} 又说：\Quote{不要按愚昧回答愚昧，免得你也像他一样。}\BibleRef{箴 26:4} 此外，我们也被吩咐要将圣物保持在自己认知范围内，不要暴露给猪和狗践踏，因为主说：\Quote{不要把圣物给狗，也不要把你们的珍珠丢在猪前，免得它们用脚践踏，转过来撕裂你们。}\BibleRef{玛 7:6} 因为你常常来找我，其目的是反驳而非学习，宁愿无礼地坚持你自己的观点，大声嚷嚷，而不愿耐心听我的见解，我觉得与你争辩是愚蠢的；因为用呼喊平息汹涌大海的怒涛，比用论据遏制你的疯狂更容易、更轻微。确实，向盲人提供光明，向聋子说话，或向牲畜提供智慧，都是徒劳而无效的劳作；因为牲畜不能领悟，盲人不能接受光明，聋子不能听见。\par

2. 鉴于这一点，我常常保持沉默，以耐心胜过不耐烦的人；因为我既不能教导一个不可教的人，也不能以宗教阻止一个不敬虔的人，也不能以温和约束一个疯狂的人。但是，当你说很多人抱怨是我们导致了战争更频繁地发生、瘟疫和饥荒肆虐、长期干旱使雨水停止时，我不应再保持沉默，免得我的沉默开始被归因于不信任而非谦逊；在我轻视这些不实指控时，我可能显得是在承认罪行。因此，我答复你，\Person{德梅特里亚努斯}，也答复那些也许被你煽动的人，你通过散播恶意话语对我们制造仇恨，使他们成为你的党羽——他们源于你自己的根和本——但我相信他们会接受我言论的合理性；因为被谎言欺骗而趋向邪恶的人，更容易被真理的力量趋向良善。\par

3. 你说这一切都是由我们造成的，现在世界所遭受的动荡和苦难应归咎于我们，因为我们的神不被你们敬拜。在这方面，由于你不懂得神圣的知识，与真理陌生，你首先必须知道：世界现在已经衰老，不再具有它过去所拥有的力量；它也不再拥有从前那样的活力和力量。这一点，即使我们保持沉默，不从圣经和神圣的宣言中提出任何证据，世界本身现在也在宣告，并以其日渐衰败的证词见证自己的衰落。冬季不再有那样充足的雨水滋养种子；夏季的太阳不再有那样多的热量催熟庄稼；春季的麦田不再那样喜乐；秋季的果实也不再那样丰盛多叶。从被掏空和疲惫的山脉中开采出的大理石层更少；黄金和白银的减少暗示着金属的早期耗尽，贫瘠的矿脉日益受限和缩减；农夫在田野中衰退，海员在海上衰退，士兵在军营中衰退，市场上的清白，法庭上的正义，友谊中的和谐，技艺中的熟练，道德中的纪律。你认为一个正在衰老的事物的实质还能像它年轻时、在崭新和充满活力时那样健壮吗？任何趋向衰败、接近终结的事物必然会减弱。因此，太阳在落山时放射的光芒不那么明亮和炽烈；因此，月亮在逐渐消逝的轨道上以耗尽的双角亏缺；先前青翠多产的树木，随着枝条干枯，渐渐在荒芜的老年中变得畸形；从前从溢出血管中喷涌而出的泉水，因衰老而枯竭，几乎涓滴不流。这就是对世界的判决，这就是天主的律法：凡有开始者必会消逝，凡生长者必会衰老，强者变弱，大者变小，当它们变得衰弱和减少时，就会走到尽头。\par

4. 你们归咎于基督徒，说万物随着世界衰老而朽坏。倘若老年人也归咎于基督徒，说他们在老年时力气减弱；他们不再像从前那样拥有同样的便利——耳聪、足捷、目明、体健、器官活力充沛、肢体丰满——尽管古时人的寿命曾超过八九百年，如今却几乎不能到达百岁。我们在孩童身上看见白发——头发尚未生长就已脱落；生命不是在老年才终止，而是从老年开始。因此，即使是在初始阶段，出生也奔向终结；\BibleRef{智 5:13}因此，凡是现在诞生的，都伴随着世界本身的衰老而衰败，以致无人应当惊奇，万物在世界里开始朽坏，而整个世界本身已处于朽坏之中，临近终结。\par

5. 此外，战争频繁肆虐，死亡与饥荒加重忧虑，健康因猛烈的疾病而受损，人类因瘟疫的荒废而消耗——要知道，这些早已被预言：在末期，祸患将增多，灾难将多样；随着审判之日临近，愤怒的天主为了鞭策人类，其谴责将越发激起。因为这些事情的发生，并非如你们虚假的抱怨以及对真理无知的缺乏经验所断言和重复的那样，是由于我们不敬拜你们的神祇，而是由于你们不敬拜天主。既然祂是世界的\OriginalTerm{主}{Lord}与统治者，万事都凭祂的旨意和指引运行，除非祂自己已经成就或允许的，便不能成就任何事，那么，当那些显示被冒犯的天主之怒的事情发生时，它们并非因我们这些敬拜天主的人而出现，而是因你们的罪孽和应得惩罚所招致——你们既不寻求天主，也不敬畏祂，因为你们没有抛弃虚妄的迷信，也没有认识真宗教，以致那唯一统御万有的天主独自受到敬拜和祈求。\par

6. 最后，请听祂亲自说话；祂以神圣的声音同时教导和警告我们：\Quote{你要敬拜上主你的天主，}祂说，\Quote{唯独事奉祂。}\BibleRef{申 6:13}又说：\Quote{除我以外，你不可有别的神。}\BibleRef{出 29:3}又说：\Quote{不可随从别的神，事奉敬拜它们，不可因你们手所做的惹我发怒，以致我毁灭你们。}\BibleRef{耶 25:6}此外，那充满圣神的先知也为天主之怒作证并宣告说：\Quote{万军的上主这样说：因为我的殿宇荒凉，你们各人却奔向自己的房屋，所以天应抑制甘露，地应停止出产；我必要使刀剑临于大地、五谷、新酒、油、人类、牲畜，以及他们手所作的一切。}\BibleRef{盖 1:9}又有另一位先知重复说：\Quote{我必降雨在一座城，在另一座城却不降雨；一块田地得雨，另一块未得雨的田地必干枯。两三个城的人要聚到一个城去喝水，却不得满足；你们仍不归向我——上主说。}\BibleRef{亚 4:7}\par

7. 看哪，上主发怒忿恨，发出威胁，因为你们不归向祂。而你们在这顽梗和藐视中，竟对以下情形感到惊奇或抱怨：雨水罕见地缺乏，土地因尘土腐朽而荒芜；贫瘠的田地几乎长不出一把纤弱苍白的小草；毁灭性的冰雹摧残葡萄树；铺天盖地的旋风连根拔起橄榄树；干旱使泉水枯竭；疫风败坏空气；疾病的虚弱耗损人的生命——尽管这一切都是因为引诱它们而来的罪恶所导致，而天主更加愤怒的是，如此重大众多的灾祸竟毫无功效！因为这些事情的发生，或是为惩戒顽梗者，或是为惩罚恶人，同一位天主在圣经中宣告说：\Quote{我白白地击打了你们的儿女，他们却没有受教。}那位将自己奉献给天主的先知也用同样的语调回应这些话，说：\Quote{祢击打他们，他们却不忧伤；祢责打他们，他们却拒绝受教。}\BibleRef{耶 5:3}看哪，来自天主的鞭刑降临，却没有对天主的敬畏；看哪，来自高天的打击和责打并不缺乏，却没有战栗，没有畏惧。倘若连这种斥责都不干预人间事务，那将会怎样？如果人类在犯罪中因不受惩罚而安全，他们的狂妄还会更大多少呢！\par

8. 你们抱怨说，如今泉水对你们不再充足，和风不再有益健康，频繁的降雨和肥沃的土地对你们的帮助不如从前；诸元素不再像古时那样服务于你们的用途和享乐。但你们自己是否事奉天主？一切事物都由祂安排为你们服务；你们是否听从祂？万事都凭祂的美意服务于你们。你们从你们的奴隶那里要求服务；虽然你们是人，却迫使你们同类的人屈服并服从你们；尽管你们在出生方面有同样的命运，在死亡方面有同样的境况，尽管你们有同样的肉身和共同的灵魂秩序，尽管你们以同样的权利和同样的法律进入这个世界，并在一定时期后离去，然而，除非奴隶按你们的喜好服侍你们，除非奴隶按你们的意愿服从你们，你们作为傲慢而过分要求服务的人，就鞭打和责打他：你们用饥饿、干渴和赤身露体折磨他，甚至常常用刀剑和监禁来迫害他。你们这些可怜的人啊！你们自己如此施行统治，难道不承认你们的天主是主吗？\par

9. 因此，在这些发生的灾祸中，天主的鞭打和惩治并非欠缺；既然它们在现今的事上毫无益处，也不能以这种毁灭的恐怖使人归向天主，那么最终仍留下永久的黑暗、不灭的火和永远的刑罚；在那里，哀求者的呻吟将不被垂听，因为在此处，发怒的天主藉着先知呼喊说：\BibleQuote{你们以色列子民，要听上主的话：因为上主要与这地的居民争辩；因为在这地上没有仁慈，没有真理，也没有认识天主的知识。只有诅咒、谎言、凶杀、偷窃、奸淫，纷争遍地，他们以血还血。因此，这地要悲哀，凡住在其上的，田野的走兽、爬行的昆虫和天空的飞鸟，都必衰残；海中的鱼也必消灭，以致无人判断，无人责备。}\BibleRef{欧 4:1} 天主说祂是愤怒且震怒的，因为在这地上没有人承认天主，天主既不为人所知，也不为人所畏惧。天主斥责并责备说谎、情欲、欺诈、残忍、不敬和愤怒的罪，却无人归向纯真。看，那些先前藉天主的话所预言的事正在发生；却没有人因相信眼前的事而受劝诫，去思想将来的事。就在那些灵魂被紧紧捆绑、封闭、几乎无法呼吸的灾难中，人仍有作恶的机会，在如此巨大的危险中，他们更多地判断别人而非自己。你愤慨于天主的愤怒，仿佛你邪恶的生活还配得什么好处，仿佛所发生的这一切事情与你的罪相比，都是无限微小、不值一提的。\par

10. 你这判断别人的人，也请暂时作自己案件的判断者；查看你自己良心的隐秘处；不，既然如今你在罪中毫无羞耻，你如此邪恶，仿佛正是邪恶本身取悦了你，你，被所有其他人看得清楚赤裸，也该自己看一看自己。因为你或是傲慢膨胀，或是贪婪吝啬，或是愤怒残忍，或是赌博挥霍，或是醉酒脸红，或是嫉妒猜忌，或是邪淫不洁，或是强暴残酷；而你竟奇怪天主的愤怒在惩罚人类时加剧，而受罚的罪恶每天在增加？你抱怨敌人兴起，仿佛即使没有敌人，你在安定的法袍中也能有平安。你抱怨敌人兴起，仿佛即使外来的兵器和来自蛮族的危险被压制，那来自权势公民的毁谤和冤屈的家庭攻击的武器，在内里不会更凶残、更严酷地使用。你抱怨荒芜和饥荒，仿佛干旱造成的饥荒大于贪婪，仿佛匮乏的严酷不会因囤积年产的涨价和积压而更加可怕。你抱怨天关闭降雨，同样地，地上的谷仓也关闭。你抱怨现在出产减少，仿佛已经出产的都给了穷人。你责备瘟疫和疾病，而瘟疫和疾病本身却揭露或增加了个人的罪行，对弱者没有表示怜悯，而贪婪和掠夺张口等待死人。同样的人在仁爱的职责上胆怯，却在追求无耻利益上鲁莽；躲避垂死者的死亡，却渴望死者的掠夺，以致看起来像是可怜人在病中被抛弃，正是为了不让他们被治愈而逃脱：因为那急切进入垂死者产业的人，\Emph{或许}希望那病人死去。\par

11. 如此巨大的毁灭恐怖不能教导纯洁；在因持续浩劫而死去的人民中间，没有人想到自己也会死。到处是抢夺、攫取、占有；对劫掠没有掩饰，也没有延误。仿佛这一切都是合法的，仿佛这一切都是合理的，仿佛不抢劫的人就在亏损、浪费自己的财产，因此每个人都急忙去掠夺。盗贼之间尚有一些犯罪时的羞耻。他们喜欢荒僻的峡谷和无人居住的荒野；他们行恶的方式，仍然使作恶者的罪被黑暗和黑夜掩盖。\Emph{然而}贪婪却公开肆虐，并且因着自己的大胆而安全，在市场的光天化日之下暴露出其急切的贪欲武器。于是骗子、毒害者、刺客，在城中央，对邪恶的渴望如同他们行恶而免罚。罪由犯罪者所犯，却找不到无罪的惩治者。没有对原告或法官的惧怕：恶人获得免罚，而谦逊者保持沉默；同谋者害怕，而审判者却可以买卖。因此，藉先知的口，神圣的精神和本能提出了事情的真相：它以一种确定而明显的方式表明，天主能够阻止不利之事，但对罪人的邪恶行为阻止祂施以救援。\BibleQuote{上主的手，}他说，\BibleQuote{岂是缩短而不能拯救？或是耳朵沉重而不能听见？但你们的罪恶使你们与天主隔绝；因为你们的罪恶，祂掩面不理你们，不肯怜悯。}\BibleRef{依 59:1} 因此，让你的罪过和过犯被数算清楚；让你良心的伤口被考虑；并且让每一个人停止抱怨天主，或者抱怨我们，如果他应该意识到自己所受的是应得的。\par

12. 请看，这正是我们主要谈论的事情——你们骚扰我们，尽管我们是无辜的；你们藐视天主，攻击并压迫天主的仆人。按\Emph{你们的说法}，这不算什么：你们的生命染上各种粗鄙的恶习、致命罪行的不义、一切血腥掠夺的总结；真正的宗教被虚假的\OriginalTerm{迷信}{superstitions}推翻；天主既不被寻求，也不被敬畏；然而除此之外，你们还用不义的迫害折磨天主的仆人，就是那些奉献于祂的威严和祂的名的人。你们自己不禁拜天主，这还不够，你们还用\OriginalTerm{亵渎神圣的敌意}{sacrilegious hostility}迫害那些敬拜祂的人。你们既不敬拜天主，也完全不允许祂被敬拜；而那些不仅敬拜人手所造的愚蠢偶像和图像，甚至敬拜怪异和怪物的人，却令你们喜悦；唯独天主的敬拜者令你们不悦。祭品的灰烬和牲畜的堆堆在你们的殿里处处冒烟，而天主的祭台却无处可寻，或被隐藏。鳄鱼、猿猴、石头、蛇被你们敬拜；唯有天主在地上不被敬拜。或者，若被敬拜，也不能不受惩罚。你们剥夺无辜、义人、天主所爱的人的家园；你们掠夺他们的财产，给他们戴上锁链，把他们关进监狱，用刀剑、野兽、火焰惩罚他们。的确，你们不满足于我们短暂的受苦和简单迅速的痛苦耗尽。你们通过撕裂我们的身体施加持久的酷刑；通过撕裂我们的内脏增加无数的惩罚；你们的残暴和凶猛不能满足于通常的酷刑；你们精巧的残酷发明了新的痛苦。\par

13. 这种嗜血的疯狂是何种无法满足的欲望？这种残忍的贪欲是何种无止境的渴求？不如从两个选项中选一个：成为基督徒要么是\OriginalTerm{罪行}{crime}，要么不是。如果是罪行，为什么你们不处死那\OriginalTerm{宣认}{confession}的人？如果不是罪行，为什么你们迫害无辜的人？因为我若否认，才应该受酷刑。如果我因害怕你们的惩罚，用欺骗的谎言隐瞒我从前是什么，以及没有敬拜你们的神的事实，那么我才该受折磨，我才该被强迫通过受苦的力量认罪，如同在其他审讯中，被指控犯有罪行而否认的罪犯，为了从其身体痛苦中逼出罪行真相的招认（因为告密的声音拒绝说出）而受酷刑。但现在，当我自愿地宣认，呼喊，并用同样的话再三作证我是基督徒时，你们为什么对承认的人施加酷刑？你们的神不是被我在隐秘和秘密的地方，而是公开地、当众地、就在市场里、在你们的官长和总督面前摧毁；因此，尽管你们先前指责我的事是轻微的，但那你们更应仇恨和惩罚的事却增加了，因为我在人多的场所，在百姓环绕下，公开宣告自己是基督徒，从而以公开的宣告使你们和你们的神都困惑？\par

14. 为什么你们把注意力转向我们身体的软弱？为什么你们与这属土的肉身的脆弱争斗？不如与\OriginalTerm{理智}{mind}的力量争战，打破灵魂的能力，摧毁我们的\OriginalTerm{信德}{faith}，如果你们能通过讨论征服，通过理性战胜；或者，如果你们的神有任何神性和能力，就让他们起来为自己辩护，让他们以自己的威严自卫。但他们能为崇拜者带来什么益处呢？如果他们不能向不崇拜他们的人报仇？因为如果报仇者比被报仇者更大，那么你们就比你们的神更大。如果你们比你们所敬拜的更大，你们就不该敬拜他们，反而该被他们当作主来敬拜和畏惧。你们的保护在神受辱时为他们辩护，正如你们的看顾在神被囚时保护他们免于灭亡。你们应当羞于敬拜那些你们自己保护的神；你们应当羞于向那些你们自己保护的神寻求保护。\par

15. 哦，但愿你们能听见并看见，当他们被我们驱魔时，被\OriginalTerm{属灵的鞭打}{spiritual scourges}折磨，并被话语的折磨从\OriginalTerm{附魔的}{possessed}身体被驱逐出去，他们在人的声音和天主的大能面前嚎哭呻吟，感受鞭打和击打，承认将来的审判！来吧，承认我们所说的是真的；既然你们说你们这样敬拜神，就相信你们所敬拜的。或者，如果你们甚至相信自己，那么他——即那魔鬼——现在已经占据了你们的心胸，用无知的黑暗蒙蔽了你们的思想，将在你们耳中谈论你们自己。你们会看到，那些你们恳求的，反来恳求我们；那些你们惧怕的，反来惧怕我们，就是你们所崇拜的。你们会看到，在我们手下，他们被捆绑，如同俘虏战栗，而你们却尊崇他们为主人：确实，即使这样，你们也会在那些错误中感到困惑，当你们看见并听见你们的神，在我们质问下立刻暴露他们是什么，甚至在你们面前也无法隐藏他们的欺骗和诡计。\par

16. 那么，心智的迟钝是什么？不，愚昧者盲目的、愚蠢的疯狂是什么？他们不愿从黑暗出来进入光明，不愿在永生死亡的束缚中接受永生的希望，也不畏惧天主威胁说：\BibleQuote{祭祀任何神祇，唯独不祭祀上主的，应被铲除。} \BibleRef{出 22:20} 又说：\BibleQuote{他们崇拜自己手指所造的；卑微人屈身，尊贵人自谦，我必不宽恕他们。} \BibleRef{依 2:8} 你为何向假神卑躬屈膝？为何在愚昧的偶像和地上的受造物前俯首为奴？天主造你是正直的；其他动物向下看，弯曲着身体朝向大地，而你的姿态是高昂的；你的面容朝上仰望天空和天主。向那里看，举目仰望，在至高之处寻求天主，好使你摆脱下面的事物；把你的心提升，依靠高天的事。你为何要与你所崇拜的蛇一同跌入死亡的毁灭？为何你藉着魔鬼的手段和与它一同堕入魔鬼的毁灭？持守你生来所拥有的崇高地位。保持天主所造的你本来的样子。使你的灵魂与你的面容和身体的姿态一致。为能认识天主，先认识你自己。弃绝人类错误所发明的偶像。归向天主，你若呼求祂，祂必帮助你。相信基督，圣父派遣祂来使我们复活并恢复。停止用你们的迫害伤害天主和基督的仆人，因为当他们受伤害时，天主的复仇会保护他们。\par

17. 因此，我们中没有人被捕时抵抗，也没有人报复你们的不义暴行，尽管我们的人数众多而丰富。我们对紧随其后的复仇的确信使我们忍耐。无辜者给有罪者让路；无害者顺从刑罚和苦刑，确信并相信我们所受的任何苦难都不会不被报复，而且我们受迫害的不义越大，对这些迫害所施行的复仇的正义和严厉也越大。不敬神者的邪恶也从未在我们所持的名字兴起，而不立即受到来自高天的惩罚。且不提古时的记忆，也不以言语的纪念重复提及为天主的敬拜者所施行的频繁复仇，最近的一件事例足以证明我们的辩护如此迅速，且在其迅速中如此有力，随后在事物的废墟、财富的毁灭、士兵的损失和堡垒的减少中。且不要有人认为这是偶然发生的，或以为是巧合，因为圣经早已规定并说：\BibleQuote{复仇在我，我必报应，主说。} \BibleRef{罗 12:19} 圣神又预告说：\BibleQuote{不要说，我要向我的敌人报仇，而要等候主，祂必帮助你。} \BibleRef{箴 20:22} 由此显而易见，这一切从天主的愤怒而来的事，不是藉着我们，而是为了我们而发生的。\par

18. 且不要让任何人认为，基督徒因这些正在发生的事而没有得到复仇，理由是他们自己似乎也受到了这些事的打击。一个人感受到世俗灾祸的惩罚，是因为他所有的喜乐和荣耀都在世界上。他若在此生遭遇不幸，便悲伤叹息，因为在此生之后他不能有福，他生命的全部果实都在这里收取，他的一切安慰在此结束，他短暂而转瞬即逝的生命在此享受一些甜蜜和快乐，但当它离开这里后，留给他的只有加上悲伤的惩罚。而那些对未来好事有信心的人，则不受当前邪恶攻击之苦。事实上，我们从未因逆境而沮丧，也未被打倒，也不在任何外在的不幸或身体的软弱中悲伤或抱怨：我们依靠圣神而非血肉生活，以心智的力量克服身体的软弱。正是那些折磨和困扰我们的事，我们知道并相信我们因此被证明和加强。\par

19. 你们以为我们与你们一样遭受逆境吗？当你们看见同样的逆境在我们和你们身上承受得不一样时。在你们中间总是有喧闹和抱怨的不耐烦；我们却有坚强而虔诚的忍耐，总是安静并常感谢天主。它在此世不为自己要求任何快乐或顺利的事，而是温良、柔和、稳固地面对这动荡世界的所有风暴，等待天主应许的时刻；因为只要这身体存在，它必然与他人有共同命运，其身体状态也是共同的。任何人类成员都不能彼此分离，除非离开此世。在此期间，我们所有人，善与恶，都包含在同一家室中。凡屋内所发生的事，我们都以同等的命运承受，直到现世生命的终结，我们将被分配到永死或永生的家室中。因此，我们与你们不在同一层次，也不平等，因为我们处在此世和此肉身中，与你们同样遭受世界和肉身的烦恼；既然痛苦的感觉就是一切惩罚，显然，你若看见他不与你们同受痛苦，他就不是你们惩罚的分享者。\par

20. 在我们中间，望德的力量与信德的坚定兴盛。在这日益衰残的世界废墟中，我们的灵魂被高举，我们的勇气不动摇；我们的忍耐始终是喜乐的，心灵常安处于天主，正如圣神藉先知所说，并用天上的言语坚固我们希望与信德的坚定：\Quote{无花果树，}祂说，\Quote{必不结果，葡萄树上必无花蕾。橄榄树的工作必失效，田地也不再出产粮食。羊群必从圈中被剪除，牛棚中也无牲畜。但我必因上主而喜乐，因我救恩的天主而欢欣。} \BibleRef{哈 3:17} 祂说，天主的人和崇拜天主的人，依靠其希望的真理，并坚固于其信德的坚定，不为这世界和今生的攻击所动摇。纵使葡萄树枯干，橄榄树欺哄，田地因草枯干而焦裂，这对基督徒又算什么呢？对天主的仆人又算什么呢？天堂正邀请他们，天主之国的全部恩宠和丰盈正等候他们。他们常因上主而欢欣，因他们的天主而喜乐；他们勇敢地忍受世间的邪恶与磨难，因为他们期待将来的恩赐与福乐：因为我们已经脱去地上的出生，现在由圣神所创造和重生，不再为世界而活，乃为天主而活，直到我们到达天主面前，才能领受天主的恩赐与应许。然而，我们常祈求仇敌的退却，求得甘霖，或祈求磨难的消除或缓和；我们倾心祈祷，求天主息怒，日夜不停地、迫切地、为你（们）的平安与救恩而恳求。\par

21. 然而，没有人应当自欺，以为因肉身的平等，我们与不敬神的人、天主的敬拜者与天主的反对者，在世俗苦难中有共同的境遇，就由此认为所有这些事并非由你们招致；因为藉天主自己的宣告和先知的见证，早已预言了天主的忿怒要临到不义的人，而那些在人性上会伤害我们的迫害也不会缺少；不仅如此，那以天主的保护来保护受害者的报复，也必伴随他们。\par

22. 此外，在这一点上，此时此刻为我们发生的事又是何等重大！有些事被赐下作为鉴戒，使复仇之天主的忿怒得以显明。但那审判的日子尚未来临，正如圣经所宣告的：\Quote{哀号吧！因为上主的日子临近了，毁灭从天主而来；看哪，上主的日子来到，残酷、忿怒、烈怒，要使大地荒凉，从其中除灭罪人。} \BibleRef{依 43:6} 又说：\Quote{看哪，上主的日子来到，如同火炉；所有外邦人和所有行恶的人必如碎秸，那来到的日子必烧尽他们，上主说。} \BibleRef{拉 4:1} 上主预言外邦人必被烧尽灭绝；外邦人就是那些与神圣族类疏远、不敬神的人，他们没有在圣神里重生，也没有成为天主的子女。因为在另一处，当上主差遣天使去毁灭世界和杀死人类时，祂更严厉地恐吓末时的情景说：\Quote{你去击打，不要顾惜你的眼。不要怜悯老少，要杀尽处女、孩童和妇女，使他们彻底灭亡。但任何人身上有印记的，都不可触摸。} \BibleRef{则 9:5} 因为只有那些重生并印上基督记号的人才能逃脱。此外，这印记是什么，印在身体上何处，天主在另一处说明：\Quote{你走遍耶路撒冷城中，在那些为城中所有可憎之事叹息哀哭的人额上画一个记号。} \BibleRef{则 9:4} 这记号与基督的苦难和血有关，凡在这记号中被寻见的人都得保全、不受伤害，这一点也由天主的证言所证明：\Quote{这血要在你们所住的房屋上作记号；我一见这血，就保护你们，当我去击打埃及地的时候，那毁灭的灾祸必不临到你们身上。} \BibleRef{出 12:13} 先前在被杀羔羊中以预表所预示的，后来藉着真理基督应验了。因此，如同埃及受击打时，犹太百姓只能借着羔羊的血和记号得逃脱；同样，当世界开始被荒废和击打时，只有被寻见在基督的血和记号中的人才能逃脱。\par

23. 因此，趁着还有时间，要留意那真实而永恒的救恩；既然如今世界的终结临近，就当将你们的心志转向天主，存着敬畏天主的心；也不要让那在世上对正义和温和之人无能的、虚妄的统治令你们喜悦，因为在田中，甚至在改良和结果子的谷物中，莠子和毒麦也占统治地位。不要说厄运发生是因为我们不敬拜你们的神明；要明白这是天主怒气的审判，使那因祂恩惠而不被承认的，至少藉着祂的审判而被承认。即使迟了，也要寻求上主；因为很久以前，天主藉祂的先知警告并劝告说：\Quote{寻求上主，你们的灵魂就必存活。} \BibleRef{亚 5:6} 即使迟了，也要认识天主；因为基督在祂来临时教导并教训这件事，说：\Quote{永生就是：认识祢，惟一的真天主，和祢所派遣的耶稣基督。} \BibleRef{若 17:3} 要相信那绝不欺骗人的那一位。要相信那预言这一切事必将发生的那一位。要相信那赐给所有相信的人永生奖赏的那一位。要相信那要叫不信的人承受地狱之火中永远的刑罚的那一位。\par

24. 那时，信德的光荣将是什么？无信的惩罚又是什么？审判之日来临，信徒的喜乐、不信者的悲哀将是什么？他们在此不愿相信，如今却无法回头去信！永不熄灭的\OriginalTerm{地狱}{Gehenna}将焚烧被定罪者，以活火吞噬的刑罚；不会有任何来源能让他们在折磨中得喘息或终结。灵魂与身体将被保留在无尽的痛苦中受苦。如此，那曾在此一时注视我们的人，将永远被我们看见；那些残忍目光在迫害我们时所获得的短暂喜乐，将被永恒的景观所补偿，按照圣经的真理所说：\BibleQuote{他们的虫不死，他们的火不灭；他们将成为所有血肉的观瞻。}\BibleRef{依 66:24} 又说：\BibleQuote{那时，义人将大义凛然地站在那些曾经磨难过他们、剥夺他们劳苦的人们面前。他们一见，必惊慌失措，害怕这意想不到的救恩；他们痛悔哀叹，由于心神的苦痛，彼此说：这就是我们曾一度嘲笑过、讥诮的谚语；我们愚昧的人曾将他们的生活视为疯狂，将他们的结局视为耻辱。怎么他们列在天主子女中，他们的分是在圣者中？因此我们离弃了真理的道路，义德的光没有照耀过我们，太阳没有为我们升起。我们在邪恶与毁灭的路上感到疲劳；我们走过了无路的荒野，却没有认识上主的道路。骄傲为我们有什么益处？财富的夸耀又为我们有什么好处？这一切都像阴影逝去。}\BibleRef{智 5:1} 惩罚的痛苦将不再有忏悔的果实；哭泣无用，祈祷无效。他们不相信永生，如今却太迟地相信永罚。\par

25. 因此，趁着还有能力，为你的安全与生命做好准备。我们向你提供我们心智与劝告的助益。因为我们不可怀恨，我们不回报恶行更令天主喜悦，我们劝告你，趁你还有力量，趁你生命中尚存些许，向天主补赎，从黑暗迷信的深渊中出来，进入真正宗教的光明。我们并不嫉妒你的安逸，也不隐瞒天主的恩赐。我们以善意回报你的仇恨；为那些加诸我们的折磨与刑罚，我们向你指出救恩的途径。相信而生活，你们这些迫害我们的人，要在永恒中与我们一同喜乐。一旦你离去，那里不再有悔改之处，也不可能补赎。在此生命要么丧失，要么得救；永恒的安全由对天主的崇拜和信德的果实提供。任何人都不要因自己的罪过或年岁而退缩，不来获得救恩。对仍留在此世界的人，没有太迟的悔改。接近天主的仁慈是敞开的，寻求和领悟真理的人易于进入。你要为你的罪祈求，即使在生命的尽头、时间的太阳西沉之时；以宣认和信仰承认祂，向独一真神天主恳求，告解的人获得宽恕，因天主的良善将救恩的仁慈赐予信徒，不死之路向死亡本身敞开。这恩宠由基督赐予；祂藉\OriginalTerm{十字架的胜利标记}{trophy of the cross}战胜死亡，以祂血的代价救赎信徒，使人与天主圣父和好，以天上的重生使我们必死的本性活跃起来。若可能，让我们都跟随祂；让我们登记在祂的圣事和标记中。祂为我们开启生命之路；祂领我们回到乐园；祂引领我们进入天国。由祂成为天主的子女，我们将与祂永远生活；与祂永远喜乐，因祂自己的血而恢复。我们基督徒将与基督一同荣耀，蒙天主圣父祝福，在天主面前永远以永恒的喜乐欢欣，并永远感谢天主。因为，那曾一度服从死亡、如今已在拥有永生中得安稳的人，不会不时刻喜乐与感恩。\par

\SourceRecord{Translated by Robert Ernest Wallis.；Ante-Nicene Fathers；Vol. 5.；Edited by Alexander Roberts, James Donaldson, and A. Cleveland Coxe.；Buffalo, NY: Christian Literature Publishing Co.,；1886.；<http://www.newadvent.org/fathers/050705.htm>.}

% structure: p0001 source=h1 target=section reason=page-title

\section{第六论}

论偶像的虚妄\par

指出偶像并非神，天主是唯一，且藉着基督将救恩赐予信徒。\par

\Emph{论证：本标题包括这篇论著的三个主要部分。作者首先指出，为他们建立庙宇、雕刻雕像、献祭牺牲、庆祝节日的人，是君王和人，而非神；因此，对他们的崇拜无论对外乡人还是罗马人都毫无用处，而且罗马帝国的力量应归因于命运而非他们，因为它是靠某种好运兴起的，并对自己的起源感到羞愧。}\par

1. 凡平民所敬拜的，并非神，从这一点可知。他们曾是君王，因着对君王的纪念，他们的百姓甚至在死后也开始敬拜他们。于是为他们建立庙宇，为他们雕刻塑像，以相似留存亡者的面容；人们献祭牺牲，庆祝节期，以给予他们尊荣。于是，这些起初作为安慰而采纳的礼仪，传至后代便成为神圣。现在让我们看看，这一真理是否在每个具体事例中得以证实。\par

2. \Person{墨利刻忒斯}和\Person{琉科忒亚}被投入海中，随后成为海神。\Person{卡斯托耳兄弟}轮流死亡，为能存活。\Person{埃斯库拉庇俄斯}被闪电击中，为能升为神。\Person{赫拉克勒斯}在\Place{俄塔山}的大火中被烧尽，为能脱去人身。\Person{阿波罗}牧养\Person{阿德墨托斯}的羊群；\Person{涅普顿}为\Person{拉俄墨冬}建造城墙，并——不幸的建造者——未得工价。\Person{朱庇特}的洞穴在\Place{克里特}可见，他的坟墓也被展示；显然\Person{萨图尔努斯}被他赶走，并且\Place{拉丁姆}由此得名，作为他的藏身之处。他首先教人刻印文字；他首先在\Place{意大利}教人铸造钱币，从此金库被称为萨图尔努斯金库。他也是田园生活的耕耘者，因此他被画成一位手持镰刀的老人。\Person{雅努斯}在他被赶走时款待了他，\Place{雅尼库鲁姆山}因而得名，一月也被设立。他自己被描绘成两张面孔，因为处于中间，他似乎同等地看向年初和年末。摩尔人确实明显敬拜君王，并不以任何伪装掩饰其名。\par

3. 由此，众神的宗教在各民族和各省中变化各异，因为没有人敬拜同一位神，而是每个人保留对自身祖先的独特崇拜。\Person{亚历山大大帝}在他写给母亲的那部非凡卷册中证实了这一点：因惧怕他的权势，一位祭司向他揭示了曾秘而不宣的神即是人的道理，即被保持的实为对祖先和君王的纪念，而崇拜和献祭的礼仪由此而生。但若神曾在某个时候诞生，为何现在不再诞生呢？——除非，\Person{朱庇特}确实已太老，或\Person{朱诺}已失去生育能力。\par

4. 但为何你以为众神能对罗马人有益，当你看到他们对自己的崇拜者在对抗罗马军队时无能无力？因为我们知道罗马的神是本地的。\Person{罗慕路斯}凭\Person{普洛库鲁斯}的伪誓而为神，\Person{皮库斯}、\Person{提贝里努斯}、\Person{皮卢姆努斯}和\Person{孔苏斯}，罗慕路斯要人敬拜一位背信之神——孔苏斯，仿佛他是谋略之神，因为他的背信导致了抢掠萨宾妇女。\Person{塔提乌斯}也发明并敬拜女神\Person{克洛阿基娜}；\Person{霍斯提利乌斯}发明并敬拜恐惧和苍白。后来，不知由谁奉献了热病，以及娼妓\Person{阿卡}和\Person{弗洛拉}。这些就是罗马的神。但\Person{玛尔斯}是色雷斯人，\Person{朱庇特}是克里特人，\Person{朱诺}或是阿耳戈斯人、萨摩斯人或迦太基人，\Person{狄安娜}是\Place{陶鲁斯}的，众神之母是\Place{伊达山}的；还有埃及怪物，并非神明，它们若有能力，必会保全自己和本族的王国。当然在罗马人中也有被征服的家神，是逃亡的\Person{埃涅阿斯}引入那里的。还有秃头\Person{维纳斯}——因她在\Place{罗马}的秃头而受的羞辱远甚于在\Person{荷马}史诗中受伤的羞辱。\par

5. 王国并非凭功绩兴起至霸权，而是因机遇而变。帝国先前由亚述人、米底人和波斯人统治；我们也知道希腊人和埃及人都曾拥有统治权。因此，在权力的兴衰变迁中，帝国时期也像其他民族一样临到了罗马人。但若你回顾其起源，你必然脸红。一群放荡者和罪犯聚集在一起，通过建立庇护所，使罪行免罚而人数大增；并且为使其王本身在罪行上更有优势，\Person{罗慕路斯}成了弑兄者；为促进婚姻，他以不和开启了那件和好的事。他们偷窃、施暴、欺骗，以增加国家的人口；他们的婚姻建立在违背待客之盟约和与岳父们的残酷战争之上。此外，执政官职位是罗马荣誉的极致，然而我们看到，执政官的开始甚至与王国相同。\Person{布鲁图斯}处死自己的儿子，以使对其凶恶的认可增加对他尊严的称许。因此，罗马王国并非从宗教的神圣、也非从征兆和占卜中成长，而是在确定的界限内持守其指定的时间。再者，\Person{雷古鲁斯}观察了征兆，仍被俘；\Person{曼基努斯}遵守了宗教义务，仍被遣送轭门下；\Person{保卢斯}的圣鸡进食了，他仍死于\Place{坎尼}；\Person{盖乌斯·恺撒}藐视反对他在冬季前将船队派往\Place{非洲}的占卜和征兆，然而他却更顺利地航行了并取得了胜利。\par

6. 然而，这一切的根源是相同的，它误导并欺骗，用蒙蔽真理的伎俩，将轻信愚昧的民众引入歧途。这些是不洁和游荡的灵，他们在尘世的恶习中浸透之后，因大地的沾染而丧失了属天的活力，自身堕落之后，仍不停寻求他人的毁灭；自身卑劣之后，又将自身堕落的错误灌输给他人。诗人们也承认这些魔鬼，苏格拉底声称自己受一个魔鬼的指示和支配；由此，贤士们便有了作恶或嘲弄的能力，其中为首的霍斯塔内斯既说无法看见真天主的形象，又宣告真天使环绕祂的宝座。柏拉图在此同一原则上亦表赞同，他主张独一的天主，而称其余的为天使或魔鬼。此外，赫尔墨斯·特里斯墨吉斯忒斯也论及独一的天主，并承认祂是不可理解的，超乎我们的估量。\par

7. 因此，这些灵潜伏在雕像和祝圣过的像中；它们以灵感灌注自己先知的心胸，使内脏的纤维跳动，指引鸟类的飞翔，掌握签运，使神谕生效，总是在真中掺杂假，因为它们既受欺骗，也欺骗人；它们扰乱人的生命，惊扰人的睡眠；它们的灵也潜入人的身体，暗中恐吓人的心智，扭曲人的四肢，损害人的健康，激起疾病以迫使人敬拜它们自己，这样，当它们饱享祭台的烟气与牲畜的堆堆时，便可解开它们所捆绑的，从而显得完成了医治。从它们而来的唯一补救是当它们自己的恶行停止之时；它们别无他求，只想将人从天主那里引开，使人转离对真宗教的认识，转向对它们自身的迷信；而由于它们自身正受惩罚，（它们希望）为自己寻找惩罚中的同伴，藉其误导使人成为其罪过的共犯。然而，当我们藉着真天主之名驱赶它们时，它们立刻屈服并承认，被迫离开附身者的身体。你们可以看见，因我们的声音和隐藏威能的作用，它们被鞭打、被火烧、被日益增长的刑罚所折磨，嚎叫、呻吟、恳求、承认从何而来、何时离去，甚至在那些敬拜它们的人耳中也是如此，要么立即跳出，要么渐渐消失，正如受苦者的信德相助，或医治者的恩宠所成就。因此，它们煽动民众憎恨我们的名字，使人在认识我们之前就开始恨我们，免得他们若认识我们，要么效仿我们，要么无法定我们的罪。\par

8. 因此，万有的唯一主是天主。因为那崇高不可能有任何同等者，只有它独自拥有一切权能。此外，让我们从大地借一个比喻来看天主的治理。何时王权的联盟曾以信实开始，或未曾以流血告终？忒拜人的兄弟情谊破裂了，不和在死后仍存于他们分离的灰烬中。一个王国不能容纳罗马的双胞胎，尽管同一子宫的庇护曾容纳过他们。庞培和凯撒是亲戚，然而在嫉妒的权力中，他们未能保持亲戚关系的纽带。对于人，你不应为此惊奇，因为整个自然界都同意这一点。蜜蜂有一个王，羊群有一个领袖，牛群有一个统治者。世界的统治者更是独一无二的；祂用自己的话命令万有，以祂的智慧安排万有，以祂的能力完成万有。\par

9. 祂不能被看见——祂太明亮，无法目睹；也不能被理解——祂太纯洁，无法辨识；也不能被估量——祂太伟大，无法感知；因此，当我们说祂不可理解时，我们才恰当地估量了祂。天主能有怎样的圣殿呢？全世界都是祂的圣殿。而人遍居各地时，难道我要将如此伟大威能关在一座小建筑里吗？祂必须在我们的心意中被奉献；在我们的胸怀中祂必须被祝圣。你也不应询问天主的名字。天主就是祂的名字。在那些需要名字的地方，是为了通过称谓的适当特征来区分众多。对于独一的天主，天主的整体名字属于祂；因此祂是独一的，祂以整体弥漫于各处。因为甚至平民百姓在许多事上自然地承认天主，当他们的心思和灵魂受到其创造者和根源的提醒时。我们常听到说：\Quote{啊，天主}和\Quote{天主看见}和\Quote{我把这事交托给天主}和\Quote{愿天主赐给你}和\Quote{如天主所愿}和\Quote{若天主应允}。这正是罪孽的极致：拒绝承认你不得不认识的那一位。\par

10. 但基督存在，以及救恩如何通过祂临到我们，其计划如下，其方式如下。首先，天主的恩宠赐给了犹太人。因此他们古时是义人；他们的祖先顺服于宗教义务。由此，他们的统治之崇高繁荣，其种族的伟大增长。但后来他们变得疏忽纪律，骄傲自满，倚靠他们的祖先，轻视神圣的诫命，失去了所赐的恩宠。他们的生活变得多么亵渎，他们因违反宗教而招致何等冒犯，连他们自己也作证，因为尽管他们以口沉默，却以结局承认。他们分散流离，漂泊不定；从自己的土地和气候中被驱逐，寄居在陌生人之中。\par

11. 此外，天主先前曾预言这事必会发生：随着时代的流逝，世界的终局临近，天主将从每个民族、人民和地方聚集更服从、更信德坚强的敬拜者到自己面前，他们将从神圣的恩赐中获得犹太人因藐视自己的宗教规条而领受又失去的慈悲。因此，为了这慈悲和恩宠，圣言和天主子被派遣为分施者和导师，古代所有先知都曾预告祂是人类的光明者和教师。祂是天主的德能，祂是理性，祂是天主的智慧和荣耀；祂进入一位童贞女；作为圣神，祂穿上血肉；天主与人结合。这就是我们的天主，这就是基督，祂作为两者的中保，穿上人性，为将人引领到圣父面前。人是什么，基督甘愿成为，为使人也能成为基督所是的。\par

12. 犹太人也知道基督必要来临，因为先知们的警告一直向他们宣告。但祂的来临被指明为双重的——一次是履行人的职分和榜样，另一次是宣认祂为天主——他们没有理解那先前的第一次来临，即隐藏在祂的苦难中，却只相信那将显于权能的唯一一次。而犹太民族不能理解这事，是他们罪的报应。他们因智慧的盲目和悟性的迟钝而受到惩罚，以致那些不配得生命的人，生命就在他们眼前，他们却看不见。\par

13. 因此，当基督耶稣按照先知先前所预言的，用祂的话语从人身上驱逐恶魔，用祂的声音的命令使瘫痪者恢复力量，洁净癞病人，开启盲人的眼睛，赐给瘸子行动的能力，使死人复活，迫使元素像仆人般服从祂，风服侍祂，海听从祂，阴府向祂屈服时；犹太人因祂肉体和身体的卑微只相信祂是人，因祂权能的权威而视祂为巫师。他们的首领和教师们——就是那些被祂在学识和智慧上征服的人——被怒火点燃，被愤慨刺激，最终抓住祂，将祂交给那时代表罗马人担任叙利亚总督的\Person{般雀·比拉多}，并以暴烈而固执的紧迫要求祂被钉十字架并死亡。\par

14. 他们要做这事，祂自己也早已预言；所有先前的先知们也同样预先作证：祂必须受苦，不是为感受死亡，而是为战胜死亡，并且当祂受苦之后，祂要重返天堂，以显示神圣威严的权能。因此事件的进程应验了预言。因为当被钉十字架时，祂抢先于刽子手的职分，自愿交出了自己的灵魂，并在第三天自由地从死者中复活。祂向门徒显现，如同祂原先的样子。祂让看见祂的人认出祂，与他们一同往来；因祂身体存在的实质而显明，祂停留了四十天，为使他们接受祂关于生命规条的教育，并学习他们应教导什么。然后祂在围绕祂的云彩中被提升到天上，为作为征服者将祂所爱、所披戴、所保护免于死亡的人——人——带到圣父面前；不久将从天降来，惩罚魔鬼并审判人类，以报复者的力量和审判者的权能；而门徒分散到世界各地，遵照他们的导师和天主的命令，传布祂的救恩规条，引导人从黑暗的迷途走向光明之路，并给盲目无知的人开启眼睛以认识真理。\par

15. 并且为使这证明不致缺少确据，为使对基督的宣认不是出于喜好，他们受酷刑、被钉十字架、受各种刑罚的考验。痛苦，即真理的考验，被施加，为使人所信赖的为赐人生命而赐予人的天主圣子基督，不仅藉声音的传报被宣告，也藉受苦的见证被宣告。因此，我们陪伴祂，我们跟随祂，我们以祂为我们道路的向导、光明的源头、救恩的创始者，向寻求并相信的人应许圣父及天堂。基督是什么，我们基督徒也将是，如果我们效法基督。\par

\SourceRecord{Translated by Robert Ernest Wallis.；Ante-Nicene Fathers；Vol. 5.；Edited by Alexander Roberts, James Donaldson, and A. Cleveland Coxe.；Buffalo, NY: Christian Literature Publishing Co.,；1886.；<http://www.newadvent.org/fathers/050706.htm>.}

% structure: p0001 source=h1 target=section reason=page-title

\section{论著七}

执事庞提乌以寥寥数语揭示本论著在其《西彼廉传》中的要旨。首先，他指出此类磨难已由基督预言，继而说明不应惧怕瘟疫，因它引人得永生；因此，不渴望更好世界的人，是因为信德欠缺。基督徒与外邦人同享今世苦难，这并不足为奇，因为他们在世界上必须受更多的苦；此后，应效法\Person{约伯}和\Person{多俾亚}的榜样，需要无怨言的忍耐。因为若非先有争斗，胜利便无从到来；无论疾病如何公平地临于义人与恶人，死亡对二者并不相同，因为义人蒙召受安慰，而不义的人被召受惩罚。\par

1. 虽然，我至爱的弟兄们，在你们当中许多人怀有坚定的心意、坚固的信德和虔诚的精神，并不因眼前频发的瘟疫而动摇，反而如坚固稳定的磐石，击碎世间狂暴的冲击和时代的汹涌波涛，自身却不受损毁，反在考验中得以磨炼而非失败；然而，由于我发现在众人中，有些或因心灵软弱，或因信德衰败，或因贪恋此世生活，或因女性之柔弱，或更甚者，因背离真理，站立不稳，未能发挥其心中神圣不可战胜的力量，故此不可掩饰或缄默不语，而应尽我微薄之力，以全部力量，并借由从主教训中采集的讲论，约束那放纵懈怠之本性，使那起初已是属天主和属基督的人，堪当成为天主和基督所悦纳的。\par

2. 我至爱的弟兄们，凡为天主作战的人，应当承认自己已置身天军营垒，寄望于天上之事，使我们不因世间的风暴与旋风而战栗，也不因主曾预言这些事将临而惶惑。主以他预见之言的劝勉，教训、教导、预备并坚强他教会的人民承受未来的一切考验，他预言并说：各处将有战争、饥荒、地震和瘟疫；以免突如其来的灾祸之惊骇动摇我们，他预先警告我们，在末期苦难将日益增多。看，这些预言已然应验；既然预言之事一一发生，所应许的一切也必随之而来；正如主亲自应许说：\BibleQuote{但当你们看见这一切发生，便知道天主的国近了}\BibleRef{路 21:31}。天主的国，亲爱的弟兄们，已经开始临近；生命的赏报、永恒救恩的喜乐、以及失乐园后新得的永恒福乐和拥有，现今随着世界的逝去而到来；天上之物已将地上之物取代，伟大之物取代微小之物，永恒之物取代衰朽之物。在此还有何忧虑和挂虑的余地？在这诸事之中，除了没有希望和信德的人，谁还会战栗悲伤？因那不愿走向基督的人才会惧怕死亡。那不愿走向基督的人，是因为他不相信自己将与基督一同为王。\par

3. 经上记载：义人因信德而生活。你若是义人，且因信德而生活，若你真实相信基督，为何既然将要与基督同在，且确信主的应许，却不欣然接纳你蒙召归于基督的确据，反而因脱离魔鬼而喜乐？诚然，那位义人\Person{西默盎}，他是真正的义人，以全备的信德遵守天主的诫命，当上天应许他未死之前必得看见基督。及至基督婴孩与母亲进入圣殿时，他在神中认出基督已经降生，就是那先前对他预言的那位；一见了他，便知自己不久将离世。因此，他欢喜自己临近的死亡，确信自己即将被召，遂将婴孩抱在怀中，赞美主，高呼说：\BibleQuote{主啊，现在可照你的话，放你的仆人平安去了！因为我亲眼看见了你的救援}\BibleRef{路 2:29}。这确凿地证明并见证，天主的仆人唯有在脱离此世狂风巨浪、抵达我们永居的港湾和永恒安全之时，才得享平安、自由和宁静的安息；唯有在此死亡完成之后，我们才来到永生。因为那就是我们的平安，我们忠信的安息，我们坚固、持久而永恒的保障。\par

4. 此外，在这世上，除了天天与魔鬼作战、不断与其箭矢和武器搏斗之外，还有什么呢？我们的争战是与贪婪、不贞、愤怒、野心；我们辛勤劳苦的角力是与肉性的恶习、世俗的诱惑。人的心灵被围困，四面八方遭受魔鬼的袭击，几乎处处难以招架，难以抵抗。若贪婪被击倒，情欲便兴起；若情欲被克服，野心取而代之；若野心被鄙视，愤怒便激化，骄傲膨胀，酗酒诱惑，嫉妒破坏和谐，猜忌割断友谊；你被迫咒骂，这是神圣律法所禁止的；你被迫发誓，这是不合法的。\par

5. 灵魂每天遭受如此多的迫害，心灵在如此多的危险中疲惫不堪，然而它却喜欢长久停留于魔鬼的武器之中，尽管我们更应渴望并希望藉着更快的死亡奔向基督；正如祂亲自教导我们说：\newline{}\BibleQuote{我实实在在告诉你们：你们要痛哭哀号，而世界却要欢乐；你们要忧愁，但你们的忧愁要变为喜乐。}\newline{}\BibleRef{若 16:20}谁不愿摆脱忧愁？谁不愿达到喜乐？但我们的忧愁何时变为喜乐，主亲自又说明，祂说：\newline{}\BibleQuote{我要再见你们，你们的心要欢乐；你们的喜乐没有人能夺去。}\newline{}\BibleRef{若 16:22}因此，看见基督就是喜乐，除非我们看见基督，否则我们不能有喜乐；那么，爱慕世界的磨难、惩罚和眼泪，而不奔向那永不失去的喜乐，是何等的心智盲目或愚昧呢？\par

6. 但是，亲爱的弟兄们，情况之所以如此，是因为缺乏信德，因为没有人相信天主所应许的事是真的，尽管祂是真实的，祂对信徒的话是永恒不变的。如果一个庄重而可敬的人应许你什么，你必定会相信那应许者，并且不会认为你会被他欺骗，因你知道他言行一致。如今天主正在对你说话；而你竟以不信的心不信地动摇？天主应许你，在你离开这个世界时，获得不死与永恒；你竟怀疑？这完全是不认识天主；这是以不信的罪冒犯基督、信徒的导师；这是身为教会中的人，却在信德之家没有信德。\par

7. 离开世界是多么大的益处，基督自身，我们救恩和善行的导师，向我们显示。当祂的门徒因祂说即将离去而忧愁时，祂对他们说：\newline{}\BibleQuote{如果你们爱我，你们必定喜乐，因为我往父那里去。}\newline{}\BibleRef{若 16:28}以此教导并表明，当我们所爱的亲者离开世界时，我们理应喜乐而非悲伤。记住这真理，蒙福的宗徒保禄在书信中规定说：\newline{}\BibleQuote{对我来说，生活原是基督，死亡乃是利益。}\newline{}\BibleRef{斐 1:21}认为最大的利益是：不再被此世界的网罗束缚，不再受制于肉身的罪恶和恶习，而从痛苦的磨难中被提去，从魔鬼的毒齿中得释放，因基督的召唤，前往永恒救恩的喜乐。\par

8. 但尽管如此，这疾病的力量攻击我们的人民，与攻击异教徒一样，使一些人感到困扰，好像基督徒信主的目的就是为了享受世界和今生，免于灾祸的触染；而不是作为那在此经受一切逆境，并为将来的喜乐而被保留的人。同样令一些人不安的是，这必死性对我们与他人是共通的；然而在这世界上，有什么不是我们与他人共通的呢？只要我们的肉身仍存，按照我们首次出生的规律，它与他人是共通的？只要我们在此世界，我们在肉身上与人类平等相连，但在精神上分离。因此，直到这可朽的穿上不朽，这必死的穿上不死，并且圣神引领我们到天主父那里，肉身的一切不利对于人类都是共通的。因此，当地土因歉收而荒芜时，饥荒不作区分；当仇敌入侵，城池被攻占时，被掳立即荒凉一切；当平静的云层不降雨时，干旱临到所有人；当嶙峋的岩石撞破船只时，船难毫无例外地临到船上所有人；眼疾、发热的侵袭、四肢的软弱，对我们与他人都是共通的，只要我们在世上带着我们这共通的肉身。\par

9. 再者，如果基督徒知道并坚守他所信奉的条件和法则，他就会意识到他必须在世上比别人受更多的苦，因为他必须更多地与魔鬼的攻击争战。圣经教导并预先警告说：\newline{}\Quote{我儿，当你来到天主的服务中时，要站在正义和敬畏中，预备你的灵魂受诱惑。}\newline{}又说：\newline{}\BibleQuote{在痛苦中忍耐，在你的谦卑中要有忍耐；因为金银在火中试炼，而蒙悦纳的人在屈辱的熔炉中。}\newline{}\BibleRef{德 2:5}\par

10. 因此，\Person{约伯}在丧失财产、子女死亡之后，且又遭受疮痍与蛆虫的严重折磨，并未被征服，而是受到了试炼；因为在他自身的挣扎与痛苦中，显示出一颗虔诚之心的忍耐，他说：\BibleQuote{我赤身出于母胎，也必赤身归于地下；主赐予，主收回；按主所视为适宜的，就这样成就了。愿主的名受赞美。}当他的妻子也因他疼痛剧烈而不耐烦，以抱怨与嫉妒的声音怂恿他说些反对天主的话时，他回答说：\BibleQuote{你说话如同愚顽的妇人一样。若我们从主手中得福，为何不能受苦呢？}在这一切临到他身上的事中，\Person{约伯}在嘴唇上并没有犯罪得罪主。（见：\BibleRef{约 1:8}）因此主天主为他作证说：\BibleQuote{你曾留意我的仆人\Person{约伯}吗？世上没有像他这样的人，纯良正直，敬畏天主，远离邪恶。}（见：\BibleRef{约 2:10}）\Person{多俾亚}在作了许多善工、并有许多光辉的仁慈行为之后，遭受失明，在逆境中敬畏并赞美天主，藉着身体的痛苦反而增加了赞美；甚至他的妻子也试图使他偏离正路，说：\BibleQuote{你的义德在哪里？看，你所受的苦。}（见：\BibleRef{多 2:14}）但他因敬畏天主而坚定稳固，并因宗教信德而武装以承受一切痛苦，在患难中没有屈服于软弱妻子的诱惑，反而因更大的忍耐更蒙天主悦纳；后来天使\Person{辣法耳}称赞他说：\BibleQuote{显扬并承认为天主的工作是可敬的。因为当你和你的儿媳\Person{撒辣}祈祷时，我在天主荣耀面前献上了你们的祈祷记录。当你以纯朴之心埋葬死者，且没有延迟起来离开你的晚餐，前去埋葬死者时，我被差来试探你。天主又派我来医治你和你的儿媳\Person{撒辣}。我是\Person{辣法耳}，是侍立在主荣耀面前出入的七位圣天使之一。}（见：\BibleRef{多 12:11}）\par

11. 义人向来拥有这种忍耐。宗徒们从主的法律中坚守这一训诫，不在逆境中抱怨，而是勇敢且忍耐地接受世间发生的一切；因为犹太民族在这件事上常常得罪天主，他们不断抱怨天主，就如主天主在\WorkTitle{户籍纪}中所作证的，说：\BibleQuote{让他们的抱怨远离我，他们就不会死亡。}（见：\BibleRef{户 17:10}）我们不可在逆境中抱怨，亲爱的弟兄们，而必须忍耐且勇敢地承受所发生的一切，因为经上记载：\Quote{天主所悦纳的祭献是破碎的精神；天主绝不轻视痛悔和谦卑的心。}此外，在\WorkTitle{申命纪}中，圣神也藉\Person{梅瑟}警告说：\BibleQuote{上主你的天主将磨炼你，以饥饿考验你，为要知道你心中是否遵守祂的诫命。}（见：\BibleRef{申 8:2}）又说：\BibleQuote{上主你的天主考验你，为要知道你是否全心全灵爱上主你的天主。}（见：\BibleRef{申 13:3}）\par

12. \Person{亚巴郎}如此中悦了天主，他为了中悦天主，甚至不惜失去他的儿子，或做出\OriginalTerm{弑亲}{parricide}的行为。你们，因着死亡的律法与命运而无法忍受失去儿子，如果被命令杀死自己的儿子，你们会怎么做呢？敬畏天主和信德应当使你们准备好面对一切，无论是私人财产的损失，还是因剧烈病痛对肢体的持续残酷折磨，或是从妻子、儿女、离世的亲人那里令人悲痛的生离死别；不要让这些事成为你们的绊脚石，而要视为战斗：不要让它们削弱或摧毁基督徒的信德，反而要在斗争中显示其力量，因为当今所有患难所加的痛苦，在确信未来祝福时都当被轻视。除非战斗先发生，否则不可能有胜利：当战斗开始时取得胜利，那时冠冕也赐给胜利者。因为舵手在风暴中被认出；士兵在战争中经受考验。没有危险时只是虚张声势。逆境中的挣扎是对真理的考验。根系深厚的树不会被风吹动，坚固木材建造的船虽被波浪拍打却不会碎裂；当打谷场扬出谷粒时，强壮饱满的谷粒藐视风力，而空穗则被吹来的阵风带走。\par

13. 此外，宗徒\Person{保禄}在遭遇海难、鞭打、肉身与身体多次严重折磨之后，他说他并未因此忧伤，反而因逆境受益，为要在他受重苦时更真实地经受试炼。他说：\BibleQuote{有一根刺加在我身上，是撒殚的使者来击打我，免得我过于高举自己。为此，我三次求主，叫它离开我；祂对我说：\InnerQuote{我的恩宠够你用，因为我的能力在软弱中得以成全。}}（见：\BibleRef{格后 12:7}）因此，当软弱、无能和任何毁灭抓住我们时，我们的能力得以成全；那时，我们的信德若在试炼中坚定不移，就会得到冠冕；正如经上所记载的：\BibleQuote{炉火试炼陶匠的器皿，患难的试探考验义人。}（见：\BibleRef{德 27:5}）简而言之，这就是我们与不认识天主之人的区别：他们在不幸中抱怨和埋怨，而逆境不会使我们偏离美德与信德的真理，反而藉痛苦坚强我们。\par

14. 这考验，即现今肠松弛不止，泻痢不停，耗尽体力；骨髓生火，化为咽喉的溃烂；肠腹因不断呕吐而震颤；双眼因注入之血而灼烧；在某些情况下，足部或四肢的某些部分因腐烂病气的传染而被切除；因躯体伤残和损失所致的虚弱，或步履蹒跚，或听觉受阻，或视觉昏暗——这一切都有益于作为信德的证明。以坚定不移的心灵之全部力量，与如此多毁灭和死亡的攻击搏斗，这是何等伟大的精神！在人类受蹂躏之中屹立不倒，不与那些在天主内无希望的人一同仆倒，反而喜乐，拥抱时机的益处；如此勇敢地彰显我们的信德，并借着所受的苦难，沿着基督所走过的窄路走向基督，我们就能按照祂自己的审判，得到祂生命和信德的赏报！确实，凡未由水和圣神重生，而被交付于革哈拿（火狱）之火的人，会怕死；凡未登记在基督的十字架与苦难中的人，会怕死；凡从今世之死转入第二次死亡的人，会怕死；凡在离开此世时，将受永恒火焰永无穷尽之刑罚折磨的人，会怕死；凡因拖延而获得好处，在此期间其呻吟和痛苦被推迟的人，会怕死。\par

15. 在这场瘟疫中，我们的许多人死去，也就是说，我们的许多人从这个世界获得了解放。这场瘟疫，对犹太人和外邦人以及基督的仇敌是灾祸，对天主的仆人却是得救的离世。义人和不义的人死而无别，这一事实不应使你以为善者与恶者的死亡相同。义人蒙召前往安息之所，不义的人被拖去受罚；信者更快地获得安全，不信者更快地获得惩罚。亲爱的弟兄们，我们对天主的恩惠是疏忽和忘恩的，不承认赐予我们的恩赐。看，贞女们平安离世，带着光荣安全无恙，不惧怕即将来临的假基督的威胁及其败坏和淫乱。男孩们避开了他们不稳定年龄的危险，幸福地获得了节制和纯洁的赏报。优雅的少妇不再惧怕酷刑；因为她通过快速的死亡逃脱了对迫害的恐惧以及刽子手的手和折磨。因对瘟疫和时日的恐惧，冷淡者被激励，松懈者被振作，懒惰者受到刺激，背弃者被迫回头，外教人被约束而相信，古老的忠信会众被召安息，新生的大军以更勇敢的精力聚集到战场，在即将来临的战斗中不惧死亡而战，因为战斗发生在瘟疫之时。\par

16. 再者，亲爱的弟兄们，这看似可怕致命的瘟疫和灾祸，是何等的事，何等的大事，何等贴切，何等必要！它探查出每个人的义德，审察人类的心思，看看健康的人是否照料病人；亲属是否深情地爱其骨肉；主人是否怜悯其垂病的仆人；医生是否不抛弃恳求的病人；凶暴者是否抑制其暴戾；贪婪者是否能因对死亡的恐惧而平息其永无餍足的炽热贪欲；高傲者是否低头；邪恶者是否软化其蛮横；当亲人死去时，富人是否甚至在那时仍有所施舍，并在无继承人时要死时仍给予。即使这场瘟疫没有带来别的，它也已为基督徒和天主的仆人带来这样的益处：我们开始欣然渴望殉道，因为我们学会了不怕死亡。这些是我们的训练，而非死亡：它们赋予心灵刚毅的光荣；借着轻视死亡，它们为荣冠预备。\par

17. 但或许有人会反对说：\Quote{那么，正是这一件事使我在当前的瘟疫中悲伤，即我本已预备好宣认，并全心以极大的勇气献身于忍受苦难，却被死亡抢先而失去了殉道。}首先，殉道不在于你的能力，而在于天主的俯允；你不能说你已失去了你尚不知道是否配得的东西。然后，此外，天主是洞察肺腑心肠者，是隐秘之事的探察者和知悉者，祂看见你，赞美你并认可你；祂看见你的德能已在你内预备好，必将给你的德能酬报。加音在向天主献祭时，难道已经杀死了他的兄弟吗？然而天主预知他心中所怀的杀弟之念，就预先定其罪。正如在那事例中，预见的天主预见了邪恶的念头和有害的意图，同样，在天主的仆人身上，当他们心中立志宣认和存有殉道之意时，这献身于善的意向就被天主审判者所加冕。心志欠缺殉道是一回事，殉道欠缺心志是另一回事。主召唤你时发现你是什么样，祂就照样审判你；因为祂亲自作证说：\BibleQuote{所有教会都将知道我是洞察肺腑心肠者。}\BibleRef{默 2:23}因为天主不要求我们的血，而是要求我们的信德。因为亚巴辣罕、依撒格、雅各伯都没有被杀；然而，他们因信德和义德的功绩而受尊崇，成了圣祖中的首位，凡被找到是忠信、正义、可称赞的人，都被聚集到他们的宴席中。\par

18. 我们应当记得，我们不应寻求自己的旨意，而应寻求天主的旨意，正如我们的主教导我们每日祈祷的那样。当我们祈求天主的旨意得以承行时，若天主从这个世界召唤并召叫我们，我们却不立即服从祂旨意的命令，这是多么荒谬和不合理！我们挣扎抵抗，如同倔强的仆人，被悲伤和痛苦拖到主面前，出于必然的束缚而非自由意志的服从离开此世；我们渴望从那位我们不愿归向的主那里获得天上的赏报。那么，如果尘世的束缚令我们喜悦，我们又为何祈祷并祈求天国来临呢？如果我们更大的渴望和更强烈的愿望是在此世服从魔鬼，而非与基督一同为王，又为何以频繁重复的祈祷恳求祂的国度之日加快到来呢？\par

19. 此外，为使天主上智安排的迹象更加明显，证明主预知未来，为祂子民的真救恩筹谋，当我们的一位同僚兼司铎弟兄因疾病疲惫不堪，因死亡临近而忧虑，祈求为自己暂缓时，在他祈祷时，当他已濒临死亡时，一位年轻人站在他身旁，可敬而威严，身材高大，面容发光，凡胎肉眼几乎无法以肉眼看清楚他，只有那即将离世的人才能看见这样一位存在。那年轻人以略带责备的语调和声音斥责他说：\Quote{你害怕受苦，你不愿离去；我该对你做什么？}这是责备和警告的话语，当人们为迫害忧虑、对召唤漠不关心时，祂不赞同他们当下的愿望，而是为未来考虑。我们临终的弟兄和同僚听到了他该告诉别人的话。因为他临终时听到，正是为了能够转述；他听到不是为了自己，而是为了我们。对于即将离世的人，他为自己能学到什么呢？是的，无疑他是为我们这些留下的人而学，以便当我们发现祈求延迟的司铎受到责备时，能认识到什么是对众人有益的。\par

20. 至于我，最卑微和最后的，天主屈尊就卑，多少次向我启示，何等频繁而明显地命令我，应当勤勉见证并公开宣告：那些因主的召唤而脱离此世的弟兄不应被哀悼，因为我们知道他们并未丧失，而是被先送了去；他们离开我们，如同旅行者、如同航海者那样先行；我们应怀念他们，但不应悲恸；当他们已在彼方穿上白衣时，我们不应在此地穿上黑衣；不应给外邦人可乘之机，让他们理所当然地责备我们：我们为那些据我们所说与天主同在的人哀悼，仿佛他们已灭绝丧失；我们用言语和话语表达信仰，却不用内心和胸中的见证来证实。我们是我们希望和信仰的背弃者：我们的言语显得虚假、伪装、不真实。用言辞彰显美德，却用行为毁灭真理，毫无益处。\par

21. 最后，宗徒保禄斥责、责备、指责那些为朋友离世而悲伤的人。他说：\BibleQuote{弟兄们，关于安眠的人，我们不愿意你们不知道，免得你们忧伤，像那些没有希望的人一样。因为我们相信耶稣死了，也复活了，同样，那些在耶稣内安眠的人，天主也必把他们与耶稣一同带来。}他说在朋友离世时忧伤的人是没有希望的。但我们生活在希望中，相信天主，相信基督为我们受苦并复活了，我们住在基督内，并且通过祂并在祂内复活，为何我们自己不愿离开此生，或为朋友离世而悲恸哀哭，仿佛他们丧失了呢？而基督自己，我们的主和天主，鼓励我们说：\BibleQuote{我就是复活，就是生命；信从我的，即使死了，仍要活着；凡活着而信从我的人，必永远不死。}\BibleRef{若 11:25} 如果我们相信基督，就让我们相信祂的话语和应许；既然我们不会永远死亡，就让我们怀着喜乐的确信奔向基督，与祂一同获胜，永远为王。\par

22. 在此期间我们死去，乃是藉着死亡过渡到不朽；永生除非我们离开此生，否则无法跟随。那不是终结，而是迁移，是渡过时间的旅程而进入永恒。谁不渴望更好的事物？谁不渴望改变并更新为基督的肖像，更快地抵达天上荣耀的尊位？正如宗徒保禄宣告说：\BibleQuote{我们的家乡是在天上，我们也等待着主耶稣基督从那里降来；祂将改变我们卑贱的身体，使之相似祂光荣的身体。}\BibleRef{斐 3:21} 主基督也应许我们将成为那样，为了使我们与祂同在，与祂生活在永恒的居所，并在天国中喜乐，祂为我们向父祈求说：\BibleQuote{父啊！我愿那些祢赐给我的人，也与我在哪里，好使他们看见祢赐给我的光荣，因为在创世之前，祢已经爱了我。}\BibleRef{若 17:24} 那将抵达基督的宝座、天国荣耀的人，不应哀悼哭泣，反而应按照主的应许，按照对真理的信德，为他这次的离去和迁移而喜乐。\par

23. 因此，我们同样发现\Person{哈诺客}也被接去了，他中悦了天主，正如在\BibleBook{}{创世纪}中圣经所见证的，说：\Quote{哈诺客中悦了天主，以后他就找不着了，因为天主把他接去了。}\BibleRef{创 5:24} 因此，在天主眼中蒙悦纳，就是配得从这世界的污染中被接去。此外，圣神也通过\Person{撒罗满}教导，中悦天主的人更早被带离，更快获得释放，免得他们在世上耽搁太久而被世界的污染玷污。他说：\Quote{他被接去，}他说，\Quote{免得邪恶改变他的悟性。因为他的灵魂中悦了天主，所以天主急忙将他从邪恶中接去。}\BibleRef{智 4:11} 同样，在\BibleBook{}{圣咏}中，以属灵的信德将自己献于天主的灵魂急忙向主说：\Quote{万军的天主，你的居所多么可爱！我的灵魂渴慕，切愿进入天主的庭院。}\par

24. 只有被世界取悦、被这谄媚而欺骗的生命以世俗快乐的诱惑所邀请的人，才愿意长久留在世上。再者，既然世界仇恨基督徒，你为什么爱那恨你的呢？为什么不更跟随那既救赎你又爱你的基督呢？\Person{若望}在他的书信中呼喊并劝告我们，不要顺从肉体的欲望，也不要爱世界。他说：\Quote{不要爱世界，也不要爱世界上的事。谁若爱世界，父的爱就不在他内。因为凡世界上的事，就是肉体的贪婪、眼目的贪婪、生活的骄矜，都不是出于父，而是出于世界的贪婪。世界和它的贪婪都要过去；但那履行天主旨意的，却永远存留，如同天主永远存留一样。}\BibleRef{若一 2:15} 亲爱的弟兄，我们应以健全的心智、坚定的信德、刚强的德性，准备好承行天主的全部旨意；放下对死亡的恐惧，思念那随之而来的永生。借此，我们要证明自己是相信的，即不为亲人的离世而悲伤，当那召唤的日子来到时，我们应毫不拖延、毫不抗拒地到召唤我们的主那里去。\par

25. 这一点，正如天主仆人常应做的，如今更应如此——如今世界正在崩溃，被邪恶灾祸的风暴所压迫；好使我们这些看见可怕之事已经开始、并知道更可怕之事即将来临的人，视尽快离开世界为最大的益处。如果你的住所墙壁因年久摇摇欲坠，屋顶在你上方颤抖，房屋因年久失修而濒临崩溃，你不是会赶紧离开吗？如果你在航行中，愤怒狂烈的风暴通过汹涌波涛预示即将到来的船难，你不是会迅速寻找港口吗？看，世界正在改变和消逝，如今不是因它的年代，而是因万物的终结来见证它的毁灭。你难道不感谢天主，不庆幸自己因早日离世而被接去，从迫在眉睫的船难和灾难中得救吗？\par

26. 我们应当思量，亲爱的弟兄们——我们应当现在且常常反省，我们已经弃绝了世界，在此暂时寄居为客旅。让我们迎接那指定我们各人归家的日子，那日子将我们从世上的罗网中夺去并释放，使我们回归乐园和天国。身处异乡的人，谁不急于返回本国？急于回到朋友身边的人，谁不渴望顺风，以便能更快拥抱他所爱的人？我们视乐园为我们的家乡——我们已开始视先祖为我们的父母：为何我们不急忙奔跑，好看见我们的家乡，问候我们的父母？那里有我们许多亲人在等候，一大群父母、兄弟、儿女在切望我们，他们已确知自己得救，仍为我们的得救而担忧。到达他们面前、进入他们的怀抱，对他们和我们共同是多大的喜乐！在天国里，没有死亡的恐惧，是何等的快乐；与永生共存，是何等崇高而永恒的幸福！那里有荣耀的宗徒团体——那里有欢喜的先知队伍——那里有因斗争和苦难得胜而加冕的无数殉道者——那里有得胜的贞女，她们以贞洁的力量制服了肉体和身体的贪欲——那里有慈悲的人得赏报，他们因喂养和帮助穷人而行了义德——他们遵守主的诫命，将世间的产业转移到了天上的宝库。亲爱的弟兄，让我们以热切的渴望奔向这些；让我们渴望尽快与他们同在，尽快来到基督面前。愿天主看到我们这份热切的渴望；愿主基督看到我们心灵和信德这意向，祂必将更大荣耀的赏报赐给那些对祂怀有更大渴望的人！\par

\SourceRecord{Translated by Robert Ernest Wallis.；Ante-Nicene Fathers；Vol. 5.；Edited by Alexander Roberts, James Donaldson, and A. Cleveland Coxe.；Buffalo, NY: Christian Literature Publishing Co.,；1886.；<http://www.newadvent.org/fathers/050707.htm>.}

% structure: p0001 source=h1 target=section reason=page-title

\section{论文第八篇}

\Emph{论行为与施舍。}\par

\Emph{论点：他极力劝勉以行为彰显信德，并强调向教会献仪及济贫的智慧，此乃基督徒家产的最佳投资。他以多处圣经证明这一点。}\par

1. 亲爱的弟兄们，天主圣父及基督那丰盛而广大的慈悲在祂的恩惠中，为我们得救所操劳并一直操劳的，实在是又多又大的：圣父派遣圣子来保全我们并赐予我们生命，为的是恢复我们；圣子甘愿被派遣并成为人子，为使我们能成为天主的子女；祂自谦自卑，为扶起先前跌倒的百姓；祂受伤，为医治我们的创伤；祂服侍，为解救受捆绑的人得自由；祂经历死亡，为向凡人显示不朽。这些是天主怜悯的诸多伟大恩惠。此外，还有一种眷顾和极大的仁慈：藉着一项救恩计划，为我们预备了更丰富的照顾，以保存那已被赎回的人！因为主在降临时，已治愈了亚当所受的创伤，并医治了古蛇的旧毒，祂便给健康的人颁赐法律，命他不要再犯罪，免得那犯罪的人遭遇更糟的事。我们曾被限制并囿于无罪的诫命之中。除非天主的仁慈再次施助，藉着指出公义和慈善的行为来打开一条确保得救的道路，好使我们藉着施舍能洗去后来沾染的一切污秽，否则人性的软弱与缺陷将毫无出路。\par

2. 圣神在圣经中发言说：\Quote{藉着施舍和信德，罪恶得以净化。}这当然不是指先前已经犯下的罪，因为那些罪已因基督的血和圣化而被净化了。此外，祂又说：\BibleQuote{如水灭火，施舍亦能消除罪恶。} \BibleRef{德 3:30}这里也表明并证明，正如在拯救之水的水池中，革赫那之火得以熄灭，同样藉着施舍和正义的行为，罪过的火焰也被制服了。因为在圣洗中，罪过只一次性获得赦免，但持续不断的努力，效法圣洗的样式，再次赐予天主的怜悯。主在福音中也这样教导。因为当门徒们被指出不先洗手就吃饭时，祂回答说：\BibleQuote{那造里面的，也造了外面的。但你们施舍吧，看，一切对你们都洁净了。} \BibleRef{路 11:41}藉此教训并表明：不是手要洗，而是心要洗，内在的污秽比外在的更要清除；但谁若洁净了内在，也洁净了外在；如果心灵洁净了，人在皮肤和身体上也开始洁净了。此外，祂劝告并指出我们何以能洁净和净化，祂补充说必须施舍。那慈悲的主教训并警告我们，必须显示仁慈；因为祂寻求拯救那些祂以高价赎回的人，祂教训那些在圣洗恩宠之后变得污秽的人，可以再次得到洁净。\par

3. 亲爱的弟兄们，让我们承认天主仁慈的救恩恩赐吧；我们既然不能没有良心上的某种创伤，就当用属灵的药物来医治我们的创伤，以洁净和净化我们的罪过。谁也不可因纯洁无瑕之心而自我陶醉，竟依仗自己的无辜，以为不需要给伤口敷药；因为经上记载：\BibleQuote{谁能说，我保持了心的洁净，我是纯洁无罪的？} \BibleRef{箴 20:9}再者，若望在他的书信中规定并说：\Quote{如果我们说我们没有罪，便是欺骗自己，真理也不在我们内。}但如果没有人能无罪，而谁若说自己没有过错，便是骄傲或愚蠢，那么天主的仁慈是多么必要、多么慈祥！它知道在那已痊愈的人身上，甚至在痊愈之后，仍可找到一些创伤，便赐予了有益的医药，以便重新医治和治愈他们的创伤！\par

4. 最后，亲爱的弟兄们，圣经中天主的劝诫，无论旧约与新约，从未停止，从未缄默，总是处处敦促天主的子民从事仁慈的行为；在圣神的激励和劝诫之下，每一位怀抱天国希望而受教的人都奉命施舍。天主命令并吩咐依撒意亚说：\BibleQuote{你要大声呼喊，不要停止，提高你的声音如同号角，向我的百姓宣告他们的过犯，向雅各伯家宣告他们的罪恶。} \BibleRef{依 58:1}当祂命令将他们的罪过归咎于他们，并满含义怒陈明他们的不义之后，祂说，即使他们恳求、祈祷、禁食，也不能赎罪；即使他们披上苦衣和灰尘，也不能平息天主的愤怒；然而在最后部分，祂表明单凭施舍就能平息天主，于是补充说：\BibleQuote{把你的饼分给饥饿的人，将无家可归的穷人领到你家里。若见赤身露体的人，给他衣服穿，不要轻视你同族的家人。那时，你的光明将在适当的时候射出，你的衣裳将迅速升起；公义必在你前面行，天主的荣耀将环绕你。那时你呼喊，天主必应允；你还在说话，祂必说：我在这里。} \BibleRef{依 58:1}\par

5. 平息天主的补救措施由天主亲口的话语所赐予；神圣的教训教导了罪人所当行的事：藉着正义的行为，天主得到满足；藉着慈悲的功德，罪恶得以洗净。在撒罗满的书中我们读到：\Quote{你要把施舍存放在穷人的心中，这施舍会为你从一切凶恶中转求。}\BibleRef{德 22:12} 又说：\Quote{谁掩耳不听弱者的哀求，他本人呼求天主时，也必无人应允。}\BibleRef{箴 21:13} 因为那自己不曾慈悲的人，不能堪当获得主的慈悲；那在穷人祈祷时不曾仁慈相待的人，也必无法从神圣的怜悯中求得任何恩惠。圣神也在圣咏集中宣讲并证明此事说：\Quote{怜悯穷困贫乏者的人是有福的；在邪恶的日子，主必拯救他。}达尼尔记念这些诫命，当拿步高王因恶梦而惊恐不安时，达尼尔为消除祸患，赐予他一项获得天主助佑的补救办法，说：\Quote{所以，大王，请接纳我的忠告：藉着施舍赎你的罪，藉着怜悯穷人赎你的不义，如此天主必容忍你的罪恶。}\BibleRef{达 4:27} 但因王未听从，便遭受了他所预见的灾祸与不幸；这些本可避免，若他曾藉施舍赎罪。辣法耳天使也作同样见证，劝人当慷慨、宽大地施舍，说：\Quote{祈祷与禁食和施舍都是好的；因为施舍能救你脱离死亡，并涤除罪恶。}\BibleRef{多 12:8} 他显示：我们的祈祷和禁食，若不辅以施舍，便无多大功效；仅凭恳求，力量微弱，不足以求得所求，除非加上善行与善工。天使揭示、表明并证实：我们的祈求因施舍而变得有效；生命因施舍而脱离危险；灵魂因施舍而免于死亡。\par

6. 亲爱的弟兄们，我们提出这些事，并非不愿以真理的见证来证明辣法耳天使所说的话。在\WorkTitle{宗徒大事录}中，事实的信德得以确立；藉着一件已完成之事的明证，我们发现灵魂不仅因施舍得救免于第二次死亡，也免于第一次死亡。塔彼达素来广行善事、广施哀矜，患病而死；伯多禄被召到她尸首前，他以宗徒般的急切赶来，身旁站着哭泣恳求的寡妇，她们展示先前所得的斗篷、外衣和所有衣物，不是用言语，而是用她自己的善行为死者祈求。伯多禄感到这样的祈求必可获得；既然主自身因寡妇们的衣物而得着衣饰，祂必不缺少对恳求者的助佑。于是，他跪下祈祷——他是寡妇和穷人的合宜代言人——将所受托的祈求呈献给主，然后转向那已洗过安放在床上的遗体，说：\Quote{塔彼达，因耶稣基督之名，起来！}\BibleRef{宗 9:40} 那曾在福音中说过，凡因祂的名所求的一切都要赐予的主，并未缺少对伯多禄的帮助。于是死亡被延缓，灵魂得以恢复；在众人的惊异与赞叹中，复活的躯体再次苏醒并进入此世的光明。慈悲的功德如此有效，正义的行为如此有力！她曾为受苦的寡妇提供生活所需的帮助，如今因寡妇们的祈祷，她堪当被召回生命。\par

7. 因此，在福音中，主——我们生命的导师与永生救恩的君王——苏醒信众的集会，并在他们苏醒后永远眷顾；在祂神圣的命令与天上的训诲中，最频密地命令并规定我们当致力于施舍，不要倚赖地上的财富，反而要积攒天上的宝藏。祂说：\Quote{变卖你们的财物，施舍给穷人。}\BibleRef{路 12:33} 又说：\Quote{不要在地上积蓄财宝，那里有虫蛀、锈蚀，也有贼挖洞偷窃；但要在天上积蓄财宝，那里没有虫蛀、锈蚀，也没有贼挖洞偷窃。因为你的财宝在哪里，你的心也必在哪里。}\BibleRef{玛 6:19} 当祂愿藉遵守法律来展示一个完美而完全的人时，祂说：\Quote{你若愿意成为成全的，去，变卖你所有的，施舍给穷人，你必有财宝在天上；然后来跟随我。}\BibleRef{玛 19:21} 此外，祂在另一处说，一位属天恩宠的商人、一位赚取永生救恩的人，当以基督宝血的代价——即从他全部的遗产中——舍去他所有财富，购买那宝贵的珍珠，即永生。祂说：\Quote{天国好像一个商人，寻找好珍珠；他找到一颗珍贵的珍珠，就去变卖他所有的一切，买了它。}\BibleRef{玛 13:45}\par

8. 最后，祂称那些勤勉地帮助和供养穷人者为亚巴郎的子孙。因为当匝凯说：\BibleQuote{我把一半的财物施舍给穷人；我若欺骗过谁，就偿还四倍。}耶稣回答说：\BibleQuote{今天救恩来到了这一家，因为他也是亚巴郎的子孙。}\BibleRef{路 19:8} 因为如果亚巴郎信了天主，这就算为他的正义，那么按照天主的诫命施舍的人，确是信天主的；保持信德真理的人，维持对天主的敬畏；此外，维持对天主的敬畏的人在怜悯穷人时顾念天主。因为他如此劳苦，是因为他相信——因为他知道天主圣言所预言的确实无误，圣经不会说谎——不结果子的树，即不结果子的人，被砍下投入火中，而怜悯的人被召进天国。祂在另一处也称勤劳结果子的人为忠信的；却把不结果子的、贫瘠的人说成没有信德，说：\BibleQuote{你们在不义的钱财上不忠信，谁还把真实的钱财托付给你们呢？如果你们在别人的财物上不忠信，谁还把属于你们的交给你们呢？}\BibleRef{路 16:11}\par

9. 你若恐惧害怕，怕你开始这样慷慨施舍后，你的家产因你的慷慨而耗尽，你也许会陷入贫困；那么在这方面要鼓起勇气，不要忧虑：那供应基督事工、并举行天上工作的地方是不会耗尽的。我这样说并非凭我自己的权威，而是凭圣经的信德和天主恩许的权威。圣神通过撒罗满说话：\BibleQuote{施舍穷人的，决不会匮乏；转目不看的，必受许多贫困。}显示怜悯人和行善的人不会匮乏，反而吝啬和贫瘠的人日后必致匮乏。此外，蒙主灵感恩宠充满的圣保禄宗徒说：\BibleQuote{那供给播种者种子，又供给食物作你们食粮的，必要使你们的种子增多，并增加你们正义果实的生长，使你们在一切事上富足起来。}\BibleRef{格后 9:10} 又说：\BibleQuote{这供应的事，不但补足了圣徒的缺乏，而且因着许多感谢天主而更加丰盛。}\BibleRef{格后 9:12} 因为当我们因施舍和劳苦而向天主感谢时，藉着穷人的祈祷，行善者的财富因天主的赏报而增加。主在福音中已经看透了这种人的心，并以预言之声斥责不忠不信的人，作证说：\BibleQuote{不要忧虑说：我们吃什么？喝什么？穿什么？这都是外邦人所求的。你们的天父知道你们需要这一切。你们先该寻求天主的国和祂的义德，这一切自会加给你们。}\BibleRef{玛 6:31} 祂说这一切都要加给那些寻求天主的国和祂的义德的人。因为主说，审判的日子来临时，那些在祂的教会中劳苦的人将被接纳进入天国。\par

10. 你害怕你的产业可能会因你开始慷慨行事而减少；可悲的人啊，你不知道当你害怕家产耗尽时，生命本身和救恩正在耗尽；当你担心你的财富减少时，你没有看到你自己正在减少，因为你爱钱财胜于爱自己的灵魂；当你害怕因你失去产业时，你自己却因产业而灭亡。因此宗徒大声疾呼说：\BibleQuote{我们什么都没有带到世界上，也什么都不能带走。所以我们有吃有穿，就当知足。至于那些想发财的人，却陷入诱惑和罗网，以及许多无知而有害的欲望，这欲望把人淹没在败坏和灭亡中。因为贪爱钱财是一切罪恶的根源，有些人因贪求而离弃了信德，使自己被许多痛苦刺透。}\BibleRef{弟前 6:7}\par

11. 你害怕你的产业也许因你开始慷慨行事而短缺？然而，对义人来说，什么时候曾发生过资源缺乏的事呢？因为经上记载：\BibleQuote{上主不使义人的灵魂忍饥挨饿。}\BibleRef{箴 10:3} 厄里亚在旷野中靠乌鸦叼来的食物生活；达尼尔被君王下令封闭在狮子坑中，却有从天上预备的膳食；而你却担心食物会缺乏，尽管你勤劳苦干，对主有功德，而主自己在福音中为责备那些心意不定、信德微小的人作证说：\BibleQuote{你们看天空的飞鸟，不播种，不收割，也不在仓里积存；你们的天父尚且养活它们；你们不比它们更贵重吗？}\BibleRef{玛 6:26} 天主养活飞鸟，麻雀每日得食；对于没有神圣感觉的受造之物，饮食并不缺乏。你认为对基督徒——对主的仆人——对专心行善的人——对主所爱的人——会有什么缺乏吗？\par

12. 除非你以为那喂养基督的人自己不被基督喂养，或者那接受天上神性事物的人会缺少地上的事物，不然这不信的念头从何而来，这邪恶而亵渎的想法从何而来？在信德之家中，一颗不信的心做什么？那不完全信赖基督的人，为何被称为基督徒？法利塞人的名字更适合你。因为在福音中，主谈论施舍时，忠信而有益地警告我们要通过智慧的善行，用不义的钱财结交朋友，使他们后来能接我们进入永久的帐幕。圣经接着说道：\BibleQuote{但法利塞人听见这一切，他们原是贪财的，就嗤笑他。}\BibleRef{路 16:14} 现在我们在教会中看到类似的人，他们的耳朵闭塞，心蒙黑暗，不接受属灵和拯救的警告。我们不必奇怪他们轻视仆人的讲论，因为我们看见主自己也受到这样的人藐视。\par

13. 那么，你为什么在那些虚妄而愚蠢的念头中自夸，仿佛因对未来的恐惧和忧虑而被阻止行善？你为什么摆出一些虚荣的借口阴影和预兆？不，要承认事实；既然你不能欺骗那些知道真相的人，就说出你内心的秘密和隐藏之事。荒芜的黑暗已经包围了你的心灵；真理之光已离开那里，而贪财的深厚黑暗已蒙蔽了你肉性的心。你是钱财的俘虏和奴隶；你被贪婪的锁链和绳索捆绑；基督曾释放了你，你又重新被捆绑。你保存你的钱财，但保存它并不能保存你。你积累一份产业，它用它的重量压迫你；你不记得天主对那富人回答的话，他愚蠢地夸耀他丰盛收成的富足：\BibleQuote{愚昧的人！}祂说，\BibleQuote{今夜就要索回你的灵魂；你所预备的一切，将归谁呢？}\BibleRef{路 12:20} 你为什么孤独地看守你的财富？为什么为了惩罚自己而堆积产业的负担，使你在这世上越富足，在天主面前越贫穷？把你的收入与主你的天主分享；与基督分享你的收益；使基督成为你地上财产的合伙人，这样祂也会使你在祂天上国度中成为祂的共同继承人。\par

14. 你错了，被蒙骗了，无论你是谁，自认为在这世上富有。听你在 \BibleBook{}{默示录} 中的主的声音，祂以公义的责备斥责像你这样的人：\BibleQuote{你说：我是富足的，已经发了财，什么都不需要；却不知道你是那困苦、可怜、贫穷、瞎眼、赤身的。我劝你向我买火炼的金子，叫你富足；又买白衣穿上，叫你赤身的羞耻不露出来；又用眼药擦你的眼睛，使你能看见。}\BibleRef{默 3:17} 因此，你，那富有而富裕的人，要从基督那里为自己买火炼的金子，这样你就可以成为纯金，像被火烧掉污秽一样，如果你通过施舍和正义的行为被炼净。给自己买白衣，使那按亚当本是赤身的你，先前可怕而不体面，可以穿上基督的白衣。而你，在基督教会中富有而富裕的贵妇，用眼药擦你的眼睛，不是用魔鬼的眼药，而是用基督的眼药，这样你就能达到看见天主，因善行和品格而蒙天主悦纳。\par

15. 但像你这样，不能在教会中劳作。因为你的眼睛被黑暗的阴云笼罩，在黑夜中看不见穷人和贫困者。你富有而富裕，你以为你庆祝主的晚餐，却根本不考虑奉献。你来参加主的晚餐却没有祭献，却分享了穷人献上的祭献。想一想福音中那记念天上诫命的寡妇，即使在贫困的困难和窘迫中仍行善，将两个小钱，即她所有的一切，投入银库。主看见并观察，她奉献不在于数量，而在于心意，不是给了多少，而是从\Emph{从}多少中给的。祂回答说：\BibleQuote{我实在告诉你们，这穷寡妇投入的，比众人投入的更多。因为众人都是拿自己有余的投入天主奉献；但这寡妇却从自己的不足中，把她所有的一切生活费都投进去了。}\BibleRef{路 21:3} 极其有福和光荣的妇人，在审判之日以前就配得从审判者口中受称赞！让富人在他们的荒芜和不信中羞愧吧。这寡妇，这贫穷的寡妇，在善行上却是富足的。尽管一切施舍都给予寡妇和孤儿，她却给予，而本该接受，这样我们就知道，那荒芜的富人将受到何等惩罚，因为连穷人都应努力行善。为了让我们明白他们的劳苦是献给天主的，并且凡行善的人都蒙主悦纳，基督称此为\Quote{天主的奉献}，并表明那寡妇将两个小钱投入了天主的奉献，从而更充分显明，怜悯穷人就是借给天主。\par

16. 但是，最亲爱的弟兄们，愿那考虑不要阻止和劝阻基督徒行公义与良善的事，免得有人以为可以因自己子女的利益而得到谅解；因为在属灵的支出上，我们应当想到基督，祂已宣告祂接纳他们；并且不可将我们的同仆、而应将主置于我们的子女之上，因为祂亲自教导并警告我们说：\BibleQuote{爱父亲或母亲超过我的人，不配属于我；爱儿子或女儿超过我的人，不配属于我。}\BibleRef{玛 10:37} 同样在 \BibleBook{}{申命记} 中，为坚固信德和爱天主，也写着类似的话：\BibleQuote{他论到自己的父母说：我没有见过你们；他也不承认自己的兄弟，也不认识自己的儿女；但这些遵守了祢的话，持守了祢的盟约。}\BibleRef{申 33:9} 因为如果我们全心爱天主，就不应将父母或子女置于天主之上。\Person{若望}也在他的书信中写道：那不肯为穷人劳苦的人，天主的爱就不在他们内。\BibleQuote{谁若有这世上的财物，看见弟兄有需要，却关闭了自己的恻隐之心，天主的爱怎能存在他内呢？}\BibleRef{若一 3:17} 因为若我们借施舍给穷人而借给天主——且当施舍给最小者时，就是施舍给基督——就没有任何理由偏爱属地的事胜过属天的事，也不应将人的事置于神圣的事之前。\par

17. 同样，在 \WorkTitle{列王纪} 中，那位寡妇在干旱和饥荒中，用完了一切，用剩下的少许面粉和油在灰上做了一张饼，吃完后就要与她的孩子们一同死去时，\Person{厄里亚}来了，要求先给他一些吃的，然后她与她的孩子再吃剩下的。她毫不犹豫地服从了；母亲在饥饿和贫穷中也没有将孩子置于\Person{厄里亚}之上。是的，在天主眼中做了令天主喜悦的事：所要求的被迅速而慷慨地献上。这不是从丰裕中取出一份，而是从微少中拿出全部，并且另一个比她自己饥饿的孩子先被喂养；在贫困和匮乏中，食物被置于仁慈之后；这样，在救恩的工作中，肉身的生命被轻视，而属神的灵魂得以保存。因此，\Person{厄里亚}作为基督的预像，表明祂按自己的仁慈回报各人，回答说：\BibleQuote{上主这样说：罐里的面不会用完，瓶里的油不会减少，直到上主使雨降在地上的日子。}\BibleRef{列上 17:14} 按照她对天主应许的信德，她所给予的东西为那寡妇倍增并堆积起来；她的义行和仁慈的功劳不断增益和增长，面粉和油的罐子都被充满。母亲并没有从她孩子那里拿走过给\Person{厄里亚}的东西，反而因她的仁爱与虔诚而使孩子们受益。那时她还不知道基督；她尚未听到祂的诫命；她也没有因被祂的十字架和苦难救赎而用血肉和饮料回报祂。因此，由此可见，在教会中，那把自己和子女置于基督之上、保存财富而不同等分给贫穷之人的罪有多大。\par

18. 此外，你（说）：家里有许多孩子；孩子的众多使你无法自由地行善工。然而，正因如此，你更应当劳苦，因为你是许多托付的父亲。有更多的人，你必须为他们祈求天主。许多人的罪需要被赎，许多人的良心需要被洁净，许多人的灵魂需要被解放。正如在这现世生活中，在抚养和教育孩子的事上，孩子越多，花费也越大；同样，在属神和属天的生活中，你的孩子越多，你劳苦的支出也应当越大。\Person{约伯}也如此为他的子女献上许多祭物；他家中的托付有多少，献给天主的牺牲也有多少。既然在天主面前每天不可能没有罪，就没有一天缺少可清除罪的牺牲。圣经证明此事说：\Quote{\Person{约伯}，一个真诚而正义的人，有七个儿子和三个女儿，他洁净他们，按他们的数目为他们献祭给天主，并为他们的罪献上一只牛犊。}如果你真爱你的子女，如果你向他们显明全然、父爱的甜蜜，你就应当更加慷慨，好藉你的义行将你的子女托付给天主。\par

19. 你不可认为那既善变又软弱的人是你们儿女的父亲，而应获得那永恒不变、属神儿女的父亲。将你为继承人所积蓄的财富托付给祂。让祂作你儿女的监护人；让祂作他们的受托人；让祂作他们的保护者，以祂神圣的尊威，免受一切世俗的伤害。国家不会夺走托付给天主的财产，国库不会侵扰它，任何法庭的诬告都不能推翻它。那在天主监护下的遗产是安稳的。这是为亲爱的抵押品预备未来；这是以慈父之心照顾未来的继承人，依照圣经的信德，经上说：\Quote{我年少，现已年老，从未见过义人被遗弃，也未见他的后裔求乞。他终日施恩，借贷，他的后裔必蒙祝福。}又说：\Quote{行为无瑕、正直的人，必将留下有福的后裔。}\BibleRef{箴20:7}因此，除非你忠实地为你的儿女筹划，除非你期望他们在宗教和真虔敬中得以保存，你就是一个不公义、背信的父亲。你更关心他们地上的产业而非天上的产业，宁愿将儿女交托给魔鬼而非基督，这是双重的罪，犯下双倍的罪行：既没有为你的儿女提供他们天父的助佑，又教导你的儿女爱他们的财产胜过基督。\par

20. 倒要像托彼特那样做你儿女的父亲。将有益而救命的训诫交给你的抵押品，正如他给他儿子的那样；命令你的儿女，他也命令他的儿子说：\Quote{如今，我儿，我吩咐你：要诚心事奉天主，在祂面前做祂所喜悦的事；并命令你的儿子们，要他们行义施舍，记念天主，时时赞美祂的圣名。}\BibleRef{多14:10}又说：\Quote{我亲爱的儿子，在你一生所有的日子里，要把天主放在心中，不要愿意违犯祂的诫命。要一生行义，不要愿意走不义的路；因为如果你行事诚实，你的工作必受尊重。要由你的财物中施舍，不要转脸不顾任何穷人。这样，天主的圣容也不会转离你。我儿，你若有，就这样做。你若财物丰裕，就多施舍；你若只有少许，也要分出那少许。施舍时不要害怕；因为你为自己积存了美好的酬报，以备急需之日；因为施舍能救人脱离死亡，不让人进入阴府。施舍，对一切施舍的人，在至高天主面前，是一件美好的礼物。}\BibleRef{多4:5}\par

21. 亲爱的弟兄们，这是何等样的礼物，其陈列在天主眼前受到赞颂？若在外邦人的礼物中，有总督或皇帝出席，似乎是一件伟大荣耀的事，送礼者中筹备和展示也更大，为了取悦上层阶级；那么，有天主和基督作为礼物的观看者，这荣耀岂不更加辉煌和伟大？在那种情况下，展示的筹备和费用岂不更加奢华和慷慨？当天上的万军聚集观看，所有天使都来齐时：在那里，为送礼者所求的，不是四驾马车或执政官职位，而是永生；所追求的也不是民众空洞易逝的青睐，而是天国永久的赏报！\par

22. 为了使懒惰和不结果实的人，以及那些因贪财而在自己救恩的果实上无所作为的人更加羞愧，并使羞辱和耻辱的绯红更加打击他们污秽的良心，让每个人将魔鬼及其仆人，即丧亡与死亡之子，放在眼前，他们冲入人群，以对比的考验激怒基督的子民——基督亲自在场审判——用这些话：\Quote{我，为了你们在我这里所看到的人，既没有遭受耳光，也没有承受鞭打，更没有忍受十字架，没有流我的血，也没有以我的苦难和血的代价救赎我的家族；但我也不应许他们天国，也不召唤他们返回乐园，重新赐予他们不朽。然而他们为我预备了何等珍贵、何等丰厚的礼物！以何等过度而费力的劳苦得来！而且，以最奢华的策划，或抵押或变卖他们的财产来获得礼物！除非有相应的展示随之而来，他们就被嘲笑和嘘声赶走，有时几乎被群众的狂怒用石头打死！}基督啊，请展示您的这些送礼者——那些富人，那些财富充裕的人——在您所主持和垂顾的教会中，他们是否也陈列了那种礼物——抵押或散尽他们的财富，是的，通过将他们的财产转变为天上的宝藏而得到改善！在我那些败坏而属世的表演中，没有人得喂养，没有人得穿衣，没有人因任何饮食的安慰而得以维持。一切，在表演者的疯狂和观众的迷误之间，都在奢靡而愚蠢的欺骗性享乐虚荣中败亡。在那里，在您的穷人身上，您得到穿衣和喂养；您应许永生给那些为您劳苦的人；而您的子民几乎不能与我那些灭亡的子民相等，尽管他们为您的神圣报酬和天上的赏赐所尊荣。\par

23. 我们如何回答这些事，最亲爱的弟兄们？以什么理由来为富人的心灵辩护呢？他们被一种世俗的荒芜和一种黑暗的夜所压倒。我们用什么借口来免他们的罪呢？因为我们比魔鬼的仆人更不如，甚至不能适度地以基督的苦难和血的代价来回报祂。祂给了我们诫命；教导了祂的仆人所当行的；向行爱德的人应许赏报，向不结果实的人警告惩罚。祂宣告了祂的判决。祂预先宣布了祂将如何审判。对懒惰者有什么借口？对不结果实者有什么辩护？但当仆人不执行命令时，主将执行祂所威胁的，因为祂说：\BibleQuote{当人子在自己的光荣中，与众天使一同降来时，那时，祂要坐在光荣的宝座上；万民都要聚集在祂面前；祂要把他们彼此分开，如同牧人分开绵羊和山羊一样：把绵羊放在右边，山羊放在左边。然后君王要对那些在祂右边的人说：来吧，我父所祝福的，承受从创世以来为你们预备的国度吧！因为我饿了，你们给了我吃的；我渴了，你们给了我喝的；我作客，你们收留了我；我赤身露体，你们给了我穿的；我患病，你们看顾了我；我在监里，你们来探望了我。那时义人要回答祂说：主啊，我们什么时候见你饿了给你吃？渴了给你喝？什么时候见你作客收留了你？赤身露体给了你穿？或者什么时候见你患病或在监里，来探望了你？君王要回答他们说：我实在告诉你们：凡你们对我这些最小兄弟中的一个所做的，就是对我做的。然后祂也要对那些在左边的说：离开我，可咒骂的，进入那为魔鬼及其使者预备的永火里吧！因为我饿了，你们没有给我吃的；我渴了，你们没有给我喝的；我作客，你们没有收留我；我赤身露体，你们没有给我穿的；我患病或在监里，你们没有来看顾我。那时他们也要回答说：主啊，我们什么时候见你饿了，或渴了，或作客，或赤身露体，或患病，或在监里，而没有服事你呢？祂要回答他们说：我实在告诉你们：凡你们没有对这些最小中的一个做的，就是没有对我做。这些人要进入永罚，而义人却要进入永生。}\BibleRef{玛 25:31-46} 基督还能向我们宣布什么？祂还能如何更激发我们行义和怜悯的作为，以至于说凡给贫困者的就是给祂自己，并说如果贫困者没有得供应，祂就感到受了亏待？这样，在教会中不被兄弟情谊感动的人，仍可能因默观基督而感动；不想到弟兄在痛苦和贫穷中的人，仍可能想到他的主，而这主恰恰住在他所轻视的那个人身上。\par

24. 因此，最亲爱的弟兄们，你们既敬畏天主，已轻视并践踏了世界，将心灵提升到天上和神圣的事物，让我们以全备的信德、虔诚的心灵和不断的劳苦，服从主，以获得主的喜悦。让我们给基督世上的衣裳，好领受天上的衣袍；让我们给祂这世界的饮食，好能与亚巴郎、依撒格和雅各伯一同参与天国的筵席。为了不多收，让我们多播种。趁着还有时间，让我们顾念我们的安全与永恒的救恩，遵奉宗徒保禄的劝告，他说：\BibleQuote{因此，我们趁着还有时间，要向众人行善，尤其向有信德的一家的人。行善不要厌倦，因为到了适当的时节，我们必获得丰收。}\BibleRef{迦 6:9-10}\par

25. 让我们思考，亲爱的弟兄们，宗徒时代信徒团体做了什么，那时起初心灵以更大的德能繁盛，信徒的信仰还以崭新的热情燃烧。那时他们卖掉房屋和田产，慷慨地、自愿地将所得交给宗徒，分施给穷人；他们将世上的产业出售和转让，将他们的土地转移到那里，以便获得永恒所有的果实，并在那里预备居所，开始永久的住所。那时，劳动的丰盛，正如爱德的一致，我们在\BibleBook{宗徒}{大事录}中读到：\BibleQuote{众信徒都一心一意；在他们中，没有人说某物是属于自己的，他们一切都归公用。}\BibleRef{宗 4:32} 这真是因属灵重生而成为天主的子女；这是以天国的法律效法天主圣父的公义。因为凡属于天主的，在我们使用中都是共有的；没有人被排除在祂的恩惠和恩赐之外，以致阻止全人类平等享有天主的良善和慷慨。同样，日光平等照耀，太阳发出光辉，雨水滋润，风吹拂，睡眠对睡眠者是一样的，星辰和月亮的光辉是共有的。在这种平等的榜样中，凡作为大地的拥有者，与弟兄分享其收益和果实的人，在其无偿的慷慨中既公平又公正，就是天主圣父的效法者。\par

26. 最亲爱的弟兄们，那些以爱德劳苦工作的人，将有何等光荣？——当主开始清点祂的子民，并根据我们的功德和善工分施所应许的赏报，以天上之物换取地上之物，以永恒之物换取暂时之物，以伟大之物换取微小之物时，那喜乐将是何等伟大崇高！祂将把我们呈献给圣父，因祂藉圣化使我们归向圣父；将把不朽和永生赐予我们，因祂藉祂宝血的活力更新了我们；将再次带我们进入乐园，开启天国，这一切都基于祂应许的信实和真理！愿这些事坚定地存留在我们的感知中，愿它们被完满的信德所理解，愿它们被我们的全心所热爱，愿它们藉着我们日益增长的劳作的慷慨而被获取。最亲爱的弟兄们，爱德的救恩劳苦真是一件光辉而神圣的事：它是信徒的大安慰，是我们安全的健康护卫，是望德的保护，是信德的保障，是罪的补救，是行善者能力内的事，既伟大又容易，是不受迫害风险的和平冠冕；是天主真实而最大的恩赐，为软弱者所必需，为强壮者所光荣，基督徒藉此恩赐成就属神恩宠，堪得审判者基督的欢心，视天主为他的债务人。为这救恩之工的棕榈枝，让我们欣然且爽快地奋斗；让我们众人在正义的斗争中，在天主和基督的注视下奔跑；让我们这些已经开始超越今生和此世的人，不因对今生和此世的任何贪恋而松懈我们的进程。若那日子——无论是赏报的日子还是迫害的日子——发现我们已预备好，若我们敏捷，若我们在爱德这竞赛中奔跑，主必不会不按我们的功德赐予赏报：在和平中，祂将赐给我们这些得胜者一顶白色的冠冕，作为我们劳作的赏报；在迫害中，祂将加上一顶紫色的冠冕，作为我们苦难的赏报。\par

\SourceRecord{Translated by Robert Ernest Wallis.；Ante-Nicene Fathers；Vol. 5.；Edited by Alexander Roberts, James Donaldson, and A. Cleveland Coxe.；Buffalo, NY: Christian Literature Publishing Co.,；1886.；<http://www.newadvent.org/fathers/050708.htm>.}

% structure: p0001 source=h1 target=section reason=page-title

\section{第九论}

\Emph{内容摘要——西彼廉本人在他致儒巴亚努斯书的结尾扼要陈述了本篇论著的缘起如下：\Quote{爱德的精神、我们团体的荣誉、信德的联系与司铎间的和谐，都藉我们的忍耐与温和得以保持。为此，我们还蒙主的许可与默感，尽我们微薄的能力，写了一本论忍耐之益的小册子，\InnerQuote{为彼此爱德的缘故，我们寄给了你们。}} 公元256年。}\par

1. 亲爱的弟兄们，我正要谈论忍耐，并要说明它的优点与益处，我应从何讲起呢？岂不更应从这一点入手：我看见此刻，你们听我讲论时，也需要忍耐，因为你们甚至连听讲和学习这一职责，若没有忍耐也无法完成吗？因为健康有益的谈话与推理，只有耐心聆听所说的话，才能有效地学到。亲爱的弟兄们，我发现，在众多天主治训的方式中，我们的望德与信德之路朝着获得天主赏报的方向前行，其中没有任何事物比那怀着敬畏与虔诚的服从、在主的诫命上劳苦的人，更应以全然的警醒，谨慎持守忍耐更为有益——无论是更利于生活，或是更有助於光荣。\par

2. 哲学家们也声称追求此德性；但在他们那里，忍耐之虚假一如他们的智慧。因为一个人若既不曾认识天主的智慧也不认识天主的忍耐，他又怎能是智慧或忍耐的呢？因为主亲自警告我们，论到那些在今世自以为有智慧的人，说：\BibleQuote{我要摧毁智慧人的智慧，我要废除明智人的明智。}\BibleRef{依 29:14} 此外，蒙福的宗徒保禄，充满圣神，被派遣去召唤和训诲外邦人，也作证并教训我们说：\BibleQuote{你们要谨慎，免得有人用哲学和虚妄的欺诈，按人的传统，按世俗的元素，而不是按基督，将你们掳去；因为在天主内住有全部圆满的神性。}\BibleRef{哥 2:8} 在另一处他说：\BibleQuote{谁也不要自欺；你们中若有人自以为今世有智慧，最好变为愚昧，好成为有智慧的。因为这世界的智慧在天主面前是愚妄。因为经上记载：我要在智者自己的诡计中捉住他们。} 又说：\BibleQuote{上主知道智慧人的思念，都是虚妄。}\BibleRef{格前 3:18} 因此，若他们当中的智慧不真，忍耐也就不真。因为如果一个人是谦卑而温良的，他就是智慧的——但我们却看不到哲学家们是谦卑或温良的，反而极其自满，且正因他们自满，便使天主不悦——那么在那些有虚饰自由的傲慢放肆，以及袒胸露怀的不知羞耻的自我夸耀之处，他们当中的忍耐便显然不真实了。\par

3. 但对于我们，亲爱的弟兄们，我们不是言语上，而是行为上的哲学家；不是在服饰上，而是在真理上展示我们的智慧；我们是更认识良心的省察，而非德性的夸耀；我们不讲论伟大的事，而是活出它们——让我们作为天主的仆人与朝拜者，在属灵的服从上，显露出我们从天主的训导学得的忍耐。因为我们与天主共有此德性。忍耐起于祂；它的荣耀与尊贵也源于祂。忍耐的起源与伟大出自天主为其创作者。人应当爱天主所珍爱的事物；上主的威严所爱之美善，它便予以称赞。如果天主是我们的主和父亲，就让我们效法我们主和我们父亲的忍耐；因为仆人应当服从，同样儿子也不可堕落到不肖。\par

4. 但是，天主的耐心是何等伟大！祂极其耐心地容忍世俗的庙宇、地上的偶像以及人为的亵渎礼仪，尽管这些是对祂威严和荣耀的藐视。祂使白日开始，使太阳的光照临善人和恶人；祂用雨水浇灌大地，没有人被排除在祂的恩惠之外，祂同样降雨给义人和不义之人。我们看到，在祂的命令下，季节以无差别的耐心服务于有罪者和无罪者、虔诚者和不敬者、感谢者和不感谢者；元素为他们服务，风为他们吹，泉源为他们流，丰收的丰盛增长，葡萄园的果实成熟，树上挂满苹果，树林披上绿叶，草地披上翠绿。尽管天主因频繁、甚至不断的过犯而被激怒，祂却缓和自己的愤怒，以耐心等待那一次定下的报复之日；虽然祂有报复的能力，却宁愿长期保持耐心，也就是说，仁慈地忍耐和推迟，以便可能的话，长期延续的祸患终会改变，而陷入错误和罪行污染的人，即使为时已晚，也能归向天主，正如祂自己警告并说：\Quote{我不愿那死去的人死亡，而愿他回头并生活。}又说：回来归向我吧，上主说。又说：\Quote{归向上主你们的天主吧，因为祂是慈悲、宽仁、忍耐、大发慈悲的，并且祂将审判转向那施加的恶。} \BibleRef{岳 2:13}此外，蒙福的宗徒在提及并召唤罪人悔改时，也指出并说：\Quote{或者你藐视祂丰富的恩宠、容忍和坚忍，不知道天主的耐心和恩宠引领你悔改吗？但因你的顽固和不悔改的心，你在愤怒之日、天主正义审判的启示之日，为自己积蓄愤怒，祂要按照各人的行为报应各人。} \BibleRef{罗 2:4}他说天主的审判是公义的，因为它是迟延的，因为它长久而大大地推迟，以便藉着天主的长期耐心，人可以受益于永生。惩罚是在罪人和不敬者身上执行的，那时对罪的悔改已不再有用。\par

5. 而且，为使我们更充分地理解，亲爱的弟兄们，耐心是天主的事，凡是温柔、耐心、温良的人，就是天主圣父的效仿者。当主在福音中赐予救恩的诫命，发出神圣的警告，教导祂的门徒达于成全时，祂规定并说：\Quote{你们听说过，\InnerQuote{你应爱你的近人，恨你的仇人。}但我对你们说：爱你们的仇人，为迫害你们的人祈祷，好使你们成为你们在天之父的子女，因为祂使自己的太阳升起，光照善人和恶人，降雨给义人和不义的人。如果你们只爱那爱你们的人，你们有什么赏报呢？税吏不也这样做吗？如果你们只问候你们的弟兄，你们做了什么特别的事呢？外邦人不也这样做吗？所以你们应当是成全的，如同你们的天父是成全的一样。} \BibleRef{玛 5:43}祂说天主的子女这样就会成为成全的。祂表明他们就是这样完成的，并教导他们藉着天上的重生而恢复，如果天主我们的父的耐心居住在我们内——如果那因亚当犯罪而失去的天主肖像在我们的行动中彰显并照耀出来。成为像天主一样，是何等的荣耀！在我们的德行中拥有那可与神圣赞美相提并论的，是何等伟大的福气！\par

6. 而且，亲爱的弟兄们，我们的天主和主耶稣基督不仅用言语教导这些，还用行动实践了。因为祂曾说过祂降下正是为此，为承行祂父的旨意；在其他彰显天主性尊威的德行奇迹中，祂也在祂忍耐的恒心中保持了祂父的耐心。最后，祂所有的行动，甚至从祂降临时起，都以耐心为伴侣；首先，从那天上的崇高降下到地上的事物，天主子并不轻视穿上人的肉体，虽然祂自己不是罪人，却担当了别人的罪。祂暂时放下祂的不朽，使自己成为可朽的，以便无罪者可以为有罪者的得救而死。主由仆人受洗；将要赐予罪赦的祂，自己却不轻视在重生池中洗祂的身体。祂禁食四十天，而祂却使别人饱飨。祂饥饿，忍受饥荒，以便那些曾饥饿于圣言和恩宠的人，能以天上的食粮得到饱足。祂与诱惑祂的魔鬼搏斗；并且仅满足于战胜敌人，不再进一步争斗，只用言语。祂统治祂的门徒，不是像主人有权能统治仆人；而是仁慈温和，以兄弟之爱爱他们。祂甚至屈尊洗宗徒的脚，以便主在祂的仆人中间如此，藉着祂的榜样教导一个同仆在同伴和同等的人中应该怎样。在顺服的人中间祂显如此，并不奇怪，因为祂能以长久的耐心容忍犹大直到最后——能与敌人一同吃饭——能认识家内的敌人而不公开指出，也不拒绝叛徒的吻。此外，在忍受犹太人时，祂有何等大的镇静和耐心：用劝说使不信者转向信仰，用让步抚慰忘恩者，温柔地回答反对者，宽仁地忍受骄傲者，谦卑地向迫害者让步，愿意聚集杀害先知和一直反抗天主的人，甚至直到祂十字架和受难的时刻！\par

7. 而且，在他的苦难与十字架上，在尚未达到死亡的残酷与流血之前，祂耐心地听到了何等羞辱的指责，忍受了何等嘲弄的凌辱，以至于祂接受了侮辱者的唾沫——\Emph{祂}曾用祂的唾沫为盲人造了眼睛；如今在祂的名下，魔鬼及其使者被祂的仆人鞭打，而祂自己却遭受了鞭打！祂被荆棘加冕，而祂却以永恒的花朵为殉道者加冕。祂的脸被手掌击打，而祂却将真正的棕榈枝赐予得胜者。祂被剥夺了地上的衣裳，而祂却以不朽的衣袍遮盖他人。祂被喂以苦胆，而祂却赐予天上的食粮。祂被给予醋喝，而祂却设立了救恩之杯。那无罪、那义者——不，祂就是无辜本身、公义本身——被列于罪犯之中，真理被假见证压制。那将要审判的被审判；天主的圣言被默默地带去宰杀。在主的十字架旁，星辰混乱，元素动荡，大地震动，黑夜遮蔽白日，太阳为了避免被迫观看犹太人的罪行，收回了它的光线和眼睛——祂却不言语，不动摇，甚至在自己的苦难中也不彰显祂的威严。直到最后，一切都被坚忍地、恒久地承受，为的是在基督内，圆满而完美的忍耐得以完成。\par

8. 在这些事之后，祂仍然接纳那些杀害祂的人，只要他们愿意悔改并归向祂；祂以救赎的忍耐，为保全人而关闭祂的教会，不拒绝任何人。那些仇敌、那些亵渎者、那些始终敌视祂名的人，如果他们为罪悔改，承认所犯的罪行，祂不仅接纳他们获得罪的赦免，而且还赐予天国的赏报。还有什么能比这更忍耐、更慈悲呢？连流基督血的人也因基督的血而活过来。基督的忍耐是如此伟大；若非如此伟大，教会就永远不会拥有保禄作为宗徒。\par

9. 但是，亲爱的弟兄们，如果我们是在基督内，如果我们穿上祂，如果祂是我们救恩的道路，我们就当跟随基督的足迹，以基督为榜样行走，正如宗徒若望教导我们说：\BibleQuote{那说自己住在祂内的，就应当照祂所行的去行。} \BibleRef{若一 2:6} 伯多禄——主屈尊将教会建立在他身上——也在他的书信中规定说：\BibleQuote{基督为我们受苦，给你们留下了榜样，叫你们跟随祂的足迹，祂没有犯过罪，口中也从未出过诡诈；祂被辱骂，却不还骂；受难时，也不恐吓，反而将自己交托给那不公义审判祂的。}\par

10. 最后，我们发现祖辈与先知，以及所有在他们先前的相似中预表基督的义人，在德行的赞美中，最留意保持的就是以坚强而稳定的平静保全忍耐。因此，亚伯尔首先开创并奉献了殉道的起源和义人的苦难，没有抵抗也没有与杀害他的兄弟争斗，而是以谦卑和温顺耐心地被杀。因此，亚巴郎相信天主，最先奠定了信德的基础和根基，在关于他儿子的试验中，毫不迟疑或拖延，以敬虔的所有忍耐服从了天主的命令。依撒格，作为主祭品的预像，当被他父亲献为燔祭时，被发现有忍耐。雅各伯，被他的兄弟从故乡驱逐，以忍耐离开；后来以更大的忍耐，他用和平的礼物谦卑地使他恢复和睦，尽管他更加不虔敬和迫害。若瑟，被他的兄弟们出卖并送走，不仅以忍耐宽恕了他们，甚至在他们来到他面前时，慷慨而仁慈地赐予了他们无偿的粮食。梅瑟屡次被忘恩负义、不信的百姓藐视，几乎被石头打死；然而他以温和与忍耐为那百姓向主恳求。但在达味身上——按肉身，基督的诞生从他而来——忍耐是多么伟大、奇妙和基督徒啊！他时常有机会能杀死那迫害他、想要杀他的撒乌耳王，却宁愿在他被交到手中时救他，不报复他的敌人，反而在他被杀后为他伸冤！总之，如此多的先知被杀害，如此多的殉道者以光荣的死亡受尊崇——他们都因忍耐的赞美而获得了天上的冠冕。因为忧伤和苦难的冠冕，除非先有忧伤和苦难中的忍耐，否则无法获得。\par

11. 但是，为了更明显地、更充分地认识忍耐是多么有用和必要，亲爱的弟兄们，让我们默想天主的审判——在世界的起初、人类之初，亚当忘了诫命，违反所赐法律时，他所接受的审判。这样我们就会知道，我们在这生命中应当多么忍耐——我们生于这样的状态，在这里因苦难和争战而辛劳。\BibleQuote{因为，}祂说\BibleQuote{你听从了你妻子的声音，吃了我唯独吩咐你不可吃的那树上的果子，地为你的缘故受咒诅：你一生日日劳苦，才能从地里得吃。地要给你生出荆棘和蒺藜，你也要吃田间的菜蔬。你必须汗流满面，才有饭吃，直到你归于你从那里出来的土地：因为你是尘土，还要归于尘土。} \BibleRef{创 3:17} 我们都因这判决的锁链被束缚捆绑，直到死亡被涂抹、我们脱离这生命。我们一生日日必须忧愁和叹息：我们必须以汗水和劳苦吃我们的饭食。\par

12. 因此，我们每个人，当他出生并被接到这世界的旅店时，都是从眼泪开始的；尽管仍无意识且对一切无知，他在那最初的诞生中除了哭泣以外什么都不知道。凭着天然的预见，未经训练的灵魂为必死生命的焦虑和劳苦而哀叹，甚至在开始时就通过它的哀号和呻吟来见证它即将进入的世界风暴。因为额上的汗水和劳作是生命持续的条件。除了耐心之外，无法为那些流汗辛劳的人提供任何安慰；这些安慰，虽然在今世对所有的人都合适且必要，但对我们尤其如此，因为我们更受到魔鬼围攻的摇撼，我们每日站在战场上，与一个老练而狡猾的敌人角力而疲惫；我们除了各种不断的诱惑战斗外，还必须在迫害的争战中抛弃家产，忍受监禁，背负锁链，舍弃生命，承受刀剑、野兽、火刑、十字架——总之，各种折磨和刑罚，都要在信德和耐心的勇气中忍受；正如主亲自教导我们说：\Quote{我给你们讲了这些事，是要使你们在我内得到平安。但在世界上你们要受苦难；然而你们放心，我已经战胜了世界。} \BibleRef{若 16:33} 如果我们这些弃绝了魔鬼和世界的人，更频繁、更猛烈地遭受魔鬼和世界的苦难与灾害，那么我们岂不更应该持守耐心，以它作为帮助者和同盟，使我们能承担一切危害之事！\par

13. 这是我们的主和师傅的健康训诫：\Quote{那坚持到底的，} 他说，\Quote{必要得救；} \BibleRef{玛 10:22} 又说：\Quote{如果你们存留，} 他说，\Quote{在我的话内，你们才真是我的门徒；你们将认识真理，而真理必会使你们获得自由。} \BibleRef{若 8:31} 我们必须坚持和忍耐，亲爱的弟兄们，以便被接纳入真理和自由希望后，我们能达致真理和自由本身；我们身为基督徒这一事实正是信德和希望的实质。但为使希望和信德达到它们的结果，需要耐心。因为我们所追求的不是现世的荣耀，而是未来的，正如宗徒保禄也告诫我们说：\Quote{我们得救是在于希望，但看得见的希望不是希望：因为一个人看见的，为什么还希望呢？但我们若希望那看不见的，我们就以耐心等待它。} \BibleRef{罗 8:24} 因此，等待和耐心是必需的，使我们能完成我们已开始成为的，并能获得我们根据天主自己的启示所相信和希望的。此外，在另一处，同一位宗徒教导正义和行善的人，以及那些以神圣高利贷的增加在天上为自己积蓄宝藏的人，他们也应当忍耐；并教导他们说：\Quote{因此，我们一有机会，就当向众人行善，尤其向那些有同样信德的家人。但我们不可行善而懈怠，因为我们到了时候，必要收获。} 他劝诫没有人应该因不耐烦而在劳动中懈怠，没有人应该被诱惑所引离或击败，而在赞美中途中和荣耀的道路上放弃；以至过去的事败坏，而那些已经开始的事也不得成全；正如经上所记载：\Quote{义人的义德，在他任何犯罪的日子，不能救他；} \BibleRef{则 33:12} 又说：\Quote{你要持守你所有的，免得别人夺取你的冠冕。} \BibleRef{默 3:11} 这话语劝勉我们以耐心和勇气坚持，以便那怀着临近的赞美奔向冠冕的人，因持续的忍耐而得冠冕。\par

14. 但是耐心，亲爱的弟兄们，不仅守护着善，也抵御着恶。它与圣神和谐一致，并与属天和神圣的事物相连，以其力量之防御，对抗那攻击和俘虏灵魂的肉身和身体的行为。让我们从许多中略观几样，以便从几样中也能理解其余的。奸淫、欺诈、凶杀，都是大罪。让耐心在心中坚强而坚定；那么，圣化了的身体和天主的殿宇就不会被奸淫玷污，献给正义的无辜就也不会被欺诈的感染所污染；而且，在领了圣体之后，手也不会被剑和血所玷污。\par

15. 爱德是兄弟情谊的纽带，和平的基础，合一的支柱和保障，它比希望和信德都大，它超越善工和殉道，它将永远与我们同在，与天主一起在天国永恒。从它那里拿走耐心；失去了耐心，它就不能持久。从它那里拿走忍耐和忍受的实质，它就毫无根基和力量继续。最后，宗徒在谈论爱德时，将忍耐和耐心与它联结起来。他说：\Quote{爱德是宽宏大量的；爱德是慈祥的；爱德不嫉妒，不自大，不动怒，不思想恶；爱一切，信一切，望一切，忍一切。} \BibleRef{格前 13:4} 于是他表明它能够坚忍地坚持，因为它知道如何忍受一切。在另一处：\Quote{互相宽恕，} 他说，\Quote{在爱中，尽一切努力保持心神在和平联系中的合一。} \BibleRef{弗 4:2} 他证明，除非弟兄们以相互容忍彼此珍惜，并通过耐心的介入保持和谐的纽带，否则合一与和平都不能保持。\par

16. 此外还有什么呢？——就是你们不可起誓，也不可咒骂；当你们的财物被夺去时，不可追讨；当你们挨打时，不但要把另一面颊转给打你的人，而且要宽恕得罪你们的弟兄，不止七次，而是七十个七次，并且宽恕他所有的罪过；要爱你们的仇敌，要为反对和迫害你们的人祈祷。除非你们坚持忍耐和坚忍的稳固，你们能做成这些事吗？我们在\Person{斯德望}身上看到这一点，当他被犹太人用暴力和石头打死时，他没有为自己求报复，反而为杀他的人求赦免，说：\BibleQuote{主，不要将这罪归于他们。}\BibleRef{宗 7:60}基督的第一位殉道者理当如此，他以光荣的死亡为将跟随他的殉道者作先行者，他不仅是主苦难的宣讲者，也是祂最忍耐之温良的效仿者。对于忿怒、纷争、争吵，这些事在基督徒身上本不该有，我还能说什么呢？如果心中有忍耐，这些事就无法在那里立足；或者它们试图进入，也会很快被排除而离去，好使平安的居所在心中持续，天主的平安乐意居住在那里。最后，宗徒告诫我们并教导说：\BibleQuote{不要使天主的圣神忧郁，你们是在祂内受了印，以待得救的日子。一切苦毒、忿怒、恼恨、喊叫、亵渎，都当从你们中间除掉。}\BibleRef{弗 4:30}因为如果基督徒已经脱离了怒气和肉体的争战，如同脱离了海上的飓风，并开始在基督的港湾中变得宁静和温良，他就不该在心中容纳忿怒或纷争，因为他既不可以恶报恶，也不可怀恨。\par

17. 此外，对于肉身的各种病痛，以及身体频繁而剧烈的折磨——人类日复一日因此感到疲惫和困扰——忍耐也是必要的。因为自从第一次违命之后，身体的强壮与不死性一同消失，软弱与死亡一同来临——而强壮只有在不死性也获得之后才能重新领受——因此，我们在这肉身的脆弱和软弱中，必须常常争战和搏斗。而这场争战和拼搏，除非借着忍耐的力量，是不能持守的。但为了我们被考验和查究，引入了各种苦难；各种试炼借着财产的损失、高烧的热度、伤口的疼痛、亲人的离世而加增。没有什么比在苦难中更能区分不义的人和义人了：不义的人不耐烦地抱怨和亵渎，而义人却因忍耐而得到证明，正如经上所记：\BibleQuote{在痛苦中要忍耐，在卑微中要有耐心；因为金银要在火中锻炼。}\BibleRef{德 2:4}\par

18. 就这样，\Person{约伯}被查究和考验，并因忍耐的德性被提升到赞美的最高顶峰。魔鬼向\Person{约伯}射出了怎样的箭矢！使用了怎样的酷刑！财产的损失加给他，众多儿女的断绝被命定给他。那位在财产上富有的主人，在儿女上更富有的父亲，瞬间既不是主人也不是父亲！伤口的溃烂又加了上去，更有侵蚀的虫豸吞噬着他化脓和腐烂的四肢。为了使\Person{约伯}在他的试炼中一无所缺，魔鬼也武装了他的妻子，使用了他那邪恶的老计谋，仿佛他能像在创世之初一样，借着女人欺骗和误导所有人。然而，\Person{约伯}并没有被这些严重而反复的冲突所击垮，也没有因忍耐的胜利而停止在那些困难和试炼中宣告天主的祝福。\Person{托彼特}也是如此，他在施行了崇高的正义和仁慈之工后，被夺去了双眼；他越是在忍耐中承担失明，就越是借着忍耐的赞美而大大地蒙受天主的恩宠。\par

19. 而且，亲爱的弟兄们，为使忍耐的益处更加显明，让我们反过来思考不忍耐会带来怎样的祸害。因为正如忍耐是基督的益处，反过来，不忍耐是魔鬼的祸害；正如基督居住和常驻在其中的人被发现有忍耐，同样，那心思被魔鬼的邪恶所占据的人总是显得不忍耐。让我们简要地看看最初的开始。魔鬼不忍耐人按天主的肖像被造，因此他首先丧亡并毁灭了他人。\Person{亚当}违背了关于死亡食物的天上命令，因不忍耐而陷入死亡；他也没有在忍耐的守护下持守从天主领受的恩宠。而为了\Person{加音}杀死他的弟弟，是他不能忍耐自己的祭品和礼物；在\Person{厄撒乌}从长子的权利降到次子的权利时，他因急于喝那红豆汤而失去了长子的优先权。为什么犹太人在天主的恩惠面前不忠和忘恩负义？难道不是不忍耐的罪行使他们首先离开了天主？他们不能忍受与天主交谈的\Person{梅瑟}的迟延，竟敢要求拜邪神，好让牛头像和泥土像作他们行军的领袖；他们从未停止不忍耐，直到他们始终不耐于受教和天主的训诫，杀死了他们的先知和所有的义人，甚至陷入钉死主和流祂血的罪行。此外，不忍耐在教会中制造异端者，并仿效犹太人，像叛逆者一样驱使他们对抗基督的平安与爱德，陷入敌意的和狂怒的仇恨。而且，不再逐一列举，凡是忍耐借着其作为建立为光荣的，不忍耐都彻底投奔毁灭。\par

20. 因此，亲爱的弟兄们，既已殷勤思量忍耐的裨益与不忍耐的祸害，就让我们以全然的警醒持守我们赖以在基督内安住的忍耐，好使我们与基督一同抵达天主。这忍耐丰盛而多样，不受狭窄的界限所约束，亦不为狭隘的边界所限。忍耐之德广泛显明，其丰饶与慷慨虽源自同一个名称的泉源，却借着从荣耀的多条道路上涌出的浩荡溪流而扩散；任何行为若未从中获得其成全的实质，便无助于赞美的圆满。正是那忍耐把我们交托并保守于天主。也是那忍耐平息愤怒、约束口舌、治理心神、守护和平、管束纪律、打破情欲的力量、压制骄傲的猛烈、熄灭敌意的火焰、抑制富人的权势、安抚穷人的匮乏、保护贞女有福的纯洁、寡妇小心的贞洁、对已婚者单一的恩爱。它使人在顺境中谦卑，在逆境中勇敢，对冤屈和轻蔑温和。它教导我们迅速宽恕伤害我们的人；若你自己做了错事，则要长久而恳切地哀求。它抵抗诱惑，忍受迫害，成就苦难与殉道。正是那忍耐坚固地巩固我们信德的基础。正是它高举我们望德的增长。正是它引导我们的行为，使我们行走在基督的忍耐中而能持守基督的道路。正是它使我们作为天主的子女恒心持守，因为我们效法天父的忍耐。\par

21. 但由于我知道，亲爱的弟兄们，许多人因着负担或伤痛之苦的刺痛，急于报复那些对他们粗暴行事和发怒的人，我们绝不可隐瞒这一事实，即我们置身于这冲突世界的风暴之中，且此外还有犹太人、外邦人和异端者的迫害，我们可耐心等待（天主）报复的日子，不可因着怨诉的仓促而急忙报复我们的痛苦，因为经上记载：\Quote{等候我，主说，直到我起来作证的那一天；因为我已判定聚集万民，聚集列国，将我的忿怒倾在他们身上。}\BibleRef{索 3:8}主命令我们等候，并勇敢地忍耐未来报复的日子；祂也在\WorkTitle{默示录}中说：\Quote{不可封了这书上的预言，因为时候近了。叫那不义的仍旧不义，污秽的仍旧污秽，为义的仍旧行义，圣洁的仍旧圣洁。看哪，我必快来！赏罚在我，要照各人所行的报应他。}\BibleRef{默 22:10}为此，那些呐喊并因悲伤爆发而急于报复的殉道者，也被吩咐仍要等候，并以忍耐等待时日的满足和殉道者的圆满。\Quote{当祂揭开}他说，\Quote{第五印，我看见在祭台底下，有为天主的道并为作证而被杀之人的灵魂。他们大声喊着说：\InnerQuote{圣洁真实的主啊，你不审判住在地上的人，给我们伸流血的冤，要等到几时呢？}于是给了他们每人一件白衣，并有话对他们说：\InnerQuote{还要安息片时，等着一同作仆人的和他们的弟兄也像他们被杀，满足了数目。}}\BibleRef{默 6:9}\par

22. 但义人之血的报复何时来临，圣神借\Person{玛拉基亚}先知宣告说：\Quote{看哪，那日临近，势如烧着的火炉，凡外邦人和一切作恶的都要如碎秸；那来临的日子必将他们烧尽。}\BibleRef{拉 4:1}这事我们在\BibleBook{}{圣咏}中也读到，那里宣告天主审判者的来临，因其审判的威严而可敬畏：\Quote{天主必显现，我们的天主，且不缄默；有火在祂面前烧起，有暴风在祂四周旋转。祂呼唤上天下地，为要审判祂的民。招聚祂的圣民到祂那里，那些用祭献与祂立约的人。诸天要宣告祂的公义，因为天主自己是审判者。}\Quote{因看！上主必像火降临，祂的战车如旋风，以烈怒施行报复；因为上主的火要施行审判，祂的刀要使他们受伤。}\BibleRef{依 66:15}又说：\Quote{上主万军的天主必出战，将战争粉碎；祂要激动战事，并向祂的敌人强力呼喊：\InnerQuote{我缄默已久，难道我要永远缄默吗？}}\BibleRef{依 42:13}\par

23. 但那位先前说祂缄默却不会永远缄默的是谁呢？诚然就是那被牵去宰杀的羊，又如羔羊在剪毛者前无声，祂始终不开口。诚然就是那在街上不喊叫、不扬声的。诚然就是不悖逆、不抗辩的那一位，当祂将背给鞭打，将颊给掌掴，也不转面躲避唾污。诚然就是那在司铎和长老控告时一言不答，令\Person{比拉多}惊奇而保持最忍耐的沉默的。这就是那位，虽然祂在苦难中默然无声，但不久在报复中将不再沉默。这就是我们的天主，即不是所有人的天主，而是信者与信徒的天主；当祂在第二次来临中显现时，将不再沉默。因为虽然祂第一次来临是披着谦卑的帷幕，但祂将要在权能中显现。\par

24. 让我们等候祂，亲爱的弟兄们，我们的审判者和复仇者，祂必要与祂自己一同为祂教会的集会以及从世界起初以来所有义人的数目复仇。那急于复仇、缺乏耐心的人，应当思想连祂自己——那复仇者——尚未得复仇。天主圣父命令祂的圣子受敬拜；宗徒\Person{保禄}谨记神圣命令，立下规矩说：\Quote{天主极其举扬祂，赐给祂一个名字，超越一切名字，为使因耶稣之名，天上、地上和地下的万膝都屈拜。} \BibleRef{斐 2:9} 在\BibleBook{}{默示录}中，天使阻止\Person{若望}想要敬拜他，说：\Quote{不可这样；我是你的同仆，和你弟兄们的同仆。你要敬拜主耶稣。} 主耶稣何等伟大，祂的忍耐何等伟大，那在天上受敬拜的，在地上尚未得复仇！亲爱的弟兄们，让我们在我们的迫害和苦难中思念祂的忍耐；让我们以充满期待的顺服迎候祂的来临；我们身为仆人，不要以不虔敬和不庄重的急切催促在我们主前得辩护。反而让我们奋勇向前，劳苦，全心警醒，恒心忍耐一切，遵守主的诫命；好使那忿怒和报复的日子来到时，我们不与不虔敬的人和罪人一同受罚，而与义人和敬畏天主的人一同得荣耀。\par

\SourceRecord{Translated by Robert Ernest Wallis.；Ante-Nicene Fathers；Vol. 5.；Edited by Alexander Roberts, James Donaldson, and A. Cleveland Coxe.；Buffalo, NY: Christian Literature Publishing Co.,；1886.；<http://www.newadvent.org/fathers/050709.htm>.}

% structure: p0001 source=h1 target=section reason=page-title

\section{第十论}

\Emph{论嫉妒与妒忌。}\par

\Emph{内容提要}——在指出嫉妒或妒忌因其邪恶的隐蔽性而更为可憎，且其根源可追溯至魔鬼之后，他举出旧约中妒忌的例证，并藉着提及特殊的恶习，指出妒忌是一切邪恶的根源。因此，基督及其宗徒不止一处禁止兄弟间的仇恨是有理由的。最后，他以天主的榜样劝勉爱仇，藉着敦促爱德所预备的赏报，劝阻妒忌之罪。\par

1. 亲爱的弟兄们，在有些人眼中，对你所见的美好事物心存嫉妒，对优于你的人心怀妒忌，似乎是一件轻微而微不足道的过错；并且，由于被视作琐碎小事而不受重视，便不令人惧怕；因不惧怕而被轻视；因被轻视而不易避之；于是它变成一种幽暗隐蔽的祸害，因其不被明智之人所察觉防备，而暗自折磨不加提防的心灵。然而，主吩咐我们要明智，并嘱咐我们以谨慎的忧虑警醒，免得那常备不懈、常潜伏伺机的仇敌悄然潜入我们胸中，从火星点燃烈焰，将小事扩大为大事；它用温和的气息与柔和的微风安抚那毫无防备与粗心大意的人，却掀起风暴与旋风，导致信德的毁坏与救恩和生命的沉沦。因此，亲爱的弟兄们，我们必须警惕，竭尽全力以谨慎而全备的警醒抵挡那仇敌；他四处咆哮，将箭矢瞄准我们身体每一处可能被击伤的部位——正如宗徒伯多禄在他的书信中所预言和教导的：\Quote{你们要节制，要警醒，因为你们的仇敌魔鬼如同咆哮的狮子巡游，寻找可吞吃的人。}\BibleRef{伯前 5:8}\par

2. 他潜伏在我们每个人周围；如同围困城中居民的仇敌，他察探城墙，试探是否有哪处城墙不够坚固、不够可靠，可以从那里潜入内部。他向眼睛呈现诱惑眼目的形象和安逸的享乐，以摧毁贞洁；他以和谐的音乐诱惑耳朵，藉着悦耳的声音松懈和削弱基督徒的刚毅；他以辱骂激怒舌头；他以令人愤慨的伤害刺激手，使之不惜谋杀；他呈现不义的财富，诱人行骗；他堆聚有害的宝藏，以钱财俘虏灵魂；他应许地上的尊荣，以剥夺天上的尊荣；他显弄虚假的事物，以偷窃真实的；当他不能隐蔽地欺骗时，他便公开而明显地威胁，以骚乱的迫害的恐惧来征服天主的仆人——他始终不安，始终敌对，在和平中诡诈，在迫害中凶暴。\par

3. 因此，亲爱的弟兄们，为对抗魔鬼一切欺骗的圈套或公开的恐吓，心灵应当整装待发，随时准备抵御，正如仇敌时时常备进攻一样。既然那些悄然潜入的箭矢更为频繁，他更为隐蔽而秘密的投掷对我们的伤害也更为严重而频繁（因其越不为人所察觉），我们就当同样警醒，以识别并抵御这些箭矢，其中就包括了嫉妒与妒忌之恶。若有人细察此事，便会发现，基督徒最应提防、最应谨慎留意的，莫过于被嫉妒与恶意所俘虏；免得有人陷入那欺骗之敌的盲目罗网——当弟兄因嫉妒而转为仇恨弟兄时，自己竟不知不觉被自己的剑所毁灭。为能更充分收集并更清楚察觉这事，让我们追溯到它的源泉与起源。让我们思量嫉妒从何而生，以及它何时、如何开始。因为如此有害的邪恶，一旦我们知道了它的源头与庞大，就更容易避开。\par

4. 从这源头，就在世界的起初，魔鬼是第一个既（自身）丧亡又（使他人）丧亡的。他原被托付于天使的尊威，被天主所接纳与喜爱，当他看见按天主的肖像所造的人时，便因恶毒的妒忌而爆发嫉妒——并非由于他嫉妒的本能先将他击倒，他自身才先被嫉妒所击倒；他是先被俘虏，然后才俘虏他人；先丧亡，然后才使他人丧亡。当他在嫉妒的煽动下，剥夺了人那已赐予的不朽恩宠时，他自己却失去了先前所拥有的。亲爱的弟兄们，那使得天使堕落、使得那崇高辉煌的尊荣可能受骗与倾覆、使得那欺骗者自身反受欺骗的，是何等巨大的邪恶啊！从此，妒忌在地上肆虐，因那即将因嫉妒而丧亡的人，服从了他毁灭的根源，在妒忌中效仿魔鬼；正如经上所载：\Quote{但因魔鬼的嫉妒，死亡进入了世界。}因此，凡属他那一边的人，都效仿他。\par

5. 于是，最终，新的兄弟情谊最初的仇恨开始了，于是，可憎的兄弟相杀发生了：不义的\Person{加音}嫉妒正义的\Person{亚伯尔}，恶人因嫉妒和妒忌迫害善人。嫉妒的愤怒如此盛行，以至于达到了那邪恶行为的顶峰，既不考虑兄弟之爱，也不考虑罪行的巨大，既不敬畏天主，也不考虑罪的惩罚。那首先显明正义的人反而被不义地击打；那不知如何仇恨的人却忍受了仇恨；那至死不抵抗的人却被不敬地杀害。\Person{厄撒乌}对兄弟\Person{雅各伯}的敌意也源于嫉妒。因为后者接受了父亲的祝福，前者因嫉妒的火把而被点燃，产生了迫害的仇恨。\Person{若瑟}被他的兄弟们出卖，出卖他的原因也来自嫉妒。当他以纯朴之心，如兄弟对兄弟，向他们讲述他在异象中看到的福乐时，他们恶毒的性情爆发成了嫉妒。此外，\Person{撒乌耳}王恨\Person{达味}，屡次迫害想要杀他——\Person{达味}无辜、仁慈、温柔、以温良忍耐——除了嫉妒的刺激之外，还有什么原因呢？因为当\Person{哥肋雅}被杀，藉天主的帮助和俯就如此强大的敌人被击溃时，惊叹的民众以欢呼的赞同爆发出来赞美\Person{达味}，\Person{撒乌耳}因嫉妒而产生了敌意和迫害的愤怒。而且，为了不一一列举，让我们观察那一次灭亡的百姓。犹太人岂不是因为这个原因而灭亡吗？因为他们宁愿嫉妒\Person{基督}也不愿相信祂？他们贬低祂所行的伟大工作，被盲目的嫉妒所欺骗，无法睁开他们心中的眼睛认识神圣的事物。\par

6. 亲爱的弟兄们，考虑到这些事，让我们以警惕和勇气，用对抗如此毁灭性邪恶的方式，坚固我们已奉献给天主的内心。让他人的死亡有益于我们的安全；让愚昧之人的惩罚给予明智之人健康。此外，没有人可以认为那种邪恶仅限于一种形式，或局限于狭窄界限内的短暂范围。嫉妒的祸害是多种多样的、多产的，广泛蔓延。它是万恶之源，灾难之泉，罪恶之苗圃，过犯之材料。由此产生仇恨，由此产生敌意。嫉妒点燃贪婪，因为一个人不能满足于自己所有的，当他看到另一个人更富有时。嫉妒激起野心，当一个人看到另一个人在荣誉上更被高举时。当嫉妒遮蔽我们的感知，并将心灵的隐秘活动置于其统治之下时，对天主的敬畏被藐视，\Person{基督}的教导被忽视，审判的日子不被预期。骄傲膨胀，残忍苦涩，不忠狡诈，急躁激动，纷争肆虐，愤怒炽热；那已服从外在权威的人不能再约束或控制自己。由此，主的和平纽带被打破；由此，兄弟之爱被侵犯；由此，真理被掺杂，合一被分裂；当司铎被贬低，主教被嫉妒，一个人抱怨自己没有而非被按立，或不愿忍受另一个人被置于他之上时，人们陷入异端和分裂。因此，那因嫉妒而傲慢、因嫉妒而乖僻的人踢踹，他反叛，在愤怒和恶意中反对的不是人，而是荣誉。\par

7. 但那是什么灵魂的啃噬之虫，什么思想的瘟疫斑点，什么心灵的锈迹——嫉妒另一个人，无论是他的德行还是他的幸福；也就是说，恨他身上的自己的功绩或天主的恩惠，将他人的益处转为自己的祸害，因杰出之人的福乐而受折磨，将他人的荣耀作为自己的惩罚，并且仿佛给自己的胸膛施加某种刽子手，将折磨者带进自己的思想和情感中，使他们以内在的痛苦撕裂我们，以恶意的蹄子击打心灵隐秘的处所。对这样的人，没有食物是愉快的，没有饮料能令人欢畅。他们总是叹息、呻吟和悲伤；因为嫉妒从未被嫉妒者放下，被占据的心灵日夜无间断地撕裂。其他邪恶有其限度；无论什么错误，都以罪行的完成为界限。在奸淫者身上，当侵犯发生时，罪行就停止了；在强盗的情况下，当凶杀完成时，罪恶就安息了；赃物的占有结束了小偷的贪婪；完成的欺骗限制了欺诈者的过错。嫉妒没有界限；它是持续不断的邪恶，是无尽的罪。被嫉妒者越有更大的成功优势，嫉妒者就越以嫉妒之火燃烧得更热。\par

8. 因此，有威胁的面容，低垂的脸色，脸上的苍白，嘴唇的颤抖，牙齿的摩擦，疯狂的言语，无节制的辱骂，手准备暴力的屠杀；即使暂时没有刀剑，也以狂怒之心的仇恨武装。因此，圣神在圣咏中说：\BibleQuote{不要嫉妒那在自己道路上顺利行走的人。} 又说：\BibleQuote{恶人窥伺义人，向他咬牙切齿。但天主笑他，因为见他日子来到。} 真福宗徒保禄指明并指出这些，他说：\BibleQuote{他们的嘴唇下有蛇毒，满口诅咒和苦毒；他们的脚飞快流血，毁灭和悲惨在他们的路上，他们不认识和平的道路；天主的敬畏不在他们眼前。}\par

9. 当四肢被剑所伤时，损害要轻微得多，危险也更小。伤口明显时，治疗容易；药膏敷上后，可见的疮伤很快痊愈。嫉妒的创伤却是隐藏而秘密的；它们不接受医治的补救，因为它们将自身封闭在良心暗处的盲目痛苦中。无论你是谁，若你嫉妒而怀恶意，请观察你对所恨之人是何等狡猾、恶毒和可憎。然而，你危害他人福祉的程度，远不及危害自身。你因嫉妒所逼迫的人，或可逃避你，但你无法逃避自己。无论你在何处，你的对头总与你同在；你的仇敌常在你自己的胸中；你的祸患封闭在内；你被无法挣脱的锁链捆绑束缚；你在嫉妒的暴政下成为俘虏；任何慰藉都无济于事。逼迫一个属于天主恩宠的人，是持续的恶行。仇恨幸福的人，是不可救药的灾难。\par

10. 因此，亲爱的弟兄们，主顾念这个危险，以免有人因嫉妒弟兄而陷入死亡的网罗。当祂的门徒问祂，他们中谁最大时，祂说：\Quote{你们中谁若成为最小的，那人就是最大的。}祂以回答断绝了一切嫉妒。祂拔除并撕去了所有啃噬嫉妒的缘由和素材。基督的门徒不可嫉妒，不可羡慕。在我们中间，不可有争高低的竞争；我们由谦卑成长到最高的境界；我们学会了如何中悦天主。最后，宗徒保禄教导并劝诫我们这些被基督的光照亮、逃离了黑夜行为的黑暗的人，应当行走在光明的事工和行为中，他写道说：\BibleQuote{黑夜已逝，白昼临近：所以让我们抛弃黑暗的行为，穿上光明的铠甲。让我们行事端正，如同在白昼；不在于狂欢醉酒，不在于纵欲放荡，不在于纷争嫉妒。}\BibleRef{罗 13:12} 如果黑暗已离开你的胸膛，如果黑夜已被驱散，如果阴霾已被赶走，如果白昼的光明已经照亮你的感官，如果你已开始成为光明之子，就做属于基督的事，因为基督就是光明和白天。\par

11. 你为什么冲入嫉妒的黑暗？为什么将自己裹入恶意的云中？为什么在嫉妒的盲目中熄灭平安和爱德的一切光明？为什么回到你曾弃绝的魔鬼那里？为什么你像\Person{加音}一样站立？因为那嫉妒弟兄、怀恨他的，就负了杀人的罪债，宗徒若望在他的书信中宣告说：\Quote{凡恨自己弟兄的，就是杀人犯；你们知道，凡杀人犯，没有永生存留在内。}又说：\BibleQuote{那说自己在光明中，却恨自己弟兄的，直到如今仍在黑暗中，且在黑暗中行走，不知道往那里去，因为黑暗蒙蔽了他的眼睛。}\BibleRef{若一 2:9} 凡恨自己弟兄的，他说，是在黑暗中行走，且不知往那里去。因为他无知而盲目地走向\OriginalTerm{革哈纳}{Gehenna}；他正匆忙奔向刑罚，也就是说，离开了基督的光，基督警告并说：\BibleQuote{我是世界的光；跟随我的，决不在黑暗中行走，必有生命的光。}\BibleRef{若 8:12} 但谁跟随基督，谁就坚守祂的诫命，行走在祂教导的道路上，追随祂的足迹和途径，效法基督所行和所教导的；这与伯多禄所劝诫和警告的相符，说：\BibleQuote{基督为我们受了苦，给你们留下了榜样，叫你们跟随祂的足迹。}\BibleRef{伯前 2:21}\par

12. 我们应当记住基督用什么名字称呼祂的子民，用什么称号命名祂的羊群。祂称他们为羊，好使他们的基督徒纯真如同羊的纯真；祂称他们为羔羊，好使他们心灵的纯朴效法羔羊的纯朴天性。为什么狼披着羊的外衣潜伏？为什么那虚假自称基督徒的人，羞辱基督的羊群？戴上基督的名，却不行在基督的道路上，这岂不是对神圣之名的嘲弄，岂不是对救恩之路的背离？因为祂亲自教导并说：遵守祂诫命的，必得生命；听见并实行祂话的，就是明智的；而且，那这样实行并教导的，在天国里被称为最大的师傅。那么，如果宣讲者口中所说的，被随后的事迹所成就，他所妥善而有益地宣讲的，对他岂不有益？但是，主更频繁地灌输给祂门徒的，更严厉地嘱咐他们在祂救恩的忠告和天上的规诫中要守护和遵守的，是什么呢？不就是以祂爱门徒同样的爱，我们也当彼此相爱吗？他既不能平和也不能有爱，当嫉妒侵入时，又如何能保持主的平安或爱呢？\par

13. 同样，\Person{保禄}宗徒在极力主张和平与爱德的功绩时，在坚决断言并教导说，无论是信德还是赒济，甚至\OriginalTerm{宣信者}{confessor}和殉道者的苦难本身，都对他无益，除非他完整无缺地保守爱德的要求时，又补充说：\Quote{爱德是宽宏的，爱德是仁慈的，爱德不嫉妒；}\BibleRef{格前 13:4}无疑是在教导和指明，凡是宽宏、仁慈、远离嫉妒和怨恨的人，就能持守爱德。此外，在另一处，当他劝告那已充满圣神、藉天上重生成为天主子的人，只应注视属神和属天的事物时，他规定并说：\Quote{弟兄们，我从前不能向你们说话如同对属神的人，只能如同对属血肉的人，如同对在基督内的婴孩。我给你们喝的是奶，不是饭；因为那时你们还不能吃，就是现在你们还是不能。因为你们仍是属血肉的；你们中间既有嫉妒、纷争和争吵，岂不仍是属血肉的，按照人的方式行事吗？}\BibleRef{格前 3:1}\par

14. 必须践踏恶习和属血肉的罪，亲爱的弟兄们，并以属神的活力踏碎属世身体的腐化瘟疫，免得我们转回到旧人的行为中，被致命的网罗缠住，正如宗徒以远见和健全的智慧预先警告我们这件事，并说：\Quote{因此，弟兄们，我们不要随从肉体生活；因为如果你们随从肉体生活，你们将要开始死；但如果你们藉着圣神治死肉体的行为，你们必要生活。因为凡受天主圣神引导的，都是天主的子女。}\BibleRef{罗 8:12}如果我们已是天主的子女，如果已开始成为祂的宫殿，如果领受了圣神而度圣洁属神的生活，如果已将目光由地上升至天上，如果已将祂和基督充满的心提升到天上和神圣的事物，那么我们就只应做配得上天主和基督的事，正如宗徒激励劝勉我们时说：\Quote{如果你们已与基督一同复活，就该追求天上的事，那里有基督坐在天主的右边；你们该思念天上的事，不该思念地上的事。因为你们已经死了，你们的生命已与基督一同藏在天主内。但当基督——你们的生命显现时，你们也要与祂一同在光荣中显现。}\BibleRef{哥 3:1}那么，我们这些在圣洗中，就旧人的属血肉的罪而言，已经死了并埋葬了的人，在天上的重生中与基督一同复活了的人，就该思念并实行基督的事，正如同一宗徒再次教导劝诫说：\Quote{第一个人出于地，属于土；第二个人出于天。那属于土的怎样，凡属于土的也怎样；那属于天的怎样，凡属于天的也怎样。我们既然带有了那属于土的肖像，也要带有了那属于天的肖像。}\BibleRef{格前 15:47}但我们不能带有天上的肖像，除非在我们已开始成为的那种状态中，我们显露出基督的相似。\par

15. 因为这正是改变你过去所是的，开始成为你过去所不是的，好使神圣的诞生在你内闪耀，好使虔敬的规矩回应天主圣父，好使在生活的荣耀和赞美中，天主在人内受光荣；正如祂亲自劝诫、警告并应许给光荣祂的人相应的赏报，说：\Quote{尊重我的，我必尊重他；轻视我的，他必受轻视。}\BibleRef{撒上 2:30}为了这光荣，主塑造并预备我们，天主子灌输天主圣父的肖像，在祂的福音中说：\Quote{你们曾听说：\InnerQuote{你应爱你的近人，恨你的仇人。}我却对你们说：爱你们的仇人，为迫害你们的人祈祷，好使你们成为你们在天之父的子女，因为祂使太阳上升，光照恶人，也光照善人；降雨给义人，也给不义的人。}\BibleRef{玛 5:43}如果人们有像自己的子女是一件喜乐和光荣的事——尤其当所生的后代以同样的特征响应父母时更为喜悦——那么在天主圣父内，当某人如此属神地诞生，以致在他的行为和赞美中宣告了种族的神圣崇高，这喜乐将是何等的大！何等正义的棕榈枝，何等的冠冕，成为这样的人，使主不致对你说：\Quote{我养育了子女，使他们成长，他们却背弃了我！}\BibleRef{依 1:2}但愿基督反为你鼓掌，并邀请你获得赏报，说：\Quote{来吧，我父所祝福的，承受从创世以来为你们预备的国度。}\BibleRef{玛 25:34}\par

16. 亲爱的弟兄们，必须藉这些默想来坚强心智。藉这类操练，心智必须坚定以抵御魔鬼的一切箭矢。让双手持守读经，心中思念主；让不断的祈祷永不停止；让救人的劳作恒久持续。我们要常忙于属灵的行动，这样，每当仇敌逼近，无论他多少次试图靠近，他都会发现我们的心怀向他关闭并武装起来。因为基督徒的冠冕不仅是在迫害时期领受的：和平也有其冠冕，胜利者在各种多样的争战中将仇敌打倒制服后，得戴这些冠冕。战胜私欲是节制的棕榈枝。抵抗愤怒和伤害是忍耐的冠冕。蔑视钱财是战胜贪婪。因信靠未来而忍受世间的逆境是信德的赞美。在顺境中不骄傲的人，为自己的谦卑获得荣耀；倾向于怜悯穷人的人，获得天上宝藏的报酬；不知嫉妒、一心一意、温良地爱弟兄的人，得享仁爱与和平的赏报。我们每日在这些德行中奔跑；我们无间歇地获得这些正义的棕榈枝和冠冕。\par

17. 为使你们这些曾被嫉妒和怨恨所占据的人也能到达这些赏报，要抛弃先前束缚你们的一切恶意，改过迁善，追随救恩的足迹走上永生之路。从你们心中拔出荆棘和蒺藜，使主的种子以丰饶的果实使你们富足，使神圣属灵的麦田在丰收的丰盛中充盈。吐出苦胆的毒，吐出纷争的毒液。让那被蛇的恶意所感染的心灵得以净化；让那已定居在心中的一切苦毒，藉基督的甘饴得以软化。如果你们从十字架的圣事领受饮食，就让那在\Place{玛辣}藉着预像曾用以甜化滋味的木头，在现实中为你们所用，以抚慰你们软化的心灵；你们便不必为日益增长的健康寻求药物。藉那使你们受伤的得以治愈。要爱那些你们先前憎恨的人；要善待那些你们曾以不公正的诋毁所嫉妒的人。效法善人，如果你们能追随他们；但如果不能追随，至少要与他们一同喜乐，并祝贺那些比你们更好的人。以合一的爱使自己成为他们的分享者；以爱德和兄弟情谊的纽带使自己成为他们的伙伴。当你们自己宽恕时，你们的债务也必被豁免。当你们怀着平安到天主面前时，你们的祭献必蒙悦纳。当你们思想那些神圣和正义的事时，你们的思想和行为必由高处指引，正如经上记载：\Quote{愿人的心思考正义的事，使他的脚步由主引导。}\par

18. 你们还有许多事要考虑。要思想那\Person{加音}因嫉妒杀害兄弟而不能进入的乐园。要思想那主只接纳一心一意之人的天国。要思想只有缔造和平的人才能称为天主的子女，他们因天上的出生和天主的法律合而为一，回应天主圣父和基督的肖像。要思想我们正站在天主的眼前，在天主亲自鉴察和审判下进行我们品行和生活的历程，这样，我们若以我们的行为取悦那鉴察我们的祂，若我们显示自己配得祂的恩惠和宽容，若我们那常在祂国中令祂喜悦的人，先在这世上令祂喜悦，我们便终能达致瞻仰祂的结果。\par

\SourceRecord{Translated by Robert Ernest Wallis.；Ante-Nicene Fathers；Vol. 5.；Edited by Alexander Roberts, James Donaldson, and A. Cleveland Coxe.；Buffalo, NY: Christian Literature Publishing Co.,；1886.；<http://www.newadvent.org/fathers/050710.htm>.}

% structure: p0001 source=h1 target=section reason=page-title

\section{第十一篇论文}

\Emph{致\Person{福图纳托}的殉道劝勉。}\par

前言。\par

\newline{}1. 你所渴望的，亲爱的\Person{福图纳托}，既然迫害与苦难的重担正压在我们身上，且在世界终结与完成之际，假基督的可憎时期已开始临近，我便愿从神圣的圣经中收集一些劝勉，以准备并坚固弟兄们的心志，借此激励基督的战士们参与天上的属灵竞赛。我被你那如此必要的愿望所约束，以至于在我有限的能力之所能及，并藉着神圣灵感的助佑所教导的范围内，要为即将作战的弟兄们从主的诫命中取出一些武器和防御。因为仅凭我们声音的号角唤醒天主的子民是不够的，除非我们藉着神圣的诵读来坚定信徒的信德，以及他们奉献并献身于天主的勇毅。\par

\newline{}2. 然而，有什么比藉着不断的劝勉，为那神圣地托付给我的子民和建立于天营中的军队，预备他们对抗魔鬼的投枪与武器，更合适或更完全地符合我自身的关怀与忧虑呢？因为未经战场操练的人不能成为合宜的战士；那寻求竞赛荣冠的人，除非先考虑自己能力的使用与技巧，否则也不会在赛道上得奖。我们所与之争战的是一个古老的仇敌，一个旧的敌人：自魔鬼首次攻击人以来，六千年已近乎完满。他藉着长年累月的实践，学会了各种推翻人的诱惑、诡计与陷阱。若他发现基督的士兵毫无准备，毫无技巧，不尽心谨慎警醒，他便乘其无知而诱骗，趁其疏忽而蒙蔽，欺其无经验而欺诈。但若人遵守主的诫命，勇敢地依附基督，与他对抗，则他必被击败，因为基督——那人所宣认的——是常胜的。\par

\newline{}3. 亲爱的弟兄，为使我的论述不致过于冗长，且不让过于散漫的风格使听者或读者疲倦，我作了一个摘要：这样，我先列出每人皆应知晓并铭记于心的标题，然后加上主言的节段，并以神圣教导的权威确立我所提议的，如此我似乎并非将自己的论文送给你，而是为他人提供讲论的材料；这做法将使个别的人获益更大。因为若我给人一件现成的衣服，那便是另一个人在使用我的衣服，而且为他人所制之物可能不适合他的身材体格。但现在，我已送你那救赎并赐予我们生命的羔羊的羊毛与紫红布；你领受后，将可按自己的意愿为自己缝制一件外衣，并且更因其为你私人的专属衣物而欢喜。你也将把我们送来的展示给他人，使他们也能按自己的意愿完成它；如此，那古老的赤身被遮盖后，众人皆穿上基督的衣裳，戴上天上恩宠的圣化。\par

\newline{}4. 此外，亲爱的弟兄们，我认为在一篇如此必要的劝勉——即造就殉道者——中，采取一项有益且健康的计划，即去除言语中所有迟延与缓慢，舍弃人类论述的曲折，只陈列天主所说的话，即基督用以劝勉其仆人走向殉道的训诫。这些神圣的诫命本身必须被提供，犹如战斗者的武器。让它们成为战争号角的激励，让它们成为战士的嘹亮号声。愿耳朵被它们唤醒，愿心智被它们预备，愿灵魂与身体的力量被坚固以承受一切苦难。唯有我们这些藉主的许可，为信徒施予第一次圣洗的人，也当为第二次圣洗预备每一个人；敦促并教导：这是一个在恩宠上更伟大、在权能上更高超、在尊荣上更贵重的圣洗——其中有天使施洗的圣洗——天主与祂的基督在其中欢跃的圣洗——此后无人再犯罪的圣洗——这圣洗完满我们信德的增长——当我们离开世界时，这圣洗立刻使我们与天主结合。在水洗中领受罪的赦免，在血洗中领受诸德的荣冠。这事当被拥抱、被渴望，并在我们祈求的一切恳求中求取，使我们这些天主的仆人也能成为祂的朋友。\par

% structure: p0008 source=h2 target=subsection reason=relative-heading-rank; base=h2

\subsection{一、偶像不是神，元素也不当被当作神崇拜。}

在\BibleBook{}{圣咏}第一百一十三篇中记载：\BibleQuote{外邦人的偶像不过是金银，是人手的制作。有口而不能言，有眼而不能见。有耳而不能听，鼻孔中毫无气息。愿制造它们的人同它们一样。}\newline{}在\BibleBook{}{智慧篇}中也记载：\BibleQuote{他们竟将外邦人的一切偶像都看作神，但这些偶像有眼不能看，有鼻不能呼吸，有耳不能听，有手不能摸，脚也迟于行走。因为是人制造了它们，是那借了气息的人塑造了它们；但没有人能造出一个与自己相似的神。既然人是可朽的，他竟用邪恶的手造出死物；他本人比他所崇拜的更好，因为他至少活过，而它们从未活过。}\BibleRef{智 15:15}\newline{}在\BibleBook{}{出谷纪}中也记载：\BibleQuote{你不可为自己制造偶像，亦不可制造任何形状。}\BibleRef{出 20:4}\newline{}此外，在\BibleBook{}{智慧篇}中关于元素记载：\BibleQuote{他们既不从观察受造物而认识那位工匠，却将火、风、流动的空气、星宿的光环、汹涌的水、太阳或月亮奉为神。若因它们的美丽而这样想，就让他们知道，主比它们更美多少。若他们钦佩它们的力量和运行，就让他们由此明白，那位创造这些强大之物者比它们更强大。}\BibleRef{智 13:1}\par

% structure: p0010 source=h2 target=subsection reason=relative-heading-rank; base=h2

\subsection{二、唯独应崇拜天主}

正如经上所记载：\BibleQuote{你要崇拜主，你的天主，唯独事奉祂。}\newline{}在\BibleBook{}{出谷纪}中也记载：\BibleQuote{你不可有别的神在我面前。}\BibleRef{出 20:3}\newline{}在\BibleBook{}{申命纪}中也记载：\BibleQuote{看，看，我就是祂，除我以外没有别的神。我使人死，也使人活；我击伤，也治愈；无人能从我手中救出。}\BibleRef{申 32:39}\newline{}在\BibleBook{}{默示录}中还记载：\BibleQuote{我看见另一位天使飞翔在天空中，拿着永久的福音，要传给地上的人和各国、各支派、各语言、各民族，大声说：\InnerQuote{你们要敬畏天主，将光荣归于祂，因为祂审判的时辰到了；你们要崇拜那创造天地、海和其中万物的天主。}}\BibleRef{默 14:6-7}\newline{}主也在祂的\BibleBook{}{福音}中提及第一条和第二条诫命，说：\BibleQuote{以色列，请听：主，你的天主，是唯一的天主；}\BibleRef{谷 12:29}\newline{}并说：\BibleQuote{你当全心、全灵、全力爱上主你的天主。这是第一条；第二条与此相似：你当爱你的近人如你自己。全部法律和先知都系于这两条诫命。}\BibleRef{玛 22:37-40}\newline{}又说：\BibleQuote{永生就是：认识祢，唯一的真天主，和祢所派遣的耶稣基督。}\BibleRef{若 17:3}\par

% structure: p0012 source=h2 target=subsection reason=relative-heading-rank; base=h2

\subsection{三、天主对向偶像献祭的人有何威胁？}

在\BibleBook{}{出谷纪}中：\BibleQuote{凡向别的神献祭，而不只向主献祭的，应被铲除。}\BibleRef{出 22:20}\newline{}在\BibleBook{}{申命纪}中也记载：\BibleQuote{他们向邪魔献祭，而不是向天主。}\BibleRef{申 32:17}\newline{}在\BibleBook{}{依撒意亚}中也记载：\BibleQuote{他们崇拜自己手指所造之物；卑微的人屈膝，尊贵的人降卑；我必不宽恕他们。}\BibleRef{依 2:8-9}\newline{}又说：\BibleQuote{你们向他们奠了酒，向他们献了祭。因此，我岂能不愤怒？——上主说。}\BibleRef{依 57:6}\newline{}在\BibleBook{}{耶肋米亚}中也记载：\BibleQuote{不要随从别的神，事奉他们，不要崇拜他们，不要以你们手中的作品惹我发怒，免得我毁灭你们。}\BibleRef{耶 7:6}\newline{}在\BibleBook{}{默示录}中也记载：\BibleQuote{若有人崇拜那兽和它的像，并在额上或手上接受了它的印记，他也要喝天主愤怒的酒，即纯酒，调和在他愤怒的杯中；他要在圣天使和羔羊面前，受火与硫磺的刑罚；他们受刑的烟上升，直到永远；那些崇拜兽和兽像的，昼夜不得安宁。}\BibleRef{默 14:9-11}\par

% structure: p0014 source=h2 target=subsection reason=relative-heading-rank; base=h2

\subsection{四、天主不轻易宽恕崇拜偶像者}

梅瑟在\BibleBook{}{出谷纪}中为民众祈祷，却未获准许，他说：\BibleQuote{我恳求祢，上主，这百姓犯了重罪，为自己制造了金神。现在，如果祢赦免他们的罪，就请赦免；不然，就请从祢所写的册上涂抹我的名字。上主对梅瑟说：谁得罪了我，我必从我的册上涂抹谁。}\BibleRef{出 32:31-33}\newline{}此外，当耶肋米亚为民众求情时，上主对他说：\BibleQuote{你不要为这百姓祈祷，不要为他们呼求祷告；因为在他们遭难时呼求我的时刻，我必不听。}\BibleRef{耶 7:16}\newline{}厄则克耳也宣告了天主对得罪祂之人的同样愤怒，说：\BibleQuote{上主的话传给我说：人子，若某地犯罪得罪我，行事不忠，我必伸手攻击它，断绝它的粮源，使饥荒临到，从它那里灭尽人和牲畜。即使诺厄、达尼尔和约伯这三人在其中，他们也不能救自己的子女，只能救自己。}\BibleRef{则 14:12-14}\newline{}同样在\BibleBook{}{撒慕尔纪上}：\BibleQuote{若人得罪人，尚可为他恳求上主；但若人得罪天主，谁能为他求情呢？}\BibleRef{撒上 2:25}\par

% structure: p0016 source=h2 target=subsection reason=relative-heading-rank; base=h2

\subsection{五、天主对崇拜偶像如此愤怒，甚至命令将那些劝说别人献祭和事奉偶像的人处死}

在\WorkTitle{申命纪}中：\Quote{如果你的兄弟，或你的儿子，或你的女儿，或你怀中的妻子，或你性命相投的朋友，暗中怂恿你说：我们去事奉别的神，别邦的神，你不可依从他，也不可听从他，你的眼不可顾惜他，不可袒护他，务要举报他。你的手要先加在他身上，把他处死，然后众人的手也加在他身上；他们要用石头砸死他，因为他想要引你离弃上主你的天主。} \BibleRef{申 13:6} 主又发言说，即使全城同意拜偶像，也不得饶恕那城：\Quote{如果你在你们的天主上主赐给你居住的一座城里，听说：我们去事奉别的神，你们素不认识的神，你就该用刀剑杀死那城里的所有居民，将城完全烧毁，使它永远成为废墟，不得再重建；这样，上主就会转消他的烈怒，向你施恩，怜悯你，使你人数增多，只要你听从上主你的天主的声音，遵守他的诫命。} \BibleRef{申 13:12} 想起这条诫命及其效力，\Person{玛塔提雅}杀了那个走近祭坛要献祭的人。但如果基督来临之前，这些关于敬拜天主和蔑视偶像的诫命已被遵守，那么自基督来临后，它们更应被遵守；因为祂来时，不仅用言语劝勉我们，也用行动，而且在受尽一切冤屈和凌辱之后，还受苦并被钉十字架，为的是以祂的榜样教导我们受苦和死亡，使人在为祂受苦时毫无借口，因为祂已为我们受苦；而且，既然祂为别人的罪受苦，每个人更应为自己本身的罪受苦。因此，主在\WorkTitle{福音}中警告说：\Quote{凡在人前承认我的，我在我天上的父前也必承认他；但谁在人前否认我的，我在我天上的父前也必否认他。} \BibleRef{玛 10:32} 宗徒保禄也说：\Quote{如果我们与他同死，也必与他同生；如果我们受苦，也必与他一同为王；如果我们否认他，他也必否认我们。} \BibleRef{弟后 2:11} 若望也说：\Quote{凡否认子的，也没有父；凡明认子的，也有父也有子。} \BibleRef{若一 2:23} 因此，主劝勉并加强我们轻视死亡，说：\Quote{你们不要害怕那杀害肉身而不能杀害灵魂的；但更要害怕那能使灵魂和肉身一同被投入地狱中的。} \BibleRef{玛 10:28} 又说：\Quote{爱惜自己生命的，必丧失生命；在现世憎恨自己生命的，必保全生命直到永生。} \BibleRef{若 12:25}\par

% structure: p0018 source=h2 target=subsection reason=relative-heading-rank; base=h2

\subsection{六、既然被基督的血所救赎和赐予生命，我们不应将任何事物置于基督之上。}

在\WorkTitle{福音}中，主发言说：\Quote{谁爱父亲或母亲超过我，不配是我的人；谁爱儿子或女儿超过我，不配是我的人；谁不背起自己的十字架跟随我，也不配是我的人。} \BibleRef{玛 10:37} 同样，在\WorkTitle{申命纪}中也记载：\Quote{他们对自己的父亲和母亲说：我不认识你们，也不承认自己的儿女；这些人遵守了你的诫命，持守了你的盟约。} \BibleRef{申 33:9} 此外，宗徒保禄说：\Quote{谁能使我们与基督的爱隔绝呢？是困苦吗？是窘迫吗？是迫害吗？是饥饿吗？是赤身露体吗？是危险吗？是刀剑吗？正如所记载的：为了你，我们终日被杀死，被视为待宰的羊。然而，靠着那爱我们的主，我们在这一切事上大获全胜。} \BibleRef{罗 8:35} 又说：\Quote{你们不是自己的，因为你们是被重价买来的。你们要在身体上光荣天主，并承载天主。} \BibleRef{格前 6:20} 又说：\Quote{基督替众人死了，为的是使活着的人不再为自己生活，而是为那替他们死而复活的主生活。} \BibleRef{格后 5:15}\par

% structure: p0020 source=h2 target=subsection reason=relative-heading-rank; base=h2

\subsection{七、那些从魔鬼的爪牙中被夺回、从这世界的罗网中被救出的人，不应再回到世界，以免失去退避的益处。}

在\WorkTitle{出谷纪}中，犹太民族——作为我们的预像和影子——当他们有天主作他们的保护者和复仇者，逃脱了法郎和埃及——即魔鬼和世界——最残酷的奴役时，却对天主不忠不感恩，抱怨梅瑟，回顾沙漠和劳苦的艰难；而不明白自由和救恩的神圣恩惠，他们寻求回到埃及——即他们被领出的世界——的奴役中，而他们本应信赖并相信天主，因为那位将祂的子民从魔鬼和世界手中救出的，也保护已得救的子民。\Quote{你为什么这样对待我们，将我们从埃及领出来呢？我们服侍埃及人比死在这旷野里还好。}他们说，\Quote{梅瑟对百姓说：你们要信赖，站稳，观看上主今天为你们施行的救恩。上主必亲自为你们作战，你们要静默。} \BibleRef{出 14:11} 主在\WorkTitle{福音}中以此劝诫我们，教导我们不应再回到我们已经弃绝并逃离的魔鬼和世界，说：\Quote{手扶着犁向后看的，不适于天主的国。} \BibleRef{路 9:62} 又说：\Quote{在田野里的人，不要返回。要记得罗特的妻子。} 为免任何人因贪恋财富或亲人的吸引而延缓跟随基督，祂又补充说：\Quote{谁不舍弃一切所有，不能做我的门徒。} \BibleRef{路 14:33}\par

% structure: p0022 source=h2 target=subsection reason=relative-heading-rank; base=h2

\subsection{八、应当努力前进，持守信德与德行，并完成天上的属神恩宠，好能获得棕榈枝与花冠。}

在《编年纪》里：\Quote{上主与你们同在，只要你们与祂同在；但如果你们离弃祂，祂也必离弃你们。}\BibleRef{编下 15:2} 在《厄则克耳》里也说：\Quote{义人的义德在他犯罪的那一天不能拯救他。}\BibleRef{则 33:12} 此外，主在福音中发言说：\Quote{唯独坚持到底的，才可得救。}\BibleRef{玛 10:22} 又说：\Quote{如果你们存留在我的话内，你们就真是我的门徒；你们必会认识真理，而真理必会使你们获得自由。}\BibleRef{若 8:31} 此外，祂预先警告我们，应当常常准备，坚定地装备武装，并补充说：\Quote{你们要把腰束好，把灯点着，你们自己像等候自己主人从婚宴回来的人一样，好叫他来到敲门时，立刻给他开门。主人来到时，发现那些仆人醒着，他们才是有福的。}\BibleRef{路 12:35} 蒙福的宗徒保禄也鼓励我们，为使我们的信德前进、增长，并达到最高点，说：\Quote{你们不知道在运动场上赛跑的，大家都跑，但只有一个得到奖赏吗？你们也该这样跑，好能夺得奖赏。他们这样跑，是为了获得那可朽坏的花冠；而我们是为获得那不可朽坏的花冠。}\BibleRef{格前 9:24} 又说：\Quote{没有一个为天主作战的人，会将自己纠缠于世事的牵挂里，好使自己能取悦于他所中悦的人。此外，如果有人参加竞赛，除非按规矩竞赛，他是得不到花冠的。}\BibleRef{弟后 2:4} 又说：\Quote{弟兄们，我因天主的仁慈恳求你们，将你们的身体献上，当作生活、圣洁、悦乐天主的祭品。不可与此世同化，反而应以更新的心思变化自己，为使你们能辨别什么是天主的旨意，什么是善事、悦乐祂的事和完美的事。}\BibleRef{罗 12:1} 又说：\Quote{我们是天主的子女：如果是子女，便是继承人，是天主的继承人，与基督同为继承人，只要我们与祂一同受苦，也必与祂一同受光荣。}\BibleRef{罗 8:16} 在《默示录》中，同样神圣宣讲的劝勉说话：\Quote{你要持守你所有的，免得别人拿走你的花冠。}\BibleRef{默 3:11} 这个恒心与坚持的榜样在《出谷纪》中指明，当时梅瑟为了击败阿玛肋克——他预表魔鬼——举起双手形成十字架的标记与神圣记号，除非他一直坚持举起双手，坚定地持续这一记号，他不能战胜敌人。经上说：\Quote{当梅瑟举起手时，以色列就占优势；当他放下手时，阿玛肋克就增强。他们拿了一块石头放在他下面，他就坐在上面。亚郎和胡尔一边一个扶着他的手，他的手就稳定直到日落。耶稣击败了阿玛肋克和他的百姓。上主对梅瑟说：你要将这事记录下来，作为纪念，传到若苏厄耳中；因为我要完全消灭天下对阿玛肋克的纪念。}\BibleRef{出 17:11}\par

% structure: p0024 source=h2 target=subsection reason=relative-heading-rank; base=h2

\subsection{九、苦难与迫害是为了考验我们而兴起。}

在《申命纪》里：\Quote{上主你们的天主考验你们，为要知道你们是否全心、全灵、全力爱上主你们的天主。}\BibleRef{申 13:3} 撒罗满又说：\Quote{火窑考验陶匠的器皿，而患难的考验考验义人。}\BibleRef{德 27:5} 保禄也作证说：\Quote{我们因希望分享天主的荣耀而欢跃。不但如此，我们在患难中也欢跃，因为我们知道患难生忍耐，忍耐生老练，老练生希望，希望不会使人蒙羞，因为天主的爱借着所赐予我们的圣神已倾注在我们心中。}\BibleRef{罗 5:2} 伯多禄在他的书信中规定说：\Quote{亲爱的，你们不要因为在你们中间发生的烈火般的考验而奇怪，好像是遭遇了什么新奇的事；反而要喜欢，因为在你们与基督一同受苦的事上你们有分，这样，在祂光荣显现时，你们也可以欢喜踊跃。如果你们为了基督的名字而受辱骂，便是有福的，因为光荣和权能的神，就是天主的神，安息在你们身上。这神在他们看来是亵渎的，但在我们看来却是光荣的。}\BibleRef{伯前 4:12}\par

% structure: p0026 source=h2 target=subsection reason=relative-heading-rank; base=h2

\subsection{十、我们不应惧怕迫害的伤害与惩罚，因为主更伟大，祂保护我们，胜于魔鬼攻击我们。}

\Person{若望}在他的书信中证明此事，说：\BibleQuote{那在你们内的，比那在世界上的更大。} \BibleRef{若一 4:4} 也在第一一七首圣咏中：\BibleQuote{我不惧怕人能对我做什么；主是我的助佑。} 又说：\BibleQuote{这些靠车，那些靠马；我们却要以我们天主上主的名夸耀。他们已被束缚，已跌倒；但我们已起来，立得正直。} 圣神更强烈地教训并显示，魔鬼的军旅不可怕，而且如果敌人向我们宣战，我们的希望反而在于那战争本身；并且通过那争斗，义人获得天主居所和永久救恩的赏报——祂在第二十六首圣咏中规定并说：\BibleQuote{虽有大军齐集攻击我，我的心决不害怕；虽有战争兴起攻击我，我仍以此寄望。我有一事求问上主，我要恳求此事：使我一生岁月，常居住在上主的殿宇中。} 也在出谷纪中，圣经宣布，我们反而因苦难而增多和加增，说：\BibleQuote{他们越压迫他们，他们就越增多，越强壮。} \BibleRef{出 1:12} 在默示录中，天主的保护应许给我们的苦难。\BibleQuote{不要怕这些事，} 它说，\BibleQuote{就是你们将要受的苦。} \BibleRef{默 2:10} 也没有别人应许我们安全和保护，只有那也曾藉先知依撒意亚说话的那一位，说：\BibleQuote{不要害怕！因为我救赎了你，曾以你的名字召叫了你：你是我的。你若渡过水，我必与你同在，河川不淹没你。你若走过火，你不会被烧，火焰不焚毁你；因为我是上主你的天主，以色列的圣者，使你安全的那一位。} \BibleRef{依 43:1} 祂也在福音中应许，天主的助佑在迫害中必不缺少天主的仆人们，说：\BibleQuote{但当人把你们交出时，不要思虑怎样说，或说什么。因为在那个时候，必赐给你们当说的话。因为不是你们说话，而是你们父的圣神在你们内说话。} \BibleRef{玛 10:19} 又说：\BibleQuote{你们心中定意，不要预先思虑如何答辩；因为我必赐给你们口才和智慧，是你们的一切仇敌所不能抵抗的。} \BibleRef{路 21:14} 如在出谷纪中，天主向迟延和战栗不敢去见百姓的梅瑟说话，说：\BibleQuote{谁给了人口？谁使人成为哑巴？谁使人成为聋子？谁使人成为明眼人和瞎子？不是我，主天主吗？现在你去，我要开启你的口，教导你当说的话。} \BibleRef{出 6:11} 对天主来说，开启一个奉献于自己的人的嘴，并激励他的宣信者在说话时具有恒心和信心，并不困难；因为在户籍纪中，祂甚至使母驴说话反对先知巴郎。因此，在迫害中，不要想魔鬼带来什么危险，而要思考天主提供什么帮助；也不要让人的恶行压倒心灵，而让天主的保护坚固信德；因为每个人，按照主的许诺和自己信德的功绩，从天主的助佑中领受多少，就如他所想领受的那样。全能者没有什么不能赐予的，除非领受者衰微的信德缺乏和退却。\par

% structure: p0028 source=h2 target=subsection reason=relative-heading-rank; base=h2

\subsection{十一、预先预言了世界将憎恨我们，并兴起迫害攻击我们，且基督徒所遭遇的并非新奇之事，因为从世界开始，善人已受苦，义人被不义之人压迫和杀害。}

主在福音中预先警告并预言说：\BibleQuote{如果世界恨你们，你们要知道，它先恨了我。如果你们属于世界，世界必爱属于自己的人；但因为你们不属于世界，而我从世界中拣选了你们，所以世界恨你们。你们要记得我对你们所说的话：仆人不能大过主人。如果他们迫害了我，也要迫害你们。}\BibleRef{若 15:18} 又说：\BibleQuote{时候要到，凡杀害你们的人，都以为是服侍天主；他们这样做，是因为不认识父，也不认识我。但我要告诉你们这些事，使你们到了时候，可以想起我告诉过你们了。}\BibleRef{若 16:2} 又说：\BibleQuote{我实实在在告诉你们：你们将要痛哭哀号，世界却要欢乐；你们将要忧愁，但你们的忧愁将变为喜乐。}\BibleRef{若 16:20} 又说：\BibleQuote{我给你们讲了这些事，为使你们在我内获得平安；但在世界上你们将遭受苦难；然而你们放心，我已经战胜了世界。}\BibleRef{若 16:33} 当祂被门徒问及有关祂来临和世界终结的征兆时，祂回答説：\BibleQuote{你们要谨慎，免得有人欺骗你们；因为将有许多人冒我的名来，说：我是基督，并要欺骗许多人。你们要开始听见战争和战争的风声；小心不要惊慌，因为这些事必须发生，但结局还不即刻来到。因为民族要起来攻击民族，国家要攻击国家，到处将有饥荒、地震和瘟疫。这一切都是灾难的开始。那时人们将把你们交出，受迫害，并杀害你们；你们为我的名将被万民憎恨。那时有许多人要跌倒，彼此出卖，彼此仇恨。将有许多假先知兴起，迷惑许多人；由于罪恶的增多，许多人的爱情将冷淡。但是，那坚持到底的人，必然得救。这天国的福音必传遍全世界，对万民作证，然后结局才会来到。所以，当你们看见达尼尔先知所说的\InnerQuote{可憎的荒凉之物}站在圣地（读者应明白），那时，在犹太的，要逃到山上；在屋顶上的，不要下来拿家里的东西；在田里的，也不要回去取衣服。但在那些日子，怀孕的和哺乳的，有祸了！你们要祈祷，免得你们的逃遁在冬天或在安息日；因为那时必有大灾难，是从世界开始至今未曾有过的，将来也不会再有。如果不减少那些日子，凡有血肉的都不会得救；但为了选民的缘故，那些日子必会减短。那时，如果有人对你们说：\InnerQuote{看，基督在这里}或\InnerQuote{在那里}，你们不要相信。因为将有许多假基督和假先知兴起，行大奇事和异兆，以至如果可能，连选民也要被迷惑。但你们要谨慎；看，我预先都告诉了你们。所以，如果他们对你们说：\InnerQuote{看，他在旷野}，你们不要去；\InnerQuote{看，他在内室}，也不要信。因为正如闪电从东方发出，直照到西方，人子的来临也要这样。无论哪里有一具尸体，鹰鸟就聚集在那里。但在那些日子的灾难以后，太阳将要昏暗，月亮不再发光，星辰要从天上坠落，天上的万有也要动摇。那时，人子的记号要出现在天上；地上所有的民族都要哀号，并要看见人子带着大能力和荣耀乘着天上的云彩降临。祂必派遣祂的天使，用大号角，从四方，从天的一边到另一边，聚集祂的选民。}\BibleRef{玛 24:4} 这些事对基督徒来说，并非新奇或突然临到；因为善良正义、在无罪的律法和真宗教的敬畏中献身于天主的人，总是借着苦难、冤屈、严酷且多样的惩罚、狭窄道路的艰辛而前进。因此，在世界之初，义人\Person{亚伯尔}最先被自己的兄弟杀害；\Person{雅各伯}被迫流亡，\Person{若瑟}被卖，\Person{撒乌耳}王迫害仁慈的\Person{达味}；\Person{阿哈布}王企图压迫坚定而勇敢地维护天主威严的\Person{厄里亚}。\Person{匝加利亚}司铎在圣殿与祭台之间被杀，使他本人成为祭献，就在他素常向天主献祭的地方。事实上，如此多义人的殉道常常被纪念；如此多信德和德行的榜样被展示给后世。三位青年：\Person{阿纳尼雅}、\Person{阿匝黎雅}和\Person{米沙耳}，年龄相同，爱心一致，信德坚定，德行恒久，胜过迫使他们的火焰和刑罚，宣布他们只服从天主，只认识祂，只敬拜祂，说：\BibleQuote{拿步高王，我们不需要在这件事上回答你。因为我们所事奉的天主，能拯救我们脱离烈火窑，并从你手中拯救我们。即使不如此，你应知道，我们绝不事奉你的神，也不敬拜你所立的金像。}\BibleRef{达 3:16} 忠信于天主并充满圣神的\Person{达尼尔}高呼说：\Quote{我只敬拜我的主天主，即创造天地的主。} \Person{多俾亚}虽然在王权和暴政的奴役下，却在心灵和精神上自由，向天主保持宣认，并崇高宣告天主的威能和威严，说：\BibleQuote{在我被掳的地方，我承认祂，并向罪恶的民族显示祂的权威。}\BibleRef{多 13:6} 在\BibleBook{玛加伯}{书}中，我们读到七位兄弟，他们在出生和德行上命运相同，以完美的圆满所象征的数目七，完成了殉道。七位兄弟就这样在殉道中结合。正如神圣的秩序中第一个七日包含七千年，正如七神和七位天使侍立在天主面前，出出入入，又如会幕中的七盏灯台，\BibleBook{}{默示录}中的七盏金灯台，以及\Person{撒罗满}那里智慧建造房屋的七根柱子；同样，这里七位兄弟的数字，以其数量涵盖了七教会，正如在\BibleBook{列王}{纪上}中我们读到不生育的生了七子。\BibleBook{依撒意亚}{先知书}中七个妇女抓住一个男人，要求用他的名字称呼她们。\Person{保禄}宗徒，提到这合法而确定的数字，写信给七教会。在\BibleBook{}{默示录}中，主向七教会及其天使发出神圣与天上的训令，这个数字现在在这七位兄弟身上实现，为使合法圆满得以完成。与这七个孩子明显相关的还有他们的根源和根苗——母亲，她后来生出了七教会，她自己先由主的声音建立在磐石上，并且是唯一的。同样并非无意义：在受苦中只有母亲与她的孩子们在一起。因为作证自己是天主子在受苦中的殉道者，现在不再算为除天主外的任何父亲的孩子，正如主在福音中所教导的，说：\BibleQuote{不要称地上的人为你们的父；因为你们的父只有一位，就是天上的父。}\BibleRef{玛 23:9}\par

但他们所宣布的宣认言词是何等的光荣！他们提供了何等辉煌、何等伟大的信德证明！他们的敌人\Person{安提约古}王——是的，在\Person{安提约古}身上预表了假基督——试图用猪肉的污染来玷污那些在宣认精神上光荣而不可战胜的殉道者的口；当他用鞭子严厉地鞭打他们，却毫无成效时，他命令烧热铁板，当铁板被烧得炽热后，他命令将那个首先发言、并以恒心和信德更激怒国王的人带上来烤炙，先拔掉并割断他那宣认了天主的舌头；这事对殉道者发生了更光荣的结果。因为那宣认了天主之名的舌头，应当首先归于天主。然后，对第二个，他设计了更剧烈的痛苦，在折磨其他肢体之前，他连头发一起剥下了他头上的皮，无疑出于仇恨的目的。因为基督是人的头，天主又是基督的头，所以那在殉道者身上撕裂头的人，是在那头中迫害天主和基督。但他信赖自己的殉道，并因天主的报应而预许自己复活的赏报，就大声喊道说：\BibleQuote{你们确实无能地将我们从这现世生命毁灭；但世界的君王必将我们这些为祂法律而死的人复活，进入永生。}第三个被挑战时，迅速伸出舌头；因为他从兄弟那里学会了藐视割舌的刑罚。而且，他坚定地伸出手来被砍掉，对这种惩罚方式极为高兴，因为他得以伸出手来模仿他主受难的形状。第四个也以同样的德行，藐视酷刑，并以属天的声音答言以制止国王，大声喊道说：\BibleQuote{那些被人置于死地的人，更好是期待从天主而来的希望，将被祂再度复活而获永生。至于你，决没有复生得生命的可能。}\BibleRef{加下 7:14}第五个，除了以信德的力量践踏国王的酷刑及其严厉而多样的折磨外，还被神性的圣神激励，具有对未来的预见和知识，向国王预告了天主的义怒和即将迅速来临的报复。\BibleQuote{你有权力，}他说，\BibleQuote{在人中，虽然你是可朽的，你任意而行。但不要以为我们的民族已被天主遗弃。你且等着，观看祂的大能，祂将怎样折磨你和你的后裔。}\BibleRef{加下 7:16}这对殉道者是何等的安慰！在他所受的痛苦中，何等实在的慰藉，不是考虑自己的折磨，而是预言折磨者的刑罚！但在第六个身上，不仅要陈明他的勇敢，还有他的谦卑；即殉道者不为自己要求什么，甚至不以骄傲的言辞谈论自己宣认的光荣，反而将遭受国王迫害归因于自己的罪，而将后来得到报仇归于天主。他教导说，殉道者是谦逊的，他们对报仇满怀信心，但在受苦中不夸耀什么。\Quote{不要，}他说，\Quote{无谓地犯错；因为我们是为了自己而受这些苦，因为我们得罪了我们的天主。但不要以为你敢于对抗天主，会不受惩罚。}还有那位可敬佩的母亲，她既未被性别的软弱所压倒，也未因多次丧子而动摇，她以喜乐的心情看着自己死去的孩子，并不把这些视为她亲爱的子女的惩罚，而是光荣，她以眼目的德行给了天主如同她的子女以肢体的折磨和痛苦所给的同样伟大的见证；当六个儿子受刑被杀后，只剩下一个兄弟，国王向他许以财富、权力和许多事物，以便通过制服一人而平息其残忍和凶暴，并请求母亲劝她儿子与她一同屈服；她劝说了，但却是以殉道者之母应有的方式——是那铭记法律和天主的人应有的方式——是那以勇敢而非娇惯的方式爱儿子的人应有的方式。因为她劝说的是要儿子宣认天主。她劝说的是兄弟不要在赞美和光荣的联盟中与他的兄弟分离；那时她只将自己视为七个儿子的母亲，如果她有幸生育了七个儿子，不是为世界，而是为天主。因此她武装他，坚强他，以更蒙福的生产怀抱着她的儿子，说：\Quote{儿子，可怜我吧，我在腹中怀了你十个月，用乳汁喂养了你三年，并养育你到这般年纪；我恳求你，儿子，你仰观天和地，并思量其中所有的一切，你就明白这一切以及人类都是天主从无中造出来的。因此，儿子，不要害怕那刽子手；但要做配得你兄弟的人，接受死亡，好使我在同样的仁慈中与你的兄弟一同接纳你。}母亲在鼓励德行方面受到极大的赞扬，但在敬畏天主和信德真理方面更受赞扬，因为她没有从那六位殉道者的光荣中为自己或儿子指望什么，也不相信兄弟们的祈祷能够为一位背弃者带来救恩，反而劝说他成为他们痛苦中的一份子，使他在审判之日能与他的兄弟一同被找到。之后，这位母亲也与她的孩子们一同死去；因为除此以外没有别的更合适：那生育并造就了殉道者的人，应与他们一起加入光荣的团体，她本人应跟随她先遣到天主那里去的子女。此外，为避免有人当有机会以某种证书或类似方式欺骗时，会接受欺骗者的邪恶角色，我们也不应缄默不言，关于\Person{厄肋阿匝尔}：当国王的仆从给他机会，让他接受可食用的肉，以便欺骗国王，假装吃了那些从祭品和不洁肉食中强制给他的东西时，他不同意这种欺骗，说这与他的年龄和尊贵都不相称，去假装那种会使他人跌倒并引入错误的事情；如果他们以为九十岁的\Person{厄肋阿匝尔}已经离弃并背叛了天主的法律，转而效仿外邦人的习俗；并且，赢得短暂的现世生命却招致触怒天主的永罚，这并不值得。他受尽折磨，最终濒临死亡，在鞭打和酷刑中死去时，呻吟着说：\Quote{主啊，祢具有圣智，显然即使我得以免死，我饱受鞭笞，身体遭受最剧烈的痛苦；但我的心灵，因敬畏祢，甘愿忍受这一切。}无疑，他的信德是真诚的，他的德行是健全而极其纯洁的，不看重\Person{安提约古}王，而看重审判者天主，并且知道，如果他欺骗人，于他的得救毫无益处，因为那审判我们良心的、唯一可畏的天主，是绝不能被嘲弄或欺骗的。因此，如果我们自己也为天主分别出来并奉献而生活——如果我们追随义人古老而神圣的足迹，就让我们经历同样的痛苦考验，同样的苦难证言，并认为我们这个时代的荣耀更大，因为古代的榜样或许可以数算，但后来，当德行和信德极为丰盛时，基督徒的殉道者却不可胜数，正如\BibleBook{}{默示录}所见证和说的：\BibleQuote{这些事以后，我看见一大伙群众，没有人能够数清，是来自各邦国、各支派、各民族、各语言的，站在宝座和羔羊面前，身穿白衣，手拿棕榈枝，大声呼喊说：救恩归于那坐在宝座上的我们的天主，并归于羔羊！有一位长老回答，对我说：这些穿白衣的是谁？他们从哪里来？我对他说：我主，祢知道。他便对我说：这些人是从大灾难中出来的，曾洗净自己的衣裳，并在羔羊的血中使它们洁白。因此，他们站在天主的宝座前，并在祂的圣殿中日夜事奉祂。}但如果基督徒殉道者的集合被证明是如此庞大，那么当一个人看到殉道者的群众不可胜数时，就不要认为成为殉道者是一件艰难或困难的事。\par

% structure: p0031 source=h2 target=subsection reason=relative-heading-rank; base=h2

\subsection{十二、在此世当前的争战与苦难之后，义人与殉道者存有何种希望与赏报}

圣神藉\Person{撒罗满}显示并预言说：\BibleQuote{他们在人眼前虽受了苦楚，但他们的希望却充满不死。他们在小事上受了磨难，却要在大事上蒙受恩待，因为天主考验了他们，发现他们配作祂自己的子民。祂试验了他们，像在炉中炼黄金；又悦纳了他们，有如全燔祭；到了时候，他们必蒙眷顾。他们必闪烁发光，有如火花在芦苇丛中奔跑。他们要审判万民，统辖各民族；他们的主要永远为王。} 同样，在此也描述了我们的报复，并宣告了那些迫害和骚扰我们的人的悔改。他说：\BibleQuote{那时，义人必以极大的信心，站在那些曾折磨他们、夺去他们劳苦果实的人面前；那些人见了，必满怀恐惧，战栗不安；因这出乎意料的救恩而惊异，遂在悲痛中叹息，彼此懊悔地说：这些就是我们曾讥笑、当作耻辱比喻的人！我们愚昧地以为他们的生活是疯狂，他们的结局是耻辱。他们怎么被列在天主的子女中，他们的份与圣者同在？可见我们迷失了真理之路，正义之光未曾照耀我们，太阳也未曾为我们升起。我们在不义和灭亡的路上疲惫，走过荒芜的旷野，却不认识主的道路。骄傲为我们有何益处？财富的炫耀又带给我们什么？这一切都像影子般消逝了。} 同样，在第一一五篇圣咏中，也显示了受苦的代价和赏报：\BibleQuote{宝贵的，它说，在主的眼中，祂的圣者的死是宝贵的。} 第一二五篇圣咏也表达了争斗的悲伤和报偿的喜悦：\BibleQuote{流泪播种的，必欢乐收割。他们边走边哭，撒播种子；但回来时，必欢呼着，禾捆满怀。} 此外，在第一一八篇圣咏中：\BibleQuote{品行完备、遵行主法律的人，是有福的。那些寻求祂的证言、全心寻找祂的人，是有福的。} 再者，主在\WorkTitle{福音}中，祂自己是迫害我们的伸冤者和受苦的赏报者，说：\BibleQuote{为义而受迫害的人是有福的，因为天国是他们的。} 又说：\BibleQuote{当人为了人子的缘故恨你们、拒绝你们、驱逐你们，并辱骂你们的名字为邪恶时，你们是有福的。在那一天，你们要欢喜踊跃，因为看，你们在天上的赏报是大的。} 又说：\BibleQuote{凡为我的缘故丧失生命的，必保全生命。} 神圣许诺的赏报不仅临到那些受辱骂和被杀害的人；即使苦难本身对信友有所欠缺，但若他们的信德保持完整、不可战胜，且基督徒放弃并藐视所有财产，显示自己是跟随基督的，那么他也要被基督列在殉道者中受尊荣，正如主亲自应许和说的：\BibleQuote{凡为了天国的缘故，离开了房屋、田地、父母、兄弟、妻子、儿女的，没有不在今世得到百倍，并在来世得到永生的。} 在\WorkTitle{默示录}中，祂也说了同样的话：\BibleQuote{我看见，他说，那些为了\Person{耶稣}的名和天主的话而被杀者的灵魂。} 当祂首先提及被杀者后，又补充说：\BibleQuote{凡没有朝拜那兽的像，也没有在额上或手上接受祂的印记的；} 所有这些人都被祂一起放在同一处，祂说：\BibleQuote{他们都活了，并与基督一同为王。} 祂说，所有的人都与基督一同生活、为王，不仅是那些被杀的人，而且凡是站在信仰和敬畏天主的坚定中，没有朝拜兽的像，也没有同意祂那致命而亵渎的谕令的人。\par

% structure: p0033 source=h2 target=subsection reason=relative-heading-rank; base=h2

\subsection{十三、我们受苦的赏报远超在此世所受的苦难本身}

蒙福的宗徒保禄证明了这一点；他因天主的俯就，被提到第三层天和乐园里，见证了他听见不可言传的话语，他夸耀自己凭眼见之信德看见了耶稣基督，他以良知的更大真理承认他所学习并看见的，并且说：\BibleQuote{现今时期的苦难，与那将要显示在我们身上的未来光荣，是不能相比的。}那么，谁不竭尽全力努力达到这样的光荣，成为天主的挚友，立即与基督一同欢乐，在历经世间的折磨与刑罚之后领受天主的赏报呢？如果对于现世的士兵而言，在打败敌人后凯旋归国是光荣的，那么当战胜魔鬼后凯旋回到乐园，并且把胜利的旌旗带回那从前\Person{亚当}因犯罪而被驱逐的地方，将从前使他堕落的那个打倒之后，那是何等更卓越、更大的光荣啊！向天主献上最中悦的礼物——纯全无损的信德、坚定不移的心灵德性、辉煌的虔敬赞美；当祂来从仇敌身上报仇时陪伴祂，当祂坐着审判时站在祂身边，与基督同为承继者，与天使同等；与圣祖、宗徒、先知一同因拥有天国而喜乐！这样的思想，什么迫害能征服？什么酷刑能战胜？那基于虔诚默想而建立的勇敢而坚定的心志，坚忍不拔；当灵魂因对未来之事的确定而坚实的信德得到强化时，它面对魔鬼的一切恐怖和世界的威胁都毫不动摇。在迫害中，大地关闭，但天堂敞开；假基督威胁，但基督保护；死亡来临，但永生随之而来；世界从被杀者那里被夺走，但乐园为他恢复而展现；时间性的生命熄灭，但永恒的生命实现。这是何等的尊严，何等的保障：欣然离开此地，在苦难和磨难中光荣地离世；瞬间闭上用于观看人和世界的眼睛，同时睁开它们以观看天主和基督！这样蒙福的离世，其速度是何等之大！你要突然地被从地上提起，被安置在天国里。我们理应在思想和默想中拥抱这些事，昼夜思虑这些事。如果迫害临到这样的天主战士身上，他那随时准备战斗的德能将不会被战胜。或者，如果他的蒙召先来到，他的信德也不会没有赏报，因为它已为殉道作了准备；天主的审判会毫不迟延地给予赏报。在迫害中，战斗得冠冕；在和平中，良心的纯洁得冠冕。\par

\SourceRecord{Translated by Robert Ernest Wallis.；Ante-Nicene Fathers；Vol. 5.；Edited by Alexander Roberts, James Donaldson, and A. Cleveland Coxe.；Buffalo, NY: Christian Literature Publishing Co.,；1886.；<http://www.newadvent.org/fathers/050711.htm>.}

% structure: p0001 source=h1 target=section reason=page-title

\section{论集十二，第一书}

\Emph{三卷\WorkTitle{驳犹太人见证集}}\par

\Person{西彼廉}致其子\Person{基林诺}，安好。我亲爱的儿子，我必须顺从你的属灵渴望，你以最恳切的祈求请求那些神圣的教导，即主藉着圣经屈尊教导和训诫我们的，使我们脱离错误的黑暗，被祂纯净闪耀的光明照亮，藉着救恩的圣事持守生命之道。诚然，正如你所求，这篇论述已由我安排；并以精简的纲要编排，使我不至于将过于散漫丰富的内容散佚，而是凭着我贫乏的记忆，将一切必要的东西收集在精选而连贯的标题之下，在这些标题下，我似乎与其说是处理了主题，不如说是为他人提供了处理主题的材料。此外，对读者而言，同样的简洁亦大有裨益，因为过于冗长的论述会涣散读者的理解与领悟，而持久的记忆则以更精确的纲要保存所读的内容。但我将我的工作编成了两卷，篇幅同样适中：其中一卷我竭力表明，犹太人按照先前所预言的，已离弃了天主，并失去了天主在往昔赐予他们、未来又应许给他们的恩宠；而基督徒则继承了他们的地位，凭着信德堪当蒙主悦纳，从万族和普世而来。第二卷同样包含基督的圣事，即祂已经来临，是照圣经所宣布的，并完成和实现了所有关于祂能被认知和辨识的预言。这些事在你阅读时，可暂时有益于奠定你信德的初步轮廓。当你更全面地研读新旧圣经，通读灵性书籍的完整卷册时，将获得更大的力量，心灵的领悟也日益增进。因为如今我们已从神圣的泉源中汲取了少许，暂且寄给你。你若也靠近同一些神圣圆满的泉源，与我们一同饮用，便能更丰盛地畅饮和满足。我亲爱的儿子，愿你永远衷心康泰。\par

% structure: p0004 source=h2 target=subsection reason=relative-heading-rank; base=h2

\subsection{一、犹太人因离弃主、追随偶像而陷于天主重怒之下}

在出谷纪中，百姓对亚郎说：\BibleQuote{起来，为我们造神，在我们前面引路；因为那个领我们出埃及地的梅瑟，我们不知道他遭遇了什么事。}\BibleRef{出 32:1} 同处，梅瑟也对主说：\BibleQuote{主，我恳求祢，这百姓犯了罪！犯了大罪。他们为自己制造了金银的神。现在，祢若赦免他们的罪，就赦免；若不然，求祢从祢所写的册上涂抹我的名字。主对梅瑟说：谁得罪了我，我就从我的册上涂抹他。}\BibleRef{出 32:31} 同样在申命纪中：他们献祭给魔鬼，而不是给天主。\BibleRef{申 32:17} 在民长纪中也说：\BibleQuote{以色列子民行了主他们祖先的天主眼中视为恶的事，就是领他们出埃及地的天主，去随从四周民族的神，触犯主，离弃了天主，事奉了巴耳。}同处又说：\BibleQuote{以色列子民又行了主眼中视为恶的事，事奉巴耳和异邦的神，离弃了主，不事奉祂。}\BibleRef{民 4:1} 在玛拉基亚中：\BibleQuote{犹大被离弃，在以色列和耶路撒冷中成了可憎之物，因为犹大亵渎了主的圣洁，在祂所爱的事上，并追求异邦的神。凡行这事的人，主必剪除，在雅各伯的帐幕中他必被轻贱。}\BibleRef{拉 2:11}\par

% structure: p0006 source=h2 target=subsection reason=relative-heading-rank; base=h2

\subsection{二、又因他们不信先知并杀害先知}

在耶肋米亚中，主说：\BibleQuote{我打发我的仆人众先知到你们那里去。黎明前我打发他们（你们却不听我，也不侧耳听），说：你们各人当回转离开自己的邪道和极恶的思念，就必永远住在我赐给你们和你们祖先的地上。}又说：\BibleQuote{你们不要随从别的神，事奉敬拜它们，也不要以你们手所做的惹我发怒，以致我驱散你们；你们却没有听从我。}\BibleRef{耶 25:6} 又在列王纪上中，厄里亚对主说：\BibleQuote{我为主万军的天主忧心如焚，因为以色列子民离弃了祢，拆毁了祢的祭台，用刀杀了祢的先知；只剩下了我一人，他们还要寻索我的性命，要将它夺去。}\BibleRef{列上 19:10} 在厄斯德拉中也说：\BibleQuote{他们背叛了祢，将祢的法律拋在背后，杀了祢的先知，那些曾警戒他们，叫他们归向祢的先知。}\BibleRef{厄下 9:26}\par

% structure: p0008 source=h2 target=subsection reason=relative-heading-rank; base=h2

\subsection{三、预先预言他们将不认识主、不明白主、也不接受主}

在\BibleBook{}{依撒意亚}中：\BibleQuote{诸天哪，侧耳而听，大地啊，留心倾听：因为上主说了；我养育儿女，将他们养大，他们竟背叛了我。牛认识主人，驴认识主人的槽；以色列却不认识我，我的子民也不明白我。祸哉，犯罪的民族，罪孽深重的百姓，邪恶的种子，败坏的儿女！你们离弃了上主，惹动了以色列圣者的怒气。}\BibleRef{依 1:2}在同一卷中，上主又说：\BibleQuote{你去告诉这百姓：你们听是听见，却不明白；看是看见，却不晓得。因为这百姓的心蒙了脂油，耳朵发沉，眼睛闭着；恐怕眼睛看见，耳朵听见，心里明白，回转过来，我便医治他们。}\BibleRef{依 6:9}在\BibleBook{}{耶肋米亚}中，上主也说：\BibleQuote{他们离弃了我这活水的泉源，为自己凿出破裂不能存水的池子。}又在同一卷中：\BibleQuote{看哪，上主的话对他们成了羞辱，他们不再喜爱它。}\BibleRef{耶 6:10}在同一卷中，上主又说：\BibleQuote{鸢知道自己的时令，斑鸠、燕子，田野的麻雀也守候归来的季节；但我的百姓不知道上主的审判。你们怎么说：我们有智慧，上主的法律与我们同在？假度量衡已经作废；经师们惭愧，智者们惊惶，被捉住了，因为他们弃绝了上主的话。}\BibleRef{耶 8:7}在\Person{撒罗满}所著的书中：\BibleQuote{恶人寻找我，却寻不着；因为他们憎恨智慧，不领受上主的话。}\BibleRef{箴 1:28}在第二十七篇圣咏中：\BibleQuote{求你按他们的行为报应他们，因为他们不留心上主的作为。}在第八十一篇圣咏中：\BibleQuote{他们不知道，也不明白；他们必在黑暗中行走。}在\BibleBook{若望}{福音}中：\BibleQuote{祂到自己的人那里，自己的人却不接受祂。凡接受祂的，祂赐给他们成为天主的子女的权能，就是信祂名的人。}\BibleRef{若 1:11}\par

% structure: p0010 source=h2 target=subsection reason=relative-heading-rank; base=h2

\subsection{四、犹太人将不理解圣经，但在最后的时期，基督降临之后，圣经将可理解}

在\BibleBook{}{依撒意亚}中：\BibleQuote{这一切的话对你们就像封严的书卷的话，若交给识字的人读，他必说：我不能读，因为封严了。但到那日，聋子必听见这书卷的话，在黑暗和云雾中的人，瞎子的眼必看见。}\BibleRef{依 29:11}在\BibleBook{}{耶肋米亚}中：\BibleQuote{在末后的日子，你们必明白这些事。}\BibleRef{耶 23:20}又在\BibleBook{}{达内尔}中：\BibleQuote{你要封闭这话，封严这书，直到末期；直到多人学习，知识增长；因为当分散时，他们必明白这一切事。}\BibleRef{达 12:4}同样，在\BibleBook{格林多}{前书}中：\BibleQuote{弟兄们，我不愿意你们不知道，我们的祖先都曾在云柱下。}\BibleRef{格前 10:1}也在\BibleBook{格林多}{后书}中：\BibleQuote{他们的心思直到今日还是迟钝的，因为有这同一帕子罩住，只在基督里才被除去；而直到今日，每逢诵读旧约时，这帕子仍存留着，没有揭去，因为在基督里才废除了；而且直到今日，每逢诵读梅瑟时，帕子还是蒙在他们心上。但几时他们归向主，帕子就除去了。}在福音中，主复活后说：\BibleQuote{这就是我还与你们同在时对你们所说的话：凡梅瑟法律、先知和圣咏上所载关于我的一切，都必须应验。于是祂开启他们的明悟，使他们理解圣经；并对他们说：经上这样记载：基督必须受苦，第三天从死人中复活；并且要因祂的名，从耶路撒冷开始，直至万邦，宣讲悔改和赦罪。}\BibleRef{路 24:44}\par

% structure: p0012 source=h2 target=subsection reason=relative-heading-rank; base=h2

\subsection{五、犹太人除非先相信基督，否则无法理解圣经}

在\BibleBook{}{依撒意亚}中：\BibleQuote{如果你们不相信，你们必不能理解。}\BibleRef{依 7:9}主在福音中也说：\BibleQuote{如果你们不信我就是那一位，你们必要死在你们的罪恶中。}\BibleRef{若 8:24}此外，义德由信德而成，且在其中获得生命，这已在\BibleBook{}{哈巴谷}中预言：\BibleQuote{义人必因我的信德而生活。}\BibleRef{哈 2:4}因此，万民之父亚巴郎相信了；在\BibleBook{}{创世纪}中：\BibleQuote{亚巴郎相信了天主，天主就以此算为他的义德。}\BibleRef{创 15:6}同样，\Person{保禄}致迦拉达人说：\BibleQuote{亚巴郎相信了天主，天主就以此算为他的义德。所以，你们要知道，凡出于信德的，都是亚巴郎的子孙。圣经预见天主将由信德使外邦人成义，就预先向亚巴郎传报：万民都要因你而蒙受祝福。因此，凡出于信德的，都与有信德的亚巴郎一同蒙受祝福。}\BibleRef{迦 3:6}\par

% structure: p0014 source=h2 target=subsection reason=relative-heading-rank; base=h2

\subsection{六、犹太人将失去耶路撒冷，并离开他们曾领受的土地}

在\BibleBook{}{依撒意亚}中：\BibleQuote{你们的国土荒凉，你们的城邑被火焚烧；你们的土地，外邦人在你们眼前吞吃它；熙雍的女儿被遗弃，被外族倾覆，像葡萄园中的草棚，像黄瓜园中的茅舍，像被围困的城。若非万军的上主给我们留下一些残余，我们早已像索多玛，与哥摩辣一样了。}\BibleRef{依 1:7}主在福音中也说：\BibleQuote{耶路撒冷，耶路撒冷，你杀死先知，又用石头砸死那些派遣到你这里来的人。我多少次愿意聚集你的子女，有如母鸡聚集自己的小鸡在翅膀底下，但你却不愿意！看，你们的房屋将给你们留下一片荒凉。}\BibleRef{玛 23:37}\par

% structure: p0016 source=h2 target=subsection reason=relative-heading-rank; base=h2

\subsection{七、他们也将失去上主的光明}

在\BibleBook{依撒意亚}{书}中：\Quote{来，让我们在上主的光中行走。因为祂已遣散了祂的子民，以色列家。} \BibleRef{依 2:5} 在\BibleBook{若望}{福音}中：\Quote{那才是真光，照亮那来到这世界上的每一个人。祂在这世界上，世界是由祂造成的，世界却不认识祂。} \BibleRef{若 3:18} 此外，在同一处：\Quote{那不信的，已受了审判，因为他没有信从天主独生子的名字。审判就在于此：光明来到了世界，世人却爱黑暗甚于光明。} \BibleRef{若 3:18}\par

% structure: p0018 source=h2 target=subsection reason=relative-heading-rank; base=h2

\subsection{八、肉身的首次割损被废除，代之而应许的是属灵的第二次割损}

在\BibleBook{耶肋米亚}{书}中：\Quote{上主这样对犹大人和耶路撒冷的居民说：你们要为自己开垦新的荒地，不要在荆棘中撒种；你们要为自己的天主行割损，割损你们心中的包皮，免得我的忿怒如火烧起，焚烧你们，无人能扑灭。} \BibleRef{耶 4:3} 梅瑟也说：\Quote{在末日，天主必割损你的心和你的后裔的心，使你们爱上主你的天主。} \BibleRef{申 30:6} 在\BibleBook{若苏厄}{书}中：\Quote{上主对若苏厄说：你要制造极锋利的石刀，第二次为以色列子民行割损。} \BibleRef{苏 5:2} 在\BibleBook{哥罗森}{书}中：\Quote{你们是在基督内受了不是人手所行的割损，而是脱去肉欲的割损，就是基督的割损。} \BibleRef{哥 2:11} 此外，因为亚当最初由天主造为未受割损者，义人亚伯尔、取悦天主并被接去的哈诺客、以及当世界和人类因过犯而灭亡时唯独被选、使人类得以保存的诺厄，还有默基瑟德，那位基督按其品位被预许的司祭。然后，因为那标记对女人无效，但所有人都被主的标记所印封。\par

% structure: p0020 source=h2 target=subsection reason=relative-heading-rank; base=h2

\subsection{九、藉梅瑟所赐的旧法律将要终止}

在\BibleBook{依撒意亚}{书}中：\Quote{那时，那些封存法律、不使人学习的人必显明出来；他必说：我等候上主，祂转脸不顾雅各伯家，我要信赖祂。} \BibleRef{依 8:16} 在\BibleBook{玛窦}{福音}中：\Quote{众先知和法律讲说预言，直到若翰为止。} \BibleRef{玛 11:13}\par

% structure: p0022 source=h2 target=subsection reason=relative-heading-rank; base=h2

\subsection{十、将要赐予新法律}

在\BibleBook{米该亚}{书}中：\Quote{因为法律将出自熙雍，上主的话将出自耶路撒冷。祂将在许多民族中施行审判，祂将制伏并揭露强盛的民族。} \BibleRef{米 4:2} 在\BibleBook{依撒意亚}{书}中：\Quote{因为法律将出自熙雍，上主的话将出自耶路撒冷；祂将在列国中施行审判。} \BibleRef{依 2:3} 同样在\BibleBook{玛窦}{福音}中：\Quote{忽然有声音从云彩中说：这是我的爱子，我所喜悦的，你们要听从祂。} \BibleRef{玛 17:5}\par

% structure: p0024 source=h2 target=subsection reason=relative-heading-rank; base=h2

\subsection{十一、将要有另一种安排和一个新盟约}

在\BibleBook{耶肋米亚}{书}中：\Quote{上主说：看，时日将到，我要与以色列家和犹大家订立新约，不像我牵着他们祖先的手，领他们出离埃及地的那天，与他们所订立的约；因为他们没有存留在我这盟约内，我就离弃了他们，上主说。因为这是那些日子以后，我将与以色列家订立的盟约，上主说：我要将我的法律放在他们心中，写在他们的心思上；我要作他们的天主，他们要作我的子民。他们各人不再教导自己的兄弟说：你们该认识上主，因为从最小的到最大的，都要认识我，因为我要宽恕他们的不义，不再记忆他们的罪恶。} \BibleRef{耶 31:31}\par

% structure: p0026 source=h2 target=subsection reason=relative-heading-rank; base=h2

\subsection{十二、旧的圣洗将要终止，新的圣洗将要开始}

在\BibleBook{依撒意亚}{书}中：\Quote{因此，你们不要追念先前的事，也不要回顾古时的事。看，我要行一件新事，如今就要发生，你们要知道；我必在旷野中开出一条道路，在旱地上流出江河，给我所拣选的子民、我所获得的人民喝水，使他们宣扬我的赞美。} \BibleRef{依 43:18} 在\BibleBook{依撒意亚}{书}中：\Quote{他们若口渴，祂必引导他们经过旷野；祂必从磐石中引出水来；磐石裂开，水就涌出；我的人民必喝。} \BibleRef{依 48:21} 此外，在\BibleBook{玛窦}{福音}中，若翰说：\Quote{我固然以水给你们授洗，叫你们悔改；但在我以后来的那一位，比我更强，我连给祂提鞋也不配；祂要以圣神和火给你们授洗。} \BibleRef{玛 3:11} 在\BibleBook{若望}{福音}中：\Quote{人除非由水和圣神而生，不能进天主的国。因为由肉生的属于肉，由神生的属于神。} \BibleRef{若 3:5}\par

% structure: p0028 source=h2 target=subsection reason=relative-heading-rank; base=h2

\subsection{十三、旧的轭将要废除，新的轭将要赐下}

在\BibleBook{圣咏}{}第二篇中：\Quote{外邦人为什么嚣张，人民为什么妄想？地上的列王齐集，官长聚首，反抗上主，反抗祂的受傅者。让我们挣断他们的捆绑，从我们身上甩掉他们的轭。} 同样在\BibleBook{玛窦}{福音}中，主说：\Quote{凡劳苦和负重担的，你们都到我这里来，我要使你们安息。你们背起我的轭，跟我学吧！因为我是良善心谦的：这样你们必要找到你们灵魂的安息。因为我的轭是美好的，我的担子是轻省的。} \BibleRef{玛 11:28} 在\BibleBook{耶肋米亚}{书}中：\Quote{在那一天，我必打碎他们颈上的轭，挣断他们的绳索；他们不再为外人劳役，而要为上主天主劳役；我必给他们兴起达味作王。} \BibleRef{耶 30:8}\par

% structure: p0030 source=h2 target=subsection reason=relative-heading-rank; base=h2

\subsection{十四、旧的牧者将要终止，新的牧者将要开始}

在\BibleBook{厄则克耳}{先知书}中：\Quote{因此，主这样说：看，我高于牧者；我必从他们手中追讨我的羊，我要使他们不再牧放我的羊，他们不得再牧养它们，我要从他们口中救出我的羊，并要以公义牧养它们。}在\BibleBook{耶肋米亚}{先知书}中，主说：\Quote{我必赐给你们符合我心的牧者，他们要用地道的知识牧养你们。}\BibleRef{耶 3:15}在\BibleBook{耶肋米亚}{先知书}中，又说：\Quote{列国啊，请听主的话，并传达到远方的海岛。你们说：那分散以色列的，必集合他，并看守他，像牧者看守自己的羊群；因为主赎回了雅各伯，从他更强者手中救出了他。}\par

% structure: p0032 source=h2 target=subsection reason=relative-heading-rank; base=h2

\subsection{十五、基督应成为天主的殿宇和圣殿，古老的殿宇要停止，新的殿宇要开始}

在\BibleBook{}{列王纪下}中：\Quote{主的话临到纳堂说：你去告诉我的仆人达味，主这样说：你不可为我建造殿宇居住；但当你的日子满足，与你祖先同睡时，我必兴起你身所出的后裔，坚定他的国。他必为我名建造殿宇，我必坚定他的王位直到永远；我要作他的父亲，他要作我的儿子；他的家必坚定，他的国在我面前直到永远。}在福音中，主说：\Quote{在圣殿里将没有一块石头留在另一块上而不被拆毁。}\BibleRef{玛 24:2}又说：\Quote{三天后，另一座将不靠人手而建立。}\par

% structure: p0034 source=h2 target=subsection reason=relative-heading-rank; base=h2

\subsection{十六、古老的祭献应被废止，新的祭献应被庆祝}

在\BibleBook{依撒意亚}{先知书}中：\Quote{你们所献的众多祭品对我有何意义？主说：我饱足了；我不愿再要公羊的燔祭和羔羊的脂油，以及公牛和山羊的血；谁从你们手中要求这些？}\BibleRef{依 1:11}在\BibleBook{}{圣咏}第49篇中：\Quote{我不吃公牛的肉，也不喝山羊的血。你们要向天主献上赞颂的祭，向至高者还你的誓愿。在患难之日呼求我，我必拯救你，你也要光荣我。}在同一篇\BibleBook{}{圣咏}中又说：\Quote{赞颂的祭要光荣我，这是使我能向他显示天主救恩的道路。}在\BibleBook{}{圣咏}第4篇中：\Quote{献上正义的祭，并信赖上主。}同样在\BibleBook{玛拉基亚}{先知书}中：\Quote{我不喜悦你们，主说，也不从你们手中悦纳祭品。因为从日出之地到日落之处，我的名在万邦中受显扬，在各处都向我的名奉献馨香和洁净的祭品，因为我的名在万邦中为大，主说。}\par

% structure: p0036 source=h2 target=subsection reason=relative-heading-rank; base=h2

\subsection{十七、旧的司铎职应被终止，一位新的司铎要来，祂将永远为司铎}

在\BibleBook{}{圣咏}第109篇中：\Quote{在晨星之前我已生了你。上主发了誓，决不会反悔：你按照默基瑟德的品位，永为司铎。}在\BibleBook{}{列王纪上}中，天主对司铎厄里说：\Quote{我必为自己兴起一位忠信的司铎，他要照我心意行事；我要为他建立坚固的家，他要在我的受傅者面前行走，直到永远。凡留在你家中的，都要为一块钱和一个饼来叩拜。}\par

% structure: p0038 source=h2 target=subsection reason=relative-heading-rank; base=h2

\subsection{十八、另一位像梅瑟一样的先知被应许，即要赐予新约的那位，更应该听从祂}

在\BibleBook{}{申命纪}中，天主对梅瑟说：\Quote{主对我说：我要从他们兄弟中兴起一位像你一样的先知，将我的话放在他口中；他要将我所吩咐的，都对他们说。凡不听从那先知因我名所说的话的，我必追讨。}\BibleRef{申 18:18}关于祂，基督在\BibleBook{若望}{福音}中也说：\Quote{你们查考圣经，以为在其中获得永生；正是这些为我作证；但你们不愿到我这里来，为获得生命。不要想我要在父前控告你们；有一位控告你们的，就是你们所寄望的梅瑟。如果你们相信梅瑟，也会相信我，因为他论及我写了书。如果你们不相信他的书，怎能相信我的话？}\par

% structure: p0040 source=h2 target=subsection reason=relative-heading-rank; base=h2

\subsection{十九、两个民族被预言，年长的和年幼的，即旧的犹太民族和将由我们组成的新民族}

在\BibleBook{}{创世纪}中：\Quote{主对黎贝加说：两国在你腹中，两族从你肚子里分离；这族要强过那族，年长的要服侍年幼的。}\BibleRef{创 25:23}在\BibleBook{欧瑟亚}{先知书}中：\Quote{我要将非我子民的称为我的子民，将不蒙爱的称为蒙爱的。从前在什么地方对他们说：你们不是我的子民，将来在那里他们必要被称为永生天主的子女。}\par

% structure: p0042 source=h2 target=subsection reason=relative-heading-rank; base=h2

\subsection{二十、从前不生育的教会将从外邦人中拥有比先前的会堂更多的子女}

在\BibleBook{依撒意亚}{先知书}中：\Quote{不生育的，欢欣吧！你这不生产的，欢乐高呼吧！因为被遗弃的子女比有丈夫的还要多。因为主曾说：扩大你帐幕的地方，伸展你居所的帷幔，不要顾忌；拉长你的绳索，坚固你的橛子；向左右扩展；你的后裔将占领外邦，居住荒城。不要害怕，因为你必得胜；不要因你受诅咒而畏惧，因为你将忘记永恒的羞耻。} \BibleRef{依 54:1} 同样，对\Person{亚巴郎}也是如此，当他先前的儿子由婢女所生时，\Person{撒辣}长久不生育；直到晚年才生下她许诺的儿子\Person{依撒格}，他是基督的预像。同样，\Person{雅各伯}有两个妻子：年长的\Person{肋阿}，眼睛昏花，是会堂的预像；年轻美丽的\Person{辣黑耳}，是教会的预像，她也长久不生育，以后生下\Person{若瑟}，他本身也是基督的预像。在\BibleBook{}{列王纪上}记载，\Person{厄耳卡纳}有两个妻子：\Person{培尼纳}和她的儿子们；以及不生育的\Person{亚纳}，她生的\Person{撒慕尔}，不是按生育的自然规律，而是按天主的慈悲和应许，当她在圣殿祈祷后；\Person{撒慕尔}出生后，是基督的预像。同样在\BibleBook{}{列王纪上}：\Quote{不生育的生了七个，有许多子女的反而衰弱。} 但这七个孩子就是七个教会。因此\Person{保禄}也写信给七个教会；\BibleBook{}{默示录}也列出了七个教会，以保持七这个数字；正如天主创造世界的七天；正如在天主面前侍立和出入的七位天使，正如\BibleBook{}{多俾亚传}中天使\Person{拉法厄尔}所说的；会幕中的七灯盏；以及天主的七眼，眷顾世界；还有\Person{匝加利亚}所说的七眼石头；七神；\BibleBook{}{默示录}中的七个灯台；以及\BibleBook{}{箴言}中智慧建造房屋的七根柱子。\par

% structure: p0044 source=h2 target=subsection reason=relative-heading-rank; base=h2

\subsection{二十一、外邦人更应相信基督}

在 \BibleBook{创世纪}{} 中：\Quote{上主天主对亚巴郎说：离开你的故乡、你的家族和父家，往我指给你的地方去。我要使你成为一个大民族，我必祝福你，使你成名，你将成为祝福：祝福你的，我必祝福他；咒骂你的，我必咒骂他。地上万民都要因你获得祝福。} \BibleRef{创 11:1} 关于同一主题，在 \BibleBook{创世纪}{} 中：\Quote{依撒格祝福雅各伯说：看，我儿子的香气，像上主祝福的丰饶田地的香气。愿天主赐给你天上的甘露，土地的肥沃，五谷、美酒和油的丰裕！愿人民服侍你，民族向你跪拜！愿你做你兄弟的主人，你父亲的儿子们向你跪拜！凡咒骂你的，必受咒骂；凡祝福你的，必受祝福。} \BibleRef{创 27:27} 关于此事也在 \BibleBook{创世纪}{} 中：\Quote{若瑟见他父亲将右手放在厄弗辣因头上，心中不悦，便握住父亲的手，要从厄弗辣因头上移到默纳协头上，对父亲说：不，父亲！这才是长子，请将你的右手放在他头上。但他不肯，说：我知道，我儿，我知道：他也将成为一个民族，他也将受尊崇；但他的弟弟却要比他更大，他的后裔要成为一大民族。} 此外在 \BibleBook{创世纪}{} 中：\Quote{犹大！你的兄弟要赞扬你；你的手必按在你仇敌的背上；你父亲的儿子要向你跪拜。犹大是只幼狮：从细嫩的枝条，我儿，你上升了；你躺卧如雄狮，又如幼狮，谁能惊动他？权杖不离犹大，领袖不离他两腿之间，直到那所委托的来到；他是万民的希望。他将驴驹拴在葡萄树上，将驴的幼驹拴在美好的葡萄树上；在酒中洗他的衣服，在葡萄汁中洗他的外氅；他的双眼因酒而威严，他的牙齿比牛奶更白。} \BibleRef{创 49:8} 因此在 \BibleBook{户籍纪}{} 中论及我们的人民写道：\Quote{看，这民族要如狮般奋起。} \BibleRef{户 23:14} 在 \BibleBook{申命纪}{} 中：\Quote{你们外邦人要作头，但这不信的民族要作尾。} \BibleRef{申 28:44} 在 \BibleBook{耶肋米亚}{} 中也说：\Quote{听号角声！他们却说：我们不愿听。为此，万民要听，并在他们中牧放牲畜的人要听。} \BibleRef{耶 6:18} 在 \BibleBook{圣咏}{} 第十七篇中：\Quote{你要立我为万国的元首：我所不认识的人民都事奉了我；他们一听见我，就服从了我。} 关于此事，主在 \BibleBook{耶肋米亚}{} 中说：\Quote{我还没有在母腹内形成你以前，我已认识了你；在你还没有出离母胎以前，我已圣化了你，并立你作万民的先知。} \BibleRef{耶 1:5} 在 \BibleBook{依撒意亚}{} 中也说：\Quote{看，我已显现他，作为万民的见证，列国的领袖和统帅。} \BibleRef{依 55:4} 在同一书中又说：\Quote{不认识你的民族，要求告你；不认识你的民族，要投奔你。} \BibleRef{依 55:5} 在同一书中又说：\Quote{在那一天，叶瑟的根子要升起，统治万民；外邦人要寄望于他；他的安息之所将充满光荣。} \BibleRef{依 11:10} 在同一书中又说：\Quote{则步隆地、纳斐塔里地，沿海之路，以及你们其他居住海边的人，和约但河外的万民。在黑暗中行走的百姓，看啊，一道大光；你们住在死荫之地的人，必有光照耀你们。} \BibleRef{依 9:1} 在同一书中又说：\Quote{上主天主对我的主基督这样说：我紧握他的右手，为使万民听从；我必粉碎列王的力量，在他面前敞开城门，使城邑不再关闭。} 在同一书中又说：\Quote{我要来聚集万民和说各种语言的人；他们都要来，看见我的荣耀。我要在他们中间竖起旗帜，并派遣他们中得救的人到遥远的万民那里去，就是那些没有听过我的名，没有见过我荣耀的民族；他们要在万民中宣扬我的荣耀。} 在同一书中又说：\Quote{虽遭遇这一切，他们仍不归化；因此他必向远方的万民竖立旗帜，将他们从地极招来。} 在同一书中又说：\Quote{那些未被告知他的人将要看见，那些未曾听过的人将要明白。} 在同一书中又说：\Quote{我向那不求问我的人显现，我被那不寻找我的人寻见。我曾说：我在这里，我在这里，对那没有呼求我名的民族。} 关于同一件事，在 \BibleBook{宗徒大事录}{} 中，保禄说：\Quote{天主的圣道本当先讲给你们听；但因你们弃绝，断定自己不配得永生，看，我们转向外邦人，因为主曾藉圣经这样说：看，我已立你为万民的光，使你把救恩带到地极。}\par

% structure: p0046 source=h2 target=subsection reason=relative-heading-rank; base=h2

\subsection{二十二、犹太人将失去，而我们应领受基督的饼与杯及祂所有的恩宠，并且基督徒的新名要在世上受祝福。}

在依撒意亚书中：\BibleQuote{主这样说：看，事奉我的人要吃饭，你们却要饥饿；看，事奉我的人要喝水，你们却要口渴；看，事奉我的人要喜乐，你们却要蒙羞；主要击杀你们。但事奉我的人必得一个新名，那名字在地上必成为祝福。}同样在该处：\BibleQuote{因此，祂要向远方的万国举起旗帜，从地极招引他们；看，他们必迅速轻快地前来；他们不饥不渴。}同样在该处：\BibleQuote{看，因此，万军的上主必从犹大和耶路撒冷除掉健壮的人和强有力的人，除掉面饼的支柱和水的支柱。}同样在第三十三篇圣咏中：\BibleQuote{你们要尝尝并看看，上主是何等甘饴。投靠祂的人真是有福。上主的圣人们，你们要敬畏上主，因为敬畏祂的人一无所缺。富人竟成了赤贫，忍饥受饿；但寻求上主的人，却不缺任何美物。}此外，在若望所载的福音中，主说：\BibleQuote{我就是生命的食粮；到我这里来的，永不会饥饿；信从我的，总不会口渴。}同样，祂在该处说：\BibleQuote{谁若口渴，让他到我这里来喝。信从我的人，就如经上说的，从他的腹中要流出活水的江河。}此外，祂在同一处说：\BibleQuote{你们若不吃人子的肉，不喝他的血，在你们内就没有生命。}\par

% structure: p0048 source=h2 target=subsection reason=relative-heading-rank; base=h2

\subsection{二十三、外邦人而非犹太人将获得天国}

在福音中，主说：\BibleQuote{有许多人要从东方和西方来，与\Person{亚巴郎}、\Person{依撒格}和\Person{雅各伯}一同坐席于天国；但本国的子民反要被驱逐到外面的黑暗中，那里要有哀号和切齿。}\par

% structure: p0050 source=h2 target=subsection reason=relative-heading-rank; base=h2

\subsection{二十四、犹太人唯有以此方式——即在祂的圣洗中洗去被杀的基督的血，并加入祂的教会，服从祂的诫命——才能获得罪的赦免}

在依撒意亚书中，主说：\BibleQuote{现在我不再赦免你们的罪。当你们伸出手时，我必掩面不看你们；即使你们多多祈祷，我也不听：因为你们的手沾满了血。你们要洗涤，自洁；从我眼前除掉你们的恶行；停止作恶，学习行善；寻求正义；保护受压迫的人；为孤儿伸冤，为寡妇辩屈。然后来，让我们一同辩论，主说：你们的罪虽似朱红，必变成雪一样洁白；虽红如丹颜，必白如羊毛。如果你们愿意听从我，就必享用地上的美物；但如果你们拒绝不听，刀剑必吞灭你们，因为上主亲口说了。}\par

\SourceRecord{Translated by Robert Ernest Wallis.；Ante-Nicene Fathers；Vol. 5.；Edited by Alexander Roberts, James Donaldson, and A. Cleveland Coxe.；Buffalo, NY: Christian Literature Publishing Co.,；1886.；<http://www.newadvent.org/fathers/050712a.htm>.}

% structure: p0001 source=h1 target=section reason=page-title

\section{论集十二，第二卷}

% structure: p0002 source=h2 target=subsection reason=relative-heading-rank; base=h2

\subsection{一、基督是首生者，祂是天主的智慧，万物藉祂而造成}

在撒罗满的\BibleBook{}{箴言}中记载：\Quote{主在造化之初，在祂的作为之先，就设立了我在祂的道路之始。在太初，在祂造地以前，在祂立定深渊以前，在泉水涌出以前，在群山奠定以前，在丘陵之前，主已生了我。祂造了大地、无人居住之处和天下无人之境。当祂预备诸天时，我已在祂身旁；当祂安置祂的宝座时，当祂在高处造了浓云，当祂在天下安置了坚固的泉源，当祂立定大地的根基时，我在祂身侧，调度一切：我是祂所喜悦的；并且我天天在祂面前欢跃，当祂因所造的世界而喜悦时。}（\BibleRef{箴 8:22-31}）又在同一位的\BibleBook{}{德训篇}中：\Quote{我出自至高者之口，是首生于一切受造物之前：我使不灭的光在天空升起，又以云彩遮盖大地：我居住在至高之处，我的宝座在云柱上：我环绕天穹，又深入深渊，行走在海浪上，并站在全地上；在万民万族中我居首位，以我自己的力量踏遍了所有高贵和卑微者的心：在我内有一切生命和德行的希望：凡渴望我的人，都到我这里来吧。}（\BibleRef{德 24:3-7}）又在第八十八篇\BibleBook{}{圣咏}中：\Quote{我要立祂为首生，为地上诸王中最高的。我必为祂永远保存我的仁慈，和祂订立忠信的盟约；并要建立祂的后裔直到永远。如果祂的子孙离弃我的法律，不遵守我的典章；如果他们亵渎我的条例，不遵守我的诫命，我必用杖责罚他们的过犯，用鞭责罚他们的罪恶；但我的仁慈必不向祂收回。}（\BibleRef{咏 88:28-34}）又在\BibleBook{若望}{福音}中，主说：\Quote{永生就是：认识你，唯一的真天主，和你所派遣的耶稣基督。我在地上已光荣了你，完成了你所托付我做的工作。现在，父啊，请以我在未有世界以前，与你同享的光荣光荣我吧。}（\BibleRef{若 17:3-5}）又保禄致\BibleBook{}{哥罗森人书}：\Quote{祂是不可见的天主的肖像，是一切受造物的首生者。}（\BibleRef{哥 1:15}）又同一处：\Quote{从死者中首生的，为使祂在一切事上成为元首。}（\BibleRef{哥 1:18}）在\BibleBook{}{默示录}中也是如此：\Quote{我是\InnerQuote{阿耳法}和\InnerQuote{敖默加}，最初的和最末的。我要把生命之泉的水白白赐给口渴的人。}（\BibleRef{默 21:6}）祂也是天主的智慧和德能，保禄在\BibleBook{}{格林多前书}中证明了这一点：\Quote{因为犹太人要求的是神迹，希腊人寻求的是智慧；而我们所宣讲的，却是被钉在十字架上的基督：这为犹太人固然是绊脚石，为外邦人是愚妄；但为那些蒙召的，不拘是犹太人或希腊人，基督却是天主的德能和天主的智慧。}（\BibleRef{格前 1:22-24}）\par

% structure: p0004 source=h2 target=subsection reason=relative-heading-rank; base=h2

\subsection{二、基督是天主的智慧；并论祂降生成人及受难的圣事，论及杯和祭台，以及被派遣宣讲的宗徒}

在撒罗满的\BibleBook{}{箴言}中记载：\Quote{智慧建造了房屋，凿了七根柱子；宰杀了牲畜，调好了美酒，铺设了筵席，派遣了她的丫鬟，在高处呼喊，向杯说：谁是无知的，请转向我；并向缺少心智的人说：你们来，吃我的饼，喝我调和的酒。你们要离弃无知，寻求智慧，以智识改正你们的明悟。}（\BibleRef{箴 9:1-6}）\par

% structure: p0006 source=h2 target=subsection reason=relative-heading-rank; base=h2

\subsection{三、同一基督是天主的圣言}

在第四十四篇\BibleBook{}{圣咏}中：\Quote{我的心涌出美善的圣言。我要向君王陈述我的作为。}（\BibleRef{咏 45:2}）又在第三十二篇\BibleBook{}{圣咏}中：\Quote{因天主的圣言，诸天得以稳固；祂口中之气，形成万有的力量。}（\BibleRef{咏 33:6}）又在\BibleBook{}{依撒意亚}中：\Quote{一句总括而缩短的言语，因正义而完成，因为天主必在普世完成一句缩短的言语。}（\BibleRef{依 10:22-23}）又在第一百零六篇\BibleBook{}{圣咏}中：\Quote{祂派遣自己的圣言，医好了他们。}（\BibleRef{咏 107:20}）此外，在\BibleBook{若望}{福音}中：\Quote{在起初已有圣言，圣言与天主同在，圣言就是天主。这圣言在起初就与天主同在。万物是藉着祂而造成的；凡受造的，没有一样不是藉着祂造成的。在祂内有生命，这生命是人的光。光在黑暗中照耀，黑暗不能胜过这光。}（\BibleRef{若 1:1-5}）又在\BibleBook{}{默示录}中：\Quote{我看见天开了，看，有一匹白马；那骑马的称为\InnerQuote{忠信}和\InnerQuote{真实}，祂按正义审判和作战。祂穿着一件浸过血的衣服，祂的名字被称为\InnerQuote{天主的圣言}。}（\BibleRef{默 19:11-13}）\par

% structure: p0008 source=h2 target=subsection reason=relative-heading-rank; base=h2

\subsection{四、基督是天主的手和臂}

在\BibleBook{}{依撒意亚}中：\BibleQuote{天主的手岂不坚强足以拯救？或者祂使自己的耳朵发沉，以致不能听见？但你们的罪恶在你们与天主之间造成隔绝，因你们的罪恶，祂转脸不看你们，以致不怜悯你们。因为你们的手染满了血，你们的手指沾满了罪恶；你们的嘴唇说了恶言，你们的舌头默想不义。无人说真话，也没有真实的判断；他们信赖虚无，说虚空的话，怀孕痛苦，产生邪恶。}\newline{}同样在那里：\BibleQuote{上主，谁相信了我们的报告？天主的臂向谁启示了呢？}\newline{}同样在那里：\BibleQuote{上主这样说：诸天是我的宝座，大地是我脚的脚凳。你们要为我建筑什么殿宇？或者什么是我的安息之所？这一切都是我的手所造的。}\newline{}同样在那里：\BibleQuote{上主天主，你的臂高举，他们却不认识；但当他们认识时，必蒙羞。}\newline{}同样在那里：\BibleQuote{上主在万民眼前显示了他的圣臂；万民，甚至地极，都要看见天主的救恩。}\newline{}同样在那里：\BibleQuote{看，我使你像打谷车的新轮子，向后翻转；你要打碎山岭，将山丘敲碎，使他们像糠秕，并簸扬它们；风要卷走它们，旋风要吹散它们：但你要因以色列的圣者而喜乐；贫穷困苦的人要欢腾。因为他们寻找水，却没有水；他们的舌头因口渴而干燥。我上主天主，以色列的天主，必应允他们，不遗弃他们；但我要在山上开辟河流，在田野中开出泉源。我要使旷野变为湿润的树林，使干旱之地成为水流渠道。我要在干旱之地栽植香柏树、黄杨树、桃金娘树、柏树、榆树和白杨树，好叫他们看见、明白、知道并一同相信，这是上主的手所做的，是以色列的圣者所显示的。}\par

% structure: p0010 source=h2 target=subsection reason=relative-heading-rank; base=h2

\subsection{五、基督同时是天使与天主}

在\BibleBook{}{创世纪}中，对亚巴郎说：\BibleQuote{上主的天使从天上呼唤他，对他说：亚巴郎，亚巴郎！他回答说：我在这里。天使说：不可在这孩子身上下手，不要伤害他。现在我知道你敬畏你的天主，因为你为了我的缘故，没有怜惜你的儿子，你所爱的独子。}\newline{}同样在那里，对雅各伯说：\BibleQuote{上主的天使在梦中对我说话：我是天主，你曾在神的地方看见我，在那里你为我竖立了石柱，并向我许了愿。}\newline{}同样在\BibleBook{}{出谷纪}中：\BibleQuote{白天，天主在云柱中走在他们前面，给他们领路；夜间在火柱中与他们同行。}\newline{}之后在同一处：\BibleQuote{天主的天使向前移动，走在以色列子民军队的前面。}\newline{}同样在那里：\BibleQuote{看，我派遣我的天使在你前面，在路上保护你，领你进入我所预备的地方。你们要留心他，听从他的话，不可违抗他，因为他必不放弃你。因为我的名在他内。}\newline{}因此主自己在福音中说：\BibleQuote{我因我父的名而来，你们却不接待我；若有人因自己的名而来，你们倒要接待他。}\newline{}又在\WorkTitle{圣咏}第一一八篇中说：\BibleQuote{奉上主之名而来的，当受赞美。}\newline{}同样在\BibleBook{}{玛拉基亚}中：\BibleQuote{我与肋未的生命与平安的盟约；我赐给他敬畏，使他敬畏我，因我的名而战栗。真理的法律在他口中，不义未见于他唇上。他用和平的言语纠正，与我们同行，使许多人远离不义。因为司铎的嘴唇应保存知识，人们要从他的口中寻求法律，因为他是全能者的天使。}\par

% structure: p0012 source=h2 target=subsection reason=relative-heading-rank; base=h2

\subsection{六、基督是天主}

在创世纪中：\Quote{天主对\Person{雅各伯}说：起来，上\Place{贝特耳}去，住在那里；在那里给曾在你逃避你兄弟\Person{厄撒乌}时显现给你的天主筑一座祭台。} 同样在\BibleBook{依撒意亚}{先知书}中：\Quote{万军的上主这样说：埃及劳碌，\Person{雇士}的货物，以及\Person{色巴}的伟人，都要归于你，做你的仆役；他们要带着锁链行走在你后面，要俯伏朝拜你，向你祈祷，因为天主在你中间，除你以外没有别的神。的确，你是天主，我们却不认识，以色列的天主，我们的救主。凡反对你的，都必蒙羞恐惧，陷入混乱。} 同样在本书中：\Quote{有声音呼号说：在旷野中预备上主的道路，修直我们天主的路。一切深谷要填满，一切山岳丘陵要削平，弯曲的要修直，崎岖的要开坦；上主的荣耀将显示，凡有血肉的都要看见天主的救恩，因为上主亲口说了。} 此外，在\BibleBook{耶肋米亚}{先知书}中：\Quote{这是我们的天主，没有别的可与他相比；他寻得了智慧的一切道路，将它赐给他的儿子\Person{雅各伯}和他所爱的\Person{以色列}。此后，他在地上出现，与人交谈。} 又在\BibleBook{匝加利亚}{先知书}中天主说：\Quote{他们要渡过狭海，击打海中的波浪，使河深的全部干涸；亚述的骄傲必受屈辱，埃及的权杖必被夺去。我必以他们的主天主巩固他们，他们必因他的名而夸耀，上主说。} 又在\BibleBook{欧瑟亚}{先知书}中主说：\Quote{我不再按我的盛怒行事，我不允许\Person{厄弗辣因}灭亡：因为我是天主，在你中间没有圣人；我不进入城中；我将追随天主。} 又在第四十四篇圣咏中：\Quote{天主，你的御座永远常存；你王国的权杖是正义的权杖。你喜爱正义，憎恶邪恶；因此天主，你的天主，用喜乐之油膏你，胜过你的同伴。} 同样在第四十五篇圣咏中：\Quote{你们要安静，知道我是天主。我要在万民中被高举，我要在地上被高举。} 又在第八十一篇圣咏中：\Quote{他们仍不知道，也不明白，在黑暗中行走。} 又在第六十七篇圣咏中：\Quote{你们应向天主歌唱，歌颂他的名；为那驾云者备路，天主是他的名。} 又在\BibleBook{若望}{福音}中：\Quote{在起初已有\OriginalTerm{圣言}{Word}，圣言与天主同在，圣言就是天主。} 同样在该福音中：\Quote{主对\Person{多默}说：伸出你的指头，看看我的手；不要无信，而要相信。\Person{多默}回答他说：我主！我天主！\Person{耶稣}对他说：你因看见了我才相信；那些没有看见而相信的，才是有福的。} 又\Person{保禄}致\Place{罗马人}书：\Quote{为了我的弟兄、我的骨肉之亲，就是被诅咒，与基督隔绝，我也愿意。他们是以色列人：义子的名分、光荣、盟约、法律的颁布、礼仪（天主的）以及恩许，都是他们的；圣祖也是他们的，并且基督按血统说，也是从他们来的，他是在万有之上，永远受赞美的天主。} 又在\WorkTitle{默示录}中：\Quote{我是\OriginalTerm{阿耳法和敖默加}{Alpha and Omega}，元始和终末。我要把生命的水白白赐给口渴的人。胜利者将拥有这些，并承受产业；我要做他的天主，他要做我的儿子。} 又在第八十一篇圣咏中：\Quote{天主站在众神会议中，在众神中施行审判。} 同样在该处：\Quote{我曾说：你们是神，都是至高者的儿子；但你们要像人一样死去。} 但如果那些义人，遵守了神圣诫命的人，可以被称为神，那么\Person{基督}、天主子，岂不更是天主！因此他自己在\BibleBook{若望}{福音}中说：\Quote{法律上不是记载着：\InnerQuote{我说过：你们是神}吗？如果那些承受天主圣言的人尚且被称为神——而圣经是不能废弃的——那么，父所祝圣并派遣到世界来的，你们就说：你说亵渎的话，因为我说了\InnerQuote{我是天主子}吗？但我若不做我父的工作，你们就不必信我；但我若做了，你们纵然不信我，也该信这些工作，从而知道父在我内，我在父内。} 又在\BibleBook{玛窦}{福音}中：\Quote{你要给他起名叫\OriginalTerm{厄玛奴耳}{Emmanuel}，意思是：天主与我们同在。}\par

% structure: p0014 source=h2 target=subsection reason=relative-heading-rank; base=h2

\subsection{七、基督、我们的天主，应来临，作为人类的照亮者和救主}

在\BibleBook{依撒意亚}{先知书}中：\BibleQuote{你们要安慰，你们软弱的手；你们无力的膝，要坚强。你们这些胆怯的心，不要害怕。我们的主必施行审判，祂必亲自来拯救我们。那时瞎子的眼要睁开，聋子的耳要听见。那时瘸子要跳跃如鹿，哑巴的舌头要说话；因为在旷野中有水涌出，在干地中有溪流。}同处又说：\Quote{不是长老，也不是天使，而是主亲自拯救他们，因为祂要爱他们，要怜恤他们，祂要亲自救赎他们。}又在同一处：\BibleQuote{我上主天主凭公义召叫了你，我必握住你的手，保护你；我要使你作人民的盟约，作外邦人的光明；为开启盲人的眼，从锁链中释放被囚的人，从黑暗的监狱中释放坐在黑暗中的人。我是上主天主，这是我的名。我决不将我的光荣让给另一个，也不将我的能力让给雕刻的偶像。}又在\BibleBook{}{圣咏集}第二十四篇中：\BibleQuote{上主，求你使我认识你的道路，教导我你的途径，并引我进入你的真理，教导我，因为你是我的救主的天主。}因此，在\BibleBook{若望}{福音}中，主说：\BibleQuote{我是世界的光；跟随我的，决不在黑暗中行走，必有生命的光。}此外，在\BibleBook{玛窦}{福音}中，天使加俾额尔对若瑟说：\BibleQuote{若瑟，达味之子，不要怕娶你的妻子玛利亚，因为那在她内受孕的是出于圣神。她要生一个儿子，你要给他起名叫耶稣，因为他要把自己的民族从罪恶中拯救出来。}也在\BibleBook{路加}{福音}中：\BibleQuote{匝加利亚充满了圣神，预言说：上主以色列的天主应受赞美，因祂眷顾救赎了祂的子民，并在祂的仆人达味家中为我们兴起了大能的救主。}在同一处，天使对牧羊人说：\BibleQuote{不要害怕！看，我给你们报告一个为全民族的大喜讯：今天在达味城里，为你们诞生了一位救世主，就是基督耶稣。}\par

% structure: p0016 source=h2 target=subsection reason=relative-heading-rank; base=h2

\subsection{八、虽然祂从起初就是天主子，但仍须按肉身再次受生}

在\BibleBook{}{圣咏集}第二篇中：\BibleQuote{上主对我说：\InnerQuote{你是我的儿子，我今日生了你。你向我请求，我便将万民赐你作产业，将地极赐你作基业。}}在\BibleBook{路加}{福音}中：\BibleQuote{当依撒伯尔听到玛利亚的问安，胎儿在她腹中跳跃，她便充满了圣神，大声呼喊说：\InnerQuote{你在妇女中是有福的，你的胎儿也是有福的。我主的母亲到我这里来，这事怎能临到我呢？}}保禄在\BibleBook{迦拉达}{书}中：\BibleQuote{但时期一满，天主就派遣了自己的儿子，生于女人。}在\BibleBook{若望}{一书}中：\BibleQuote{凡明认耶稣基督是在肉身内降世的神，便是出于天主。凡否认耶稣是在肉身内降世的，就不是出于天主，而是出于假基督的精神。}\par

% structure: p0018 source=h2 target=subsection reason=relative-heading-rank; base=h2

\subsection{九、这将成为祂诞生的标记：祂将由童贞女所生——既是人又是天主——人子与天主子}

在\BibleBook{依撒意亚}{先知书}中：\BibleQuote{上主又对阿哈次说：\InnerQuote{你向上主你的天主要一个征兆，或求诸高处，或求诸深处。}阿哈次说：\InnerQuote{我不求，我不试探上主。}依撒意亚说：\InnerQuote{达味家族，你们听着！你们使人厌烦还算小事，还要使我的天主厌烦吗？因此，我主要亲自给你们一个征兆：看，有位贞女要怀孕生子，给他起名叫厄玛奴耳。他要吃奶酪和蜂蜜，直到他知道弃恶择善的时候。}}天主曾预言那要踏碎魔鬼头颅的女人后裔。在\BibleBook{}{创世纪}中：\BibleQuote{天主对蛇说：\InnerQuote{因你做了这事，你在一切牲畜和野兽中是可咒骂的；你要用肚子爬行，毕生日日吃土。我要把仇恨放在你和女人之间，你的后裔和她的后裔之间；他的后裔要踏碎你的头颅，你要咬伤他的脚跟。}}\par

% structure: p0020 source=h2 target=subsection reason=relative-heading-rank; base=h2

\subsection{十、基督既是人又是天主，兼备双重本性，好成为我们与圣父之间的中保}

在\BibleBook{耶肋米亚}{先知书}中：\Quote{祂是人，谁能认识祂？}在\BibleBook{}{户籍纪}中：\BibleQuote{由雅各伯将出现一颗星，由以色列将兴起一位君王。}同处：\BibleQuote{由他的后裔将出现一位君王，他要统治许多民族；他的王国要升高如哥格，他的王国要扩大；天主要将他从埃及领出。他的荣耀有如独角兽，他要吞吃他仇敌的民族，要吸取他们肥髓的骨髓，用箭射穿他的仇敌。他伏卧如雄狮，有如幼狮；谁能使他起立？祝福你的，必受祝福；诅咒你的，必受诅咒。}在\BibleBook{依撒意亚}{先知书}中：\BibleQuote{上主的神临于我身上，因为祂给我傅了油，派遣我向贫穷人传报喜讯，向心灵破碎的人宣告释放，向俘虏宣告自由，向盲人宣告复明，宣布上主恩慈之年，和我们天主复仇的日子。}因此，在\BibleBook{路加}{福音}中，加俾额尔对玛利亚说：\BibleQuote{天使回答她说：\InnerQuote{圣神要临于你，至高者的能力要庇荫你，因此那要诞生的圣者，将称为天主子。}}在\BibleBook{格林多}{前书}中：\BibleQuote{第一个人出于地，属于土；第二个人出于天。那属于土的怎样，凡是属于土的也怎样；那属于天的怎样，凡是属于天的也怎样。我们既然肖似那属于土的，也要肖似那属于天的。}\par

% structure: p0022 source=h2 target=subsection reason=relative-heading-rank; base=h2

\subsection{十一、基督应按肉身生于达味的后裔}

在\BibleBook{}{列王纪下}中：\BibleQuote{主的话传于\Person{纳堂}说：你去告诉我的仆人\Person{达味}，上主这样说：你不可为我建造殿宇居住；但当你的日子满期，你与你的祖先同眠时，我必兴起你腰中所出的后裔，我要坚定他的国。他要为我建造一座殿宇，我要永远坚定他的宝座；我要作他的父亲，他要作我的子；他的家将要坚定，他的国在我面前直到永远。} 在\BibleBook{}{依撒意亚}中也如此说：\BibleQuote{从\Person{叶瑟}的根茎要发出一条嫩枝，从他的根上将生出一朵花；上主的神要住在他身上，就是智慧和明达的神，超见和刚毅的神，知识和孝爱的神；敬畏上主的神要充满他。} 在第一百三十一篇\BibleBook{}{圣咏}中：\BibleQuote{天主向\Person{达味}起了誓，决不反悔；我要从你的腹中，立你的后裔，坐我的宝座。} 在\BibleBook{路加}{福音}中：\BibleQuote{天使对她说：\Person{玛利亚}，不要害怕！因为你在天主前获得了宠幸。看，你将怀孕生子，并要给他起名叫\Person{耶稣}。他将要伟大，并被称为至高者的儿子；上主天主要把他祖先\Person{达味}的宝座赐给他，他要为王统治\Person{雅各伯}家直到永远，他的王权没有终结。} 在\BibleBook{}{默示录}中：\BibleQuote{我看见坐在宝座上的天主右手中有书卷，里面和背面都写着字，用七个印密封着；我又看见一位强壮的天使大声呼喊：谁配得拿起那书卷，并揭开它的印？在天上、地上、地下，没有谁能打开书卷，也不能观看它。我就大哭，因为没有人配得打开书卷，也不能观看它。长老中有一位对我说：不要哭！看，犹大支派的狮子，\Person{达味}的根，已得胜，能打开那书卷和它的七个印。}\par

% structure: p0024 source=h2 target=subsection reason=relative-heading-rank; base=h2

\subsection{十二、基督应在白冷诞生}

在\BibleBook{}{米该亚}中：\BibleQuote{你，\Place{白冷}——\Place{厄弗辣大}！你在\Place{犹大}各城中并非最小，因为将有一位领袖从你那里出来，来统治\Place{以色列}；他的根源是从亘古，从太初就有。} 在\BibleBook{玛窦}{福音}中：\BibleQuote{当\Person{耶稣}在\Place{犹大}的\Place{白冷}诞生时，在\Person{黑落德}王的日子里，看，贤士从东方来到\Place{耶路撒冷}，说：那生下来作犹太人的君王在哪里？我们在东方看见了他的星，特来朝拜他。}\par

% structure: p0026 source=h2 target=subsection reason=relative-heading-rank; base=h2

\subsection{十三、基督在首次来临时应处于卑微状态}

在\BibleBook{}{依撒意亚}中：\BibleQuote{主啊，谁相信了我们的报道？上主的臂膀向谁启示了？我们在祂面前如同儿童，像旱地里的根苗。祂没有威仪，也没有光荣；我们看见祂，祂没有形貌，也没有美丽；但祂的形貌毫无尊荣，比众人更缺乏。祂是一个身患灾病的人，知道如何忍受软弱；因为祂的脸被转开，祂被轻视，不被重视。祂担当了我们的罪，为我们忧愁；我们却以为祂是在受苦、受打击、受磨难。但祂是为了我们的过犯而被刺伤，为了我们的罪恶而被压伤；祂受的惩罚使我们获得和平，因祂的创伤我们得了痊愈。我们都像羊一样迷了路，各人偏行己路。天主却将祂的罪归于祂；祂因为受苦，却不开口。} 在同一卷书中：\BibleQuote{我没有悖逆，也没有反抗。我把我的背给鞭打的人，我的颊给掌击的人；我没有把我的脸躲避唾污，天主是我的援助。} 在同一卷书中：\BibleQuote{祂不喧嚷，也不扬声，在街市上也听不见祂的声音。压伤的芦苇祂不折断，将熄的灯心祂不吹灭；但祂要真实地施行公理。祂要发光，并不动摇，直到祂在世上设立了公理，海岛都要仰望祂的名。} 在\BibleBook{}{圣咏}第二十一篇中：\BibleQuote{但我是虫，不是人；是人的耻辱，被百姓所蔑视。凡看见我的都嘲笑我，他们撇嘴摇头说：他依靠上主，让上主救他吧！既然他喜悦上主，就让上主拯救他吧！} 在同一处：\BibleQuote{我的力量如同瓦片被晒干，我的舌头紧贴着我的牙床。} 在\BibleBook{}{匝加利亚}中：\BibleQuote{上主使我看见大司祭\Person{耶稣}，站在上主的天使面前，魔鬼站在他的右边控告他。\Person{耶稣}穿着污秽的衣服，站在天使面前；天使回答那些站在他面前的说：脱去他污秽的衣服。且对他说：看，我已消除了你的罪恶。并给他穿上礼服，将洁净的冠冕戴在他头上。} \Person{保禄}在\BibleBook{斐理伯}{书}中说：\BibleQuote{祂虽具有天主的形象，却不以自己与天主同等为抢夺的，反而空虚自己，取了奴仆的形象，成为人的模样；且以人的形状出现，祂谦卑自下，服从至死，甚至死在十字架上。为此，天主极其举扬祂，赐给祂一个名字，超越一切名字，使得天上、地上和阴间的一切，因\Person{耶稣}的名字而屈膝，并要众口宣认\Person{耶稣基督}是主，以光荣天主圣父。}\par

% structure: p0028 source=h2 target=subsection reason=relative-heading-rank; base=h2

\subsection{十四、祂是犹太人要处死的义人}

在\BibleBook{}{智慧篇}中：\BibleQuote{让我们抓住义人，因为祂令我们厌烦，与我们的行为对立，并指责我们违反法律。祂声称认识天主，并自称是天主子；祂成为我们思想的揭露者；连看祂都令我们痛苦，因为祂的生活与别人不同，祂的行径改变了。我们被祂视为轻浮，祂回避我们的道路，如同避开不洁；祂赞美义人的最终结局，并夸耀有天主为祂的父。那么，我们来看祂的话是否真实，试试祂的结局如何。我们要以羞辱和酷刑审问祂，好知道祂的敬畏，考验祂的忍耐。我们要以最可耻的死亡定祂的罪。他们这样想，就错了。因为他们的恶意蒙蔽了他们，他们不认识天主的圣事。}同样在\BibleBook{依撒意亚}{先知书}中：\BibleQuote{看，义人怎样灭亡，没有人理解；义人被取去，没有人留意。因为义人从不义面前被取去，他的安葬将在平安中。}关于这事，早在\BibleBook{}{出谷纪}中已预言：\BibleQuote{不可杀害无辜和义人。}同样在\BibleBook{}{福音}中：\BibleQuote{犹大被后悔所驱使，对司祭和长老说：我出卖了无辜的血，犯了罪。}\par

% structure: p0030 source=h2 target=subsection reason=relative-heading-rank; base=h2

\subsection{十五、基督被称为将被宰杀的绵羊和羔羊，以及论及苦难的\OriginalTerm{圣事（奥迹）}{sacrament (mystery)}}

在\BibleBook{依撒意亚}{先知书}中：\BibleQuote{祂被牵走，如同绵羊被牵去宰杀，又像羔羊在剪毛者前默不作声，祂没有开口。在祂的卑微中，祂的审判被夺去：谁能述说祂的世代？因为祂的生命从地上被夺去。因我百姓的过犯，祂被引到死亡；我使恶人作祂的坟墓，使富人作祂的死亡；因为祂没有行过强暴，口中也没有诡诈。因此祂要使许多人获得，并分得强者的战利品；因为祂将灵魂交付于死亡，被列于罪犯之中。祂承担了多人的罪，并为他们的过犯被交出。}同样在\BibleBook{耶肋米亚}{先知书}中：\BibleQuote{主，求祢赐我知识，我就知道；随即我看见他们的计谋。我好像一只驯良的羔羊被牵去宰杀；他们设计谋害我，说：来，我们把树木扔在祂的面饼上，将祂从活人之地剪除，使祂的名字不再被纪念。}同样在\BibleBook{}{出谷纪}中，天主对梅瑟说：\BibleQuote{让他们各家取一只羊，按家族，一只无瑕疵、完全的、公的、一岁的羊羔。从绵羊和山羊中取来，以色列子民全会众要在黄昏时宰杀；取一些血，涂在房屋的两边门框和门楣上，在他们吃羔羊的房屋里。当夜要吃肉，用火烤；与无酵饼和苦菜同吃。不可吃生的，也不可水煮，只可火烤；头与腿和内脏一起烤。不可留到早晨；骨头一根也不可折断。剩下的到早晨要用火烧。你们要这样吃：腰间束带，脚上穿鞋，手中拿杖，快快地吃：这是主的逾越节。}同样在\BibleBook{}{默示录}中：\BibleQuote{我看见在宝座和四个活物中间，并在长老中间，有一只羔羊站立，好像被杀过的，有七角七眼，就是天主的七神，奉差遣往普天下去。祂前来，从坐宝座者的右手中取了书卷。祂取了书卷，四个活物和二十四位长老就俯伏在羔羊面前，各拿着琴和盛满了香的金炉，这香就是圣人们的祈祷。他们唱新歌，说：主，祢配取书卷，揭开它的印；因为祢曾被宰杀，用祢的血从各支派、各方言、各民族、各邦国中赎回了我们，使我们成为天主的国度，并使我们成为司祭，他们要在地上为王。}同样在\BibleBook{若望}{福音}中：\BibleQuote{第二天，若翰看见耶稣向他走来，就说：看，天主的羔羊，除去世人罪的！}\par

% structure: p0032 source=h2 target=subsection reason=relative-heading-rank; base=h2

\subsection{十六、基督也被称为石头}

在\BibleBook{依撒意亚}{}中：\BibleQuote{上主这样说：看，我要在\Place{熙雍}的地基上安放一块宝贵的石头，精选的，主要的，角石，尊贵的；信赖祂的人必不蒙羞。} 也在第一一八篇圣咏中：\BibleQuote{匠人弃而不用的石头，反而成了屋角的基石；这是上主的作为，在我们眼中神妙莫测。这是上主所安排的一天，我们要欢欣鼓舞。上主，求祢救助！上主，求祢赐福！奉上主之名而来的，应受赞颂。} 也在\BibleBook{匝加利亚}{}中：\BibleQuote{看，我要引出我的仆人，\OriginalTerm{苗芽}{Orient}是他的名字，因为我要在\Person{约叔亚}面前安置一块石头，这块石头有七只眼睛。} 也在\BibleBook{申命纪}{}中：\BibleQuote{你要将这一切法律极其清楚地写在石头上。} 也在\BibleBook{若苏厄}{}中：\BibleQuote{他拿来一块大石头，立在上主面前；\Person{若苏厄}对人民说：看，这块石头将为你们作证，因为它听到了上主今天向你们所说的所有话；它将在末日为你们作证，当你们离开你们的天主时。} 也在\BibleBook{宗徒大事录}{}中，\Person{伯多禄}说：\BibleQuote{诸位百姓的首领和以色列的长老，请听：我们今日因为向一个跛子行了善事，他怎样痊愈了，被你们盘问。你们和全体以色列百姓都应知道：凭\Place{纳匝肋}人\Person{耶稣基督}的名，即你们所钉死，天主从死者中所复活的那位，站在你们面前痊愈了。除祂以外，别无救援，因为在天下人间，没有赐予别的名字，使我们赖以得救的。这石头就是你们匠人所弃而不用的，反而成了屋角的基石。} 这就是\BibleBook{创世纪}{}中\Person{雅各伯}放在头下的石头，因为人的头是\Person{基督}；当他睡觉时，梦见一架梯子直达天上，上主立在梯子上，天使们上去下来。\Person{雅各伯}将这石头祝圣并敷以傅油圣事，象征基督。这就是\BibleBook{出谷纪}{}中\Person{梅瑟}坐在山顶上的石头，当纳维的儿子\Person{若苏厄}与\Place{阿玛肋克}作战时；借着石头的圣事和\Person{梅瑟}坐姿的稳固，\Place{阿玛肋克}被\Person{若苏厄}打败，即魔鬼被\Person{基督}打败。这就是\BibleBook{撒慕尔纪上}{}中的大石头，当牛用车子将约柜送回来时，外邦人送还后，约柜放在那石头上。这也是\BibleBook{撒慕尔纪上}{}中\Person{达味}用以击杀\Person{哥肋雅}的额头并将其杀死的石头；象征魔鬼及其仆役因此被击倒——即他们尚未被盖印的那部分头被战胜。借着这印记，我们也永远安全存活。这就是以色列打败外邦人后，\Person{撒慕尔}所立并称其为\OriginalTerm{厄本厄则尔}{Ebenezer}的石头；即帮助之石。\par

% structure: p0034 source=h2 target=subsection reason=relative-heading-rank; base=h2

\subsection{十七、后来这块石头将成为一座山，并充满全地}

在\BibleBook{达尼尔}{}中：\BibleQuote{看，有一尊巨大的像；这像的形状非常可怕，它直立在你面前；它的头是纯金的，胸和臂是银的，腹和腿是铜的，脚一部分是铁，一部分是泥，直到有一块石头从山中凿出，没有凿石者的手，击打在像的泥铁脚上，把它们砸得粉碎。于是铁、泥、铜、银、金都混在一起，变得像夏天禾场上的糠秕，被风吹散，什么也没有留下。那块击打像的石头变成了一座大山，充满全地。}\par

% structure: p0036 source=h2 target=subsection reason=relative-heading-rank; base=h2

\subsection{十八、在末世，这同一座山将要显现，外邦人将要来到其上，所有义人也将登上它}

在\BibleBook{依撒意亚}{}中：\BibleQuote{在末世，上主的山必要显出，天主的殿宇将在群山之巅；它将被高举于丘陵之上，万民都要涌向它，许多人要行走并说：来，我们登上上主的山，来到\Person{雅各伯}天主的殿；祂必将祂的道路指教我们，我们必行在祂的路上。因为法律将从\Place{熙雍}发出，上主的话将从\Place{耶路撒冷}而来；祂将治理万民，指责许多民族；他们要把刀剑打成犁头，把枪矛打成镰刀，不再学习战斗。} 也在第二十三篇圣咏中：\BibleQuote{谁能登上上主的山？谁能站在祂的圣所？那手洁心清的人，不慕虚幻，不发虚伪誓言。他必蒙上主赐福，由拯救他的天主获得宽仁。这就是寻求祂的族类，寻求\Person{雅各伯}天主面容的人。}\par

% structure: p0038 source=h2 target=subsection reason=relative-heading-rank; base=h2

\subsection{十九、基督是新郎，以教会为新娘，由此要生出属灵的子女}

在 \BibleBook{岳厄尔}{} 中：\BibleQuote{你们要在熙雍吹号角，制定斋戒，召集盛会；聚集民众，圣化教会，召集长老，聚集吃奶的婴孩；愿新郎从洞房出来，新娘从内室出来。} 也在 \BibleBook{耶肋米亚}{} 中：\BibleQuote{我必从犹大城邑和耶路撒冷街市上，除掉欢乐的声音、喜乐的声音、新郎的声音和新娘的声音。} 也在第十八篇 \BibleBook{}{圣咏} 中：\BibleQuote{他如同新郎走出洞房，又像壮士欣然奔路。他从天这边出现，绕到天那边，没有什么能躲过他的热力。} 也在 \BibleBook{默示录}{} 中：\BibleQuote{你来，我要把新妇，就是羔羊的妻子，指给你看。他就在圣神内带我到了一座大山上，又指给我看从天上由天主那里降下的圣城耶路撒冷，具有天主的荣耀。} 也在 \BibleBook{若望}{福音} 中：\BibleQuote{你们为我作证，我曾对从耶路撒冷派到我那里的人说过：我不是基督，我只是被派在他前面。那有新娘的是新郎；但新郎的朋友，站着听他，因新郎的声音而欢喜快乐。} 这事的神秘意义曾在农的儿子\Person{若苏厄}身上显示出来：当时命令他脱下鞋，无疑是因为他自己不是新郎。因为在法律中，凡拒绝婚姻的人应脱下鞋，而要作新郎的人则应穿鞋：\BibleQuote{当若苏厄在耶里哥时，他举目观看，见一个人站在他面前，手里拿着一把拔出的剑，对他说：你是帮助我们，还是帮助我们敌人？那人说：我是上主军队的统帅，现在我来。若苏厄便俯伏在地，对他说：我主，你要你仆人做什么？上主军队的统帅对若苏厄说：将你脚上的鞋脱下，因为你所站的地方是圣地。} 也在 \BibleBook{}{出谷纪} 中，梅瑟被命令脱下鞋，因为他也不是新郎：\BibleQuote{上主的使者从荆棘丛中的火焰里显现给他。他看见荆棘被火烧着，却没有烧毁。梅瑟说：我要过去看这大异象，为什么荆棘没有烧毁。上主天主见他走来观看，便从荆棘中叫他，说：梅瑟，梅瑟。他说：我在这里。他说：不可到这边来，除非你将脚上的鞋脱下，因为你所站的地方是圣地。又对他说：我是你父亲的天主，亚巴郎的天主，依撒格的天主，雅各伯的天主。} 这也在 \BibleBook{若望}{福音} 中显明：\BibleQuote{若翰回答他们说：我用水施洗，但有一位站在你们中间，你们却不认识他；他是在我以后来的，我却在他以前，我连解他的鞋带也不配。} 也在 \BibleBook{路加}{福音} 中：\BibleQuote{你们要束上腰，点上灯，如同仆人等待主人从婚筵回来，他来到叩门时，就立刻给他开门。主人来到时，看见那些仆人警醒，他们就有福了。} 也在 \BibleBook{默示录}{} 中：\BibleQuote{上主全能天主为王了！让我们欢喜快乐，将荣耀归于他，因为羔羊的婚期到了，他的新娘也准备好了。}\par

% structure: p0040 source=h2 target=subsection reason=relative-heading-rank; base=h2

\subsection{二十、犹太人将把基督钉在十字架上}

在 \BibleBook{依撒意亚}{} 中：\BibleQuote{我整天向悖逆、反抗我的百姓伸出我的手，他们行走不善之道，随从自己的罪恶。} 也在 \BibleBook{耶肋米亚}{} 中：\BibleQuote{来，让我们把树投到他的饼里，从地上除灭他的生命。} 也在 \BibleBook{}{申命纪} 中：\BibleQuote{你的生命必在你眼前悬而不定；你昼夜恐惧，不敢自信能活下去。} 也在第二十一篇 \BibleBook{}{圣咏} 中：\BibleQuote{他们刺穿了我的手和我的脚，数尽了我的骸骨。他们注视着我，看见了我，分了我的衣裳，为我的外衣拈阄。但是上主，求你不要远离我，快来援助我。求你从刀下救出我的灵魂，从狗爪下救出我唯一的生命。求你救我脱离狮子的口，并使我从野牛角下获得卑微。我要向我的弟兄宣扬你的名，在教会中我要赞美你。} 也在第一一八篇 \BibleBook{}{圣咏} 中：\BibleQuote{因敬畏你，用钉子穿透我的肉。} 也在第一四〇篇 \BibleBook{}{圣咏} 中：\BibleQuote{我的双手高举，如同晚祭。} 关于这祭献，索福尼亚说：\BibleQuote{在上主天主面前恐惧吧！因为他的日子临近了，上主预备了祭献，圣化了他所拣选的人。} 也在 \BibleBook{匝加利亚}{} 中：\BibleQuote{他们要仰望那被他们刺透的我。} 也在第八十七篇 \BibleBook{}{圣咏} 中：\BibleQuote{上主，我终日呼求你，向你伸出我的双手。} 也在 \BibleBook{}{户籍纪} 中：\BibleQuote{天主并非人，不会被悬吊，人子也不受威胁。} 因此，主在福音中说：\BibleQuote{梅瑟怎样在旷野里举起了蛇，人子也必照样被举起来，使凡信子的人得永生。}\par

% structure: p0042 source=h2 target=subsection reason=relative-heading-rank; base=h2

\subsection{二十一、在苦难与十字架的标记中蕴含一切德能与权能}

在\BibleBook{哈巴谷}{先知书}中：\BibleQuote{祂的威能遮蔽了诸天，大地充满了对祂的赞美，祂的光辉必如光明；祂手中必有角。在那里，祂荣耀的威能得以建立，祂奠定了祂坚强的爱。在祂面前，圣言必前行，并依祂的步履，往平原上去。} 在\BibleBook{依撒意亚}{先知书}中也说：\BibleQuote{看，有一个婴孩为我们诞生了，赐给了我们一个儿子，权柄必担在祂的肩上；祂的名要称为奇谋的使者。} 藉此十字圣号，若苏厄藉梅瑟征服了阿玛肋克。在\BibleBook{}{出谷纪}中，梅瑟对若苏厄说：\BibleQuote{你从人中选择出征的人，明天出去与阿玛肋克交战。看，我要站在山顶上，手中拿着天主的棍杖。梅瑟一举手，以色列就获胜；但梅瑟一放手，阿玛肋克就占优势。梅瑟的手沉重了，他们就搬了一块石头放在他下面，他坐在上面；亚郎和胡尔一边一个托着他的手，直到日落。若苏厄击溃了阿玛肋克和他的人民。上主对梅瑟说：你要将这事写在书上作纪念，也叫若苏厄听见，因为我要将阿玛肋克的纪念从天下完全抹去。}\par

% structure: p0044 source=h2 target=subsection reason=relative-heading-rank; base=h2

\subsection{二十二、在此十字圣号，为所有额上印有这记号的人，获得救恩}

在\BibleBook{厄则克耳}{先知书}中，上主说：\BibleQuote{你要走遍耶路撒冷，在那些为城中一切可憎之事而叹息哀哭的人额上划一个\Quote{I}记号。} 又在同处说：\BibleQuote{你们去击杀，不要顾惜你们的眼睛。不要怜恤老人、少年、处女、小孩和妇女，将他们完全毁灭。但任何身上有记号的人，你们不可触碰，而要从我的圣所开始。} 在\BibleBook{}{出谷纪}中，天主对梅瑟说：\BibleQuote{这血要在你们所住的房屋上作记号；我一见这血，就必保护你们。当我打击埃及地时，你们中间就不会有毁灭的灾殃。} 在\BibleBook{}{默示录}中：\BibleQuote{我看见羔羊站在熙雍山上，同祂一起的有十四万四千人，都有祂的名号和祂父的名号写在额上。} 又在同处说：\BibleQuote{我是阿耳法和敖默加，最初的和最末的，元始和终结。那些遵守祂诫命的人是有福的，为使他们有权柄到生命树那里。}\par

% structure: p0046 source=h2 target=subsection reason=relative-heading-rank; base=h2

\subsection{二十三、在祂受难时，正午应有黑暗}

在\BibleBook{亚毛斯}{先知书}中：\BibleQuote{到那一天，上主说：太阳要在正午落下，光明之日变为黑暗；我要将你们的庆节变为悲哀，将你们的歌曲变为哀歌。} 在\BibleBook{耶肋米亚}{先知书}中也说：\BibleQuote{她惊慌失措，因为她生了儿女，她的灵魂疲乏。她的太阳在正午落下，她蒙受羞辱和诅咒；我要将其余的人交在敌人面前，使他们被刀剑所杀。} 在福音中：\BibleQuote{从第六时辰起，遍地都黑暗了，直到第九时辰。}\par

% structure: p0048 source=h2 target=subsection reason=relative-heading-rank; base=h2

\subsection{二十四、祂不应被死亡所胜，也不应留在阴府}

在\BibleBook{}{圣咏}第二十九篇中：\BibleQuote{上主，祢从阴府中把我的灵魂救回。} 又在\BibleBook{}{圣咏}第十五篇中：\BibleQuote{祢不将我的灵魂遗弃在阴府，也不让祢的圣者见到腐朽。} 又在\BibleBook{}{圣咏}第三篇中：\BibleQuote{我躺下安睡，又醒起来，因为上主扶持了我。} 又据\BibleBook{若望}{福音}：\BibleQuote{没有人夺去我的生命，是我自己舍弃的。我有权舍弃，也有权再取回。这是我由我父所接受的命令。}\par

% structure: p0050 source=h2 target=subsection reason=relative-heading-rank; base=h2

\subsection{二十五、祂应在第三日从死者中复活}

在\BibleBook{欧瑟亚}{先知书}中：\BibleQuote{过两天祂必使我们复生，第三天我们必兴起。} 在\BibleBook{}{出谷纪}中也说：\BibleQuote{上主对梅瑟说：你下去警告百姓，今天明天使他们圣洁，洗净他们的衣服，预备后天，因为第三天，上主要在西乃山上降临。} 在福音中：\BibleQuote{邪恶淫乱的世代寻求征兆，除了先知约纳的征兆外，再没有征兆给他们。正如约纳在大鱼腹中三天三夜，人子也要在地心里三天三夜。}\par

% structure: p0052 source=h2 target=subsection reason=relative-heading-rank; base=h2

\subsection{二十六、在祂复活后，祂应从祂父手中接受一切权能，而祂的权能应永远长存}

在\BibleBook{达尼尔}{}书中：\BibleQuote{我在夜间见一神视，看见一位相似人子者，乘着天上的云彩而来，直到万古常存者前，被引到他面前。那些站在他旁边的人，把他引到他面前：他便被赐予王权，一切地上万民都要事奉他；他的王权是永远的，永不消逝，他的国度永不灭亡。}同样在\BibleBook{依撒意亚}{}书中：\BibleQuote{现在我要起来，上主说；现在我要受光荣，现在我要被高举，现在你们要看见，现在你们要明白，现在你们要羞愧。你们精神的力量是虚妄的：火要吞灭你们。}同样在第一百零九篇圣咏中：\BibleQuote{上主对我主说：你坐在我右边，等我使你的仇敌，作你的脚凳。上主从熙雍伸出你权力的棍杖：你要在仇敌中统治。}同样在\BibleBook{默示录}{}中：\BibleQuote{我转过身来，要看那同我说话的声音。我看见七个金灯台，在灯台中间有一位相似人子者，身穿长袍，胸间束着金带。他的头和头发像白羊毛、像雪；他的眼睛像火焰；他的脚像在火炉中锻炼过的铜；他的声音像大水的响声。他右手拿着七颗星；从他口中射出一把双刃的利剑；他的面容好像烈日发光。我一看见他，就仆倒在他脚前，像死人一样。他把右手按在我身上说：不要害怕！我是元始，我是终末，我是生活的，我曾死过，看，我永远活着，我拿着死亡和地狱的钥匙。}同样在福音中，主复活后对门徒说：\BibleQuote{天上地下的一切权柄都给了我，所以你们要去使万民成为门徒，因父及子及圣神之名给他们授洗，教训他们遵守我所吩咐你们的一切。}\par

% structure: p0054 source=h2 target=subsection reason=relative-heading-rank; base=h2

\subsection{二十七、除非经过祂的圣子耶稣基督，不可能到达天主父}

在福音中：\BibleQuote{我是道路、真理、生命，除非经过我，谁也不能到父那里去。}同样在同一处：\BibleQuote{我是门；谁若经过我进来，必将得救。}同样在同一处：\BibleQuote{许多先知和义人，曾希望看见你们所看见的，而没有看见；希望听见你们所听见的，而没有听见。}同样在同一处：\BibleQuote{信子的人，有永生；不信从子的人，不会有生命，反之天主的义怒常在他身上。}同样保禄致\BibleBook{厄弗所}{书}：\BibleQuote{他来了，就向你们远处的人传了和平的福音，也向近处的人传了和平的福音，因为藉着他，我们双方在一个圣神内，得以进到父面前。}同样致\BibleBook{罗马}{书}：\BibleQuote{因为所有的人都犯了罪，缺乏天主的荣耀，但他们赖他的恩赐和恩宠成义，藉着在基督耶稣内的救赎。}同样在\BibleBook{伯多禄前}{书}中：\BibleQuote{基督一次为了罪而死，义人代替不义的人，为把我们领到天主面前。}同样在同一处：\BibleQuote{为此，福音也曾传给死者，好使他们虽在人肉体方面受了审判，却在神方面靠天主生活。}同样在\BibleBook{若望一}{书}中：\BibleQuote{凡否认子的，也否认父；凡宣认子的，也拥有父和子。}\par

% structure: p0056 source=h2 target=subsection reason=relative-heading-rank; base=h2

\subsection{二十八、耶稣基督将来要作为审判者而来}

在\BibleBook{玛拉基亚}{}书中：\BibleQuote{看，主的日子来到，如烈火熊熊；所有骄傲的人和所有作恶的人，都要像碎秸；那日子一到，必要将他们烧尽，万军的上主说。}同样在第四十九篇（或第五十篇）圣咏中：\BibleQuote{上主，万神之神，已发言，从日出之地到日落之处，从熙雍，光辉的美饰，天主来临。我们的天主要来临，决不缄默；在他面前有烈火焚烧，在他四周有风暴旋转。他呼唤了上天下地，为审判他自己的百姓。你们应集合他的圣民，就是那些以祭献与他立约的人。诸天要宣布他的正义，因为审判者是天主。}同样在\BibleBook{依撒意亚}{}书中：\BibleQuote{万军的上主必出征，将战争粉碎；他必激起争斗，向他的敌人强力喊叫。我沉默了，难道我永远沉默？}同样在第六十七篇圣咏中：\BibleQuote{愿天主兴起，使他的仇敌四散，愿恨他的人由他面前逃跑。如烟消散，愿他们消散；如蜡从火前熔化，愿恶人在天主面前灭亡。但愿义人在天主面前欢乐，踊跃喜乐。请歌颂天主，赞美他的圣名；为那驾云者铺路。雅威是他的名。在他面前，他们必混乱，他是孤儿的慈父，是寡妇的审判者。天主在他的圣所内：天主使孤独者安居在家中，领被囚者出来，获享自由；唯有悖逆者住在荒芜之地。天主，当你率领百姓在旷野前行时。}同样在第八十一篇圣咏中：\BibleQuote{天主，求你起来，审判大地，因为你要在万国中消灭一切。}同样在\BibleBook{玛窦}{福音}中：\BibleQuote{我们与你有什么相干？达味之子！你为什么在这里提前来惩罚我们？}同样在\BibleBook{若望}{福音}中：\BibleQuote{父不审判任何人，而是把审判权全交给了子，为叫众人尊敬子如同尊敬父。不尊敬子的，就是不尊敬派遣他的父。}同样在保禄致\BibleBook{格林多后}{书}中：\BibleQuote{我们都必须出现在基督的审判台前，为使各人按他肉身所行的，或善或恶，领取相当的报应。}\par

% structure: p0058 source=h2 target=subsection reason=relative-heading-rank; base=h2

\subsection{二十九、祂将作为君王永远为王}

在匝加利亚中记载：\BibleQuote{你要向熙雍女子说：看，你的君王来到你面前：正义的，获得救恩的；谦逊的，骑在一头未驯服的驴上。} 同样在依撒意亚中：\BibleQuote{谁能向你宣示那永恒的地方？那行走在正义中，阻止自己的手接受礼物；掩住耳朵不听流血的事；闭上眼睛不看邪恶的人：这人必住在高崖的磐石洞穴中，必赐给他食粮，他的水必稳妥。你必亲眼看见光荣的君王。} 同样在玛拉基亚中：\BibleQuote{我是伟大的君王，上主说，我的名在万民中是显赫的。} 同样在第二篇圣咏中：\BibleQuote{但我由他建立在熙雍他的圣山上，被立为君王，宣布他的统治。} 同样在第二十一篇圣咏中：\BibleQuote{大地四极都要想起，并归向上主；万邦的众民族都要在你面前朝拜。因为王国属于上主，他要统治万民。} 同样在第二十三篇圣咏中：\BibleQuote{城门，请抬起你们的门楣！古老的门户，请高高抬起！光荣的君王要进来。谁是这位光荣的君王？是英勇有力的上主，是战场上的英雄。城门，请抬起你们的门楣！古老的门户，请高高抬起！光荣的君王要进来。谁是这位光荣的君王？万军的上主，他就是光荣的君王。} 同样在第四十四篇圣咏中：\BibleQuote{我的心涌现优美的言词：我向君王陈述我的作品；我的舌头像书法家灵巧的笔。你在世人中最为美丽，你的口唇充满恩宠，因此天主永远祝福你。英雄啊，请把剑佩在腰间！为了你的尊威和荣耀，请向前迈进，为真道、谦逊和正义而统治。} 同样在第五篇圣咏中：\BibleQuote{我的君王，我的天主，因为我向你祈祷。上主，早晨你必听见我的声音；早晨我站在你面前，仰望你。} 同样在第九十六篇圣咏中：\BibleQuote{上主为王，愿大地踊跃，愿无数岛屿欢乐！} 此外，在第四十四篇圣咏中：\BibleQuote{王后身穿金衣站在你右边，她服饰华丽，色彩缤纷。女儿，请听，请看，请侧耳！忘记你的人民和你的父家，因为君王恋慕你的美丽，他是你的上主天主。} 同样在第七十三篇圣咏中：\BibleQuote{但天主自古以来是我们的君王，他在大地之中施行了救恩。} 同样在玛窦福音中：\BibleQuote{当耶稣在犹大的白冷，在黑落德王的时候诞生后，看，有贤士从东方来到耶路撒冷，说：那生来作犹太人的君王在哪里？我们在东方看见了他的星，特来朝拜他。} 此外，根据若望福音，耶稣说：\BibleQuote{我的国不属于这世界；如果我的国属于这世界，我的臣民就会为我作战，使我不致被交给犹太人；但我的国不是从这来的。比拉多对他说：那么，你是君王了？耶稣回答说：你说的是，我是君王。我为此而生，也为此而来到世界上，为给真理作证；凡属于真理的，必听我的声音。}\par

% structure: p0060 source=h2 target=subsection reason=relative-heading-rank; base=h2

\subsection{三十、他自己就是审判者和君王}

在第七十一篇圣咏中：\BibleQuote{天主，请将祢的审判权赐给君王，将祢的正义赐给君王的儿子，使祂以正义审判祢的百姓。}在\WorkTitle{默示录}中：\BibleQuote{我看见天开了，看，有一匹白马；那坐在上面的，称为\OriginalTerm{忠信而真实}{Faithful and True}；祂凭正义审判和作战。祂的眼睛有如火焰，祂头上戴着许多冠冕；祂写有一个名号，除祂以外没有人知道；祂穿着溅了血的衣服，祂的名字称为\OriginalTerm{圣言}{the Word of God}。天上的军队也骑着白马，穿着洁白的细麻衣跟随祂。从祂口中出来一把双刃利剑，用来打击万民；祂要用铁杖牧放他们，并践踏全能天主烈怒的酒醡。在祂的衣服和大腿上写着名号：\OriginalTerm{万王之王，万主之主}{King of kings, and Lord of lords}。}同样在\WorkTitle{福音}中：\BibleQuote{当人子在自己的光荣中，与所有天使一同降来时，那时祂要坐在光荣的宝座上；万民都要聚集在祂面前，祂要把他们彼此分开，如同牧人分开绵羊和山羊一样；把绵羊放在右边，山羊放在左边。那时君王要对那些在祂右边的说：来吧，我父所祝福的，承受那自创世以来为你们预备的国度吧！因为我饿了，你们给了我吃的；我渴了，你们给了我喝的；我做客，你们收留了我；我赤身露体，你们给了我穿的；我患病，你们看顾了我；我在监里，你们来探望了我。那时义人要回答说：\InnerQuote{主啊！我们什么时候见祢饿了而给祢吃的？或渴了而给祢喝的？什么时候见祢做客而收留了祢？或赤身露体而给祢穿的？什么时候见祢患病或在监里而来探望过祢？}君王要回答他们说：我实在告诉你们：凡你们对我这些最小兄弟中的一个所做的，就是对我做的。然后祂又对那在左边的说：\InnerQuote{可咒骂的，离开我，到那为魔鬼及其使者预备的永火里去吧！因为我饿了，你们没有给我吃的；我渴了，你们没有给我喝的；我做客，你们没有收留我；我赤身露体，你们没有给我穿的；我患病或在监里，你们没有来探望我。}那时他们也要回答说：\InnerQuote{主啊！我们几时见祢饿了或渴了，或做客，或赤身露体，或患病，或在监里，而没有给祢效劳？}祂要回答他们说：我实在告诉你们：凡你们没有给这些最小中的一个所做的，就是没有给我做。这些人要进入永罚，而义人却要进入永生。}\par

\SourceRecord{Translated by Robert Ernest Wallis.；Ante-Nicene Fathers；Vol. 5.；Edited by Alexander Roberts, James Donaldson, and A. Cleveland Coxe.；Buffalo, NY: Christian Literature Publishing Co.,；1886.；<http://www.newadvent.org/fathers/050712b.htm>.}

% structure: p0001 source=h1 target=section reason=page-title

\section{\WorkTitle{第十二论，第三卷}}

西彼廉向他子基林诺问安。因你向主天主所显的信德与虔诚，爱子啊，你请求我从圣经中为你摘录一些有关我们学派宗教教导的要目，作为训诲；寻求一条简明的圣学课程，好使你的心志，既已顺从于天主，不致因冗长而众多的书卷感到厌倦，却能受天上箴规的概要所教导，拥有健全而丰富的纲要来滋养记忆。因我亏欠你丰厚而爱戴的服从，我做了你所愿的。我劳苦一次，使你不再恒常劳苦。因此，尽我微薄的能力所能包含，我收集了某些主的训诫与神圣教导，它们对读者既容易又有用，因少许事物浓缩于短小篇幅，既易于通读，又常可重温。爱子，我祝你永远全心安康。\par

% structure: p0003 source=h2 target=subsection reason=relative-heading-rank; base=h2

\subsection{一、论善工与慈惠的益处}

在\BibleBook{依撒意亚}{}中：祂说：\BibleQuote{你要大声呼喊，不要止息；扬起你的声音如同号角；向我的百姓宣告他们的过犯，向雅各伯家宣告他们的邪恶。他们天天寻求我，乐意明白我的道，好像行义的国民，未曾离弃他们天主的判断。他们向我求问公义，乐意亲近天主，说：\InnerQuote{什么？我们禁食，你未见之；我们刻苦己心，你不知之。}其实在你们禁食的日子，你们仍寻求自己的意愿；或者你们欺压属下的人，或者你们禁食却起争竞和恼怒，或者用拳头打邻舍。你们这样禁食，岂能使你们的声音听闻于上呢？这禁食并非我所拣选，除非人谦卑自己的心灵。你们若弯颈如环，铺下麻布和灰，也不可称为可悦纳的禁食。我所拣选的禁食，不是这样，——主说——而是松开不义之绳，解下轭上之索，使被欺压的得自由，折断一切不义的轭。将你的饼分给饥饿的人，将飘流的穷人接到你家中。见赤身的给他衣服遮体，不可嫌弃你骨肉之亲。那时，你的光明将如曙光初现，你的衣服也必迅速兴起；公义必在你前面行，主的荣耀必作你的后盾。那时你呼求，主必应允；你尚在说话，祂必说：我在这里。}\newline{}关于同一件事，在\BibleBook{约伯传}{}中：\BibleQuote{我拯救了穷乏人脱离强者的手；我帮助了孤儿，无人帮助他。寡妇的口曾祝福我，因我是盲人的眼，也是瘸子的脚，又是软弱人的父。}\newline{}同样在\BibleBook{多俾亚传}{}中：\BibleQuote{我对多俾亚说：我儿，你去从我们弟兄中寻一个全心记念天主的穷人，带他来，他必与我同享午餐。我儿，我等你直到你来。}\newline{}同处又说：\BibleQuote{我儿，在你一生所有的日子，要记念天主，不要违背祂的诫命。你要一生行义，不要走不义的道路；因为你若行事真诚，你的工作必蒙悦纳。要由你的财产中施舍，不要转脸不顾任何穷人。这样，天主的脸也必不转离你。我儿，你按照所有而行：你若多有，就多施舍；你若少有，就连那少的分施。施舍时不要害怕：你为自己积攒美好的报应，以备需用之日；因为施舍救人脱免死亡，且不容人进入黑暗。施舍为所有在至高天主面前行这事的人，是美好的善功。}\newline{}同一主题，在\BibleBook{箴言}{}中：\BibleQuote{怜悯穷人的，就是借给主。}\newline{}同处又说：\BibleQuote{周济穷人的，必不缺乏；转目不看的，必多受贫穷。}\newline{}同处又说：\BibleQuote{施舍与信德能涤除罪恶。}\newline{}同处又说：\BibleQuote{若你的仇人饿了，你要给他吃；若他渴了，你要给他喝：因为你这样行，就是把炭火堆在他头上。}\newline{}同处又说：\BibleQuote{水能灭火，施舍能除罪。}\newline{}同处又说：\BibleQuote{不要说：去，明天再来，我必给你；因为你现在能行善。因为你不知明日将发生何事。}\newline{}同处又说：\BibleQuote{塞耳不听穷人哀求的，他将来呼求天主，也无人听。}\newline{}同处又说：\BibleQuote{行为无瑕疵、行公义的人，遗下蒙福的儿女。}\newline{}在\BibleBook{德训篇}{}中：\BibleQuote{我儿，你若有余，当行善，并向主奉献相称的祭品；要记念死亡并不迟延。}\newline{}同处又说：\BibleQuote{要把施舍藏在穷人心中，这要为你恳求脱离一切灾祸。}\newline{}关于此事，在\BibleBook{圣咏}{第三十六篇}中，仁慈也惠及后代：\BibleQuote{我从前年幼，现在年老，从未见过义人被弃，也未见过他的后裔讨饭。他终日恩待人，借贷给人；他的后裔也蒙福。}\newline{}同一事，在\BibleBook{圣咏}{第四十篇}中：\BibleQuote{眷顾贫穷和困乏人的有福了：在患难之日，主必救他。}\newline{}又在\BibleBook{圣咏}{第一百一十一篇}中：\BibleQuote{他施舍钱财，周济穷人；他的仁义存到永远。}\newline{}同一事，在\BibleBook{欧瑟亚}{}中：\BibleQuote{我喜欢仁慈胜过祭献，喜欢认识天主胜过全燔祭。}\newline{}同一事，在\BibleBook{玛窦}{福音}中：\BibleQuote{饥渴慕义的人是有福的，因为他们要得饱饫。}\newline{}同处又说：\BibleQuote{怜悯人的人是有福的，因为他们要受怜悯。}\newline{}同处又说：\BibleQuote{要为自己积蓄财宝在天上，那里没有虫蛀、锈坏，也没有贼挖窟窿来偷；因为你的财宝在哪里，你的心也在那里。}\newline{}同处又说：\BibleQuote{天国好像一个商人，寻找好珠子；遇见一颗宝贵的珠子，就去变卖他一切所有的，买了它。}\newline{}即使微小的工作也有益处，同处又说：\BibleQuote{凡因门徒的名，只把一杯凉水给这些小子中的一个喝，我实在告诉你们，他的赏赐绝不会失去。}\newline{}施舍不可拒绝任何人，同处又说：\BibleQuote{凡求你的，就给他；有人想要向你借贷，不可推辞。}\newline{}同处又说：\BibleQuote{你若要进入永生，就当遵守诫命。他说：什么诫命？耶稣说：不可杀人，不可奸淫，不可作假见证，孝敬父母，又当爱人如己。那少年人说：这一切我都遵守了，还缺少什么呢？耶稣说：你若愿意成全，就去变卖你所有的，分给穷人，就必有财宝在天上；你还要来跟从我。}\newline{}同处又说：\BibleQuote{当人子在自己的光荣中，与众天使一同降来时，那时祂要坐在光荣的宝座上；万民都要聚集在祂面前；祂要把他们分别出来，好像牧羊人分别绵羊和山羊一样；把绵羊放在右边，山羊放在左边。于是君王要对那些在右边的说：你们这些蒙我父降福的，来承受自创世以来为你们预备的国度吧。因为我饿了，你们给了我吃的；我渴了，你们给了我喝的；我作客，你们收留了我；我赤身露体，你们给了我穿的；我患病，你们看顾了我；我在监里，你们来探望了我。那时义人回答祂说：主啊！我们什么时候见祢饿了而给祢吃？或渴了而给祢喝？什么时候见祢作客而收留？或赤身而给祢穿？我们什么时候见祢患病或在监里而来探望祢？君王回答他们说：我实在告诉你们，你们对我这些最小兄弟中的一个所做的，就是对我做的。然后祂要对在左边的说：你们这些可咒骂的，离开我，到那为魔鬼和他的使者预备的永火里去吧！因为我饿了，你们没有给我吃的；我渴了，你们没有给我喝的；我作客，你们没有收留我；我赤身露体，你们没有给我穿的；我患病或在监里，你们没有来探望我。那时他们也要回答：主啊！我们什么时候见祢饿了或渴了，或作客，或赤身露体，或患病，或在监里，没有服事祢呢？祂要回答他们说：我实在告诉你们，你们没有对我这些最小中的一个做的，就是没有对我做。这些人要进入永罚，而义人却要进入永生。}\newline{}关于同一事，在\BibleBook{路加}{福音}中：\BibleQuote{变卖你们所有的，施舍给人。}\newline{}同处又说：\BibleQuote{那造里面的，也造了外面。但你们施舍吧，看哪，一切对你们都是洁净的。}\newline{}同处又说：\BibleQuote{主啊，我把财产的一半给了穷人；我若欺骗过谁，就四倍偿还。耶稣对他说：今天救恩临到了这一家，因为他也是亚巴郎之子。}\newline{}同一事，在\BibleBook{格林多}{后书}中：\BibleQuote{愿你们的富余补足他们的缺乏，好使他们的富余也补足你们的缺乏，以求均平，正如所记载的：多收的没有余，少收的也没有缺。}\newline{}同处又说：\BibleQuote{少种少收，多种多收。各人要随本心所酌定的，不要作难，不要勉强，因为天主喜爱乐意奉献的人。}\newline{}同处又说：\BibleQuote{正如所记载的：他施舍钱财，周济穷人；他的仁义存到永远。}\newline{}同处又说：\BibleQuote{那赐种给撒种的，必供给粮食供你吃，且必增多你的种子，增加你仁义的果实，使你们凡事富足。}\newline{}同处又说：\BibleQuote{这供应的事，不但补足了圣徒的缺乏，而且因着许多感谢，也越发增添给天主。}\newline{}关于同一事，在\BibleBook{若望}{一书}中：\BibleQuote{凡有世上财物的，看见弟兄有需要，却关闭了慈心，天主的爱怎能在他里面呢？}\newline{}同一事，在\BibleBook{路加}{福音}中：\BibleQuote{你设宴或备晚餐，不要请你的朋友、弟兄、亲属和富有的邻舍，恐怕他们也要回请你，你就得了报答。但你摆设筵席，要请贫穷的、软弱的、瞎眼的、瘸腿的：你就有福了，因为他们无力报答你；但在义人复活时，你必得报答。}\par

% structure: p0005 source=h2 target=subsection reason=relative-heading-rank; base=h2

\subsection{二、论善工与施舍：即使因能力微小而做得较少，但意愿本身已足够}

在\BibleBook{格林多}{后书}中：\BibleQuote{若有乐意的意愿，必按人所有的蒙悦纳，不是按他所没有的；也不要叫别人轻松，你们却受累。}\par

% structure: p0007 source=h2 target=subsection reason=relative-heading-rank; base=h2

\subsection{三、论爱德与手足之情，当以虔诚和坚忍实践之}

在\BibleBook{玛拉基亚}{书}中：\BibleQuote{不是同一的天主创造了我们？我们众人不是同有一位父亲？为什么你们各人竟离弃自己的兄弟？}同样，按\BibleBook{若望}{福音}：\BibleQuote{我把平安留给你们，我将我的平安赐给你们。}同样在该处：\BibleQuote{这是我的命令：你们该彼此相爱，如同我爱了你们一样。人若为自己的朋友舍命，再没有比这更大的爱了。}同样在该处：\BibleQuote{缔造和平的人是有福的，因为他们要称为天主的子女。}同样在该处：\BibleQuote{我实在告诉你们：如果你们中有两个人在地上同心合意，无论为什么事祈祷，我在天之父必要成全他们。因为哪里有两个或三个人因我的名聚会，我就在他们中间。}同样，在\BibleBook{格林多}{前书}中：\BibleQuote{我原不能把你们当作属神的人，只得当作属血肉的人，当作在基督内的婴孩。我给你们喝的是奶，不是饭；因为那时你们还不能吃，就是如今你们还是不能，因为你们仍是属血肉的。你们中间既有嫉妒、纷争、分裂，岂不是属血肉的，并按照人的样子行事吗？}同样在该处：\BibleQuote{我若有全备的信心，甚至能移山，却没有爱德，我什么也不算。我若把我所有的一切施舍给人，并舍身被火焚烧，却没有爱德，为我毫无益处。爱德是含忍的，爱德是慈祥的，爱德不嫉妒，不夸张，不自大，不作无理的事，不求己益，不动怒，不记怨，不以不义为乐，却与真理同乐。凡事包容，凡事相信，凡事盼望，凡事忍耐。爱德永不止息。}同样，在\BibleBook{迦拉达}{书}中：\BibleQuote{你应爱你的近人如你自己。但如果你们彼此咬伤、吞食，你们要小心，免得彼此消灭。}同样，在\BibleBook{若望}{一书}中：\BibleQuote{天主的子女和魔鬼的子女，在此可以认出：凡不行义的人，不是出于天主，凡不爱自己弟兄的，也是如此。因为那不爱自己弟兄的，就是杀害弟兄的；你们知道：凡杀害弟兄的，没有永远的生命留在他内。}同样在该处：\BibleQuote{若有人说自己爱天主，却恨自己的弟兄，便是说谎者；因为那不爱自己所看见的弟兄的，怎能爱那看不见的天主呢？}同样，在\BibleBook{}{宗徒大事录}中：\BibleQuote{众信徒都一心一意，凡各人所有的，没有人说是自己的，都归公用。}同样，在\BibleBook{玛窦}{福音}中：\BibleQuote{若你到祭台前献礼物时，在那里想起你的弟兄有什么怨你的事，就把礼物留在祭台前，先去与你的弟兄和好，然后再来献你的礼物。}同样，在\BibleBook{若望}{一书}中：\BibleQuote{天主是爱，那住在爱内的，就住在天主内，天主也住在他内。}同样在该处：\BibleQuote{那说自己是在光明中，却恨自己弟兄的，是说谎者，并且至今仍在黑暗中行走。}\par

% structure: p0009 source=h2 target=subsection reason=relative-heading-rank; base=h2

\subsection{四、论不可自夸，因无物属己}

在\BibleBook{若望}{福音}中：\BibleQuote{人不能领受什么，除非是从天上赐予他的。}同样，在\BibleBook{格林多}{前书}中：\BibleQuote{你有什么不是领受的呢？既然是领受的，为什么你夸耀，好像不是领受的呢？}同样，在\BibleBook{撒慕尔}{上}中：\BibleQuote{不要夸耀，不要说高傲的话，也不要口出狂言，因为上主是洞悉一切的天主。}同样在该处：\BibleQuote{壮士的弓已被折断，衰弱者却腰缠力量。}同样，在\BibleBook{玛加伯}{下}中：\BibleQuote{服从天主本是公正的，有死的人不该想与天主相等。}同样在该处：\BibleQuote{不要害怕罪人的话，因为他的光荣要变为粪土和虫子。他今天被高举，明天却不见踪影，因为他已归于土里，他的思念也灭亡了。}\par

% structure: p0011 source=h2 target=subsection reason=relative-heading-rank; base=h2

\subsection{五、论凡事当保持谦逊与宁静}

在\BibleBook{依撒意亚}{先知书}中：\BibleQuote{上主天主这样说：上天是我的宝座，下地是我的脚凳。你们要为我建筑什么殿宇？或是我安息的地方在哪里？这一切都是我手所造的，这一切都是我的。我垂顾谁呢？除非是谦卑而安静的人，和因我的话而战栗的人。} 同一件事在\BibleBook{玛窦}{福音}中：\BibleQuote{温良的人是有福的，因为他们要承受土地。} 同样，在\BibleBook{路加}{福音}中：\BibleQuote{你们中间最小的，将成为伟大的。} 同一处又：\BibleQuote{凡高举自己的，将被贬抑；凡贬抑自己的，将被高举。} 同一事在\BibleBook{}{罗马书}中：\BibleQuote{不可心高气傲，但要战战兢兢；因为天主既然没有怜惜原来的树枝，也恐怕不怜惜你。} 同一事在\BibleBook{}{圣咏}第三十三篇中：\BibleQuote{他必拯救谦卑的人。} 又\BibleBook{}{罗马书}：\BibleQuote{向众人还清所欠的：该向谁纳税，就纳税；该向谁进贡，就进贡；该敬畏谁，就敬畏谁；该尊敬谁，就尊敬谁；除了彼此相爱外，你们不可再欠人什么。} 又在\BibleBook{玛窦}{福音}中：\BibleQuote{他们喜欢坐首席，喜爱会堂里的上座，喜爱市场上的问候，喜爱人们称他们为\InnerQuote{辣彼}。但你们不要被称为辣彼，因为你们的师傅只有一位。} 又在\BibleBook{若望}{福音}中：\BibleQuote{仆人不能大过主人，宗徒不能大过派遣他的人。你们既然知道这些事，如果实行，便是有福的。} 又在\BibleBook{}{圣咏}第八十一篇中：\BibleQuote{你们应对穷人和卑微的人施行正义。}\par

% structure: p0013 source=h2 target=subsection reason=relative-heading-rank; base=h2

\subsection{六、所有良善而正义的人所受的苦更多，但应当忍受，因为他们在受考验。}

在\BibleBook{撒罗满}{箴言}中：\BibleQuote{火炉试验陶匠的器皿，患难的试验锻炼义人。} 又在\BibleBook{}{圣咏}第五十篇中：\BibleQuote{天主，你喜爱的祭献是忏悔的精神；忏悔和谦卑的心，天主决不会轻视。} 又在\BibleBook{}{圣咏}第三十三篇中：\BibleQuote{天主亲近那些忏悔的人，他要拯救心灵谦卑的人。} 又在同一处：\BibleQuote{义人的苦难很多，但上主要从这一切中拯救他们。} 关于同一事在\BibleBook{约伯}{传}中：\BibleQuote{我赤身出自母胎，也将赤身归入地下；上主赐予，上主收回；随上主所喜悦的，就这样成就了；愿上主的名受赞美。在这一切事上，约伯在唇舌上并没有犯罪得罪上主。} 关于同一事在\BibleBook{玛窦}{福音}中：\BibleQuote{哀恸的人是有福的，因为他们要受安慰。} 又在\BibleBook{若望}{福音}中：\BibleQuote{我给你们讲了这些事，是要你们在我内有平安。但在世界上你们将受苦难；然而你们放心，我已战胜了世界。} 关于同一事在\BibleBook{格林多后}{书}中：\BibleQuote{有一根刺加在我身上，是撒殚的使者，击打我，不让我过于高举自己。为此我曾三次求主，使它离开我。但主对我说：我的恩宠为你足够了，因为我的能力在软弱中才得以完善。} 关于同一事在\BibleBook{}{罗马书}中：\BibleQuote{我们因希望天主的荣耀而欢跃。不仅这样，我们也因患难而欢跃，因为我们知道：患难生忍耐，忍耐生老练，老练生希望，希望不会使人蒙羞，因为天主的爱已藉着所赐予我们的圣神，倾注在我们心中。} 关于同一事在\BibleBook{玛窦}{福音}中：\BibleQuote{那引到死亡的街是多么宽广，路是多么宽阔，有许多人从那里进去；那引到生命的门是多么狭窄，路是多么狭小，找到他的人很少。} 又在\BibleBook{多俾亚}{传}中：\BibleQuote{你的义德在哪里？请看，你所受的是什么痛苦？} 又在\BibleBook{}{智慧篇}中：\BibleQuote{恶人猖獗的地方，义人叹息；但恶人灭亡时，义人必增多。}\par

% structure: p0015 source=h2 target=subsection reason=relative-heading-rank; base=h2

\subsection{七、我们不可使我们所领受的圣神忧郁。}

\Person{保禄}宗徒致\BibleBook{}{厄弗所书}：\BibleQuote{不要使天主的圣神忧郁，因为你们是在他内受了恩印，以待得救的日子。一切毒辣、怨恨、忿怒、争吵、毁谤，都要从你们中除掉。}\par

% structure: p0017 source=h2 target=subsection reason=relative-heading-rank; base=h2

\subsection{八、必须克服愤怒，免得它迫使我们犯罪。}

在\BibleBook{撒罗满}{箴言}中：\BibleQuote{忍耐的人胜于勇士；克制自己怒气的人，胜于攻克一座城。} 又在同一处：\BibleQuote{愚鲁的人当日就表露愤怒；精明的人却遮掩羞辱。} 关于同一事在\BibleBook{}{厄弗所书}中：\BibleQuote{你们生气，但不要犯罪；不要让太阳在你们的愤怒中落下。} 又在\BibleBook{玛窦}{福音}中：\BibleQuote{你们曾听见过对古人说：\InnerQuote{不可杀人}；谁若杀人，应受裁判。我却对你们说：凡无端向弟兄发怒的，都应受裁判。}\par

% structure: p0019 source=h2 target=subsection reason=relative-heading-rank; base=h2

\subsection{九、弟兄们应当彼此扶持。}

致\BibleBook{}{迦拉达书}：\BibleQuote{各人应存想他人，免得你们也受诱惑。你们应彼此担待重担，这样就满了基督的法律。}\par

% structure: p0021 source=h2 target=subsection reason=relative-heading-rank; base=h2

\subsection{十、我们只应信赖天主，并在他内夸耀。}

在\BibleBook{耶肋米亚}{}中：\BibleQuote{智者不要因他的智慧夸耀，壮士不要因他的勇力夸耀，富人不要因他的财富夸耀；夸耀的只应因这件事：他明白并认识我是主，我在地上施行慈悲、公道和正义，因为我喜爱这些，主说。} 关于同一事，在第五十四篇圣咏中：\BibleQuote{我寄望于主；我不害怕人能对我做什么。} 同书又言：\BibleQuote{我的灵魂只服从天主。} 又在第一百一十七篇圣咏中：\BibleQuote{我不害怕人能对我做什么；主是我的助佑。} 同书又言：\BibleQuote{信赖主胜过信赖人，寄望于主胜过寄望于王侯，这是好的。} 关于同一事，在\BibleBook{达尼尔}{}中：\BibleQuote{沙得辣客、默沙客和阿贝得乃哥回答拿步高王说：君王，关于这事，我们无需回答你。我们事奉的天主能拯救我们；他必从烈火窑中和你的手中拯救我们。即或不然，君王，你应当知道，我们决不事奉你的神，也不朝拜你所立的金像。} 同样在\BibleBook{耶肋米亚}{}中：\BibleQuote{凡信赖世人的人，是可咒骂的；凡信赖主的人，是有福的，他的希望必在于天主。} 关于同一事，在\BibleBook{申命纪}{}中：\BibleQuote{你当朝拜主你的天主，只事奉他。} 关于同一事，在\BibleBook{致罗马人}{书}中：\BibleQuote{他们朝拜事奉受造物，而离弃造物主。因此，天主任凭他们陷入可耻的情欲。} 关于同一事，也在\BibleBook{若望}{一书}中：\BibleQuote{那在你们内的，比那在世界上的更大。}\par

% structure: p0023 source=h2 target=subsection reason=relative-heading-rank; base=h2

\subsection{十一、凡已获得信赖、脱去了旧人者，只应注视天上的和属神的事物，并不留意他已弃绝的世界。}

在\BibleBook{}{依撒意亚}中：\BibleQuote{寻求主；找着了他时，就呼求他。但当他临近你时，愿恶人离弃他的道路，不义的人撇弃他的思想；愿他归向主，必蒙怜悯，因为他必多多赦免你的罪。}关于同一件事，在\BibleBook{}{训道篇}中：\BibleQuote{我见了太阳之下所行的一切事；看，全是虛空。}关于同一件事，在\BibleBook{}{出谷纪}中：\BibleQuote{你们要这样吃：束上腰，脚上穿鞋，手中拿棍，快快地吃；这是主的逾越节。}关于同一件事，在\BibleBook{玛窦}{福音}中：\BibleQuote{不要忧虑说：我们吃什么？喝什么？穿什么？因为这一切都是外邦人所求的。你们的天父知道你们需要这一切。你们先寻求天主的国和他的正义，这一切自会加给你们。}同样，在同一处：\BibleQuote{不要为明天忧虑，因为明天自有明天的忧虑；一天的难处一天当就够了。}同样，在同一处：\BibleQuote{手扶着犁向后看的人，不配进天主的国。}还有，在同一处：\BibleQuote{你们看天空的飞鸟：也不种，也不收，也不积蓄在仓里，你们的天父尚且养活它们。你们不比它们更贵重吗？}关于同一件事，按照\BibleBook{路加}{福音}：\BibleQuote{你们腰里要束上带，灯也要点着；自己好像仆人等候主人从婚筵回来；他来到叩门，就立刻给他开门。主人来到，看见仆人警醒，那仆人就有福了。}关于同一件事，在\BibleBook{玛窦}{福音}中：\BibleQuote{狐狸有洞，天空的飞鸟有窝，人子却没有枕头的地方。}还有，在同一处：\BibleQuote{凡不舍弃一切所有的，不能作我的门徒。}关于同一件事，在\BibleBook{格林多}{前书}中：\BibleQuote{你们不是自己的，因为你们是重价买来的。要在你们身上光荣天主，并承载天主。}同样，在同一处：\BibleQuote{时间是短促的。从此以后，那有妻子的要像没有妻子；哀哭的，要像不哀哭；快乐的，要像不快乐；置买的，要像无有所得；享用这世界的，要像不享用的；因为这世界的局面正在逝去。}同样，在同一处：\BibleQuote{第一个人出于地，属于土；第二个人出于天。那属于土的怎样，凡属于土的也怎样；那属于天的怎样，凡属于天的也怎样。我们既然有属土的肖像，也要有属天的肖像。}关于同一事，在\BibleBook{斐理伯}{书}中：\BibleQuote{众人只求自己的事，并不求基督耶稣的事。他们的结局是丧亡，他们的天主是自己的肚腹，以自己的羞辱为光荣，只思念地上的事。因为我们的家乡是在天上，我们也等待救主，就是主耶稣基督，他要改变我们卑贱的身体，使之相似他光荣的身体。}关于同一事，在\BibleBook{迦拉达}{书}中：\BibleQuote{但是我绝不夸耀，只夸耀我们主耶稣基督的十字架，因这十字架，世界于我已被钉死，我于世界也被钉死。}关于同一件事，在\BibleBook{弟茂德}{后书}中：\BibleQuote{凡当兵的，不让世务缠身，好使那招他当兵的人喜悦。人若在场上比武，非按规矩，就不能得冠冕。}关于同一件事，在\BibleBook{哥罗森}{书}中：\BibleQuote{你们若是与基督同死，脱离了世俗的元素，为什么仍像在世俗中活着，服从那些虛妄的事呢？}还有，关于同一件事：\BibleQuote{你们既然与基督一同复活了，就该追求天上的事，那里有基督坐在天主的右边。你们要思念天上的事，不要思念地上的事；因为你们已经死了，你们的生命与基督一同藏在天主内。当基督——你们的生命显现时，你们也要与他一同在光荣中出现。}关于同一件事，在\BibleBook{厄弗所}{书}中：\BibleQuote{脱去你们从前生活方式的旧人，这旧人是因私欲的迷惑而败坏的；但要在你们心思的灵里更新，并且穿上新人，这新人是按照天主在正义、圣洁和真理中所造的。}关于同一件事，在\BibleBook{伯多禄}{前书}中：\BibleQuote{你们要像外方人和旅客，戒绝肉身的私欲，这些私欲与灵魂交战；在外邦人中要品行端正，使那些毁谤你们是作恶的人，因看见你们的好行为，便在鉴察的日子归荣耀于天主。}关于同一件事，在\BibleBook{若望}{一书}中：\BibleQuote{那说自己住在基督内的，就该自己照主所行的去行。}同样，在同一处：\BibleQuote{不要爱世界，和世界上的事。谁若爱世界，父的爱就不在他内。因为凡世界上的事，就是肉身的私欲、眼目的私欲、以及人生的骄矜，都不是出于父，而是出于世界。世界和它的私欲都要过去；但那履行天主旨意的，却永远存在，正如天主永远存在一样。}同样，在\BibleBook{格林多}{前书}中：\BibleQuote{你们应把旧酵母除净，好使你们成为新的面团，如同你们原是無酵的一样。因为我们的逾越节羔羊基督，已被祭杀作了牺牲。所以我们过节，不可用旧酵母，也不可用奸詐和邪恶的酵母，但只用真诚和纯正的無酵饼。}\par

% structure: p0025 source=h2 target=subsection reason=relative-heading-rank; base=h2

\subsection{十二、不可发誓}

在撒罗满内：\Quote{常起誓的人必充满不义，瘟疫必不离其家；若妄发誓，他必不得成义。}\newline{}同一事，据\BibleBook{玛窦}{福音}：\Quote{（你们又听见对古人说：不可背誓，但要向主履行你们的誓愿。）我却对你们说：什么誓都不可起；（不可指着天起誓，因为它是天主的宝座；不可指着地起誓，因为它是祂的脚凳；也不可指着耶路撒冷起誓，因为它是大王的城；也不可指着你的头起誓，因为你不能使一根头发变白或变黑。）你们的话该当是：是就说是，非就说非；（因为多余的都出于邪恶。）}\newline{}同一事，在\BibleBook{}{出谷纪}：\Quote{不可妄呼上主你的天主的名。}\par

% structure: p0027 source=h2 target=subsection reason=relative-heading-rank; base=h2

\subsection{十三、不可诅咒}

在\BibleBook{}{出谷纪}：\Quote{不可诅咒，也不可辱骂你百姓的首长。}\newline{}又，在\BibleBook{}{圣咏}第三十三篇：\Quote{谁愿意生活，爱好幸福美善的日子？就应谨守口舌，不说坏话，嘴唇不欺诳。}\newline{}同一事，在\BibleBook{}{肋未纪}：\Quote{上主对梅瑟说：将那诅咒的人带到营外；所有听见的人，都按手在他头上，然后以色列全会众用石头砸死他。}\newline{}同一事，在保禄致\BibleBook{厄弗所}{书}：\Quote{一切坏话都不可出口，但只讲有助于信德建树的善言，使听者得恩宠。}\newline{}同一事，在致\BibleBook{罗马}{书}：\Quote{祝福，而不诅咒。}\newline{}同一事，在\BibleBook{玛窦}{福音}：\Quote{若向弟兄说\InnerQuote{傻瓜}，应受火狱的罚。}\newline{}同一事，据同一\BibleBook{玛窦}{福音}：\Quote{但我告诉你们：人所说的每句废话，在审判之日都要交账。因为要凭你的话定为义人，也凭你的话定为罪人。}\par

% structure: p0029 source=h2 target=subsection reason=relative-heading-rank; base=h2

\subsection{十四、总不可抱怨，关于所发生的一切事应赞美天主}

在\BibleBook{}{约伯传}：\Quote{\InnerQuote{你说句话反对天主，然后死罢！}但他注视她，说：\InnerQuote{你说话如同一个愚妇。我们从天主手中得福，难道就不能受祸吗？}在这一切所遭遇的事上，约伯嘴唇没有犯罪，在主前没有过错。}\newline{}又，同处：\Quote{你注意了我的仆人约伯么？世上没有一人像他：一个无怨言的人，真诚敬拜天主，远离一切邪恶。}\newline{}同一事，在\BibleBook{}{圣咏}第三十三篇：\Quote{我时时赞美上主，对祂的赞颂常在我口。}\newline{}同一事，在\BibleBook{}{户籍纪}：\Quote{让他们的抱怨远离我，他们就不致死亡。}\newline{}同一事，在\BibleBook{宗徒}{大事录}：\Quote{约在半夜，保禄和息拉祈祷赞美天主，囚犯都侧耳听着。}\newline{}又，在保禄致\BibleBook{斐理伯}{书}：\Quote{但一切事都因爱而作，没有抱怨和争辩，使你们无可指摘，成为天主纯洁无瑕的子女。}\par

% structure: p0031 source=h2 target=subsection reason=relative-heading-rank; base=h2

\subsection{十五、天主考验人，为使他们得到证明}

在\BibleBook{}{创世纪}：\Quote{天主试探亚巴郎，对他说：带你心爱的独生子依撒格，往高地去，在我指示你的山上，将他献为全燔祭。}\newline{}同一事，在\BibleBook{}{申命纪}：\Quote{上主你的天主试探你，要晓得你是否尽心尽性爱上主你的天主。}\newline{}同一事，在\BibleBook{}{智慧篇}：\Quote{在人看来，他们受了苦楚，但他们的希望充满不死之恩；他们在小事上受挫折，却在大事上得福乐，因为天主试验了他们，发现他们配得上自己。如金子般在炉中锻炼了他们，悦纳了他们如全燔祭。到了时候，他们必受尊崇；他们要审判万邦，统治各族；上主要永远为王。}\newline{}同一事，在\BibleBook{玛加伯}{书}：\Quote{亚巴郎在试探中不是显为忠信吗？这就算为他的义德。}\par

% structure: p0033 source=h2 target=subsection reason=relative-heading-rank; base=h2

\subsection{十六、殉道的益处}

在\BibleBook{}{箴言}中：\BibleQuote{忠信的殉道者由邪恶中拯救自己的灵魂。}\newline{}同处又载：\BibleQuote{那时，义人必大敢地站在那些磨难他们、夺去他们劳苦的人面前。当他们看见义人时，必为可怕的恐惧所扰乱，并因这出乎意料的迅速救恩而惊奇，彼此悔改并因心灵痛苦而叹息说：这些人就是我们一度曾讥笑、当作笑柄的人；我们愚人曾以为他们的生活是疯狂，他们的结局毫无尊荣。他们怎么被列在天主的儿女中，他们的份子在圣人之内呢！所以我们离开了真理的道路，正义的光没有照耀我们，太阳也没有为我们升起。我们因邪恶和丧亡的道路而疲惫，我们穿过了艰险的荒野；但我们却不认识上主的道路。骄傲为我们有什么益处呢？财富的夸耀给我们带来了什么呢？这一切都像影子一样消逝了。}\newline{}同一主题见于\BibleBook{}{圣咏}第一百一十五篇：\BibleQuote{在上主眼中，祂的圣人之死是宝贵的。}\newline{}又见于\BibleBook{}{圣咏}第一百二十五篇：\BibleQuote{含泪播种的，必欢乐收割。他们走去时，边走边哭，撒下种子；但归来时，必欢呼雀跃，怀抱禾捆。}\newline{}同一主题见于\BibleBook{若望}{福音}：\BibleQuote{爱惜自己性命的，必要丧失性命；在现世憎恨自己性命的，必要保存性命进入永生。}\newline{}同处又载：\BibleQuote{当他们把你们交出时，不要思虑你们该说什么；因为说话的不是你们，而是你们父的圣神在你们内说话。}\newline{}同处又载：\BibleQuote{时候要到，凡杀害你们的，都以为是尽恭敬天主的义务；他们这样做，是因为不认识父，也不认识我。}\newline{}同一主题见于\BibleBook{玛窦}{福音}：\BibleQuote{为义而受迫害的人是有福的，因为天国是他们的。}\newline{}同处又载：\BibleQuote{不要害怕那些杀害身体而不能杀害灵魂的；但更要害怕那能使灵魂和身体都投入地狱的。}\newline{}同处又载：\BibleQuote{凡在人前承认我的，我也必在我天上的父前承认他；凡在人前否认我的，我也必在我天上的父前否认他。并且坚持到底的，必要得救。}\newline{}同一主题见于\BibleBook{路加}{福音}：\BibleQuote{当人们憎恨你们，并与你们隔离，驱逐你们，并以你们的名字为恶名，好像歹徒一样，为人子的缘故时，你们是有福的。在那一天，你们要欢喜跳跃，因为你们的赏报在天上是大的。}\newline{}同处又载：\BibleQuote{我实在告诉你们，没有人为了天主的国而舍弃房屋、父母、兄弟、妻子或儿女，而不在现世得到多倍的报答，并在来世获得永生的。}\newline{}同一主题见于\BibleBook{}{默示录}：\BibleQuote{当羔羊揭开第五印时，我看见在祭台下面，那些为了天主的话和他们的见证而被杀者的灵魂。他们大声呼喊说：圣洁真实的主啊！你还不审判地上的人，为我们伸流血之冤，要到几时呢？于是有白衣赐给他们各人，又有话对他们说，还要安息片刻，等他们同仆和弟兄的数目满了，就是那些将要像他们一样被杀的人。}\newline{}同处又载：\BibleQuote{此后，我观看，见有一大群人，没有人能数得过来，来自各邦国、各支派、各民族、各种语言，站在宝座前和羔羊面前，身穿白衣，手拿棕榈枝。他们大声呼喊说：救恩属于坐在宝座上的我们的天主，也属于羔羊。众长老中有一位问我说：这些穿白衣的是谁？是从哪里来的？我回答说：我主，你知道。他对我说：这些人是从大灾难中出来的，曾用羔羊的血把衣裳洗白净了。所以他们在天主宝座前，昼夜在祂的圣殿中事奉祂；坐在宝座上的必住在他们中间。他们不再饥，不再渴，日头和炎热也必不伤害他们，因为宝座中的羔羊必牧养他们，领他们到生命水的泉源；天主也必擦去他们一切的眼泪。}\newline{}同处又载：\BibleQuote{得胜的，我必将生命树上的果子赐给他吃，这树在我天主的乐园中。}\newline{}同处又载：\BibleQuote{你务要至死忠信，我就赐给你那生命的冠冕。}\newline{}同处又载：\BibleQuote{那警醒并保守自己衣裳，免得赤身而行、叫人看见他羞耻的，是有福的。}\newline{}同一主题，保禄在\BibleBook{弟茂德}{后书}中说：\BibleQuote{我现在被奠祭，我离世的时候到了。那美好的仗我已经打过了，当跑的路我已经跑尽了，所信的道我已经守住了。从此以后，有公义的冠冕为我存留，就是按公义审判的主到了那日要赐给我的；不但赐给我，也赐给凡爱慕祂显现的人。}\newline{}同一主题，在\BibleBook{}{罗马书}中说：\BibleQuote{我们是天主的子女；既是子女，便是承继者，天主的承继者，与基督同为承继者；如果我们与祂一同受苦，也必与祂一同受光荣。}\newline{}同一主题见于\BibleBook{}{圣咏}第一百一十八篇：\BibleQuote{行为完全、遵行上主法律的人是有福的。遵守祂的法度、一心寻求祂的人是有福的。}\par

% structure: p0035 source=h2 target=subsection reason=relative-heading-rank; base=h2

\subsection{十七、我们在今世所受之苦，与所许之赏赐相比微不足道}

在\BibleBook{}{保禄致罗马人书}中：\BibleQuote{今时的苦楚，与将来要在我们身上显示的荣耀，是不能相比的。}论同一主题在\BibleBook{}{玛加伯书}中：\BibleQuote{主啊，祢拥有神圣的知识，显然，我本可免于一死，却因受鞭打而忍受最残酷的肉体痛苦；然而因着敬畏祢，我甘心在精神上忍受这一切。}又在同一处：\BibleQuote{你虽无能，却将我们逐出此世；但世界的君王必将我们这些为祂的律法而死的人复活，进入永生的复活。}又在同一处：\BibleQuote{被人置于死地，更好是从天主那里期待希望，由祂复活。因为对于你，是没有生命的复活。}又在同一处：\BibleQuote{你在人中有权，虽然你是有朽的，你任意而行。但不要以为我们的种族被天主遗弃。忍耐，看祂的大能怎样折磨你和你的后裔。}又在同一处：\BibleQuote{不要无故犯错；因为我们遭受这些，是咎由自取，因我们得罪了我们的天主。但不要以为你可以不受惩罚，因为你竟敢与天主争战。}\par

% structure: p0037 source=h2 target=subsection reason=relative-heading-rank; base=h2

\subsection{十八、没有什么比天主和基督的爱更应优先}

在\BibleBook{}{申命纪}中：\BibleQuote{你当全心、全灵、全力爱上主你的天主。}又据\BibleBook{玛窦}{福音}：\BibleQuote{爱父亲或母亲超过我的，不配是我的；爱儿子或女儿超过我的，不配是我的；不背起自己的十字架跟随我的，也不配是我的门徒。}又据\BibleBook{}{保禄致罗马人书}：\BibleQuote{谁能使我们与基督的爱隔绝呢？是患难吗？是困苦吗？是迫害吗？是饥饿吗？是赤身露体吗？是危险吗？是刀剑吗？正如经上所记载：\InnerQuote{为了你的缘故，我们终日被处死，被看作待宰的羊。}然而靠着那爱我们的主，我们在这一切事上已经得胜有余了。}\par

% structure: p0039 source=h2 target=subsection reason=relative-heading-rank; base=h2

\subsection{十九、我们不应服从自己的意愿，而应服从天主的旨意}

在\BibleBook{若望}{福音}中：\BibleQuote{我从天降下，不是为行我的旨意，而是为行那派遣我者的旨意。}论同一主题，据\BibleBook{玛窦}{福音}：\BibleQuote{父啊！如果可能，请让这杯离开我；但不要照我的意愿，而要照祢的意愿。}又据日常祈祷：\BibleQuote{愿祢的旨意奉行在人间，如同在天上。}又据\BibleBook{玛窦}{福音}：\BibleQuote{不是凡对我说\InnerQuote{主啊，主啊}的人，就能进入天国；而是那承行我在天之父旨意的人，才能进入天国。}又据\BibleBook{路加}{福音}：\BibleQuote{但那知道主人旨意，却不遵行的仆人，必受重罚。}在\BibleBook{若望}{书}中：\BibleQuote{但那承行天主旨意的人，永远常存，正如天主自己也永远常存。}\par

% structure: p0041 source=h2 target=subsection reason=relative-heading-rank; base=h2

\subsection{二十、希望和信德的基础与力量在于敬畏}

在第一百一十篇中：\BibleQuote{敬畏上主是智慧的开端。}同样的事在\BibleBook{}{智慧篇}中：\BibleQuote{智慧的开端是敬畏天主。}也在同一作者的\BibleBook{}{箴言}中：\BibleQuote{以敬畏之心尊重万物的人是有福的。}同样的事[在\BibleBook{}{依撒意亚}中：\BibleQuote{我还垂顾谁呢？除非是那谦卑、和平，且因我的话而战栗的人。}]同样的事在\BibleBook{}{创世纪}中：\BibleQuote{上主的天使从天上呼唤他说：\Quote{亚巴郎！亚巴郎！}他答说：\Quote{我在这里。}天使说：\Quote{不可在这孩子身上下手，不要伤害他！现在我知道你敬畏天主，因为你为了我甘愿牺牲你的独生子。}}也在第二篇圣咏中：\BibleQuote{你们要以敬畏事奉上主，战战兢兢地向祂欢呼。}也在\BibleBook{}{申命纪}中，天主对梅瑟说：\BibleQuote{你召集人民到我面前，我要让他们听我的话，好使他们在我所赐给他们生活的土地上，一生学习敬畏我。}也在\BibleBook{}{耶肋米亚}中：\BibleQuote{看，时日将到——上主断语——我愿与以色列家和犹大家订立新约，不像我同他们的祖先所立的盟约，在他们拉着他们的手领出埃及地的那天，因为他们没有在我的盟约中存留，我也厌弃了他们——上主断语——因为这是我将与以色列家订立的盟约：在那时日之后——上主断语——我将把我的法律放在他们心中，写在他们的脑海里；我要做他们的天主，他们要做我的子民。他们不再教导自己的同胞兄弟说：\Quote{认识上主！}因为从最小的到最大的，他们都要认识我，因为我要宽恕他们的过犯，不再记忆他们的罪恶。如果上天能被上举，大地被下沉——上主断语——我也不会因他们所做的一切而弃绝以色列子民。看，我将从我被愤怒、忿恨和大怒所分散的各地聚集他们，并将他们带回这地方，使他们安居；他们要做我的子民，我要做他们的天主。我要赐给他们另一条道路和另一颗心，使他们一生敬畏我，同他们的子女一起得福；我要为他们订立永久的盟约，不再离弃他们；我要将我的敬畏放在他们心中，使他们不离开我；我要照顾他们，使他们得福，并在一心一意、信赖中栽植他们在自己的土地上。}也在\BibleBook{}{默示录}中：\BibleQuote{那坐在自己宝座上的二十四位长老，在天主面前俯伏在地，朝拜天主说：\Quote{全能的上主天主，你今在昔在，我们感谢你，因为你已取得你的大能，开始为王。万民发怒，你的愤怒已到，审判死人的时候到了，你应赏赐你的仆人众先知和敬畏你名字的圣徒，无论大小，并消灭败坏大地的人。}}也在同一处：\BibleQuote{我又看见一片如同玻璃的海洋，混杂着火；那兽与羔羊一同牧养；祂名字的数目是一百四十四，那些战胜了兽和兽像以及它的名号数字的人，站在玻璃海上，拿着天主的琴，歌唱梅瑟的颂歌，即天主的仆人的歌，和羔羊的歌说：\Quote{上主全能的天主，你的作为伟大而奇！万民之王，你的道路公正又真实！谁敢不敬畏你，不光荣你的名字？惟独你是圣的！因为万民都要前来，在你面前朝拜，因为你正义的判断已显现了。}}也在同一处：\BibleQuote{我又看见另一位天使飞翔在天际，带着永久的福音向地上居住的人，向各邦国、支派和民族宣讲，大声说：\Quote{你们要敬畏天主，光荣他，因为他审判的时辰来到了；你们要朝拜那创造天地海洋和泉源的那一位。}}也在\BibleBook{}{达尼尔}中：\BibleQuote{在巴比伦住着一个人，名叫约阿金；他娶了一个妻子，名叫苏撒纳，是赫尔基雅的女儿，她很美丽，且敬畏上主。她的父母都是义人，曾依照梅瑟的法律教导他们的女儿。}又在\BibleBook{}{达尼尔}中：\BibleQuote{我们在全地今日因自己的罪而卑微，如今没有君王、先知、领袖、全燔祭、供物、牺牲、馨香和向您献祭的场所，以寻得您的垂怜。但让我在谦卑的灵魂和心灵中被接受，像公绵羊和公牛的全燔祭，又如许多最肥嫩的羔羊。但愿我们的祭献今天能在您面前蒙悦纳，他们的力量就要消耗，因为凡信赖您的必不蒙羞。现在我们全心追随您，敬畏您，寻求您的慈颜。不要使我们受辱，请按您的温和和您丰厚的仁慈对待我们，拯救我们。}也在同一处：\BibleQuote{君王非常高兴，下令将达尼尔从狮子圈里提上来；狮子没有伤害他，因为他信赖自己的天主。君王下令，把那些控告达尼尔的人带来，连同他们的妻子儿女一起投入狮子圈。他们还没有到达圈底，狮子就抓住了他们，将他们的骨头咬得粉碎。那时达理阿王写信给住在全地的各民族、各支派、各语言说：\Quote{愿你们的平安十分昌盛！我颁布命令：在我统治的全国，人人都应在达尼尔所事奉的天主面前战栗敬畏，因为他是永活的神，永远长存，他的国度不会灭亡，他的王权永无止境；他独自在天上和地下行奇迹异能，拯救了达尼尔免于狮子的爪牙。}}也在\BibleBook{}{米该亚}中：\BibleQuote{我该用什么亲近上主，屈膝敬拜高高在上的天主？要带全燔祭，带一岁的牛犊吗？上主会喜欢千万只公山羊，或万条油河吗？为我的过犯，应献我的长子，为我的灵魂的罪，应献我亲生的儿子？人啊！已通知了你什么是善，上主向你要求的是什么？无非是履行正义，爱好慈善，与你的天主谦逊同行。上主的呼声向城市呼唤，祂必拯救敬畏祂名的人。}也在\BibleBook{}{米该亚}中：\BibleQuote{求你用你的牧杖牧养你的子民，作你产业的羊群；那些独居在加尔默耳中间的，你要拔出。让他们在巴商和基肋阿得像往古一样牧放。使他们出离埃及的日子，我必显示奇迹。万民见了，必因他们的一切势力而困惑，以手掩口。他们的耳朵变聋，他们要像蛇一样舔土。他们在地上爬行，必惊慌失措，他们必向上主他们的天主战兢，并因你而害怕。哪一位神像你，赦免过犯，宽恕罪愆？}在\BibleBook{}{纳鸿}中：\BibleQuote{山岭在他面前震动，丘陵摇撼；大地在他面前崩塌，世界和所有居民也都在战栗。在他盛怒之下，谁能站得住？在他烈怒之下，谁能立得起？他的愤怒如火倾倒，岩石在他面前崩裂。上主对信赖他的人，在患难之日是良善的，并认识那些敬畏他的人。}也在\BibleBook{}{哈盖}中：\BibleQuote{则鲁巴贝耳——沙耳提耳的儿子，犹大支派，和耶叔亚——约匝达克的儿子，大司祭，以及全体遗民，都听从了上主他们天主的声音，因为上主派他来到他们当中；民众因为天主而畏惧。}也在\BibleBook{}{玛拉基亚}中：\BibleQuote{我与他们订立的盟约是生命与平安；我赐给他们敬畏，好因我的名而敬畏我。}也在第三十三篇圣咏中：\BibleQuote{敬畏上主，你们祂的圣徒；因为敬畏祂的人一无所缺。}也在第十八篇圣咏中：\BibleQuote{敬畏上主是纯洁的，永远常存。}\par

% structure: p0043 source=h2 target=subsection reason=relative-heading-rank; base=h2

\subsection{二十一、我们不可轻率判断他人}

据\WorkTitle{路加福音}记载：\BibleQuote{你们不要判断，你们就不被判断；不要定罪，你们就不被定罪。}关于同一主题，在\WorkTitle{罗马书}中：\BibleQuote{你是何人，竟判断别人的仆人？他或站或倒，自有他的主人；而且他必要站住，因为天主能够使他站住。}又说：\BibleQuote{因此，你这判断人的，无论你是谁，都是无可推诿的；因为你在判断别人的事上，正定你自己的罪，因为你所判断的，你自己也作。然而你这判断行恶的人，自己却行同样的恶，还指望能逃脱天主的审判吗？}又在\WorkTitle{格林多前书}中：\BibleQuote{所以，那自以为站得稳的，要谨慎，免得跌倒。}又说：\BibleQuote{若有人以为自己知道什么，这还不知道他应当怎样知道。}\par

% structure: p0045 source=h2 target=subsection reason=relative-heading-rank; base=h2

\subsection{二十二、我们受人亏欠时，必须赦免并宽恕}

在福音中，日常祈祷里：\BibleQuote{宽免我们的罪债，如同我们也宽免得罪我们的人。}又据\WorkTitle{马尔谷福音}记载：\BibleQuote{当你们站着祈祷时，若你们有什么怨怪别人的，就要宽恕，好叫你们在天上的父，也宽恕你们的过犯。但你们若不宽恕，你们在天上的父也不宽恕你们的过犯。}又在同一处：\BibleQuote{你们用什么尺度量给人，也要用什么尺度量给你们。}\par

% structure: p0047 source=h2 target=subsection reason=relative-heading-rank; base=h2

\subsection{二十三、不可以恶报恶}

在\WorkTitle{罗马书}中：\BibleQuote{不以恶报恶。}又在同一处：\BibleQuote{不可为恶所胜，反要以善胜恶。}关于同一件事，在\WorkTitle{默示录}中：\BibleQuote{祂对我说：\InnerQuote{你不可封住这书上预言的话，因为时期近了。让行恶的仍旧行恶，让污秽的仍旧污秽，让行义的仍旧行义，让圣洁的仍旧圣洁。看，我快要来，我的赏报在我身边，要按照各人的行为还报各人。}}\par

% structure: p0049 source=h2 target=subsection reason=relative-heading-rank; base=h2

\subsection{二十四、除祂圣子耶稣基督外，无人能到圣父那里}

据\WorkTitle{若望福音}记载：\BibleQuote{我是道路、真理、生命：除非经过我，谁也不能到父那里去。}又在同一处：\BibleQuote{我就是门：谁若经过我进来，必得救。}\par

% structure: p0051 source=h2 target=subsection reason=relative-heading-rank; base=h2

\subsection{二十五、人除非受洗并重生，不能进入天主的国}

据\WorkTitle{若望福音}记载：\BibleQuote{人除非由水和圣神重生，不能进入天主的国。因为凡由血肉生的，属于血肉；凡由圣神生的，属于神。}又在同一处：\BibleQuote{你们若不吃人子的肉，不喝祂的血，在你们内就没有生命。}\par

% structure: p0053 source=h2 target=subsection reason=relative-heading-rank; base=h2

\subsection{二十六、除非人在行为和事工上得益，否则受洗和领受圣体圣事就算不得什么}

在\WorkTitle{格林多前书}中：\BibleQuote{你们不知道在运动场上赛跑的，固然都跑，但只有一个得奖赏吗？你们也该这样跑，好能得到奖赏。凡参加比赛的，不过为得到可朽坏的花冠；我们却是为得到不朽坏的花冠。}在\WorkTitle{玛窦福音}中：\BibleQuote{凡不结好果子的树，就要砍下来，丢在火里。}又在同一处：\BibleQuote{到那天，许多人要对我说：\InnerQuote{主啊，主啊，我们不是因你的名字说过预言，因你的名字驱过魔鬼，因你的名字行过许多异能吗？}那时我要向他们声明说：\InnerQuote{我从来不认识你们；你们这些作恶的人，离开我吧！}}又在同一处：\BibleQuote{你们的光当在人前照耀，好使他们看见你们的善行，光荣你们在天之父。}又，保禄在\WorkTitle{斐理伯书}中：\BibleQuote{你们应照耀如同世上的光。}\par

% structure: p0055 source=h2 target=subsection reason=relative-heading-rank; base=h2

\subsection{二十七、即使受洗的人，若不保持纯洁，也会失去所获得的恩宠}

据\WorkTitle{若望福音}记载：\BibleQuote{看，你已经痊愈了；不要再犯罪，免得遭遇更不幸的事。}又在\WorkTitle{格林多前书}中：\BibleQuote{你们不知道你们是天主的圣殿，天主的圣神住在你们内吗？若有人毁坏天主的圣殿，天主必要毁坏他。}关于同一件事，在\WorkTitle{编年纪}中：\BibleQuote{天主与你们同在，你们若与祂同在；如果你们离弃祂，祂也必离弃你们。}\par

% structure: p0057 source=h2 target=subsection reason=relative-heading-rank; base=h2

\subsection{二十八、对于得罪天主（即圣神）的人，在教会中不能获得赦免}

据\WorkTitle{玛窦福音}记载：\BibleQuote{凡说话干犯人子的，可得赦免；但谁若说话干犯圣神，无论在今世或来世，都不得赦免。}又据\WorkTitle{马尔谷福音}记载：\BibleQuote{世人的一切罪，连一切亵渎的话，都可得赦免；但谁若亵渎了圣神，永远不得赦免，而要承担永久的罪恶。}关于同一件事，在\WorkTitle{列王纪上}中：\BibleQuote{若有人得罪别人，人可以为他向主祈祷；但若有人得罪天主，谁能为他祈祷呢？}\par

% structure: p0059 source=h2 target=subsection reason=relative-heading-rank; base=h2

\subsection{二十九、关于对那名字的憎恨，早已预言}

在\BibleBook{路加}{福音}中：\BibleQuote{你们要为我的名被众人憎恨。}又在\BibleBook{若望}{福音}中：\BibleQuote{世界若恨你们，你们该知道，它恨你们以前，已经恨了我。你们若是属于世界，世界必爱属于自己的；但因为你们不属于世界，而我拣选了你们出离世界，为此世界恨你们。你们要记得我对你们所说的话：仆人不能大于主人。如果他们迫害了我，也要迫害你们。}又在\BibleBook{巴路克}{书}中：\BibleQuote{因为时候将到，你们和你们以后的人要寻找我，要听智慧与明达的话，却找不到我。但外邦人渴望见到智者，却不得遇见；不是因为这世界的智慧要缺少，或在地上消失；但法律的话也不会从世上断绝。因为智慧只在于少数警醒、静默、安静、彼此交谈的人中；因为有些人畏惧他们，视他们为邪恶而惧怕。但有些人不相信至高者法律的话。有些人面容惊愕而不信；那些抗拒的人也将相信，并会反对和阻碍真理之神。此外，有些人将明智于谬误之神，并宣告谕令，好似出自至高者和全能者。此外，有些人是\Emph{拥有信德者}。其他人在对至高者的信德上刚强有力，憎恶外人。}\par

% structure: p0061 source=h2 target=subsection reason=relative-heading-rank; base=h2

\subsection{三十、人向天主所许的愿，必须迅速偿还}

在\BibleBook{}{箴言}中：\BibleQuote{你向天主许了愿，不可迟延偿还。}关于同一事，在\BibleBook{申命}{纪}中：\BibleQuote{你若向主你的天主许了愿，你不可迟延偿还，因为主你的天主必要向你追讨，那将成为你的罪。你口中所出的，你应遵守，并应履行你亲口所许诺的礼物。}关于同一事，在第四十九篇\BibleBook{}{圣咏}中：\BibleQuote{以赞颂的祭献向天主奉献，并向至高者偿还你的愿。在患难之日呼求我，我必拯救你，你必荣耀我。}关于同一事，在\BibleBook{宗徒}{大事录}中：\BibleQuote{为什么撒殚充满了你的心，使你欺骗圣神？你的田产原属你支配，你没有欺骗人，而是欺骗了天主。}又在\BibleBook{耶肋米亚}{书}中：\BibleQuote{懒散地做天主工作的人是可咒骂的。}\par

% structure: p0063 source=h2 target=subsection reason=relative-heading-rank; base=h2

\subsection{三十一、不信的人已经受了审判}

在\BibleBook{若望}{福音}中：\BibleQuote{不信的人已受了审判，因为他没有信从天主独生子的名字。审判就在于此：光明来到了世界，世人却爱黑暗甚于光明。}关于同一事，在第一篇\BibleBook{}{圣咏}中：\BibleQuote{为此，不虔敬的人在审判中站不住，罪人在义人的集会中也不能立足。}\par

% structure: p0065 source=h2 target=subsection reason=relative-heading-rank; base=h2

\subsection{三十二、论贞洁和节制的益处}

在\BibleBook{创世}{纪}中：\BibleQuote{我要大大增加你的忧愁和叹息，你将在痛苦中生子；你要依恋你的丈夫，他要管辖你。}关于同一事，在\BibleBook{玛窦}{福音}中：\BibleQuote{这话不是人人都能领悟的，只有那些得了恩赐的人才能领悟。因为有些阉人从母胎生来就是如此，有些阉人是被人阉的，有些阉人是为了天国而自阉的。能领悟的，就领悟吧！}又在\BibleBook{路加}{福音}中：\BibleQuote{今世之子也娶也嫁。但那些算得配得那世界、并从死者中复活的人，也不娶也不嫁；因为他们不能再死，他们相似天使，既是复活之子。至于死人复活，梅瑟在荆棘篇中指明：他说，主是亚巴郎的天主，依撒格的天主，雅各伯的天主。他不是死人的天主，而是活人的天主：因为众人为他而生活。}又在圣保禄致\BibleBook{格林多}{前书}中：\BibleQuote{男人不亲近女人是好的。但因着淫乱的危险，每男人当有自己的妻子，每女人当有自己的丈夫。丈夫对妻子当尽本分，妻子对丈夫也当如此。妻子对自己的身体没有主权，而是丈夫；同样，丈夫对自己的身体没有主权，而是妻子。你们不要彼此亏负，除非两相情愿，暂为分时，为专务祈祷；然后再回到一起，免得撒殚因你们不能自制而诱惑你们。我说这话，是宽容，不是命令。我切愿众人像我一样。但每人各有从天主而来的恩赐：有人这样，有人那样。}也在同一处：\BibleQuote{未结婚的男子挂念主的事，想怎样悦乐天主；但结了婚的男子挂念世俗的事，想怎样悦乐妻子。同样，女子和未出嫁的贞女挂念主的事，为在身体和灵魂上成为圣洁；但已出嫁的女子挂念世俗的事，想怎样悦乐丈夫。}又在\BibleBook{出谷}{纪}中，当主命令梅瑟为第三日圣化百姓时，他圣化了他们，并补充说：\BibleQuote{准备好，三天内不可亲近女人。}又在\BibleBook{列王}{纪上}：\BibleQuote{司祭回答达味说：我手中没有普通饼，只有圣饼。如果少年人没有亲近女人，他们就可以吃。}又在\BibleBook{}{默示录}中：\BibleQuote{这些人没有与女人玷污自己，因为他们保持了贞洁；这些人无论羔羊往哪里去，都跟随他。}\par

% structure: p0067 source=h2 target=subsection reason=relative-heading-rank; base=h2

\subsection{三十三、圣父不审判任何事，而是圣子审判；并且凡不荣耀圣子的，也不能荣耀圣父}

在\BibleBook{若望}{福音}中记载：\BibleQuote{圣父不审判任何人，却把一切审判交给了圣子，使众人尊敬圣子如同尊敬圣父。不尊敬圣子的，就是不尊敬派遣他的圣父。}又在\BibleBook{}{圣咏集}第七十一篇中：\BibleQuote{天主，请将你的审判权赐给君王，将你的正义赐给君王的儿子，使他以正义审判你的百姓。}又在\BibleBook{}{创世纪}中：\BibleQuote{上主从天上上主那里，降下硫磺与火，焚烧了\Place{索多玛}和\Place{哈摩辣}。}\par

% structure: p0069 source=h2 target=subsection reason=relative-heading-rank; base=h2

\subsection{三十四、信徒不应像外邦人那样生活}

在\BibleBook{}{耶肋米亚}中：\BibleQuote{上主这样说：不要随从外邦人的道路。}关于同一件事，即人应当与外邦人分离，免得成为他们罪的同伴，并分担他们的惩罚，在\BibleBook{}{默示录}中：\BibleQuote{我又听见从天上另有声音说：我的百姓，你们要从她那里出来，免得你们分担她的罪恶，也免得你遭受她的灾殃；因为她的罪恶滔天，上主天主已想起了她的不义。所以祂加倍报应她，用她调和的杯加倍给她调和；她怎样荣耀自己，怎样奢华，也照样叫她痛苦悲哀。因为她心里说：我坐王后，并不是寡妇，决不见悲哀。所以在一日之内，她的灾祸要一齐来到，就是死亡、悲哀、饥荒；她又要在火中被烧，因为审判她的上主天主是刚强的。地上的君王，凡与她行淫、同她享乐的，也都要为她哭泣哀号。}在\BibleBook{}{依撒意亚}中：\BibleQuote{你们抬上主器皿的，要从他们中间出去。}\par

% structure: p0071 source=h2 target=subsection reason=relative-heading-rank; base=h2

\subsection{三十五、天主忍耐是为使我们悔改罪过，得以更新}

在\Person{撒罗满}的\BibleBook{}{德训篇}中：\BibleQuote{不要说：我犯了罪，有什么灾祸降到我身上？因为至高者是容忍的偿还者。}在\BibleBook{}{罗马书}中：\BibleQuote{难道你藐视祂丰厚的恩慈、宽容与忍耐，不晓得天主的恩慈是领你悔改吗？你竟任着你刚硬不悔改的心，为自己积蓄忿怒，直到忿怒的日子，就是天主公义审判的日子显明的时候，祂要照各人的行为报应各人。}\par

% structure: p0073 source=h2 target=subsection reason=relative-heading-rank; base=h2

\subsection{三十六、妇女不应世俗打扮}

在\BibleBook{}{默示录}中：\BibleQuote{拿七碗的七位天使中，有一位来对我说：你来，我要将坐在众水上的大淫妇所要受的审判指给你看。地上的君王与她行淫。我看见一个女人骑在兽上。那女人穿着紫色和朱红色的衣服，用金子、宝石、珍珠为妆饰；手拿金杯，杯中盛满了可憎之物，就是她淫乱的污秽。}在致弟茂德书中：\BibleQuote{愿你们的妇女以廉耻和端庄为妆饰，不以编发、黄金、珍珠、或华贵的衣服为妆饰，但要以善行，就如那自称为虔敬天主的妇女所应有的。}关于同一件事，在伯多禄致本都人的书信中：\BibleQuote{妇女不应注重外表的装饰，如金饰、衣服，而应注重内心的装饰。}在\BibleBook{}{创世纪}中：\BibleQuote{\Person{塔玛尔}披上外衣，打扮自己；\Person{犹大}看见她，以为她是妓女。}\par

% structure: p0075 source=h2 target=subsection reason=relative-heading-rank; base=h2

\subsection{三十七、信徒除为所负之名外，不应因其他过犯受罚}

在伯多禄致本都人的书信中：\BibleQuote{你们中不要有人因作贼、或杀人、或作恶、或干预别人的事而受苦，但若是作为基督徒而受苦。}\par

% structure: p0077 source=h2 target=subsection reason=relative-heading-rank; base=h2

\subsection{三十八、天主的仆人应纯洁，以免陷入世俗的刑罚}

在\BibleBook{}{罗马书}中：\BibleQuote{你不怕权柄吗？你只管行善，就必从它得称赞。}\par

% structure: p0079 source=h2 target=subsection reason=relative-heading-rank; base=h2

\subsection{三十九、我们在基督内获得了生活的榜样}

在伯多禄致本都人的书信中：\BibleQuote{因为基督为我们受了苦，给你们留下榜样，叫你们跟随祂的脚踪行。祂没有犯罪，口里也没有诡诈；祂被骂不还口，受害不说威吓的话，只将自己交托那按公义审判的主。}在\BibleBook{}{斐理伯书}中：\BibleQuote{祂虽具有天主的形体，却没有以自己与天主同等为应当把持的，反而空虚自己，取了奴仆的形体，成了人的模样；祂既成了人，就谦卑自下，服从至死，且死在十字架上。因此，天主极其举扬祂，赐给祂一个名字，超越所有名字，使天上、地上、地下的一切，因耶稣的名字而屈膝，并一切唇舌都明认耶稣基督是主，以光荣天主父。}关于同一件事，在\BibleBook{若望}{福音}中：\BibleQuote{我是你们的主、你们的师傅，尚且洗你们的脚，你们也当彼此洗脚。我给你们作了榜样，叫你们照着我向你们所作的去作。}\par

% structure: p0081 source=h2 target=subsection reason=relative-heading-rank; base=h2

\subsection{四十、不可喧嚷自夸地行事}

在\BibleBook{玛窦}{福音}中：\BibleQuote{不要叫左手知道右手所作的，好使你的施舍隐而不露；你父在暗中看见，必要报答你。}同处：\BibleQuote{当你施舍时，不要在你前面吹号，像假善人在会堂及街市上所行的一样，为受人们的称赞；我实在告诉你们，他们已获得了他们的赏报。}\par

% structure: p0083 source=h2 target=subsection reason=relative-heading-rank; base=h2

\subsection{四十一、不可说愚妄和冒犯人的话}

在\BibleBook{}{厄弗所书}中：\BibleQuote{淫词、妄语和戏笑的话都不相宜，总不可说。}\par

% structure: p0085 source=h2 target=subsection reason=relative-heading-rank; base=h2

\subsection{四十二、信德全然有益，我们能照我们所信的行事}

在\BibleBook{}{创世纪}中：\BibleQuote{亚巴郎相信了天主，天主就以此算为他的正义。}又在\BibleBook{}{依撒意亚}中：\BibleQuote{假使你们不相信，你们必然不能理解。}又在\BibleBook{玛窦}{福音}中：\BibleQuote{小信德的人哪！你为什么怀疑？}同样在那处：\BibleQuote{如果你们有信德像芥子那样大，你们向这座山说：从这边移到那边去，它必移过去；在你们没有不可能的事。}又根据\BibleBook{马尔谷}{福音}：\BibleQuote{凡你们祈祷和所求的一切，相信你们必得，就必给你们。}同样在那处：一切事情为信的人都是可能的。在\BibleBook{}{哈巴谷}中：\BibleQuote{义人必因我的信德而生活。}又在\BibleBook{}{达尼尔}中：\BibleQuote{\Person{亚纳尼亚}、\Person{阿匝黎雅}和\Person{米沙耳}依赖天主，便得脱离火燄。}\par

% structure: p0087 source=h2 target=subsection reason=relative-heading-rank; base=h2

\subsection{四十三、相信的人能立即获得（即赦免与平安）}

在\BibleBook{宗徒}{大事录}中：\BibleQuote{看，这里有水；有甚么阻止我受洗呢？\Person{斐理伯}说：你若全心相信，便可以。}\par

% structure: p0089 source=h2 target=subsection reason=relative-heading-rank; base=h2

\subsection{四十四、彼此有争执的信徒不应向外邦法官申诉}

在\BibleBook{格林多}{前书}中：\BibleQuote{你们中间有人对别人有了争执，岂敢在不义的人面前，而不在圣人面前请求裁判吗？你们不知道圣人要审判这世界吗？}又说：\BibleQuote{你们彼此间既然有诉讼的争执，这已是你们的过失了。你们为什么不情愿受委屈？为什么不情愿吃亏？你们反倒冤枉人，亏负人，况且是亏负弟兄！你们不知道不义的人不能承继天主的国吗？}\par

% structure: p0091 source=h2 target=subsection reason=relative-heading-rank; base=h2

\subsection{四十五、望德关乎未来之事，因此我们对所应许之事的信德应当忍耐}

在\BibleBook{}{罗马书}中：\BibleQuote{我们得救是在乎望德。但是看得见的望德，不是望德；既然看见了，为甚么还盼望呢？如果我们盼望那看不见的，我们就忍耐着盼望。}\par

% structure: p0093 source=h2 target=subsection reason=relative-heading-rank; base=h2

\subsection{四十六、妇女在教会中应当沉默}

在\BibleBook{格林多}{前书}中：\BibleQuote{妇女在教会中应当缄默。若她们愿意学习甚么，在家里问自己的丈夫好了。}又对\BibleBook{弟茂德}{前书}：\BibleQuote{女人要沉静学习，事事服从。我不许女人施教，也不许她管辖男人，但要她沉静。因为先造的是\Person{亚当}，后造的是\Person{厄娃}；\Person{亚当}没有受骗，而是女人受了骗。}\par

% structure: p0095 source=h2 target=subsection reason=relative-heading-rank; base=h2

\subsection{四十七、我们受苦且未能察觉天主在万事上的帮助，是由于我们的过失和应得的惩罚}

在\BibleBook{}{欧瑟亚}中：\BibleQuote{以色列子民，请听主的话。因为主控告本地的居民：本地没有诚实，没有仁慈，也没有认识天主；只有诅咒、谎言、杀戮、偷窃、奸淫，且日增不已；他们流人血，接踵流血。因此，此地要哀悼，居其上的居民、田野的走兽、天空的飞鸟都要消逝，连海中的鱼也要灭绝；以至于没有人能裁判，没有人能谴责。}关于同一件事，在\BibleBook{}{依撒意亚}中：\BibleQuote{主的手岂是短了而不能拯救？他的耳朵岂是沉了而不能听见？而是你们的罪恶使你们与你们的天主隔绝；是你们的恶行使他掩面不理你们。因为你们的手沾满了血，你们的手指沾染了罪恶；你们的嘴唇说出谎言，你们的舌头议论邪恶。没有人谈论正义，也没有人公正判断。他们信赖虚无，口说空话，孕育痛苦，产生邪恶。}又在\BibleBook{}{索福尼亚}中：\BibleQuote{我必要彻底消灭地面上的一切，主说。我要消灭人和牲畜，空中的飞鸟和海中的鱼，并铲除恶人。}\par

% structure: p0097 source=h2 target=subsection reason=relative-heading-rank; base=h2

\subsection{四十八、我们不可取利息}

在\BibleBook{}{圣咏集}第十五篇中：\BibleQuote{他不上以利息借贷，也不接受贿赂加害无辜。这样行事的人，永不动摇。}又在\BibleBook{}{厄则克耳}中：\BibleQuote{那正义的人：不压迫人，归还债务人的抵押，不强夺，把自己的食物分给饥饿的人，给赤身者衣服穿，不放债取利。}又在\BibleBook{}{申命纪}中：\BibleQuote{不可向你兄弟索取利息，无论是金钱利息或是食物利息。}\par

% structure: p0099 source=h2 target=subsection reason=relative-heading-rank; base=h2

\subsection{四十九、连我们的仇人也必须爱}

在\BibleBook{路加}{福音}中：\BibleQuote{如果你们爱那些爱你们的人，你们还有甚么赏报呢？因为罪人也爱那些爱他们的人。}又根据\BibleBook{玛窦}{福音}：\BibleQuote{你们要爱你们的仇人，为迫害你们的人祈祷，好使你们成为你们在天之父的子女；因为他使太阳上升，光照恶人，也光照善人；降雨给义人，也给不义的人。}\par

% structure: p0101 source=h2 target=subsection reason=relative-heading-rank; base=h2

\subsection{五十、信德的圣事不可被亵渎}

在\Person{撒罗满}的\BibleBook{}{箴言}中：\BibleQuote{不要对愚人说话，怕他轻视你的明智。}又在\BibleBook{玛窦}{福音}中：\BibleQuote{不要把圣物给狗，也不要把你们的珍珠投在猪前，免得它们用脚践踏，并且转过来撕裂你们。}\par

% structure: p0103 source=h2 target=subsection reason=relative-heading-rank; base=h2

\subsection{五十一、任何人不应因劳作而自高}

在\Person{撒罗满}的\BibleBook{}{德训篇}中：\BibleQuote{不要因善事而自高。}又在\BibleBook{路加}{福音}中：\BibleQuote{你们中间谁有仆人耕地或放羊，从田里回来，便对他说：你快过来坐下吃饭呢？而不对他说：你给我准备晚饭，束上腰伺候我，等我吃喝完了，然后你才吃喝？仆人作完所吩咐的事，主人岂要感谢他吗？你们也是这样，作完所吩咐的事，要说：我们是无用的仆人，我们不过作了我们应做的事。}\par

% structure: p0105 source=h2 target=subsection reason=relative-heading-rank; base=h2

\subsection{五十二、信或不信的自由在于自由意志}

在\BibleBook{申命}{纪}中：\BibleQuote{看，我已将生命与死亡，福与祸，摆在你们面前。你要选择生命，为叫你们生存。}又在\BibleBook{依撒意亚}{}中：\BibleQuote{若是你们愿意，且听从我，你们将享用地上的美物。但若是你们拒绝，且不听从我，刀剑必吞噬你们，因为上主的口说了这话。}又在\BibleBook{路加}{福音}中：\BibleQuote{天主的国就在你们中间。}\par

% structure: p0107 source=h2 target=subsection reason=relative-heading-rank; base=h2

\subsection{五十三、论天主的奥妙不可窥透，因此我们的信德应当是纯朴的}

在\BibleBook{格林多}{前书}中：\BibleQuote{我们现在是借着镜子观看，模糊不清，到那时，就要面对面了。我现在所认识的，只是局部的，那时我就要完全认识了，如同我全被认识一样。}又在撒罗满的\BibleBook{智慧}{篇}中：\BibleQuote{以纯朴的心寻求祂。}又同一书中：\BibleQuote{以纯朴行事的人，行事稳妥。}又同一书中：\BibleQuote{不要寻求过于你的事，也不要探究强于你的事。}又在\BibleBook{撒罗满}{书}中：\BibleQuote{不要过分正义，也不要过度明智，免得你惊惶失措。}又在\BibleBook{依撒意亚}{}中：\BibleQuote{祸哉，那些在自己眼中为义的人。}又在\BibleBook{玛加伯}{}中：\BibleQuote{达内尔因着他的纯朴，从狮子的口中被救出。}又在\BibleBook{罗马}{书}中：\BibleQuote{啊，天主的富饶、上智和知识是多么高深！祂的决断是多么不可测量！祂的道路是多么难觅！谁曾知道上主的心意？或者谁做过祂的顾问？或者谁先给了祂，而该受祂的报答呢？因为万物都出于祂，依赖祂，归于祂。愿光荣归于祂至于永远！}又在\BibleBook{弟茂德}{后书}中：\BibleQuote{至于那些愚昧而粗野的辩论，务要躲避，因为知道它们会生争论。但天主的仆人不应争执，却要对众人温和。}\par

% structure: p0109 source=h2 target=subsection reason=relative-heading-rank; base=h2

\subsection{五十四、论没有一人是没有污秽、没有罪的}

在\BibleBook{约伯}{传}中：\BibleQuote{谁能从污秽中生出洁净呢？没有一人，即使他的生命在世上只有一天。}又在\BibleRef{咏 50}中：\BibleQuote{请看，我是在罪孽中生成的，我母亲在罪恶中怀了我。}又在\BibleBook{若望}{一书}中：\BibleQuote{如果我们说我们没有罪，就是欺骗自己，真理也不在我们内。}\par

% structure: p0111 source=h2 target=subsection reason=relative-heading-rank; base=h2

\subsection{五十五、论我们不应取悦人，而应取悦天主}

在\BibleRef{咏 52}中：\BibleQuote{取悦人的人，蒙受羞辱，因为天主使他们归于无有。}又在\BibleBook{迦拉达}{书}中：\BibleQuote{如果我还想取悦人，我就不是基督的仆人了。}\par

% structure: p0113 source=h2 target=subsection reason=relative-heading-rank; base=h2

\subsection{五十六、论没有一件事能向天主隐藏}

在撒罗满的\BibleBook{智慧}{篇}中：\BibleQuote{上主的眼睛在各地注视善人与恶人。}又在\BibleBook{耶肋米亚}{}中：\BibleQuote{我岂是近处的天主，而不是远处的天主？人若藏在隐密处，我岂不能看见他？我岂不充满天与地？上主说。}又在\BibleBook{撒慕尔}{上}中：\BibleQuote{人看容貌，上主却看人心。}又在\BibleBook{默示录}{}中：\BibleQuote{众教会都要知道，我是洞察肺腑心肠的；我要按照你们的行为，赏报你们每人。}又在\BibleRef{咏 18}中：\BibleQuote{谁能认清自己的过犯？求你赦免我未觉察的罪愆。}又在\BibleBook{格林多}{后书}中：\BibleQuote{我们众人都必须出现在基督的审判台前，为使各人按他肉身所行的，或善或恶，领取报应。}\par

% structure: p0115 source=h2 target=subsection reason=relative-heading-rank; base=h2

\subsection{五十七、论信徒受管教而得以保存}

在\BibleRef{咏 117}中：\BibleQuote{上主惩戒了我，却没有将我交于死亡。}又在\BibleRef{咏 88}中：\BibleQuote{我要用棍杖惩罚他们的过犯，用鞭笞责罚他们的罪孽。但我绝不将我的慈爱从他们身上收回。}又在\BibleBook{玛拉基亚}{}中：\BibleQuote{祂要坐下，如同熔炼金银者，炼净肋未的子孙。}又在\BibleBook{福音}{}中：\BibleQuote{你非到还清最后一文钱，决不能从那里出来。}\par

% structure: p0117 source=h2 target=subsection reason=relative-heading-rank; base=h2

\subsection{五十八、论不应因死亡而悲伤；因为生时是劳苦与危险，死时是平安与复活的确定}

在\BibleBook{}{创世纪}中：\newline{}\BibleQuote{那时，主对亚当说：因为你听从了你妻子的话，吃了我唯独命令你不可吃的那树上的果子，地必因你的所为而受诅咒；你一生日日劳苦，叹息吃食；地必给你生出荆棘和蒺藜，你也要吃田间的蔬菜，汗流满面，直到你归回你被取出的地土；因为你是土，你必归于土。}\newline{}又在同一处：\BibleQuote{哈诺客悦乐了天主，此后未见踪影，因为天主将他提去。}\newline{}在\BibleBook{}{依撒意亚}中：\BibleQuote{凡有血肉的尽是草，一切荣华如同野花。草枯干，花凋谢；但主的话永远常存。}\newline{}在\BibleBook{}{厄则克耳}中：\newline{}\BibleQuote{他们说：我们的骨头枯干了，我们的希望破灭了，我们已绝了命。为此你应预言，对他们说：主这样说：看，我要打开你们的坟墓，把你们从坟墓中领出来，领你们进入以色列地；我要把我的圣神放在你们身上，你们就会生活；我要把你们安置在你们的土地上；那时你们要知道，我主说了，也必实行——主说。}\newline{}又在\BibleBook{}{智慧篇}中：\BibleQuote{他被提去，免得邪恶改变他的悟性，因为他的灵魂悦乐了天主。}\newline{}又在第八十三篇\BibleBook{}{圣咏}中：\newline{}\BibleQuote{万军的主，你的居所多么可爱！我的灵魂渴慕且切望天主的庭院。}\newline{}在\BibleBook{得撒洛尼}{人书}中：\newline{}\BibleQuote{弟兄们，我们不愿意你们不知道关于安眠的人的事，免得你们忧伤，像那些没有望德的人一样。因为我们若信耶稣死了，也复活了，同样，那些在耶稣内安眠的人，天主也必领他们与他一同前来。}\newline{}又在\BibleBook{格林多}{前书}中：\newline{}\BibleQuote{愚昧的人！你所播下的，若不先死，决不会活过来。}\newline{}又说：\BibleQuote{星的光荣与星的光荣不同：复活也是这样。播种的是可朽坏的，复活起来的是不可朽坏的；播种的是羞辱的，复活起来的是光荣的；播种的是软弱的，复活起来的是强健的；播种的是属生灵的身体，复活起来的是属神的身体。}\newline{}又说：\BibleQuote{因为这可朽坏的必须穿上不可朽坏的，这可死的必须穿上不死。当这可朽坏的穿上了不可朽坏的，这可死的穿上了不死的时候，那时经上的话就应验了：死亡被吞灭而成为胜利。死亡啊！你的胜利在哪里？死亡啊！你的刺在哪里？}\newline{}在\BibleBook{若望}{福音}中：\newline{}\BibleQuote{圣父！我在那里，愿你所赐给我的人也与我在那里，使他们看见你赐给我的光荣，因为在创世之前，你已经爱了我。}\newline{}又在\BibleBook{路加}{福音}中：\newline{}\BibleQuote{主啊！现在照你的话，让你的仆人平安去吧；因为我亲眼看见了你的救恩。}\newline{}在\BibleBook{若望}{福音}中又说：\BibleQuote{如果你们爱我，就该喜欢我往圣父那里去，因为圣父比我大。}\par

% structure: p0119 source=h2 target=subsection reason=relative-heading-rank; base=h2

\subsection{五十九、外邦人视为神明的偶像}

在\BibleBook{}{智慧篇}中：\BibleQuote{外邦人所有的偶像，他们都当作神，这些偶像既不能用眼睛看，也不能用鼻孔呼吸，不能用耳朵听，不能用手上的指头触摸；它们的脚也走得缓慢。因为是人造了它们；那借了气息的人，塑造了它们。但没有人能造一个像自己一样的神。因为他是必死的，他用邪恶的手造出一个死物。但他本人比他所敬拜的更好，因为他确实活着，而它们从未活过。} \newline{}关于同一件事：\BibleQuote{那些观看受造物的人也不知道工匠是谁，却以为火、风、迅疾的空气、星辰的轨道、充沛的水、太阳和月亮，就是统治世界的众神；如果因这些的美丽而这样想，就让他们知道主比这些更美丽；如果他们赞叹它们的能力和运作，就让他们从这些事上看出，那建立了这些伟大事物的那一位比它们更强。} \newline{}也在\BibleBook{}{圣咏集}第一百三十四篇中：\BibleQuote{外邦人的偶像是银和金，是人手的作品。它们有口，却不能说话；有眼，却不能看；有耳，却不能听；口中也没有气息。愿制造它们的人变得像它们，凡信赖它们的也如此。} \newline{}也在\BibleBook{}{圣咏集}第九十五篇中：\BibleQuote{外邦人的众神都是魔鬼，但主创造了诸天。} \newline{}也在\BibleBook{}{出谷纪}中：\BibleQuote{你们不可为自己制造金银的神。} 又说：\BibleQuote{你不可为自己制造偶像，也不可制造任何事物的肖像。} \newline{}也在\BibleBook{}{耶肋米亚}中：\BibleQuote{上主这样说：不要照外邦人的道路行走，因为他们亲自惧怕那些事物，因为外邦人的合法之事是虚妄的。从森林砍伐的木头是木匠的作品，熔化的银和金被漂亮地布置：他们用锤子和钉子加固它们，它们就不会移动，因为它们是固定的。银子从\Place{塔尔史士}运来，金子从\Place{摩阿布}来。一切都是工匠的作品；他们要给它披上蓝色和紫色；抬起来，他们要扛着它们，因为它们不能前行。不要害怕它们，因为它们不作恶，它们里面也没有善。你们要这样说：那没有创造天地的神必从地上、从这天底下消灭。天因这事战栗，极其剧烈地震撼，上主说。我的人民行了这些恶事。他们离弃了活水的泉源，为自己挖了破裂的水池，不能蓄水。你的爱恋击打了你，你的邪恶要控告你。你要知道并看见，你离弃我是一件痛苦的事——上主你的天主说——你没有仰望我——你的主说。因为从古时你就厌恶我的轭，折断你的绳索，说：我不愿事奉，却要登上每一座高山，每一处高丘，每一棵茂密的树下：在那里我要因淫乱而蒙羞。他们对木头和石头说：你是我的父亲；对石头说：你生了我；他们向我转过背，而不是脸。} \newline{}在\BibleBook{}{依撒意亚}中：\BibleQuote{龙已坠落或消化；他们的雕刻作品变成了像走兽和牲畜。辛苦、饥饿、无力，你要将他们捆在颈上作为重担。} \newline{}又说：\BibleQuote{聚集在一起，他们也不能从战争中得救；但他们自己却与你一同被掳。} \newline{}又说：\BibleQuote{你们把我比作谁？你们看，要明白你们心里错了，你们从袋中倾倒金子，在天平上称银子，使重量合适。工匠用手制造了这些东西；他们俯伏朝拜它，把它扛在肩上，就这样行走。但如果他们把它放下来，它就会留在原地，不能移动；它不会听那些向它呼求的人；它不能救他们脱离灾祸。} \newline{}也在\BibleBook{}{耶肋米亚}中：\BibleQuote{上主创造了天地，用力量安排了世界，用智慧展开了天空，以及天空中的众水。他从地极带出云彩，从云中发出闪电；从宝藏中带出风。人人因自己的知识成为愚拙，每个工匠因自己的雕像而困惑；因为他铸造了虚谎：其中没有气息。封闭在它们里面的作品成为虚无；在它们被审察的时候必要灭亡。} \newline{}也在\BibleBook{若望}{默示录}中：\BibleQuote{第六位天使吹号。我听见金约柜四角中的一角，那是在天主面前的，对那吹号的第六位天使说：释放那捆绑在大\Place{幼发拉底河}上的四位天使。那四位天使就被释放了，他们预备好了，在某个时辰、某日、某月、某年，要杀死人类的三分之一；马兵的数量是二万万：我听见了他们的数目。然后我在异象中看见了马和骑马的，他们穿着火、风信子和硫磺的胸甲；马的头好像狮子的头；从它们口中出来火、烟和硫磺。藉这三种灾祸，人类的三分之一被杀死了，就是从它们口中出来的火、烟和硫磺，也在它们的尾巴上；因为它们的尾巴好像蛇；它们有头，并用头作恶。其余没有被这些灾祸杀死的人，也没有悔改他们手所作的工，以至于不再敬拜魔鬼和偶像，就是金、银、铜、石、木的像，这些既不能看，也不能走，他们也没有悔改他们的凶杀。} \newline{}也在同一处：\BibleQuote{第三位天使跟随他们，大声说：若有人朝拜那兽和它的像，并且在额上或手上接受了它的印号，他必要喝天主愤怒的酒，并要在圣天使和羔羊眼前受火与硫磺的刑罚；他们受痛苦的烟上升，直到永远。}\par

% structure: p0121 source=h2 target=subsection reason=relative-heading-rank; base=h2

\subsection{六十、不可贪求过度的饮食之欲}

在\BibleBook{}{依撒意亚}中：\BibleQuote{我们吃喝吧，因为明天我们就要死了。这罪直到你死也不得赦免。} 又在\BibleBook{}{出谷纪}中：\BibleQuote{百姓坐下吃喝，起来玩耍。} 保禄在\BibleBook{格林多}{前书}中：\BibleQuote{食物不能使我们得宠于天主；我们若吃，并不增加什么；我们若不吃，也并不减少什么。} 又说：\BibleQuote{你们聚集一同吃饭时，要彼此等待。谁饿了，就在家里吃，免得你们聚集受审判。} 在\BibleBook{}{罗马书}中：\BibleQuote{天主的国不在于吃喝，而在于正义、平安，以及在圣神内的喜乐。} 在\BibleBook{若望}{福音}中：\BibleQuote{我有你们不知道的食物。我的食物就是承行派遣我者的旨意，并完成他的工作。}\par

% structure: p0123 source=h2 target=subsection reason=relative-heading-rank; base=h2

\subsection{六十一、不可追求贪恋财物与金钱}

在撒罗满的\BibleBook{}{德训篇}中：\BibleQuote{爱银子的，不因银子知足。} 又在\BibleBook{}{箴言}中：\BibleQuote{屯积粮食的，必受人民咒骂；祝福却降在出售者头上。} 又在\BibleBook{}{依撒意亚}中：\BibleQuote{祸哉，那些以房接房，以地连地，以致不留余地的人！你们要独自住在地上吗？} 又在\BibleBook{}{索福尼亚}中：\BibleQuote{他们建造房屋，不得居住；栽种葡萄园，不得喝酒，因为上主的日子临近了。} 又在\BibleBook{路加}{福音}中：\BibleQuote{人纵然赚得了全世界，却丧失了自己，为他有什么益处呢？} 又说：\BibleQuote{主对他说：糊涂人！今夜就要索回你的灵魂，你所预备的一切，将归谁呢？} 又说：\BibleQuote{你该记得你生前享过福，而拉匝禄同样受过苦。如今他这里受安慰，而你却受痛苦。} 又在\BibleBook{}{宗徒大事录}中：\BibleQuote{伯多禄对他说：金银我没有；但我把我有的给你：因纳匝肋人耶稣基督的名字，起来行走吧！于是拉住他的右手，扶他起来。} 又在\BibleBook{弟茂德}{前书}中：\BibleQuote{我们带什么到世界上，也不能带走什么。只要有衣有食，就应当知足。但那些想发财的人，陷于诱惑，落入罗网和许多有害的欲望中，使人沉溺于败坏和灭亡。因为贪财是万恶的根源；有些人因贪求钱财，离弃了信德，使自己陷入许多痛苦。}\par

% structure: p0125 source=h2 target=subsection reason=relative-heading-rank; base=h2

\subsection{六十二、不可与外邦人联姻}

在\BibleBook{}{多俾亚传}中：\BibleQuote{从你祖宗的女儿中娶一女子，不要娶不属于你支派的外族女子。} 又在\BibleBook{}{创世纪}中：亚巴郎派仆人为儿子依撒格从本族中娶了黎贝加。又在\BibleBook{}{厄斯德拉}中：当犹太人被掳掠时，天主仍不满意，除非他们抛弃外族妻子以及她们所生的儿女。又在保禄的\BibleBook{格林多}{前书}中：\BibleQuote{妻子活着的时候是被束缚的；但如果丈夫死了，她就自由了，可以随意嫁人，只要是在主内。但若照常守节，则更为有福。} 又说：\BibleQuote{你们不知道你们的身体是基督的肢体吗？我可以把基督的肢体取来，作为娼妓的肢体吗？断乎不可！你们不知道那与娼妓结合的，便是与她成为一体吗？因为经上说：\InnerQuote{二人成为一体。}但那与主结合的，便是与主成为一灵。} 又在\BibleBook{格林多}{后书}中：\BibleQuote{不要与不信者同负一轭。正义与不义有什么相通？光明与黑暗有什么联系？} 论到撒罗满，在\BibleBook{列王纪}{上}中：\BibleQuote{外族女子使他的心偏离，追随她们的神。}\par

% structure: p0127 source=h2 target=subsection reason=relative-heading-rank; base=h2

\subsection{六十三、奸淫之罪极为严重}

在保禄的\BibleBook{格林多}{前书}中：\BibleQuote{人所犯的一切罪都在身体以外，但行淫的人却是得罪自己的身体。你们不是自己的，因为你们是重价买来的。所以要在你们身上光荣并承载主。}\par

% structure: p0129 source=h2 target=subsection reason=relative-heading-rank; base=h2

\subsection{六十四、哪些肉性之事导致死亡，哪些灵性之事引向生命}

保禄在\BibleBook{}{迦拉达书}中：\BibleQuote{肉性欲与圣神相争，圣神与肉性欲相争；二者彼此敌对，致使你们不能做你们所愿意的事。肉性的作为是显而易见的：奸淫、邪淫、不洁、放荡、崇拜偶像、行邪术、谋杀、仇恨、争斗、嫉妒、忿怒、自私、分党、异端、妒忌、醉酒、宴乐，以及类似的事。我预先警告你们：行这样事的人决不能承受天主的国。然而圣神的果实是：爱德、喜乐、平安、忍耐、良善、信德、温和、节制、贞洁。凡属基督的人，已把肉身连同邪情私欲钉在十字架上了。}\par

% structure: p0131 source=h2 target=subsection reason=relative-heading-rank; base=h2

\subsection{六十五、一切罪恶在圣洗中被洗净}

在保禄的\BibleBook{格林多}{前书}中：\BibleQuote{无论是行淫的、拜偶像的、奸淫的、作娈童的、亲男色的、偷窃的、欺诈的、醉酒的、辱骂的、勒索的，都不能承受天主的国。你们中从前也曾这样；但如今你们已经洗净了，已经圣化了，因我们主耶稣基督的名，并因我们天主的圣神。}\par

% structure: p0133 source=h2 target=subsection reason=relative-heading-rank; base=h2

\subsection{六十六、在教会规诫中应遵守天主的纪律}

在\BibleBook{耶肋米亚}{书}中：\Quote{我要将合我心意的牧者赐给你们；他们必用知识和明智牧养你们。} 又在\BibleBook{箴言}{篇}中：\Quote{我儿，不可轻视天主的管教，也不可厌烦他的谴责，因为天主谴责他所爱的。} 又在\BibleRef{咏 2}中：\Quote{你们要接受管教，免得上主发怒，你们便从正路上灭亡，因为他的怒火快要燃起。凡投靠他的人真是有福。} 又在\BibleRef{咏 49}中：\Quote{天主却对恶人说：你怎敢讲述我的诫命，你的口怎敢提及我的盟约？你憎恶管教，将我的话抛在脑后。} 又在\BibleBook{智慧}{篇}中：\Quote{抛弃管教的人是有祸的。}\par

% structure: p0135 source=h2 target=subsection reason=relative-heading-rank; base=h2

\subsection{六十七、预言世人将轻视健全的管教}

保禄在\BibleBook{弟茂德}{后书}中：\Quote{那时人不再容忍健全的道理；反而依从自己的私欲，为自己聚拢许多师傅，耳朵发痒，随从自己的情欲，并且掩耳不听真理，偏向无稽的传说。}\par

% structure: p0137 source=h2 target=subsection reason=relative-heading-rank; base=h2

\subsection{六十八、必须远离生活放荡、违反纪律的人}

保禄在\BibleBook{得撒洛尼}{后书}中：\Quote{我们曾因主耶稣基督的名命令你们，要远离一切游手好闲、不按从我们所领受的传统生活的弟兄。} 又在\BibleRef{咏 49}中：\Quote{你见了盗贼，就与他同伙，又与奸夫同分赃物。}\par

% structure: p0139 source=h2 target=subsection reason=relative-heading-rank; base=h2

\subsection{六十九、天主的国不在于世上的智慧，也不在于口才，而在于对十字架的信德和生活的德能}

保禄在\BibleBook{格林多}{前书}中：\Quote{基督派遣我，不是为施洗，而是为宣传福音，且不用智慧的言辞，免得基督的十字架失去效力。原来十字架的道理，为丧亡的人是愚妄，但为我们得救的人，却是天主的德能。因为经上记载：\InnerQuote{我要摧毁智者的智慧，废弃贤者的聪明。}智者在哪里？经师在哪里？今世的辩士在哪里？天主岂不是使今世的智慧变成了愚妄吗？因为世人凭自己的智慧，未能认识天主，天主就乐意用愚妄的道理来拯救那些相信的人。犹太人要求神迹，希腊人寻求智慧，而我们所宣讲的却是被钉在十字架上的基督：这为犹太人确是绊脚石，为外邦人是愚妄，但为那些蒙召的人，不拘是犹太人或希腊人，基督却是天主的德能和天主的智慧。} 又说：\Quote{谁也不要自欺。你们中若有人自以为在今世是聪明的，便该变为愚妄，好成为真正的智者，因为今世的智慧在天主面前是愚妄。经上记载：\InnerQuote{上主责斥智者的计谋都是枉然。}} 又说：\Quote{上主知道智者的思念，原来都是虚幻。}\par

% structure: p0141 source=h2 target=subsection reason=relative-heading-rank; base=h2

\subsection{七十、必须服从父母}

保禄在\BibleBook{厄弗所}{书}中：\Quote{孩子们，你们要服从父母，这是理所当然的。\InnerQuote{孝敬你的父亲和母亲}——这是第一条带恩许的诫命——为使你得到幸福，并在地上延年益寿。}\par

% structure: p0143 source=h2 target=subsection reason=relative-heading-rank; base=h2

\subsection{七十一、父亲也不应对儿女过于严厉}

在同一处：\Quote{你们作父亲的，不要惹你们的子女发怒，但要用主的管教和劝诫养育他们。}\par

% structure: p0145 source=h2 target=subsection reason=relative-heading-rank; base=h2

\subsection{七十二、仆人信主后应更忠心地服侍肉身的主人}

保禄在\BibleBook{厄弗所}{书}中：\Quote{你们作仆人的，要怀着恐惧和战栗，以诚实的心，听从你们肉身的主人，如同听从基督一样。不要只在人眼前服事，像是讨人喜欢的，但该像基督的仆人，从心里遵行天主的旨意。}\par

% structure: p0147 source=h2 target=subsection reason=relative-heading-rank; base=h2

\subsection{七十三、主人也应更加温和}

在同一处：\Quote{你们作主人的，也要照样对待他们，不可恐吓他们，因为你们知道，他们和你们的主人在天上，而且他不看人的情面。}\par

% structure: p0149 source=h2 target=subsection reason=relative-heading-rank; base=h2

\subsection{七十四、所有经得起考验的寡妇都当受尊敬}

保禄在\BibleBook{弟茂德}{前书}中：\Quote{要尊敬那些真正守寡的寡妇。至于那放荡的寡妇，虽生犹死。} 又说：\Quote{至于年轻的寡妇，你要拒绝，因为她们一旦情欲冲动，违背基督，就想再嫁；她们犯了罪，因为她们背弃了起初的信德。}\par

% structure: p0151 source=h2 target=subsection reason=relative-heading-rank; base=h2

\subsection{七十五、每人应先照顾自己的人，尤其是信徒}

宗徒在\BibleBook{弟茂德}{前书}中：\Quote{若有人不照顾自己的亲属，尤其是家人，便是背弃信德，比不信的人更坏。} 关于同一事，在\BibleBook{依撒意亚}{书}中：\Quote{你若看见赤身露体的，给他衣服穿；不要轻忽你的骨肉之亲。} 关于这些家人，福音中说：\Quote{若有人称家主为贝耳则步，何况他的家人呢？}\par

% structure: p0153 source=h2 target=subsection reason=relative-heading-rank; base=h2

\subsection{七十六、不可轻率控告长老}

在\BibleBook{弟茂德}{前书}中：\Quote{不可受理对长老的控告。}\par

% structure: p0155 source=h2 target=subsection reason=relative-heading-rank; base=h2

\subsection{七十七、罪人应公开受责}

保禄在\BibleBook{弟茂德}{前书}中：\Quote{你在众人面前责备那些犯罪的，使其余的人有所警戒。}\par

% structure: p0157 source=h2 target=subsection reason=relative-heading-rank; base=h2

\subsection{七十八、不可与异端者交谈}

致\BibleBook{弟铎}{书}：\Quote{对异端者，在一次警告以后，就该远离，因为你知道这种人已走入歧途，犯罪作恶，自定罪案。} 关于同一事，在\BibleBook{若望}{一书}中：\Quote{他们从我们中间出去，但不是属于我们的；如果属于我们，必会留在我们中间。} 又在\BibleBook{弟茂德}{后书}中：\Quote{他们的话如同毒疮越烂越大。}\par

% structure: p0159 source=h2 target=subsection reason=relative-heading-rank; base=h2

\subsection{七十九、清白的人能放心祈求并得到}

在\BibleBook{若望}{一书}中：\Quote{如果我们的心不责备我们，我们在天主面前便可放心大胆；无论我们求什么，必从他获得。} 又在\BibleBook{玛窦}{福音}中：\Quote{心里洁净的人是有福的，因为他们要看见天主。} 又在\BibleRef{咏 23}中：\Quote{谁能登上上主的山？谁能站在他的圣所？是那手洁心清的人。}\par

% structure: p0161 source=h2 target=subsection reason=relative-heading-rank; base=h2

\subsection{八十、魔鬼除非得天主准许，否则无权加害于人}

在\BibleBook{若望}{福音}中记载：\BibleQuote{\Person{耶稣}说：若不是从上头赐给你的，你对我毫无权柄。}又在\WorkTitle{列王纪上}中记载：\BibleQuote{\Person{天主}又激起\Person{撒殚}来攻击\Person{撒罗满}本人。}在\WorkTitle{约伯传}中，天主先准许，然后才允许魔鬼行事；在福音中，主先准许，对\Person{犹达斯}说：\BibleQuote{你要做的事，快做吧。}又在\Person{撒罗满}的\WorkTitle{箴言}中记载：\BibleQuote{君王的心在天主手中。}\par

% structure: p0163 source=h2 target=subsection reason=relative-heading-rank; base=h2

\subsection{八十一、雇工的工资应当迅速支付}

在\WorkTitle{肋未纪}中记载：\BibleQuote{你不可将雇工的工资留过夜到早晨。}\par

% structure: p0165 source=h2 target=subsection reason=relative-heading-rank; base=h2

\subsection{八十二、不可使用占卜}

在\WorkTitle{申命纪}中记载：\BibleQuote{不可使用预兆和征兆。}\par

% structure: p0167 source=h2 target=subsection reason=relative-heading-rank; base=h2

\subsection{八十三、不可在头上留绺发}

在\WorkTitle{肋未纪}中记载：\BibleQuote{不可从你的头发中留绺。}\par

% structure: p0169 source=h2 target=subsection reason=relative-heading-rank; base=h2

\subsection{八十四、不可拔胡须}

\BibleQuote{不可毁损你胡须的形状。}\par

% structure: p0171 source=h2 target=subsection reason=relative-heading-rank; base=h2

\subsection{八十五、当主教或司铎来到时，我们应当起立}

在\WorkTitle{肋未纪}中记载：\BibleQuote{在长者面前要起立，并要尊敬司铎的尊严。}\par

% structure: p0173 source=h2 target=subsection reason=relative-heading-rank; base=h2

\subsection{八十六、不可制造分裂，即使脱离者仍持守同一信德和同一传统}

在\WorkTitle{德训篇}，即\Person{撒罗满}的作品中记载：\BibleQuote{劈柴的人若铁器脱落，必因此而处于危险。}又在\WorkTitle{出谷纪}中记载：\BibleQuote{应在一间屋内吃，不可将肉块带到屋外。}又在\BibleBook{}{圣咏}中记载：\BibleQuote{看，兄弟们同居共处，多么美好，多么快乐！}又在\BibleBook{玛窦}{福音}中记载：\BibleQuote{不与我结合的，就是反对我；不与我聚集的，就是分散。}又在\Person{保禄}致\BibleBook{格林多}{前书}中记载：\BibleQuote{弟兄们，我因我们的主\Person{耶稣基督}之名，劝你们大家言谈一致，不要有分裂，但要同心合意，思想一致。}又在\BibleBook{}{圣咏}中记载：\BibleQuote{天主使同居共处的人一心一意。}\par

% structure: p0175 source=h2 target=subsection reason=relative-heading-rank; base=h2

\subsection{八十七、信徒应当纯洁，兼有机警}

在\BibleBook{玛窦}{福音}中记载：\BibleQuote{机警如蛇，纯洁如鸽。}又说：\BibleQuote{你们是地上的盐。盐若失了味，可用什么使它再咸呢？它毫无用处，只好丢在外面，任人践踏。}\par

% structure: p0177 source=h2 target=subsection reason=relative-heading-rank; base=h2

\subsection{八十八、不可欺骗弟兄}

在\Person{保禄}致\BibleBook{得撒洛尼}{前书}中记载：\BibleQuote{不可在这事上欺骗他的弟兄，因为天主是这一切的报应者。}\par

% structure: p0179 source=h2 target=subsection reason=relative-heading-rank; base=h2

\subsection{八十九、世界的结局突然来临}

宗徒说：\BibleQuote{主的日子要像夜间的盗贼一样来到。当人们说\InnerQuote{平安无事}时，毁灭就突然临到他们身上。}又在\WorkTitle{宗徒大事录}中记载：\BibleQuote{没有人能知道父以自己的权柄所定的时期和日期。}\par

% structure: p0181 source=h2 target=subsection reason=relative-heading-rank; base=h2

\subsection{九十、妻子不可离开丈夫；如果离开，应当保持独身}

在\Person{保禄}致\BibleBook{格林多}{前书}中记载：\BibleQuote{至于已婚者，我命令——其实不是我，而是主命令：妻子不可离开丈夫；如果离开，应保持独身，或与丈夫和好；丈夫也不可休妻。}\par

% structure: p0183 source=h2 target=subsection reason=relative-heading-rank; base=h2

\subsection{九十一、每人所受的诱惑不超过他能承受的}

在\Person{保禄}致\BibleBook{格林多}{前书}中记载：\BibleQuote{你们所受的诱惑，无非是普通人所能受的；天主是忠信的，他决不许你们受诱惑超过你们所能承担的，而且在受诱惑时，他必给开一条出路，使你们能忍受。}\par

% structure: p0185 source=h2 target=subsection reason=relative-heading-rank; base=h2

\subsection{九十二、并非所有合法之事都可去做}

\Person{保禄}在致\BibleBook{格林多}{前书}中记载：\BibleQuote{凡事都可行，但不全有益；凡事都可行，但不全有建树。}\par

% structure: p0187 source=h2 target=subsection reason=relative-heading-rank; base=h2

\subsection{九十三、预言将有异端兴起}

在\Person{保禄}致\BibleBook{格林多}{前书}中记载：\BibleQuote{在你们中间难免有异端，为使那些经得起考验的人显明出来。}\par

% structure: p0189 source=h2 target=subsection reason=relative-heading-rank; base=h2

\subsection{九十四、应以敬畏和尊荣领受圣体圣事}

在\WorkTitle{肋未纪}中记载：\BibleQuote{凡灵魂若吃了属于主的救恩祭肉，而他身上仍有不洁，那灵魂应从民间剪除。}又在致\BibleBook{格林多}{前书}中记载：\BibleQuote{谁若不相称地吃主的饼，喝主的杯，就是冒犯主的身体和主的血。}\par

% structure: p0191 source=h2 target=subsection reason=relative-heading-rank; base=h2

\subsection{九十五、我们应与善人共处，远离恶人}

在\Person{撒罗满}的\WorkTitle{箴言}中记载：\BibleQuote{不可将不虔敬的人带入义人的居所。}又在同一作者的作品，即\WorkTitle{德训篇}中记载：\BibleQuote{愿义人作你的宾客。}又说：\BibleQuote{忠实的朋友是生命和不朽的良药。}又在同一处记载：\BibleQuote{远离那有权杀害的人，你就不必怀疑恐惧。}又在同一处记载：\BibleQuote{找到真朋友，且向静听的耳朵说正义话的人，是有福的。}又在同一处记载：\BibleQuote{用荆棘围住你的耳朵，不要听恶舌。}又在\BibleBook{}{圣咏}中记载：\BibleQuote{你与善人，必显为善；与正直者，必显正直；与乖戾者，必显乖戾。}又在\Person{保禄}致\BibleBook{格林多}{前书}中记载：\BibleQuote{恶友败坏善德。}\par

% structure: p0193 source=h2 target=subsection reason=relative-heading-rank; base=h2

\subsection{九十六、我们应当用行动而非言语劳作}

在\Person{撒罗满}的著作中，在\BibleBook{德训}{篇}中：\Quote{你的舌头不可急躁，你的行为不可无用而懈怠。}又\Person{保禄}在\BibleBook{格林多}{前书}中：\Quote{天主的国不在于言词，而在于能力。}又\BibleBook{罗马}{书}：\Quote{不是听法律的在天主面前成为义人，而是行法律的才被成义。}又\BibleBook{玛窦}{福音}：\Quote{谁若实行并教导这样的诫命，他将被称为天国中最大的。}又同处：\Quote{凡听了我这些话而实行的，我要把他比作一个聪明人，他把房屋建在磐石上。雨降下，洪水来了，风吹来，冲击那房屋，它没有倒塌，因为它建在磐石上。凡听了我这些话而不实行的，我要把他比作一个糊涂人，他把房屋建在沙土上。雨降下，洪水来了，风吹来，撞击那房屋，它就倒塌了，且倒塌得很厉害。}\par

% structure: p0195 source=h2 target=subsection reason=relative-heading-rank; base=h2

\subsection{九十七、必须尽快趋向信德和达到目标}

在\Person{撒罗满}的著作中，在\BibleBook{德训}{篇}中：\Quote{不要迟延归向天主，也不要一日一日地拖延，因为祂的愤怒会突然来到。}\par

% structure: p0197 source=h2 target=subsection reason=relative-heading-rank; base=h2

\subsection{九十八、望教者如今不应再犯罪}

在\Person{保禄}致\BibleBook{罗马}{书}中：\Quote{让我们作恶以成就善事吧！——他们的定罪是公道的。}\par

% structure: p0199 source=h2 target=subsection reason=relative-heading-rank; base=h2

\subsection{九十九、审判将按照时代，或为法律之前的公平时代，或为梅瑟之后的法律时代}

\Person{保禄}致\BibleBook{罗马}{书}：\Quote{凡没有法律而犯罪的人，也必不带法律而丧亡；凡在法律之下犯罪的人，也必按法律受审判。}\par

% structure: p0201 source=h2 target=subsection reason=relative-heading-rank; base=h2

\subsection{100、天主的恩宠应当是无偿的}

在\BibleBook{宗徒}{大事录}中：\Quote{愿你的银子与你一同丧亡，因为你以为天主的恩宠可以用金钱获得。}又在福音中：\Quote{你们白白地领受了，也要白白地施与。}又同处：\Quote{你们使我父的殿成了商场，又使祈祷的殿成了贼窝。}又在\BibleBook{}{依撒意亚}中：\Quote{凡口渴的，请到水边来！凡没有钱的，也请来，买来喝，不用钱。}又在\BibleBook{}{默示录}中：\Quote{我是\OriginalTerm{阿耳法}{Alpha}和\OriginalTerm{敖默加}{Omega}，元始和终末。我要把生命之水白白地赐给口渴的人喝。胜利者必要承受这些产业，我要作他的天主，他要作我的儿子。}\par

% structure: p0203 source=h2 target=subsection reason=relative-heading-rank; base=h2

\subsection{101、圣神屡次以火显现}

在\BibleBook{出谷}{纪}中：\Quote{整个西乃山冒烟，因为上主在火中降临其上。}又在\BibleBook{宗徒}{大事录}中：\Quote{忽然从天上来了一阵响声，好像暴风刮来，充满了他们所在的那座房屋。又有如火的舌头显现出来，分开停留在他们每人头上，他们就都充满了圣神。}又在祭献中：凡天主所悦纳的，就有火从天上降下，焚尽了祭品。在\BibleBook{出谷}{纪}中：\Quote{上主的使者从荆棘的火焰中显现出来。}\par

% structure: p0205 source=h2 target=subsection reason=relative-heading-rank; base=h2

\subsection{102、所有善人都应乐意接受责备}

在\Person{撒罗满}的著作\BibleBook{}{箴言}中：\Quote{责备恶人的人必被憎恨；责备明智人，他必爱你。}\par

% structure: p0207 source=h2 target=subsection reason=relative-heading-rank; base=h2

\subsection{103、必须戒除多言}

在\Person{撒罗满}的著作中：\Quote{多言必不免犯罪；约束嘴唇的，便是明智。}\par

% structure: p0209 source=h2 target=subsection reason=relative-heading-rank; base=h2

\subsection{104、不可说谎}

\Quote{说谎的嘴唇为上主所憎恶。}\par

% structure: p0211 source=h2 target=subsection reason=relative-heading-rank; base=h2

\subsection{105、对在家庭本分上犯错的人应常常纠正}

在\Person{撒罗满}的著作中：\Quote{不用棍杖的，是憎恨儿子。}又：\Quote{不可停止纠正孩童。}\par

% structure: p0213 source=h2 target=subsection reason=relative-heading-rank; base=h2

\subsection{106、受伤害时应保持忍耐，将报复交给天主}

不要说：我要向我的敌人报仇；但要等候上主，祂必帮助你。又别处：\Quote{报复属于我，我必要偿还，上主说。}又在\BibleBook{}{索福尼亚}中：\Quote{你们等候我，上主说，在我起来作证的日子；因为我的判断是聚集外邦人，使我夺取列王，在他们身上倾倒我的愤怒。}\par

% structure: p0215 source=h2 target=subsection reason=relative-heading-rank; base=h2

\subsection{107、不可诽谤}

在\Person{撒罗满}的著作\BibleBook{}{箴言}中：\Quote{不要爱诽谤，免得你被除去。}又在\BibleBook{}{圣咏集}第四十九篇中：\Quote{你坐着，毁谤你的兄弟，对你母亲的儿子设置绊脚石。}又在\Person{保禄}致\BibleBook{哥罗森}{书}中：\Quote{不要毁谤任何人，也不要争讼。}\par

% structure: p0217 source=h2 target=subsection reason=relative-heading-rank; base=h2

\subsection{108、不可对邻舍设陷阱}

在\Person{撒罗满}的著作\BibleBook{}{箴言}中：\Quote{为邻舍挖坑的，自己必掉进坑中。}\par

% structure: p0219 source=h2 target=subsection reason=relative-heading-rank; base=h2

\subsection{109、应探望病人}

在\Person{撒罗满}的著作中，在\BibleBook{德训}{篇}中：\Quote{不要懈怠探望病人，因为从这些事上，你将在爱德中得以坚固。}又在福音中：\Quote{我病了，你们来探望了我；我在监里，你们来看了我。}\par

% structure: p0221 source=h2 target=subsection reason=relative-heading-rank; base=h2

\subsection{110、传舌者应受诅咒}

在\BibleBook{德训}{篇}，\Person{撒罗满}的著作中：\Quote{传舌和一口两舌的人是可诅咒的，因为他要扰乱许多和平的人。}\par

% structure: p0223 source=h2 target=subsection reason=relative-heading-rank; base=h2

\subsection{111、恶人的祭献不蒙悦纳}

同处：\Quote{至高者不悦纳不义之人的礼物。}\par

% structure: p0225 source=h2 target=subsection reason=relative-heading-rank; base=h2

\subsection{112、在今世拥有更大权柄的人，受审判更为严厉}

在\Person{撒罗满}的著作中：\Quote{掌权的人将受最严厉的审判；因为卑微的人获得怜悯，但有权势的人必受强烈的痛苦。}又在\BibleBook{}{圣咏集}第二篇：\Quote{如今，你们君王，要慎思明辨；你们审判大地的，应受惩戒。}\par

% structure: p0227 source=h2 target=subsection reason=relative-heading-rank; base=h2

\subsection{113、应保护寡妇和孤儿}

在撒罗满的著作中：\BibleQuote{你们要像父亲一样怜悯孤儿，像丈夫一样对待他们的母亲；如果你们听从，你们将成为至高者的儿子。}又是在\BibleBook{}{出谷纪}中：\BibleQuote{你们不可欺压任何寡妇和孤儿。但如果你们欺压他们，他们向我哀号呼求，我必垂听他们的哀声，并要发怒攻击你们；我要用刀剑消灭你们，使你们的妻子成为寡妇，你们的儿女成为孤儿。}又是在\BibleBook{}{依撒意亚}中：\BibleQuote{你们要替孤儿伸冤，为寡妇辩护；然后来，我们辩论，上主说。}又是在\BibleBook{}{约伯传}中：\BibleQuote{我曾从强暴者的手中拯救穷人，帮助那无依无靠的孤儿；寡妇的口曾祝福我。}又是在\BibleBook{圣咏}{第六十七篇}中：\BibleQuote{孤儿的父亲，寡妇的判官。}\par

% structure: p0229 source=h2 target=subsection reason=relative-heading-rank; base=h2

\subsection{114、人应当在有肉身时告解}

在\BibleBook{圣咏}{第五篇}中：\BibleQuote{但在坟墓中，谁会向你宣认呢？}又是在\BibleBook{圣咏}{第二十九篇}中：\BibleQuote{尘土岂能向你宣认呢？}又是在别处论及告解：\BibleQuote{我宁愿罪人悔改，而不愿他死亡。}又是在\BibleBook{}{耶肋米亚}中：\BibleQuote{上主这样说：跌倒的岂能不起来？转身离去的岂能不再归回？}\par

% structure: p0231 source=h2 target=subsection reason=relative-heading-rank; base=h2

\subsection{115、谄媚是有害的}

在\BibleBook{}{依撒意亚}中：\BibleQuote{称你为有福的人，是引领你走入歧途，扰乱你脚下的道路。}\par

% structure: p0233 source=h2 target=subsection reason=relative-heading-rank; base=h2

\subsection{116、天主更爱那在圣洗中许多罪过得到赦免的人}

在\BibleBook{路加}{福音}中：\BibleQuote{那得了多赦免的，爱得多；那得了少赦免的，爱得少。}\par

% structure: p0235 source=h2 target=subsection reason=relative-heading-rank; base=h2

\subsection{117、必须与魔鬼进行激烈的战斗，因此我们应勇敢站立，好能得胜}

在\BibleBook{}{厄弗所书}中：\BibleQuote{我们的战斗不是对抗血肉，而是对抗这个黑暗世界的掌权者和统治者，对抗天上邪恶的属灵势力。因此，要穿上天主所赐的全副武装，使你们能在邪恶的日子抵挡得住，并在完成一切之后，仍能站立得住。要以真理的福音束腰，穿上正义的甲，脚上穿着和平福音的准备；又拿起信德的盾牌，用以扑灭恶者的一切火箭；并戴上救恩的头盔，拿起圣神的剑，就是天主的话。}\par

% structure: p0237 source=h2 target=subsection reason=relative-heading-rank; base=h2

\subsection{118、论假基督，他将以人的形态出现}

在\BibleBook{}{依撒意亚}中：\BibleQuote{这就是那震动大地，惊扰列王，使全地变为荒野的人。}\par

% structure: p0239 source=h2 target=subsection reason=relative-heading-rank; base=h2

\subsection{119、法律的轭是沉重的，为我们所摆脱；主的轭是轻松的，为我们所负担}

在\BibleBook{圣咏}{第二篇}中：\BibleQuote{为什么异民激动，万民冥想虚妄的事？地上的君王起来，首领聚在一起，反对上主和祂的受傅者。让我们挣断他们的捆绑，抛弃他们的轭。}又是在\BibleBook{玛窦}{福音}中：\BibleQuote{凡劳苦和负重担的，你们都到我这里来，我要使你们安息。你们背起我的轭，跟我学习，因为我是良善心谦的；这样你们必得心安。因为我的轭是甘甜的，我的担子是轻松的。}又是在\BibleBook{}{宗徒大事录}中：\BibleQuote{圣神和我们决定，不把别的重担加给你们，只有这几件事是必须的：即要禁绝偶像崇拜、血和奸淫。你们不愿意人怎样待你们，你们也不要怎样待人。}\par

% structure: p0241 source=h2 target=subsection reason=relative-heading-rank; base=h2

\subsection{120、我们应恒切祈祷}

在\BibleBook{}{哥罗森书}中：\BibleQuote{你们要恒切祈祷，在此警醒。}又是在\BibleBook{圣咏}{第一篇}中：\BibleQuote{他喜爱上主的法律，昼夜默思祂的法律。}\par

\SourceRecord{Translated by Robert Ernest Wallis.；Ante-Nicene Fathers；Vol. 5.；Edited by Alexander Roberts, James Donaldson, and A. Cleveland Coxe.；Buffalo, NY: Christian Literature Publishing Co.,；1886.；<http://www.newadvent.org/fathers/050712c.htm>.}
